\documentclass[11pt]{amsart}
\usepackage{amsmath,amssymb,amsthm,mathtools,mathrsfs}
\usepackage[margin=1.08in]{geometry}
\usepackage{booktabs}
\usepackage{array}
\usepackage{longtable}
\usepackage{tabularx}
\usepackage[hypertexnames=false]{hyperref}
\hypersetup{hidelinks}
\usepackage[final,expansion=false]{microtype}
\usepackage{url}
\setlength{\parskip}{4pt plus 1pt}
\emergencystretch=2.5em

\newtheorem{theorem}{Theorem}[section]
\newtheorem{proposition}[theorem]{Proposition}
\newtheorem{lemma}[theorem]{Lemma}
\newtheorem{corollary}[theorem]{Corollary}
\theoremstyle{definition}
\newtheorem{definition}[theorem]{Definition}
\newtheorem{problem}[theorem]{Problem}
\theoremstyle{remark}
\newtheorem{remark}[theorem]{Remark}

\newcommand{\C}{\mathbb C}
\newcommand{\R}{\mathbb R}
\newcommand{\N}{\mathbb N}
\newcommand{\Z}{\mathbb Z}
\newcommand{\Q}{\mathbb Q}
\newcommand{\Rel}{\operatorname{Re}}
\newcommand{\Iml}{\operatorname{Im}}
\newcommand{\supp}{\operatorname{supp}}
\newcommand{\TV}{\operatorname{TV}}
\newcommand{\Kp}{\mathcal K^{\mathrm{prim}}_{2T}}
\newcommand{\Kw}{\mathcal K_{2T}}
\newcommand{\PT}{\mathcal P_T}
\newcommand{\IT}{I_T}
\newcommand{\lamP}{\lambda^{\mathrm{prim}}_T}
\newcommand{\lamW}{\lambda^{\mathrm W}_T}
\newcommand{\lamA}{\lambda_{\min}(A_T|\PT)}
\newcommand{\sT}{\sigma_T}
\newcommand{\gh}[1]{\widehat g_{#1}}
\newcommand{\nuh}{\widehat\nu}
\newcommand{\wh}{\widehat w}
\newcommand{\Fh}{\widehat F}
\newcommand{\cF}{c_F}
\newcommand{\Lev}{\mathfrak E}
\newcommand{\Wcl}{\mathcal W_T}
\newcommand{\NT}{\mathcal N_T}
\newcommand{\Aconst}{\mathsf A_T}
\newcommand{\Ainf}{\mathcal A_\infty}
\newcommand{\PiT}{\Pi_T}
\newcommand{\laminf}{\lambda^{\infty}_{\min}}
\newcommand{\laminfP}{\lambda^{\infty,\mathrm P}_{\min}}
\newcolumntype{Y}{>{\raggedright\arraybackslash}X}

\begin{document}

\title[A dual theory of Weil positivity at finite support]
{A Dual Theory of Weil Positivity at Finite Support:\\
Outside-Window Certificates, Non-Attainment,\\
and the Explicit Formula as the Dual Variable}

\author{David Alejandro Trejo Pizzo}
\thanks{Independent Researcher. \texttt{dtrejopizzo@gmail.com}. The certificate data of \S\ref{sec:cert} accompany the paper.}

\date{\today}


\begin{abstract}
Let $W_T$ be Weil's quadratic form on $L^2(-T,T)$ and $\mathfrak a_T$ its restriction to the
primitive subspace, on which the two Tate moments vanish; Weil's criterion says that the Riemann
Hypothesis is equivalent to $W_T\ge0$ for every $T$. We develop the dual side of this problem. The
archimedean multiplier is a L\'evy exponent with an explicit completely monotone density, and the
bottom of $\mathfrak a_T$ is bounded below by the value of a semi-infinite linear program whose dual
variables are finite measures supported outside $[-2T,2T]$. We prove strong duality with finitely
atomic dual sequences, show that the dual is never attained, and give tail lemmas that reduce a
certificate to a bounded interval; an explicit certificate proves $\lambda_{\min}\ge8.8577\cdot10^{-5}$
at $T=\frac12\log3$. We then show that every certificate of the complete form is, up to an arbitrarily
small correction, the explicit formula evaluated on modulated flat-top test functions, so that its
value is a lower bound for a smoothing $\sum_\rho\widehat w(\tau-\gamma_\rho)$ of the zeros; the
Riemann Hypothesis is equivalent to the existence of such a nonnegative smoothing at every window, and
to the indeterminacy of a Kre\u{\i}n extension problem at every window, while no finite set of zeros
carries one. Further results concern the limits of lattice certificates, the negative index of the
archimedean form (computed exactly for $T\le0.90$ away from two short intervals), the decay of the
margin as $T\to\infty$, the exact relation between the primitive and the complete form at a single
window, and a finite sum--product identity with a faithful Gamma-moment detector.  The remaining
moment-annihilation statement and the first-contact bridge are stated explicitly and are not proved.
\end{abstract}

\maketitle
\tableofcontents

\section{Introduction}

\subsection{The problem}
Throughout, $T>0$, $\IT=(-T,T)$, and for $F\in L^2(\IT)$ (extended by zero) we write
$\Fh(\tau)=\int F(t)e^{i\tau t}\,dt$ and $\cF(u)=\int F(t)\overline{F(t-u)}\,dt$ for the
autocorrelation, so that $\cF$ is continuous, supported in $[-2T,2T]$, and
$\widehat{\cF}=|\Fh|^2$ on $\R$. Put
\begin{equation}\label{eq:symbol}
m(\tau)=\Rel\psi\!\left(\tfrac14+\tfrac{i\tau}{2}\right)-\log\pi,\qquad
\sT(\tau)=m(\tau)-\sum_{n\in\NT}v_n\cos(\tau\log n),\qquad v_n=\frac{2\Lambda(n)}{\sqrt n},
\end{equation}
where $\NT=\{n\ge2:\ \log n<2T\}$ is the set of \emph{contacts} and
$\Aconst:=\sum_{n\in\NT}v_n$ is the total mass of the prime comb. The truncated Weil form is
\begin{equation}\label{eq:form}
\mathfrak a_T(F)=\frac1{2\pi}\int_\R\sT(\tau)\,|\Fh(\tau)|^2\,d\tau ,
\end{equation}
defined on the form domain $\{F\in L^2(\IT):\int\log(2+|\tau|)|\Fh|^2<\infty\}$; it is the
quadratic form of the operator $A_T$ of \cite{Paper42}, namely the archimedean multiplier minus
$\sum_{n\in\NT}w_n(S_{\log n}+S_{-\log n})$, $w_n=\Lambda(n)/\sqrt n$, compressed to $\IT$. The two
\emph{Tate moments} are $A(F)=\int F(t)e^{t/2}dt$ and $B(F)=\int F(t)e^{-t/2}dt$, and the
\emph{primitive subspace} is $\PT=\{F\in L^2(\IT):A(F)=B(F)=0\}$. The \emph{complete Weil form} on
the window is
\begin{equation}\label{eq:W}
W_T(F)=\mathfrak a_T(F)+2\,\Rel\big(A(F)\overline{B(F)}\big),
\end{equation}
which by the explicit formula equals $\sum_\rho\Phi_F(\rho)\Phi_F(1-\rho)$,
$\Phi_F(s)=\int F(t)e^{(s-\frac12)t}dt$, the sum over the nontrivial zeros of $\zeta$
(\cite{Weil52}, \cite[Thm.~8.4]{Paper45}; for tests that are not smooth, under absolute convergence of
the zero sum, Lemma~\ref{lem:roughEF}). Weil's criterion says that the Riemann Hypothesis (RH) is
equivalent to $W_T\ge0$ for every $T$, and primitive tests suffice (\cite[Thm.~11.3]{Paper45}; see
also \cite[Thm.~1]{Bombieri2000}):
\[
\mathrm{RH}\iff \forall T:\ W_T\ge0 \ \text{on }L^2(\IT)\iff \forall T:\ \mathfrak a_T\ge0\ \text{on }\PT .
\]
Write $\lamA$ for the bottom of $\mathfrak a_T$ on the unit sphere of $\PT$ and $\lamW$ for that
of $W_T$ on the unit sphere of $L^2(\IT)$; $\lamW\le\lamA$ since the two forms agree on $\PT$.
Positivity of the primitive form ($\lamA\ge0$) is known for $T\le\log2$ \cite{Paper42}.
Positivity of the complete form ($\lamW\ge0$) is known for $T\le\frac12\log2$
\cite{Yoshida1992,Bombieri2000} and, by a certified matrix computation, for $T\le0.8$
\cite{Zhu2026}; we use the latter result as it is stated there. We extend both slightly, to
$T\le\frac12\log5=0.80472$, by a certified computation (\S\ref{sec:log5}).

The primal side of the problem, the minimization of $\mathfrak a_T$ over unit vectors, has been
studied variationally by Bombieri \cite{Bombieri2000} and computationally by Galerkin compression
\cite{Paper42,Zhu2026,Groskin2026}. This paper is about the dual side: which objects certify
$\lamA\ge c$, what their structure is, and whether they can be exhibited.

\subsection{Results}
The archimedean multiplier is a L\'evy exponent (\S\ref{sec:levy}): $m(\tau)=C_0+\int_0^\infty
(1-\cos\tau u)K(u)\,du$ with $C_0=\psi(\frac14)-\log\pi$ and the completely monotone density
$K(u)=\frac12/\sinh(u/2)+\frac12/\cosh(u/2)$; equivalently
$\Rel\psi(\frac14+\frac{i\tau}2)-\psi(\frac14)=\varphi(\tau^2)$ with $\varphi$ a complete Bernstein
function (Theorem~\ref{thm:S}). Hence
\begin{equation}\label{eq:A}
\mathfrak a_T(F)=C_0\|F\|^2+\int_0^\infty\big(\|F\|^2-\cF(u)\big)K(u)\,du-\sum_{n\in\NT}v_n\cF(\log n)
\end{equation}
(Theorem~\ref{thm:A}), and the whole dependence on $T$ of the archimedean part is the elementary tail
$\int_{2T}^\infty K$.

Writing $\mu=|\Fh|^2d\tau/2\pi$, the primal problem is a moment problem over probability measures
whose cosine transform vanishes on $|u|\ge2T$ and which annihilate two explicit functions $\gh0,\gh1$.
Its dual (\S\ref{sec:duality}) is: maximize $\inf_\tau[\sT+\nuh+\beta_0\gh0+\beta_1\gh1]$ over finite
even signed measures $\nu$ supported in $\{|u|\ge2T\}$ and $\beta\in\R^2$. We prove strong duality,
with finitely atomic $\nu$ sufficing (Theorem~\ref{thm:D1}), so that positivity, when it holds with a
margin, is always provable by a finite certificate; the dual supremum is never attained
(Theorem~\ref{thm:O}), whereas the primal infimum is \cite[Thm.~3]{Bombieri2000}; exponential columns
exhaust the dual (Theorem~\ref{thm:E}); and $\beta=(2,0)$ is the unique choice for which a
certificate bounds the complete form, the two numbers restoring the polar terms of the explicit
formula (Proposition~\ref{prop:P}). Three tail lemmas reduce the verification of a certificate to a
bounded interval (\S\ref{sec:tails}), and at $T=\frac12\log3$, the largest window with a single
contact, an explicit certificate of $437$ atoms gives
\[
\lamA\ \ge\ \lamP\ \ge\ 8.8577\cdot10^{-5}
\]
(Theorem~\ref{thm:C}). This is an independent proof, by an explicit dual feasible point, of a
positivity statement already known at that window. The exchange between escaped primal mass and
value is exactly logarithmic (Theorem~\ref{thm:Vp}), and no pointwise envelope can lower the
verification threshold $2\pi e^{\Aconst}$, which grows doubly exponentially in $T$
(Proposition~\ref{prop:env}).

The main structural result is that the dual variable is the explicit formula (\S\ref{sec:Z}). For an
even weight $w$ with $w\equiv1$ on $[-2T,2T]$ and sufficient decay, an explicit measure $\nu_w$ outside
the window is dual feasible with $\beta=(2,0)$, and
\[
\sT+\nuh_w+2\gh0=f_w,\qquad f_w(\tau)=\sum_\rho\wh(\tau-\gamma_\rho)\quad(\rho=\tfrac12+i\gamma_\rho),
\]
the explicit formula evaluated on $u\mapsto w(u)\cos(\tau u)$ (Theorem~\ref{thm:Z}). Every
certificate of the complete form is of this kind up to an arbitrarily small uniform correction, and
$\sup_w\inf f_w=\lamW$ (Theorem~\ref{thm:Zexh}). Consequently RH is equivalent to the existence, for
every $T$, of a flat-top weight whose smoothing of the zeros is nonnegative on $\R$
(Corollary~\ref{cor:RH}). A flat-top weight never has a nonnegative Fourier transform
(Lemma~\ref{lem:Rrig}), so such a certificate always expresses a cancellation among the zeros, and no
finite set of zeros carries one (Theorem~\ref{thm:finZ}). In the same language RH is equivalent to the
indeterminacy, at every window, of a Kre\u{\i}n extension problem, which holds unconditionally for
$T\le0.8$ by \cite{Zhu2026} (Theorem~\ref{thm:KEP3}, Corollary~\ref{cor:KEP08}). A certified
computation in a Legendre basis, with a Schur and Haynsworth inertia count in ball arithmetic, gives
$\lamA\ge2.5\cdot10^{-12}$ and $\lamW\ge10^{-18}$ at $T=\frac12\log5$, the last window with contact set
$\{2,3,4\}$; hence both forms are positive and $\mathrm{KEP}(T)$ is indeterminate for
$T\le\frac12\log5$, slightly beyond the range of \cite{Zhu2026} (Theorems~\ref{thm:log5A},
\ref{thm:log5W}, Corollary~\ref{cor:KEPlog5}).

The remaining sections study the limits of the method. Lattice certificates have an explicit upper
bound, one small Toeplitz eigenvalue (Proposition~\ref{prop:TK}); these ceilings converge to $\lamP$
(Theorem~\ref{thm:R}), and interval enclosures show that no lattice of order $M\le8$ certifies at
$\frac12\log3$ and none of order $M\le24$ at $T=0.68$ or $T=\log2$ (Theorems~\ref{thm:N}, \ref{thm:Np}).
Positivity is a property of the complete initial segment of primes: no partial comb of small mass is
positive on a large window (Theorem~\ref{thm:SC}). With the primes deleted, the negative index of the
archimedean form on $\PT$ is $0$ at $\frac12\log3$ and $1$ at $\log2$, and is computed exactly for
$T\le0.90$ outside two short intervals (Theorem~\ref{thm:index}); its positivity threshold lies in
$[0.58,0.60]$ (Corollary~\ref{cor:Tarch}). At a fixed window the sign of $\lamP$ is decidable
(Proposition~\ref{prop:decide}); the margin decays, $\lamA\le4\cdot10^9Te^{-15.449T}$ for $T\ge1$
(Theorem~\ref{thm:margin}); and positivity at a fixed window is a condition on the finite vector
$(\Lambda(n))_{n<e^{2T}}$ which modifications of $\Lambda$ preserving its asymptotics, its zeros and its
pole violate at explicit windows (Theorem~\ref{thm:finmod}). The complete and the primitive form at one
window are related by a diagonal $2\times2$ matrix, the polar capacity (Theorem~\ref{thm:parityschur},
Proposition~\ref{prop:polarcap}). At $T=\log2$ the complete form is positive on an explicit
seven-dimensional subspace (Theorem~\ref{thm:nest}). The last section proves a finite identity
combining Poisson summation with $\Lambda*1=\log$, constructs a faithful Gamma-moment detector and
shows explicitly that the proposed return inequalities are equivalent to annihilating that detector.
The actual open problems are direct Gamma-moment annihilation at primitive first contact
(Problem~\ref{prob:annihilate}) and the first-contact bridge of Remark~\ref{rem:completecontact}.

Apart from the certified computation at $T\le\frac12\log5$, none of the results gives the sign of $W_T$ or
$\mathfrak a_T$ at a window $T>\log2$; the equivalences with RH are reformulations, not proofs.

\subsection{Related work}\label{subsec:related}
Positivity of the Weil functional on a fixed window goes back to Yoshida~\cite{Yoshida1992}, who
proved strict positivity of the Hermitian form on $K(a)$ at $a=\tfrac12\log 2$ by an approximate
finite-matrix computation, together with qualitative small-window and high-frequency lemmas
(\cite[Lemmas~2--3]{Yoshida1992}). Bombieri~\cite{Bombieri2000} gave the variational theory: the
infimum of the functional on the unit sphere of $L^2(E)$ is attained (\cite[Thm.~3]{Bombieri2000}),
the window minima are continuous and decreasing (\cite[Thm.~5]{Bombieri2000}), the finite
truncations never have eigenvalue $0$, and the number of negative eigenvalues equals the number of
off-line conjugate pairs (\cite[Thm.~8]{Bombieri2000}); his \cite[Thm.~12]{Bombieri2000} makes the
window $|I|<\log 2$, where the prime sum is empty, quantitative. Suzuki's work on screw functions
\cite{Suzuki2023jlms,Suzuki2023screwline,Suzuki2025hilbert} recasts the same form through Kre\u{\i}n's
class $\mathcal G_a$, with a L\'evy--Khintchine representation of the zero measure, trace-class kernels
and a de Branges model of the Weil Hilbert space; Theorem~\ref{thm:finZ} is a finite-rank counterpart
of that picture, in that a finite symmetric zero set makes the associated Kre\u{\i}n problem
determinate. Two recent preprints treat the truncated form by matrix methods: Groskin~\cite{Groskin2026}
proves an exact finite Guinand--Weil dictionary and shows that the omitted archimedean tail is a totally
positive Cauchy--Stieltjes increment, and Zhu~\cite{Zhu2026} reduces the window problem to one finite
matrix by a pointwise envelope, certifies $Q(f)\ge 8.9\cdot 10^{-18}\|f\|_2^2$ for
$\supp f\subseteq[-0.8,0.8]$ with a matching upper bound and a conjectural Landau--Widom decay law, and
proves that pointwise-envelope certificates face a doubly exponential barrier (\cite[Thm.~1.4]{Zhu2026}).
Our window is smaller than Zhu's; what is new here is the dual formulation, the non-attainment of the
dual by finite objects, and the identification of the dual variable with the explicit formula. The
Delsarte and Cohn--Elkies linear programming method \cite{CohnElkies2003} is the structural analogue
of Theorem~\ref{thm:D1}. The archimedean form alone was studied by Connes and Consani
\cite{ConnesConsani2021}.

\subsection{Conventions}
$\widehat h(\tau)=\int_\R h(u)e^{i\tau u}du$; for even $h$ this is $\int h\cos(\tau u)$. For a
finite even measure $\nu$, $\nuh(\tau)=\int\cos(\tau u)\,d\nu(u)$, and for a probability measure
$\mu$ on $\R$, $\check\mu(u)=\int\cos(\tau u)\,d\mu(\tau)$. The functions
\begin{equation}\label{eq:g}
g_0(u)=\cosh\tfrac u2\,\mathbf 1_{|u|<2T},\qquad g_1(u)=u\sinh\tfrac u2\,\mathbf 1_{|u|<2T},
\end{equation}
have transforms $\gh0,\gh1$ (closed forms in Lemma~\ref{lem:gbounds}). Galerkin eigenvalues are
upper bounds on $\lamA$ and are quoted as such; numerical values that are not certified are collected
in Appendix~\ref{app:diag}. Several results use a computer verification in interval arithmetic; the
methods are described in Appendix~\ref{app:encl}.

We write $B_*=\Aconst-C_0$, $\PiT$ for the prime form, $\Ainf$ for the archimedean form and $\Gamma_T$
for its restriction without the pole term (Proposition~\ref{prop:nodes}); $N(t)$ is the zero-counting
function; $M$ is a lattice order, $d=2T/M$ the lattice step and $P=2\pi/d$ the period. $(D)$ is the
dual program and $D(T,d,U)$ the deficit of Problem~\ref{prob:A}. Any other reuse of a letter is
local to the proof in which it occurs.



\section{The archimedean multiplier as a L\'evy exponent}\label{sec:levy}

\begin{lemma}[the L\'evy density]\label{lem:K}
For $u>0$ put $K(u)=\dfrac{2e^{3u/2}}{e^{2u}-1}=\dfrac{1}{2\sinh(u/2)}+\dfrac{1}{2\cosh(u/2)}
=2\sum_{k\ge0}e^{-(2k+\frac12)u}$. Then $K$ is positive, completely monotone,
$K(u)=u^{-1}+O(1)$ as $u\downarrow0$, $K(u)=2e^{-u/2}(1+O(e^{-u}))$ as $u\to\infty$, and
\begin{equation}\label{eq:m}
m(\tau)=\Rel\psi\!\left(\tfrac14+\tfrac{i\tau}2\right)-\log\pi
= C_0+\int_0^\infty(1-\cos\tau u)\,K(u)\,du,\quad C_0=\psi(\tfrac14)-\log\pi=-5.3721834\ldots
\end{equation}
\end{lemma}
\begin{proof}
The three expressions for $K$ agree: $\frac{1}{2\sinh(u/2)}+\frac1{2\cosh(u/2)}
=\frac{\cosh(u/2)+\sinh(u/2)}{\sinh u}=\frac{2e^{u/2}}{e^u-e^{-u}}=\frac{2e^{3u/2}}{e^{2u}-1}$, and
$\frac{2e^{3u/2}}{e^{2u}-1}=2e^{-u/2}\sum_{k\ge0}e^{-2ku}$. Complete monotonicity is that of a
convergent series of decaying exponentials with positive coefficients. For \eqref{eq:m} use
$\psi(z)=-\gamma+\int_0^\infty\frac{e^{-t}-e^{-zt}}{1-e^{-t}}dt$ ($\Rel z>0$): with $z=\frac14+\frac{i\tau}2$,
\[
\Rel\psi(z)-\psi(\tfrac14)=\int_0^\infty\frac{e^{-t/4}\big(1-\cos(\tau t/2)\big)}{1-e^{-t}}\,dt
=\int_0^\infty\frac{2e^{-u/2}(1-\cos\tau u)}{1-e^{-2u}}\,du\qquad(t=2u),
\]
and $\frac{2e^{-u/2}}{1-e^{-2u}}=\frac{2e^{3u/2}}{e^{2u}-1}=K(u)$. The integral converges since
$1-\cos\tau u=O(u^2)$ at $0$ and $K=O(e^{-u/2})$ at $\infty$.
\end{proof}

\begin{theorem}[complete Bernstein structure]\label{thm:S}
Put $a_k=k+\frac14$ and $\varphi(x)=\sum_{k\ge0}\dfrac{x}{a_k(x+4a_k^2)}$ for $x\ge0$. Then
\emph{(i)} $\Rel\psi(\frac14+\frac{i\tau}2)-\psi(\frac14)=\varphi(\tau^2)$ and $\varphi$ is a
complete Bernstein function; \emph{(ii)} the L\'evy measure of the negative-definite function
$\tau\mapsto\varphi(\tau^2)$ on $\R$ is $\frac12K(|u|)du$, which is completely monotone on each
half-line, so the associated L\'evy process is a pure-jump subordinate Brownian motion
$B_{S_t}$ with a special subordinator $S$ of Laplace exponent $\varphi$; in particular it has no
Gaussian component; \emph{(iii)} $\int_{2T}^\infty K(u)\,du=\log\coth\frac T2+\pi-2\arctan e^{T}$.
\end{theorem}
\begin{proof}
(i) From $\psi(z)=-\gamma+\sum_{k\ge0}\big[\frac1{k+1}-\frac1{k+z}\big]$,
$\Rel\big[\frac1{a_k}-\frac1{a_k+i\tau/2}\big]=\frac1{a_k}-\frac{a_k}{a_k^2+\tau^2/4}
=\frac{\tau^2}{a_k(4a_k^2+\tau^2)}$, which is the stated series at $x=\tau^2$. Each summand
$x\mapsto a_k^{-1}\,x/(x+4a_k^2)$ is a complete Bernstein function (Stieltjes representation with
a point mass at $4a_k^2$), the summands are $O(k^{-3})$ locally uniformly, and the complete
Bernstein functions form a convex cone closed under pointwise limits
\cite[Ch.~6--7]{SSV2012}. (ii) By Lemma~\ref{lem:K}, $\varphi(\tau^2)=\int_\R(1-\cos\tau u)\,
\frac12K(|u|)du$, so the L\'evy measure is $\frac12K(|u|)du$; it is completely monotone on
$(0,\infty)$ and $\varphi(x)/x\to0$, so there is no Gaussian part; the representation as a
subordinate Brownian motion with a special subordinator is \cite[Thm.~13.10, Cor.~13.12]{SSV2012}.
(iii) $\int\frac{du}{2\sinh(u/2)}=\log\tanh\frac u4$ and $\int\frac{du}{2\cosh(u/2)}=2\arctan e^{u/2}$.
\end{proof}

\begin{remark}
Statement~(i) appears in earlier unpublished work of the author; (ii) and (iii) are new, and (ii)
shows in particular that the archimedean part is not a diffusion.
\end{remark}

\begin{theorem}[the form as a L\'evy energy]\label{thm:A}
For every $F\in L^2(\IT)$, with $\Lev(\Phi):=\int_0^\infty(\Phi(0)-\Rel\Phi(u))K(u)\,du\in[0,\infty]$
(for complex $F$ the autocorrelation is Hermitian and only its even part $\Rel\cF$ enters; the
prime term below is likewise $\Rel\cF(\log n)$, and from \S\ref{sec:duality} on all measures are even),
\begin{equation}\label{eq:Aform}
\mathfrak a_T(F)=C_0\|F\|^2+\Lev(\cF)-\sum_{n\in\NT}v_n\cF(\log n),
\end{equation}
both sides being finite exactly when $\int\log(2+|\tau|)|\Fh|^2<\infty$. Consequently, with
$\Kw=\{\cF:F\in L^2(\IT),\|F\|=1\}$ and the primitive class
\begin{equation}\label{eq:Kprim}
\Kp=\Big\{\Phi\in\Kw:\ \int\Phi(u)\cosh\tfrac u2\,du=0,\ \int\Phi(u)\,u\sinh\tfrac u2\,du=0\Big\},
\end{equation}
one has
\begin{align*}
\lamP&:=\inf_{\Phi\in\Kp}\Big[C_0+\Lev(\Phi)-\sum_{n\in\NT}v_n\Phi(\log n)\Big]=C_0-\Lambda_T,\\
\Lambda_T&:=\sup_{\Phi\in\Kp}\Big[\sum_{n\in\NT}v_n\Phi(\log n)-\Lev(\Phi)\Big],
\end{align*}
and $\lamP\le\lamA$.
\end{theorem}
\begin{proof}
By \eqref{eq:m} and Tonelli (both $1-\cos\tau u$ and $K$ are nonnegative),
$\frac1{2\pi}\int m|\Fh|^2=C_0\|F\|^2+\int_0^\infty K(u)\Big[\frac1{2\pi}\int(1-\cos\tau u)|\Fh|^2d\tau\Big]du$,
and $\frac1{2\pi}\int\cos(\tau u)|\Fh(\tau)|^2d\tau=\cF(u)$ by Plancherel; likewise for the
prime terms. Finiteness: $\int_0^\infty(1-\cos\tau u)K(u)du=\log|\tau|+O(1)$ as $|\tau|\to\infty$
by Lemma~\ref{lem:K}, since $\int_0^1(1-\cos\tau u)\frac{du}u=\log|\tau|+O(1)$. For the last
statements: the two integrals in \eqref{eq:Kprim} are $\Rel(A\bar B)$ and $\Rel(A'\bar B-A\bar B')$
with $A'=\int tF(t)e^{t/2}dt$, $B'=\int tF(t)e^{-t/2}dt$ (expand $\cosh\frac{t-s}2$ and
$(t-s)\sinh\frac{t-s}2$ in exponentials), so $\{\cF:F\in\PT,\|F\|=1\}\subseteq\Kp$, whence
$\lamP\le\lamA$. In terms of $\widehat\Phi$, an entire function of exponential type $2T$ with
$\widehat\Phi(i/2)=\overline{\widehat\Phi(-i/2)}$, the two conditions say that
$\Rel\widehat\Phi(i/2)=0$ and $\Rel\big[i(\widehat\Phi'(i/2)-\widehat\Phi'(-i/2))\big]=0$; when
$F\in\PT$ one has the stronger $\widehat\Phi(\pm i/2)=\widehat\Phi'(\pm i/2)=0$, i.e.\ $\widehat\Phi$
vanishes to order $\ge2$ at $\pm i/2$ (Remark~\ref{rem:Kprim}).
\end{proof}

\begin{remark}\label{rem:Kprim}
$\Kp$ is slightly larger than the set of primitive autocorrelations: for real $F$ its two
conditions read $AB=0$ and $A'B-AB'=0$, i.e.\ $F\in\PT$, or $A=A'=0$, or $B=B'=0$. Every
certificate below bounds $\lamP$, hence $\lamA$; the equality $\lamP=\lamA$ is not needed
anywhere and is not asserted. Since $W_T(F)=\mathfrak a_T(F)$ whenever $A\bar B$ is purely
imaginary, also $\lamW\le\lamP$.
\end{remark}

\begin{corollary}[monotonicity; upper bound]\label{cor:mono}
With the contact set held fixed, $T\mapsto\Lambda_T$ is nondecreasing and $T\mapsto\lamP$
nonincreasing. Moreover
$\lamA\le C_0+\Aconst+\varphi\big((3\pi/2T)^2\big)$, which equals $1.291$ at $T=\frac12\log3$ and
$2.327$ at $T=\log 2$.
\end{corollary}
\begin{proof}
Monotonicity: $\Kp$ grows with $T$ and $\Lev$ does not depend on $T$. Upper bound: the form
$\Lev(\cF)$ is the Dirichlet form of the L\'evy process of Theorem~\ref{thm:S} killed on
leaving $\IT$; its eigenvalues are bounded by those of the subordinate killed Brownian motion,
$\lambda_n\le\varphi((n\pi/2T)^2)$ \cite[Prop.~2.3]{SongVondracek2003}, the Brownian motion being
normalized with characteristic exponent $\xi^2$; $\PT$ has codimension
$2$, so the bottom of $\Lev$ on $\PT$ is at most $\lambda_3$; and $-\sum v_n\cF(\log n)\le\Aconst\|F\|^2$.
\end{proof}

The bound of Corollary~\ref{cor:mono} is an upper bound and says nothing about the sign; at
$T=\log2$ its two crude steps lose $2.3$ out of $\Aconst-C_0=7.62$, and there the bottom of the form
is not known even numerically (Remark~\ref{rem:galerkin}). The rest of the paper is concerned with
lower bounds.

The hypothesis that the contact set is fixed can be removed from Corollary~\ref{cor:mono}, for
all three bottoms at once.

\begin{lemma}[monotonicity across contact thresholds]\label{lem:monoT}
The functions $T\mapsto\lamA$, $T\mapsto\lamP$ and $T\mapsto\lamW$ are nonincreasing on $(0,\infty)$,
including at the points $T=\frac12\log n$ where a contact enters.
\end{lemma}
\begin{proof}
Let $T<T'$ and extend $F\in L^2(\IT)$ by zero to $I_{T'}$. The norm, the Tate moments $A(F),B(F)$ and
the autocorrelation $\cF$ are unchanged, and $\cF$ vanishes on $|u|\ge2T$. A contact
$n\in\mathcal N_{T'}\setminus\NT$ has $\log n\ge2T$, so $\cF(\log n)=0$, and by \eqref{eq:Aform}
$\mathfrak a_{T'}(F)=\mathfrak a_T(F)$. Hence $\PT\subset\mathcal P_{T'}$ isometrically with the same
form values, which gives the statement for $\lamA$; the same computation with $W_T=\mathfrak a_T+2\Rel(A\bar B)$
gives it for $\lamW$, and for $\lamP$ the class $\Kp$ grows with $T$, $\Lev$ is $T$-free and the new
contacts contribute $\Phi(\log n)=0$.
\end{proof}

For the complete form, monotonicity comes with continuity; this is \cite[Thm.~5]{Bombieri2000}.
Bombieri proves that the even and odd minima $\mu_\pm$ of the complete functional
on $L^2([M^{-1},M])$ are continuous and decreasing. Here $W_T$ splits into its values on the even and
odd parts of $F$: the symbol is even, and $\Rel(A\bar B)=|C|^2-|S|^2$ with $C=\int F\cosh\frac t2$, which
sees only the even part, and $S=\int F\sinh\frac t2$, which sees only the odd part. So, after the change of
variables $x=e^t$, $\lamW=\min(\mu_+,\mu_-)$ at $M=e^T$. A direct proof in the present normalization follows.

\begin{lemma}[continuity of the complete bottom]\label{lem:contT}
The function $T\mapsto\lamW$ is continuous on $(0,\infty)$.
\end{lemma}
\begin{proof}
Fix $T$ and let $r\in[\frac12,2]$. The dilation $U_rF(t)=r^{-1/2}F(t/r)$ is unitary from $L^2(\IT)$ onto
$L^2(I_{rT})$, with $\widehat{U_rF}(\tau)=r^{1/2}\Fh(r\tau)$ and $c_{U_rF}(u)=\cF(u/r)$. It maps the form
domain onto the form domain, since $\log(2+|s|/r)\le\log(2+|s|)+\log2$. Hence
$\lambda^{\mathrm W}_{rT}=\inf\{W_{rT}(U_rF):F\text{ in the form domain},\ \|F\|=1\}$. By \eqref{eq:form}
and \eqref{eq:W},
\[
W_{rT}(U_rF)=\frac1{2\pi}\int_\R m(s/r)|\Fh(s)|^2ds-\sum_{2\le n<e^{4T}}v_n\Rel\cF(\tfrac{\log n}r)
+2\Rel\big(A(U_rF)\overline{B(U_rF)}\big),
\]
where the comb may be summed over $n<e^{4T}$ for both windows, because $\cF$ vanishes on $|u|\ge2T$. This
also takes care of contacts that enter or leave. We compare the three terms with those of $W_T(F)$.
\emph{(i) Archimedean.} $m$ is even and $m'(x)=-\frac12\Iml\psi'(\frac14+\frac{ix}2)$ with
$\psi'(z)=O(1/|z|)$ on $\Rel z=\frac14$. So $M_1:=\sup_x|xm'(x)|<\infty$ and
$|m(s/r)-m(s)|\le M_1|\log r|$, and the difference of the archimedean terms is at most
$M_1|\log r|\,\|F\|^2$.
\emph{(ii) Comb.} $|\Rel\cF(u)-\Rel\cF(u')|\le\frac1{2\pi}\int\min(2,|\tau||u-u'|)|\Fh|^2d\tau$. Put
$\delta=4T|1/r-1|$ and $\|F\|_Q^2:=\frac1{2\pi}\int\log(2+|\tau|)|\Fh|^2d\tau$. Splitting at
$|\tau|=\delta^{-1/2}$, the difference of the comb terms is at most
$S_T\big(\delta^{1/2}\|F\|^2+2\|F\|_Q^2/\log(2+\delta^{-1/2})\big)$, where $S_T=\sum_{n<e^{4T}}v_n$.
\emph{(iii) Polar.} $A(U_rF)=r^{1/2}\int F(s)e^{rs/2}ds$, so
$|A(U_rF)-A(F)|\le\|F\|\,\|r^{1/2}e^{rs/2}-e^{s/2}\|_{L^2(\IT)}$, and likewise for $B$. Both moments are
bounded by $C_T\|F\|$ uniformly in $r\in[\frac12,2]$. So the difference of the polar terms is at most
$\delta_P(r)\|F\|^2$ with $\delta_P(r)\to0$ as $r\to1$.
By Lemma~\ref{lem:mbounds}(ii) there is $C_1$ with $\log(2+|\tau|)\le m(\tau)+C_1$. By
Lemma~\ref{lem:mbounds}(iii) and $|2A\bar B|\le4\sinh T\,\|F\|^2$ (Cauchy--Schwarz), this gives
$\|F\|_Q^2\le W_T(F)+C'\|F\|^2$ with $C'=C_1+\Aconst+4\sinh T$. Collecting (i)--(iii), there is
$\eta(r)\to0$ as $r\to1$ with
\[
|W_{rT}(U_rF)-W_T(F)|\le\eta(r)\big(W_T(F)+C'\|F\|^2\big)
\]
for all $F$ in the form domain, and the right side is nonnegative. Take infima over unit $F$. This gives
$(1-\eta)\lamW-\eta C'\le\lambda^{\mathrm W}_{rT}\le(1+\eta)\lamW+\eta C'$, so
$|\lambda^{\mathrm W}_{rT}-\lamW|\le\eta(r)(|\lamW|+C')\to0$.
\end{proof}

The same proof with the comb deleted gives continuity of $\laminf$ (\S\ref{subsec:arch}). The
argument does not apply to $\lamA$ or $\lamP$, because $U_r$ does not preserve $\PT$.


\section{Primal and dual programs}\label{sec:duality}

\subsection{Primal and dual}
Fix $T$ and the weights $v=(v_n)_{n\in\NT}$ (any nonnegative weights may be used; $v_n=2\Lambda(n)/\sqrt n$
is the case of $\zeta$).

\begin{definition}\label{def:PD}
\textbf{(P)} $\ \lamP=\inf\big\{\int_\R\sT\,d\mu:\ \mu$ an even probability measure on $\R$,
$\check\mu(u)=0$ for $|u|\ge2T$, $\int\gh0\,d\mu=\int\gh1\,d\mu=0\big\}$.

\textbf{(D)} $\ D_T=\sup\big\{\inf_{\tau\in\R}\big[\sT(\tau)+\nuh(\tau)+\beta_0\gh0(\tau)+\beta_1\gh1(\tau)\big]:\
\nu$ a finite even signed Borel measure with $\supp\nu\subseteq\{|u|\ge2T\}$, $\beta\in\R^2\big\}$.
A triple $(\nu,\beta_0,\beta_1)$ with $f:=\sT+\nuh+\beta\cdot\widehat g\ge c$ on $\R$ is a
\emph{certificate of value $c$}.
\end{definition}

\begin{proposition}[the primal is the eigenvalue problem; weak duality]\label{prop:weak}
The value of \textup{(P)} is the $\lamP$ of Theorem~\ref{thm:A}. Every certificate of value $c$
satisfies $c\le\lamP\le\lamA$.
\end{proposition}
\begin{proof}
If $F$ is a unit vector, $\mu_F=|\Fh|^2d\tau/2\pi$ is feasible ($\check\mu_F=\cF$) with
$\int\sT d\mu_F=\mathfrak a_T(F)$, and $\int\gh j\,d\mu_F=\int g_j\cF$. Conversely if $\mu$ is
feasible, $\Phi=\check\mu$ is continuous, positive definite, $\Phi(0)=1$, supported in $[-2T,2T]$,
so by the Boas--Kac theorem \cite{BoasKac1945} $\Phi=\cF$ for some $F\in L^2(\IT)$, $\|F\|=1$, and
$\Phi\in\Kp$; thus the feasible set of (P) is exactly $\Kp$ read through $\mu=|\Fh|^2d\tau/2\pi$.
Weak duality: for feasible $\mu$, $\int\nuh\,d\mu=\int\check\mu\,d\nu=0$ because $\check\mu=\cF$
vanishes identically on the closed set $\{|u|\ge2T\}$ ($\cF(\pm2T)=0$ by continuity, so an atom
of $\nu$ at $\pm2T$ is harmless), and $\int\gh j\,d\mu=0$; hence
$\int\sT d\mu=\int f\,d\mu\ge c$.
\end{proof}

\begin{lemma}[bounds on the multipliers]\label{lem:gbounds}
$\gh0(\tau)=2\int_0^{2T}\cosh\frac u2\cos\tau u\,du$ and $\gh1(\tau)=2\int_0^{2T}u\sinh\frac u2\cos\tau u\,du$
are even, real, bounded, tend to $0$ at infinity, and for $\tau>0$
\[
|\gh0(\tau)|\le k_0(\tau):=\frac{2\cosh T}{\tau}+\frac{\sinh T}{\tau^2},\qquad
|\gh1(\tau)|\le k_1(\tau):=\frac{8T\sinh T}{\tau},
\]
both bounds decreasing in $\tau$. Moreover $\gh0(\tau)=\frac{2\cosh T\sin2T\tau}{\tau}+O(\tau^{-2})$,
and for any $g=y_0g_0+y_1g_1$, $\widehat g(\tau)=\frac{2g(2T^-)\sin2T\tau}{\tau}+\frac{2g'(2T^-)\cos2T\tau}{\tau^2}+O(\tau^{-3})$.
\end{lemma}
\begin{proof}
Integrate by parts: $\gh0=\frac{2\cosh T\sin2T\tau}\tau-\frac1\tau\int_0^{2T}\sinh\frac u2\sin\tau u\,du$,
and once more, $\frac1\tau\int_0^{2T}\sinh\frac u2\sin\tau u\,du=-\frac{\sinh T\cos2T\tau}{\tau^2}+\frac{\gh0(\tau)}{4\tau^2}$;
solving, $|\gh0|(1+\frac1{4\tau^2})\le\frac{2\cosh T}\tau+\frac{\sinh T}{\tau^2}$. For $\gh1$, one
integration by parts gives $\gh1=\frac{4T\sinh T\sin2T\tau}{\tau}-\frac2\tau\int_0^{2T}\big(\sinh\frac u2+\frac u2\cosh\frac u2\big)\sin\tau u\,du$,
and the integrand is the derivative of $u\sinh\frac u2$, whose integral over $[0,2T]$ is $2T\sinh T$.
The asymptotic expansion is two integrations by parts for a function real-analytic on $[0,2T]$
with $g'(0)=0$.
\end{proof}

\subsection{Strong duality}

\begin{theorem}[strong duality; finitely atomic dual sequences]\label{thm:D1}
$D_T=\lamP$. Moreover, for every $\varepsilon>0$ there are finitely many atoms
$u_1,\dots,u_N\ge2T$, reals $a_1,\dots,a_N$ and $\beta\in\R^2$ such that
$\sT(\tau)+\sum_j2a_j\cos(u_j\tau)+\beta_0\gh0(\tau)+\beta_1\gh1(\tau)>\lamP-\varepsilon$ for all $\tau\in\R$.
In particular, for every $T$ with $\lamP>0$, $\lamA>0$ is provable by a finite certificate.
\end{theorem}
\begin{proof}
\emph{(i) Finite relaxations.} For a finite set $U=\{u_1<\dots<u_N\}\subset[2T,\infty)$ let
$P(U)=\inf\{\int\sT d\mu:\ \mu\ \text{even probability},\ \check\mu(u_i)=0\ \forall i,\ \int\gh j d\mu=0\}$
and $D(U)$ the supremum in (D) restricted to $\nu$ supported on $\pm U$. Then $P(U)\le\lamP$
(the feasible set of (P) is contained in that of $P(U)$) and $D(U)\le D_T\le\lamP$.

\emph{(ii) $P(U)=D(U)$.} Write $f_0\equiv1$, $f_i(\tau)=\cos(u_i\tau)$ ($1\le i\le N$),
$f_{N+1}=\gh0$, $f_{N+2}=\gh1$, $b=(1,0,\dots,0)\in\R^{N+3}$, and let
$C=\{\sum_kt_kf(\tau_k):\ t_k\ge0\}$ be the convex cone generated by the moment curve
$\tau\mapsto f(\tau)=(f_0,\dots,f_{N+2})(\tau)$. The value function
$w(b')=\inf\{\int\sT d\mu:\ \mu\ge0,\ \int f\,d\mu=b'\}$ is convex, $>-\infty$ everywhere
($\sT\ge-\Aconst-|C_0|$ and $\int f_0d\mu=b'_0$ bounds the mass) and finite on $C$ (finite atomic
$\mu$ have finite value). Its Legendre biconjugate at $b$ is the Lagrangian dual
$w^{**}(b)=\sup\{y_0:\ \sT-\sum_ky_kf_k\ge0\text{ on }\R\}$, which is $D(U)$ after
$a_i=-y_i/2$, $\beta_j=-y_{N+1+j}$. A proper convex function agrees with its biconjugate at every
interior point of its domain \cite[Thm.~7.4, Cor.~12.1.1]{Rockafellar}; so it suffices to show
$b\in\operatorname{int}C$.

Suppose not. Since $C$ is a convex cone containing $0$, there is $y\ne0$ with
$\langle y,z\rangle\ge\langle y,b\rangle$ for all $z\in C$; taking $z=0$ gives $y_0\le0$, and
taking $z=tf(\tau)$, $t\to\infty$, gives
$G(\tau):=y_0+\sum_iy_i\cos(u_i\tau)+y_{N+1}\gh0(\tau)+y_{N+2}\gh1(\tau)\ge0$ for all $\tau$.
The Ces\`aro mean $X^{-1}\int_0^XG\to y_0$ (each $\cos(u_i\tau)$ has mean $0$ because $u_i>0$,
and $\gh j\to0$ by Lemma~\ref{lem:gbounds}); as $G\ge0$ this forces $y_0=0$, and then
$X^{-1}\int_0^X|G|\to0$. Testing against the bounded functions $\cos(u_i\tau)$ gives
$y_i/2=\lim X^{-1}\int_0^XG\cos(u_i\tau)\,d\tau=0$ for each $i$. There remains
$\widehat g\ge0$ with $g=y_{N+1}g_0+y_{N+2}g_1$. By the expansion in Lemma~\ref{lem:gbounds},
$\widehat g\ge0$ forces $g(2T^-)=0$ and then $g'(2T^-)=0$; but
$g(2T^-)=y_{N+1}\cosh T+2T\,y_{N+2}\sinh T$ and $g'(2T^-)=\frac12y_{N+1}\sinh T+y_{N+2}(\sinh T+T\cosh T)$,
whose determinant is $\sinh T\cosh T+T>0$. So $y=0$, a contradiction; hence $b\in\operatorname{int}C$.

\emph{(iii) The limit.} Choose finite sets $U_1\subset U_2\subset\cdots$ with $\bigcup U_n$ dense in
$[2T,\infty)$; $P(U_n)$ is nondecreasing and $\le\lamP$. Pick feasible $\mu_n$ with
$\int\sT d\mu_n\le P(U_n)+1/n$. Since $\sT\to+\infty$ and is bounded below, $(\mu_n)$ is tight;
let $\mu$ be a weak-$*$ limit point, a probability measure. For $u\in U_N$ and $n\ge N$,
$\check\mu_n(u)=0$ and $\cos(u\,\cdot)$ is bounded continuous, so $\check\mu(u)=0$; $\check\mu$
is continuous, so $\check\mu\equiv0$ on $[2T,\infty)$; likewise $\int\gh jd\mu=0$. Thus $\mu$ is
feasible for (P), and by lower semicontinuity of $\mu\mapsto\int\sT d\mu$ under weak convergence
($\sT$ continuous and bounded below), $\lamP\le\int\sT d\mu\le\liminf\int\sT d\mu_n=\lim P(U_n)\le\lamP$.
Hence $D_T\ge\sup_nD(U_n)=\sup_nP(U_n)=\lamP\ge D_T$, and the atomic statement follows.
\end{proof}

\begin{remark}
Theorem~\ref{thm:D1} is the Weil criterion at window $T$ in dual coordinates:
\[
\mathfrak a_T\ge0\ \text{on}\ \PT\ \Longleftarrow\ \lamP\ge0\iff
\forall\varepsilon>0\ \exists\ \text{atomic }\nu\ \text{on}\ |u|\ge2T,\ \beta\in\R^2:\
\sT+\nuh+\beta\cdot\widehat g\ge-\varepsilon\ \text{on }\R,
\]
(the first implication is $\lamP\le\lamA$),
and RH is equivalent to the right-hand side holding for every $T$. The same proof, with the two $\widehat g$-constraints deleted, gives strong duality for
the complete Weil form: $\lamW=\sup\{\inf_\tau[\sT+2\gh0+\nuh]\}$ over the same class of $\nu$
(Proposition~\ref{prop:P} explains the constant $2$).
\end{remark}

\subsection{Non-attainment}

\begin{theorem}[non-attainment]\label{thm:O}
Let $\nu$ be any finite even signed measure with $\supp\nu\subseteq\{|u|\ge2T\}$ and let
$h_0,\dots,h_J$ be finitely many continuous functions on $[-2T,2T]$ that annihilate $\Kp$
(for instance $g_0,g_1$). Then for every $\beta\in\R^{J+1}$
\[
\inf_{\tau\in\R}\Big[\sT(\tau)+\nuh(\tau)+\sum_j\beta_j\widehat h_j(\tau)\Big]<\lamP\quad\text{strictly.}
\]
The statement holds for every finite weight family, including $v=0$.
\end{theorem}
\begin{proof}
A minimizer $\Phi_\star=c_{F_\star}\in\Kp$ of the primal exists: the sublevel sets of
$\mathfrak a_T$ on the unit sphere of $L^2(\IT)$ are compact, because $m\to\infty$ gives uniform
control of $\int_{|\tau|>R}|\Fh|^2$ while $F\mapsto\Fh|_{[-R,R]}$ is Hilbert--Schmidt on $L^2(\IT)$
(equivalently: the feasible set of (P) is tight and weak-$*$ closed and $\mu\mapsto\int\sT d\mu$ is
weak-$*$ lower semicontinuous, exactly as in step (iii) of the proof of Theorem~\ref{thm:D1}; on
$\PT$ itself this is \cite[Thm.~3]{Bombieri2000}). Put $\varrho=\sT+\nuh+\sum\beta_j\widehat h_j-\lamP$, a
continuous function ($\nu$ finite). By the annihilation properties,
$\frac1{2\pi}\int\varrho|\Fh_\star|^2=\mathfrak a_T(F_\star)-\lamP=0$. If $\varrho\ge0$ then $\varrho=0$ a.e.\ on
$\{\Fh_\star\ne0\}$, whose complement is the discrete zero set of a nonzero entire function of
exponential type; so $\varrho\equiv0$ on $\R$, and by Fourier inversion, as distributions in $u$,
\[
\check\sT+\nu+\sum_j\beta_jh_j=\lamP\,\delta_0 .
\]
Restrict to the open set $0<|u|<2T$, where $\nu=0$: by \eqref{eq:m}, $\check m=-\frac12K(|u|)$
there, and the prime part of $\check\sT$ is $-\sum_{n\in\NT}\frac{v_n}2(\delta_{\log n}+\delta_{-\log n})$;
hence $\sum_j\beta_jh_j(u)=\frac12K(|u|)+\sum_{n\in\NT}\frac{v_n}2(\delta_{\log n}+\delta_{-\log n})$
on $0<|u|<2T$. The left side is a bounded continuous function; the right side is unbounded as
$u\to0$ ($K(u)\sim1/u$) and carries atoms as soon as $\NT\neq\emptyset$. Contradiction.
\end{proof}

\begin{remark}\label{rem:OvsB}
Theorem~\ref{thm:O} forbids attainment, not certificates: whenever $\lamP>0$ any dual value in
$(0,\lamP)$ suffices, and Theorem~\ref{thm:D1} says such values are reached by finite atomic
$\nu$. Every maximizing sequence of the dual must therefore degenerate (its total variation is
unbounded; Corollary~\ref{cor:T2} below quantifies the cost). On the primal side the infimum is
attained \cite[Thm.~3]{Bombieri2000}. The proof shows the reason for the asymmetry: the
singularity $K(u)\sim1/u$ of the archimedean L\'evy density at $u=0$ and the prime atoms lie inside
the window, where the dual measure is not allowed to have support.
\end{remark}

\subsection{Euler--Lagrange, and the polar multiplier}

\begin{theorem}[finite-interval Wiener--Hopf equation]\label{thm:W}
Let $\mathcal A$ denote the convolution operator on $\R$ with symbol $\sT$, acting on
$L^2(\IT)$ extended by zero and taking values in tempered distributions. A unit vector
$F\in\PT$ in the form domain is a critical point of $\mathfrak a_T$ on $\PT$ with critical value
$\lambda$ if and only if there are $\alpha,\beta\in\C$ (real for real $F$) with
\[
(\mathcal A-\lambda)F=\alpha\cosh\tfrac t2+\beta\sinh\tfrac t2\qquad\text{on }\IT,
\]
in the sense of distributions on $\IT$. The symbol $\sT-\lambda$ is real, even, tends to
$+\infty$, and changes sign; it is not elliptic, and no Wiener--Hopf index is defined for it.
\end{theorem}
\begin{proof}
Criticality means $\langle(\mathcal A-\lambda)F,G\rangle=0$ for all $G\in C_c^\infty(\IT)$ with
$\int Ge^{\pm t/2}=0$, a subspace of codimension $2$; a distribution on $\IT$ annihilating the
common kernel of two continuous linear functionals is a linear combination of them. Conversely
such an $F$ is critical. $\sT(0)=C_0-\Aconst<0$ and $\sT\to+\infty$.
\end{proof}

\begin{proposition}[the polar multiplier]\label{prop:P}
For every $F\in L^2(\IT)$ in the form domain,
$\frac1{2\pi}\int\gh0|\Fh|^2=\Rel\big(A(F)\overline{B(F)}\big)$ and hence
\[
\frac1{2\pi}\int\big(\sT+2\gh0\big)|\Fh|^2\,d\tau=W_T(F);
\]
if moreover $F$ is real and $\sum_\rho|\Phi_F(\rho)\Phi_F(1-\rho)|<\infty$ (in particular if
$F\in C_c^\infty(\IT)$ is real), then $W_T(F)=\sum_\rho\Phi_F(\rho)\Phi_F(1-\rho)$
(Lemma~\ref{lem:roughEF} below).
Consequently a certificate with $(\beta_0,\beta_1)=(2,0)$ and value $c$ proves $W_T\ge c$ on all
of $L^2(\IT)$, and $\beta_0=2$ is the unique value of $\beta_0$ with this property.
\end{proposition}
\begin{proof}
$\frac1{2\pi}\int\gh0|\Fh|^2=\int g_0\cF=\iint\cosh\frac{t-s}2F(t)\overline{F(s)}$ and
$\cosh\frac{t-s}2=\frac12e^{t/2}e^{-s/2}+\frac12e^{-t/2}e^{s/2}$ gives $\frac12(A\bar B+B\bar A)$.
The identification of $W_T$ with the zero sum is the explicit formula for the autocorrelation
$\cF$, \cite{Weil52,Bombieri2000}, \cite[Thm.~8.4]{Paper45}, whose polar term is
$2\Phi_F(0)\Phi_F(1)=2\Rel(A\bar B)$ for the real forms considered here. Uniqueness: for
$\beta_0\ne2$, $\frac1{2\pi}\int(\sT+\beta_0\gh0)|\Fh|^2=W_T(F)+(\beta_0-2)\Rel(A\bar B)$, and
$\Rel(A\bar B)$ takes both signs on the unit sphere of $L^2(\IT)$ (Lemma~\ref{lem:W} below, witnesses
$(1,\pm1,0,0)$), so no lower bound of the left side bounds $W_T$; the same argument with
$\Rel(A'\bar B-A\bar B')$ forces $\beta_1=0$.
\end{proof}

\cite[Thm.~8.4]{Paper45} is stated for real $F\in C_c^\infty$. The first equality of
Proposition~\ref{prop:P} holds on the whole form domain and is the only one the certificates use;
wherever the zero sum is used for a rough test (Corollary~\ref{cor:RH}, Theorem~\ref{thm:margin},
Remark~\ref{rem:Cheb}), it enters through the following lemma.

\begin{lemma}[the explicit formula for rough tests]\label{lem:roughEF}
Let $F\in L^2(\IT)$ be real, with \emph{(a)} $\int\log(2+|\tau|)|\Fh|^2<\infty$ and \emph{(b)}
$\sum_\rho|\Phi_F(\rho)\Phi_F(1-\rho)|<\infty$. Then
\[
\sum_\rho\Phi_F(\rho)\Phi_F(1-\rho)=W_T(F)=2\Phi_F(0)\Phi_F(1)+\mathfrak a_T(F),
\]
the sum being absolutely convergent. In particular $\mathfrak a_T(F)=\sum_\rho\Phi_F(\rho)\Phi_F(1-\rho)$
for $F\in\PT$, and a zero $\rho=\frac12+i\gamma$ on the critical line contributes $|\Fh(\gamma)|^2$.
\end{lemma}
\begin{proof}
For real $F$, $\Phi_F(1)=A(F)$ and $\Phi_F(0)=B(F)$, so $2\Phi_F(0)\Phi_F(1)=2\Rel(A\bar B)$. Let
$\eta_\varepsilon\ge0$, $\eta_\varepsilon\in C_c^\infty(-\varepsilon,\varepsilon)$, $\int\eta_\varepsilon=1$,
and $F_\varepsilon=F*\eta_\varepsilon\in C_c^\infty$, supported in $[-T-\varepsilon,T+\varepsilon]$.
Apply \cite[Thm.~8.4]{Paper45} to $F_\varepsilon$ at the window $T+\varepsilon$. Zero side:
$\Phi_{F_\varepsilon}=\Phi_F\Phi_{\eta_\varepsilon}$, $|\Phi_{\eta_\varepsilon}(s)|\le e^{\varepsilon|\Rel s-\frac12|}\le e^{\varepsilon/2}$
in the critical strip and $\Phi_{\eta_\varepsilon}\to1$ pointwise, so by (b) and dominated convergence
the zero sum (which for $F_\varepsilon$ converges absolutely, so that the symmetric order is
immaterial) tends to $\sum_\rho\Phi_F(\rho)\Phi_F(1-\rho)$. Polar term:
$\Phi_{F_\varepsilon}(0)\Phi_{F_\varepsilon}(1)\to\Phi_F(0)\Phi_F(1)$. Archimedean term:
$|\Fh_\varepsilon|^2=|\Fh|^2|\widehat{\eta_\varepsilon}|^2\le|\Fh|^2$ and $|m(\tau)|\le|C_0|+\log(1+|\tau|)$
(Lemma~\ref{lem:mbounds}(i),(ii); $m\le0$ for $|\tau|\le\tau_0$), so by (a) and dominated convergence
$\frac1{2\pi}\int m|\Fh_\varepsilon|^2\to\frac1{2\pi}\int m|\Fh|^2$. Prime term: $c_{F_\varepsilon}\to\cF$
uniformly, the sum runs over the finitely many $n<e^{2T+2\varepsilon}$, and for
$e^{2T}\le n<e^{2T+2\varepsilon}$ one has $\cF(\log n)=0$; so it tends to $-\sum_{n\in\NT}v_n\cF(\log n)$.
\end{proof}

The certificate of \S\ref{sec:cert} was computed with $\beta$ free and has $\beta_0=2.0026$,
$\beta_1=-0.0097$. We do not know whether optimal dual multipliers converge to $(2,0)$ as the
certificate class is enlarged; Section~\ref{sec:Z} explains why $(2,0)$ is the natural limit.

\section{Tail bounds and exhaustion}\label{sec:tails}

A certificate is an inequality $f\ge c$ on all of $\R$. The lemmas of this section reduce it to
a bounded interval. All of them rest on the following elementary facts.

\begin{lemma}\label{lem:mbounds}
\emph{(i)} $m$ is increasing on $[0,\infty)$. \emph{(ii)} For every $\tau>0$,
\[
\log\frac{\tau}{2\pi}-\frac1\tau-\frac1{2\tau^2}\ \le\ m(\tau)\ \le\ \log\frac{\tau}{2\pi}+\frac1\tau+\frac12\log\Big(1+\frac1{4\tau^2}\Big),
\]
and the lower bound is increasing in $\tau$. \emph{(iii)} $-\sum_{n\in\NT}v_n\cos(\tau\log n)\ge-\Aconst$, with
$\sup_\tau\sum_{n\in\NT}v_n\cos(\tau\log n)=\Aconst$.
\end{lemma}
\begin{proof}
(i) In $\Rel\psi(\frac14+iy)=-\gamma+\sum_{n\ge0}\big[\frac1{n+1}-\frac{n+\frac14}{(n+\frac14)^2+y^2}\big]$
every summand increases with $|y|$. (ii) Use $\psi(z)=\log z-\frac1{2z}-\int_0^\infty g(t)e^{-zt}dt$
for $\Rel z>0$, $g(t)=\frac1{e^t-1}-\frac1t+\frac12$ \cite[5.9.13]{DLMF}, at $z=\frac14+iy$, $y=\tau/2$.
Here $g(0^+)=0$, $g'(t)=t^{-2}-\frac14\sinh^{-2}(t/2)\in(0,\infty)$ and $g(\infty)=\frac12$, so one
integration by parts gives $|\int_0^\infty g e^{-zt}dt|=|z|^{-1}|\int_0^\infty g'e^{-zt}dt|\le\frac1{2|z|}\le\frac1{2y}$.
Also $\Rel\log z=\frac12\log(\frac1{16}+y^2)\in[\log y,\log y+\frac12\log(1+\frac1{16y^2})]$ and
$\Rel\frac1{2z}=\frac{1/8}{1/16+y^2}\in[0,\frac1{8y^2}]$. Substituting $y=\tau/2$ and $-\log\pi$
gives (ii); the derivative of the lower bound is $\tau^{-1}+\tau^{-2}+\tau^{-3}>0$. (iii) The
first is trivial. For the second, the numbers $\log p$, $p$ prime, are linearly independent over
$\Q$, so by Kronecker's theorem the vector $(\tau\log p)_{p<e^{2T}}$ is dense modulo $2\pi$; at
$\tau$ where all $\cos(\tau\log p)>1-\varepsilon$ also $\cos(\tau k\log p)>1-k^2\varepsilon$, so the
sum over all $n\in\NT$ exceeds $\Aconst-O(\varepsilon)$.
\end{proof}

\subsection{Concentration of the near-optimal feasible measures}\label{subsec:conc}

The program (P) has no minimizer (Theorem~\ref{thm:O}), so the reduction of its value to a bounded
interval cannot be routed through an extremal measure: it has to be uniform over the near-optimal
feasible set. Write
\[
\mathcal F_\varepsilon=\Big\{\mu\ \text{feasible for (P)}:\ \int\sT\,d\mu\le\lamP+\varepsilon\Big\},
\qquad B_*=\Aconst-C_0,\qquad G(\varepsilon)=\lamP+\varepsilon+B_*,
\]
which is non-empty for every $\varepsilon>0$.

\begin{theorem}[concentration of near-optimal feasible measures]\label{thm:P3}
Let $\varepsilon>0$ and $\eta\in(0,1)$.
\emph{(i)} Every $\mu\in\mathcal F_\varepsilon$ satisfies $\int(m-C_0)\,d\mu\le G(\varepsilon)$ and,
for every $R>0$, $\mu(\{|\tau|\ge R\})\le G(\varepsilon)/(m(R)-C_0)$.
\emph{(ii)} Consequently $\mu(\{|\tau|\ge R\})\le\eta$ for every $\mu\in\mathcal F_\varepsilon$ and
every $R\ge\max\{2,R(\varepsilon,\eta)\}$, where
\[
R(\varepsilon,\eta):=2\pi\exp\Big(1+C_0+\frac{G(\varepsilon)}\eta\Big).
\]
\emph{(iii)} Let $(\nu,\beta)$ be dual feasible of value $c$, put
$E_R:=\Aconst+\|\nu\|_{\TV}+|\beta_0|k_0(R)+|\beta_1|k_1(R)-c$, and let
$R\ge\max\{2,R(\varepsilon,\eta)\}$. Then, for every $\mu\in\mathcal F_\varepsilon$,
\[
0\le\lamP-c\le\int_{[-R,R]}(f-c)\,d\mu+\varepsilon_{\mathrm{tail}}(\mu,R),\qquad
\varepsilon_{\mathrm{tail}}(\mu,R)\le\int_{|\tau|>R}(m-C_0)\,d\mu+(E_R+C_0)^+\,\eta ,
\]
the first tail integral being $\le G(\varepsilon)$ always, and $\to0$ as $R\to\infty$ for each fixed
$\mu$.
\end{theorem}
\begin{proof}
(i) $m-C_0\ge0$ by Lemma~\ref{lem:mbounds}(i), and $\int\sT d\mu\ge\int m\,d\mu-\Aconst$ by
Lemma~\ref{lem:mbounds}(iii); hence $\int(m-C_0)d\mu\le\int\sT d\mu+\Aconst-C_0\le G(\varepsilon)$.
Chebyshev on $\{|\tau|\ge R\}$, where $m-C_0\ge m(R)-C_0$, gives the second bound. (ii) By
Lemma~\ref{lem:mbounds}(ii), $m(R)\ge\log(R/2\pi)-1$ for $R\ge2$, so $m(R)-C_0\ge G(\varepsilon)/\eta$
as soon as $\log(R/2\pi)\ge1+C_0+G(\varepsilon)/\eta$. (iii) By Proposition~\ref{prop:weak},
$\int f\,d\mu=\int\sT d\mu$ for feasible $\mu$, so $\int(f-c)d\mu=\int\sT d\mu-c\in[\lamP-c,\lamP+\varepsilon-c]$;
both pieces of the split at $R$ are nonnegative because $f\ge c$. On $|\tau|>R$,
$f-c\le m+E_R=(m-C_0)+(E_R+C_0)\le(m-C_0)+(E_R+C_0)^+$ by Lemmas~\ref{lem:mbounds}
and~\ref{lem:gbounds}; apply (i) to the first summand and (ii) to the second.
\end{proof}

The positive part in (iii) is needed: $C_0<0$, so $E_R+C_0$ can be negative for certificates of
small total variation, although it is positive for every certificate built in this paper
($\|\nu\|_{\TV}\approx48.7$).

\begin{corollary}[verification on a bounded interval]\label{cor:AP}
Suppose $(\nu,\beta)$ satisfies $f\ge c$ \emph{only} on $[-R,R]$ with
$R\ge\max\{2,R(\varepsilon,\eta)\}$, and put
\[
\mathsf W:=\Aconst+\|\nu\|_{\TV}+|\beta_0|\,\gh0(0)+|\beta_1|\,\gh1(0)-C_0 ,
\]
so that $f\ge-\mathsf W$ on all of $\R$. Then $c\le(\lamP+\varepsilon+\mathsf W\eta)/(1-\eta)$ for
every $\varepsilon>0$.
\end{corollary}
\begin{proof}
For $c<0$ there is nothing to prove. For $c\ge0$ take $\mu\in\mathcal F_\varepsilon$: then
$\lamP+\varepsilon\ge\int\sT d\mu=\int f\,d\mu\ge c\,\mu([-R,R])-\mathsf W\mu(|\tau|>R)\ge c(1-\eta)-\mathsf W\eta$
by Theorem~\ref{thm:P3}(ii).
\end{proof}

\begin{remark}\label{rem:P3exp}
The exponential dependence of $R(\varepsilon,\eta)$ on $1/\eta$ is forced and is not an artifact of
Chebyshev's inequality. Theorem~\ref{thm:Vp}(ii) constructs, for every $\delta>0$, feasible measures
$\mu'_{\eta,R}$ with mass $\ge\eta(1-o(1))$ beyond $R-\sqrt R$ and value
$\le\lamP+\delta+\eta(\log R+C_1)+o(\eta)$; such a $\mu'$ lies in $\mathcal F_\varepsilon$ as soon as
$\log R\lesssim(\varepsilon-\delta)/\eta$. So the radius at which mass $\eta$ is forced to have left
is $\log R\asymp\mathrm{const}/\eta$, exactly the shape of (ii); we do not know the optimal
constant, which lies between $\varepsilon$ and $G(\varepsilon)$. The radius $R(\varepsilon,\eta)$ has
the same shape as the envelope threshold of Proposition~\ref{prop:env}, with $B_*/\eta$ in place of
$\Aconst$ (numerical values in Appendix~\ref{app:diag}).
\end{remark}

\subsection{Lattice atoms: the graded periodic tail}

\begin{lemma}[graded periodic tail]\label{lem:F1}
Let $d=2T/M$, $u_j=jd$ for $M\le j\le J$ (so $u_M=2T$), $\nu=\sum_ja_j(\delta_{u_j}+\delta_{-u_j})$,
$P=2\pi/d$, so that $\nuh(\tau)=\sum_j2a_j\cos(u_j\tau)$ is $P$-periodic. Let
$f=\sT+\nuh+\beta_0\gh0+\beta_1\gh1$. Fix $\tau_1>0$, $N\ge1$, and for $x\in[\tau_1,\tau_1+P)$ put
\begin{align*}
W_k(x)&=m(x+kP)-\sum_{n\in\NT}v_n\cos\big((x+kP)\log n\big)-|\beta_0|k_0(x+kP)-|\beta_1|k_1(x+kP)\ \ (k<N),\\
W_N(x)&=m(x+NP)-\Aconst-|\beta_0|k_0(x+NP)-|\beta_1|k_1(x+NP).
\end{align*}
If $\nuh(x)+W_k(x)\ge c$ for all $x\in[\tau_1,\tau_1+P)$ and all $0\le k\le N$, then $f(\tau)\ge c$
for all $\tau\ge\tau_1$.
\end{lemma}
\begin{proof}
Every $\tau\ge\tau_1$ is uniquely $x+kP$ with $x\in[\tau_1,\tau_1+P)$, $k\ge0$, and
$\nuh(\tau)=\nuh(x)$ since $u_jP=2\pi j$. For $k<N$ the $k$-th row is $f(x+kP)$ with the exact
prime term and the two multiplier terms replaced by the lower bounds of Lemma~\ref{lem:gbounds}.
For $k\ge N$: $m(x+kP)\ge m(x+NP)$, $k_i(x+kP)\le k_i(x+NP)$, and the prime term is $\ge-\Aconst$.
\end{proof}

The prime term is never required to be periodic: it is evaluated exactly on the first $N$
copies of one period of $\nuh$ and replaced by its infimum beyond. The lemma therefore holds
for any finite set of contacts; the incommensurability of $\log2$ and $\log3$ plays no role.
It is a sufficient criterion only (the rows are conservative in $\beta$ by
$|\beta_0|k_0+|\beta_1|k_1$).

\subsection{Exponential columns: the monotone tail}

\begin{proposition}[exponential columns]\label{prop:T0}
For $s>0$ and an integer $p\ge0$ let $d\nu_{s,p}(u)=\frac12(|u|-2T)^pe^{-s(|u|-2T)}\mathbf 1_{|u|\ge2T}du$,
a positive finite measure of mass $p!/s^{p+1}$ supported outside the window. Then
\[
\nuh_{s,p}(\tau)=\Rel\Big[e^{2iT\tau}\frac{p!}{(s-i\tau)^{p+1}}\Big],\qquad
|\nuh_{s,p}(\tau)|\le\operatorname{cap}_{s,p}(\tau):=\frac{p!}{(s^2+\tau^2)^{(p+1)/2}},
\]
and $\operatorname{cap}_{s,p}$ is decreasing in $\tau\ge0$.
\end{proposition}
\begin{proof}
$\nuh_{s,p}(\tau)=\int_{2T}^\infty(u-2T)^pe^{-s(u-2T)}\cos\tau u\,du=\Rel\big[e^{2iT\tau}\int_0^\infty t^pe^{-(s-i\tau)t}dt\big]$.
\end{proof}

\begin{lemma}[monotone tail]\label{lem:T}
Let $\nu=\sum_{k=1}^Kc_k\nu_{s_k,p_k}+\sum_{j=1}^Ja_j(\delta_{u_j}+\delta_{-u_j})$ with $u_j\ge2T$, and
$f=\sT+\nuh+\beta_0\gh0+\beta_1\gh1$. Define for $\tau>0$
\[
L(\tau)=\log\frac\tau{2\pi}-\frac1\tau-\frac1{2\tau^2}-\Aconst-\sum_k|c_k|\operatorname{cap}_{s_k,p_k}(\tau)
-2\sum_j|a_j|-|\beta_0|k_0(\tau)-|\beta_1|k_1(\tau).
\]
Then $f\ge L$ on $(0,\infty)$, $L$ is nondecreasing, and $L(\tau_1)\ge c$ for one $\tau_1$ implies
$f\ge c$ on $[\tau_1,\infty)$.
\end{lemma}
\begin{proof}
Lemma~\ref{lem:mbounds}(ii),(iii), Proposition~\ref{prop:T0}, the trivial bound on the atoms and
Lemma~\ref{lem:gbounds}; every term of $L$ is nondecreasing.
\end{proof}

\begin{corollary}[cost of the tail]\label{cor:T2}
For an atom-free $\nu$ with $A:=\sum_k|c_k|p_k!$ and $A':=|\beta_0|(2\cosh T+\sinh T)+|\beta_1|8T\sinh T$,
the tail inequality $L(\tau_1)\ge c$ holds for every $\tau_1\ge\max\{A+A'+2,\ 2\pi e^{c+\Aconst+1}\}$,
so the tail costs linearly in the $\ell^1$ size of the certificate. For atoms it costs
exponentially, $\tau_1\approx2\pi e^{\Aconst+2\sum|a_j|}$, which is why Lemma~\ref{lem:F1} and not
Lemma~\ref{lem:T} is used for lattice certificates. Neither cost is uniform over the dual class:
by Theorem~\ref{thm:O}, $\|\nu\|\to\infty$ along every maximizing sequence.
\end{corollary}
\begin{proof}
For $\tau\ge1$, $\operatorname{cap}_{s,p}(\tau)\le p!/\tau$ and $|\beta_i|k_i\le A'/\tau$, so
$L(\tau)\ge\log\frac\tau{2\pi}-\Aconst-\frac{A+A'+2}{\tau}$.
\end{proof}

\begin{lemma}[hybrid tail]\label{lem:hybrid}
Let $\nu=\nu_{\mathrm{lat}}+\sum_\ell c_\ell\nu_{s_\ell,p_\ell}$ with $\nu_{\mathrm{lat}}$ as in Lemma~\ref{lem:F1}.
For $x\in[\tau_1,\tau_1+P)$ define $H_k(x)=\nuh_{\mathrm{lat}}(x)+W_k(x)+\sum_\ell c_\ell\nuh_{s_\ell,p_\ell}(x+kP)$
for $0\le k<N$, and $H_N(x)=\nuh_{\mathrm{lat}}(x)+W_N(x)-\sum_\ell|c_\ell|\operatorname{cap}_{s_\ell,p_\ell}(x+NP)$.
If $H_k\ge c$ on $[\tau_1,\tau_1+P)$ for all $0\le k\le N$ then $f\ge c$ on $[\tau_1,\infty)$.
\end{lemma}
\begin{proof}
As in Lemma~\ref{lem:F1}; on the rows $k<N$ the column term is evaluated exactly, and for
$k\ge N$ it is bounded by $-\sum_\ell|c_\ell|\operatorname{cap}(x+kP)\ge-\sum_\ell|c_\ell|\operatorname{cap}(x+NP)$
since $\operatorname{cap}$ decreases.
\end{proof}

\subsection{Exhaustion by exponential columns}

\begin{theorem}[exponential columns exhaust the dual]\label{thm:E}
Let $\mathcal E$ be the real span of $\{\nu_{s,p}:s>0,p\ge0\}$. Then
$\sup\{\inf_\tau[\sT+\nuh+\beta\cdot\widehat g]:\ \nu\in\mathcal E,\beta\in\R^2\}=\lamP$, and the same
holds with $\mathcal E$ replaced by the absolutely continuous measures with bounded density and
compact support in $\{|u|\ge2T\}$.
\end{theorem}
\begin{proof}
``$\le$'' is weak duality. For ``$\ge$'' fix $\varepsilon>0$ and, by Theorem~\ref{thm:D1}, an atomic
$\nu_0=\sum_{j\le J}a_j(\delta_{u_j}+\delta_{-u_j})$ and $\beta$ with
$c_0:=\inf[\sT+\nuh_0+\beta\cdot\widehat g]>\lamP-\varepsilon/2$. By Lemma~\ref{lem:T} there is
$\tau_1<\infty$ with $\sT+\nuh_0+\beta\cdot\widehat g\ge c_0-\varepsilon/8$ on $[\tau_1,\infty)$,
and the lemma uses only $\sum|a_j|$, so the same bound holds for every signed measure of the
same total variation supported in $\{|u|\ge2T\}$. Replace each atom $a_j\delta_{u_j}$ by
$a_j$ times the uniform probability density on $(u_j,u_j+\eta)$: since
$|\cos u\tau-\cos u_j\tau|\le|\tau|\eta$, the transform changes by at most $2\sum|a_j|\tau_1\eta$
on $[0,\tau_1]$, which is $<\varepsilon/8$ for small $\eta$; the new $\nu_1$ has bounded density
and compact support and value $\ge c_0-\varepsilon/4$. Finally $\{t^pe^{-t}\}_{p\ge0}$ is total in
$L^1(0,\infty)$ (if $g\in L^\infty$ annihilates them, its Laplace transform is holomorphic on
$\Rel z>0$ with all derivatives vanishing at $z=1$, hence $g=0$), so there is $\nu_2\in\mathcal E$
with $\|\nu_2-\nu_1\|_{\TV}<\varepsilon/8$, and $\|\nuh_2-\nuh_1\|_\infty\le\|\nu_2-\nu_1\|_{\TV}$ on all of $\R$.
\end{proof}

\begin{proposition}[the canonical column; band-limitation]\label{prop:M}
Let $d\nu_K=\frac12K(|u|)\mathbf 1_{|u|\ge2T}du$, a positive measure of mass
$c_K=\int_{2T}^\infty K$ (Theorem~\ref{thm:S}(iii)). Then for all $\tau$
\[
\sT(\tau)+\nuh_K(\tau)=\big(C_0+c_K\big)+\int_0^{2T}(1-\cos\tau u)K(u)\,du-\sum_{n\in\NT}v_n\cos(\tau\log n),
\]
a function whose distributional Fourier transform in $u$ is supported in $[-2T,2T]$ (it still grows
like $\log|\tau|$: the singularity of $K$ at $u=0$ lies inside the window). Moreover
$\nu_K=2e^{-T}\nu_{\frac12,0}(1+O(e^{-2u}))$:
the archimedean L\'evy density selects the scale $s=\frac12$ among the exponential columns.
\end{proposition}
\begin{proof}
Split $\int_0^\infty(1-\cos\tau u)K$ at $2T$ and use $\nuh_K(\tau)=\int_{2T}^\infty K\cos\tau u$.
\end{proof}

This identity is not a certificate (as it must not be, by Theorem~\ref{thm:O}): its right side at
$\tau=0$ is $C_0+c_K-\Aconst<0$. It says that the outside-window L\'evy tail removes from the symbol
every frequency outside the window; apart from the $\log|\tau|$ growth generated by $K$ near $u=0$,
what remains is carried by frequencies in $[0,2T]$. Section~\ref{sec:Z} identifies this remainder.

\section{An explicit certificate at \texorpdfstring{$T=\tfrac12\log3$}{T = (1/2) log 3}}\label{sec:cert}

At $T=\frac12\log3$ the window is $[-\log3,\log3]$, the contact set is $\NT=\{2\}$, and
$v_2=2\log2/\sqrt2=0.98025814\ldots=\Aconst$. Positivity of the primitive form here is known
\cite{Paper42}, and positivity of the complete form by \cite[Thm.~1.2]{Zhu2026}. We give an independent
proof by an explicit dual feasible point.

\begin{theorem}[explicit certificate]\label{thm:C}
Let $T=\frac12\log3$, $d=2T/24=\frac{\log3}{24}$,
$u_j=jd$ for $24\le j\le460$ (so $u_{24}=\log3$ and $u_{460}=21.0567\ldots$), let
$a_{24},\dots,a_{460}$ be the $437$ rational numbers given as exact decimal strings in the
certificate data accompanying this paper, and let
$\beta_0=2.00256170925026\ldots$, $\beta_1=-0.00965443404485\ldots$. Put
$\nu=\sum_ja_j(\delta_{u_j}+\delta_{-u_j})$, so $\|\nu\|_{\TV}=2\sum|a_j|=90.75$, $\nuh(0)=1.7337$.
Then
\[
\sT(\tau)+\nuh(\tau)+\beta_0\gh0(\tau)+\beta_1\gh1(\tau)\ \ge\ c:=8.8577\cdot10^{-5}\qquad\text{for every }\tau\in\R,
\]
and therefore $\lamA\ge\lamP\ge8.8577\cdot10^{-5}>0$. The proof includes a computer verification
in interval arithmetic (Appendix~\ref{app:encl}).
\end{theorem}

\begin{proof}
The function $f$ is even, so it suffices to take $\tau\ge0$. Split $[0,\infty)$ as $[0,\tau_1]$
with $\tau_1=30$, where $f$ itself is bounded below, and $[\tau_1,\infty)$, where
Lemma~\ref{lem:F1} with $N=2$ and $P=2\pi/d=137.2608\ldots$ reduces the inequality to the three
rows $\nuh+W_k\ge c$, $k=0,1,2$, on the single period $[30,167.2608]$.

Each of the four functions ($f$ on $[0,\tau_1]$ and the three rows) is covered by the exhaustive
second-order procedure of Proposition~\ref{prop:TAY} (Appendix~\ref{app:encl}). The second-derivative
bound is $M_2=6264.92$ for $f$, computed from the exact decimals; the three rows admit the smaller
constant $6262.99$, and the larger one is used for all four. The initial cell width is $0.02$, the
minimal width $w_{\min}=10^{-9}$; there is no cap on the number of cells and no cell is discarded.
The cover uses $2\,378\,593$ evaluations, leaves no unresolved interval, and consists of $20$
segments tiling the four intervals exactly.

The target in the cover is $8.8578\cdot10^{-5}$ minus a floating-point allowance
$\varepsilon_{\mathrm{fp}}=10^{-9}$. This allowance exceeds the a priori error of the double-precision
evaluator: argument rounding in the $437$ cosines contributes at most
$2\sum_j|a_j|\cdot u_{460}\cdot(\tau_1+NP)\cdot2^{-53}<10^{-10}$, the Taylor remainder at the minimal
width at most $M_2w_{\min}^2/8<10^{-15}$, and the summation of the $445$ terms together with the
evaluation of $\psi$ less than $10^{-12}$.

The constants $H_0,H_1$, $\psi'(\frac14)$, $k_0,k_1$ and the values of $m$ at the four minimizers
are enclosed in interval arithmetic at $40$ digits, $m$ through the recurrence
$\Rel\psi(z)=\Rel\psi(z+J)-\sum_{j<J}\Rel(z+j)^{-1}$ and the asymptotic series with Stieltjes'
remainder bound. The four minima $f(17.645)$, $h_0(37.815)$, $h_1$, $h_2(58.817)$ are strictly
positive intervals with lower ends $\ge8.8578\cdot10^{-5}$; the smallest is
$h_0(37.8153)=8.85786\cdot10^{-5}$. The certificate data and the verification statistics accompany
the paper (\S\ref{sec:repro}), and an independent replay from the decimal strings reproduces all
figures.
\end{proof}

\begin{remark}[how the certificate was found]
The atoms were placed on the lattice $u_j=j\cdot2T/M$ because Lemma~\ref{lem:F1} then controls the
tail on one period; the resolution near the window edge, rather than the range of the atoms, turned
out to be the decisive parameter. A linear program in $(a_j,\beta,c)$ with cutting planes on a
$\tau$-grid produced an approximate profile; since values on a grid are unreliable at the required
accuracy, the profile was re-optimized in $128$ blocks on a grid $10^3$ times finer (details in
Appendix~\ref{app:diag}). The result has three visible features:
\begin{enumerate}
\item $\beta_0=2.0026$ and $\beta_1\approx0$: the optimizer restores the two polar terms of the
explicit formula that the primitive class removes (Proposition~\ref{prop:P});
\item the weights $a_j$ alternate in sign at $70\%$ of consecutive lattice sites, decay smoothly away
from the window edge (mean $|a_j|$ from $0.27$ near $u=2T$ to $0.02$ at $u\approx40$), and show no
arithmetic structure (no $\log2$-spacing, no equal amplitudes);
\item the contact set $\{f-c<10^{-5}\}$ is a comb: a band inside $[0,13]$, just beyond the last
sign change of $\sT$ at $\tau=10.53$, plus one phase per tail row. This is the complementary
slackness behind Theorem~\ref{thm:O}: the contact set must grow as $c\uparrow\lamP$ and can never be
all of $\R$.
\end{enumerate}
None of these features is visible in the Feshbach certificate of \cite{Paper42}.
\end{remark}

\subsection{Comparison}\label{subsec:compare}
The value $c$ is about $9\%$ of the numerical Galerkin upper bound on $\lamA$ (Appendix~\ref{app:diag}); by
Theorem~\ref{thm:D1} the remainder is available to the same class of certificates, and by
Theorem~\ref{thm:O} it is not attained. Zhu \cite[Thm.~1.2]{Zhu2026} proves
$Q(f)\ge8.9\cdot10^{-18}\|f\|^2$ for $\supp f\subseteq[-0.8,0.8]$ (autocorrelation support
$1.6>\log3$) for the complete form $Q=W_T$; the two results are compatible and his window is larger.
His proof combines a pointwise envelope of $\sT$ with a certified Cholesky factorization of one
$200\times200$ Legendre block; ours is an explicit dual feasible point, whose ingredients (polar
multiplier, outside-window measure, contact comb) are interpretable, and whose existence for every
$c<\lamP$ is guaranteed by Theorem~\ref{thm:D1}. Zhu's barrier for envelope certificates
\cite[Thm.~1.4]{Zhu2026} and Proposition~\ref{prop:env} below describe the same phenomenon from the
primal and the dual side.

\section{Obstructions}\label{sec:negative}

\subsection{The moment map has interior}

\begin{lemma}[moment interior]\label{lem:W}
For $F\in L^2(\IT)$ real put $Q_0(F)=\int\cF\cosh\frac u2=A B$ and
$Q_1(F)=\int\cF\,u\sinh\frac u2=A'B-AB'$. The image of the set of unit vectors under
$F\mapsto(Q_0,Q_1)$ contains a neighborhood of $0$ in $\R^2$, realised by $C_c^\infty(\IT)$
functions; equivalently, $0$ is an interior point of the (convex) moment image
$\mathfrak m(\mu)=(\int\gh0d\mu,\int\gh1d\mu)$ of the feasible set of (P) without the two
constraints, and the interior is realised by measures of finite value.
\end{lemma}
\begin{proof}
The formulas for $Q_0,Q_1$ follow from expanding $\cosh\frac{t-s}2$ and $(t-s)\sinh\frac{t-s}2$
in exponentials. The functions $e^{t/2},e^{-t/2},te^{t/2},te^{-t/2}$ are linearly independent
on $\IT$, and a nontrivial combination is real-analytic, so it cannot annihilate
$C_c^\infty(\IT)$: the linear map $F\mapsto(A,B,A',B')$ is onto $\R^4$ from $C_c^\infty(\IT)$.
If $0$ were not interior, some $(t_0,t_1)\ne0$ would satisfy $t_0Q_0(F)+t_1Q_1(F)\ge0$ for all
$F$ (normalization only rescales $Q_j$ by a positive factor). Take $(A,B,A',B')=(1,0,0,\mp1)$:
$Q_0=0$, $Q_1=\pm1$, so $t_1=0$; take $(1,\pm1,0,0)$: $Q_1=0$, $Q_0=\pm1$, so $t_0=0$. A finite
convex combination of the four witnesses covers a ball, and each has finite value because
$\Fh$ is Schwartz.
\end{proof}

\subsection{Mass against value}

\begin{theorem}[logarithmic exchange]\label{thm:Vp}
Let $B_*=\Aconst-C_0$ ($=7.621$ at $T\in(\frac12\log3,\log2]$).
\emph{(i)} For every feasible $\mu$ of (P), with $\varepsilon:=\mu(\{|\tau|\ge R\})$,
$\int\sT d\mu\ge-B_*+\varepsilon\,(m(R)-m(0))$, and $m(R)-m(0)=\log R+O(1)$.
\emph{(ii)} For every $\delta>0$ there is $C_1<\infty$ such that for all large $R$ and all
$\varepsilon\in[0,1]$ there is a feasible $\mu'_{\varepsilon,R}$ with
$\mu'_{\varepsilon,R}(\{|\tau|\ge R-\sqrt R\})\ge\varepsilon(1-o(1))$ and
$\int\sT d\mu'_{\varepsilon,R}\le\lamP+\delta+\varepsilon(\log R+C_1)+o(\varepsilon)$.
\emph{(iii)} Consequently no inequality $\int\sT d\mu\ge-B_*+\varepsilon\,\omega(R)$ with
$\omega(R)/\log R\to\infty$ holds on the feasible set: the exchange rate between escaped mass and
value is $\Theta(\log R)$, and converting the depth $B_*$ into a nonnegative value through escaped
mass alone requires $R\gtrsim e^{B_*/\varepsilon}$.
\end{theorem}
\begin{proof}
(i) $\sT\ge-B_*$ everywhere and $\sT\ge m(R)-\Aconst$ on $|\tau|\ge R$; split the integral.
(ii) Pick $F_0\in C_c^\infty(\IT)$ with $c_{F_0}\in\Kp$ and $\mathfrak a_T(F_0)\le\lamP+\delta$
(such $F_0$ are dense in the primitive class), $\mu_0=\mu_{F_0}$, which has a rapidly decaying
density. Set $\mu_{\varepsilon,R}=(1-\varepsilon)\mu_0+\frac\varepsilon2(\mu_0(\cdot-R)+\mu_0(\cdot+R))$;
then $\check\mu_{\varepsilon,R}=[(1-\varepsilon)+\varepsilon\cos Ru]\check\mu_0$ is still supported in
$[-2T,2T]$, so $\mu_{\varepsilon,R}$ is feasible for the relaxation without the two moment
constraints; its moment defect is $\varepsilon\,d(R)$ with $d(R)=(\int g_0\check\mu_0\cos Ru,\int g_1\check\mu_0\cos Ru)\to0$
by Riemann--Lebesgue. By Lemma~\ref{lem:W} there are $r_0>0$, $V^*<\infty$ and, for each $w$
with $|w|\le r_0$, a relaxed-feasible $\mu_w$ of value $\le V^*$ and moment $w$. Put
$\lambda=\varepsilon|d(R)|/(r_0+\varepsilon|d(R)|)$ and choose $\mu_w$ with
$w=-\frac{1-\lambda}\lambda\varepsilon d(R)$ (of norm $r_0$); then
$\mu'_{\varepsilon,R}=(1-\lambda)\mu_{\varepsilon,R}+\lambda\mu_w$ has moment $0$, is feasible for
(P), and its value exceeds that of $\mu_{\varepsilon,R}$ by at most
$\lambda(B_*+V^*)=o(\varepsilon)$. The value of $\mu_{\varepsilon,R}$ is
$(1-\varepsilon)\mathfrak a_T(F_0)+\varepsilon\int\sT(\tau\pm R)d\mu_0$, and by
Lemma~\ref{lem:mbounds}(ii) $\int m(\tau+R)d\mu_0\le\log R+\int|\tau|d\mu_0/R+O(1)$ while the prime
part is bounded by $\Aconst$. (iii) Apply (i) to $\mu'_{\varepsilon,R}$ and let $\varepsilon=1$.
\end{proof}

\begin{remark}
Theorem~\ref{thm:Vp} excludes a rate lemma depending on $R$ alone, uniformly over the feasible
set. It does not exclude rate statements that depend on the data of a certificate class (the
lattice step $d$, the number of atoms, the $\ell^1$ size); those are the object of the finite
program of \S\ref{sec:path}.
\end{remark}

\subsection{The envelope threshold}

\begin{proposition}[envelope threshold]\label{prop:env}
\emph{(i)} Every certificate $(\nu,\beta)$ satisfies $f(\tau)\ge\log\frac\tau{2\pi}-\frac1\tau-\frac1{2\tau^2}-\Aconst-\|\nu\|_{\TV}-|\beta_0|k_0(\tau)-|\beta_1|k_1(\tau)$,
so $f\ge0$ holds automatically for $\tau\ge\tau_{\mathrm{env}}:=2\pi e^{\Aconst+\|\nu\|_{\TV}+o(1)}$.
\emph{(ii)} No bound of the form $f\ge m(\tau)-E(\tau)$ with $E$ depending on $\tau$ only through
an upper bound of the prime comb can replace $\Aconst$ by a smaller constant, because
$\sup_\tau\sum_{n\in\NT}v_n\cos(\tau\log n)=\Aconst$ (Lemma~\ref{lem:mbounds}(iii)).
\emph{(iii)} $\Aconst=\sum_{n<e^{2T}}2\Lambda(n)/\sqrt n=(4+o(1))e^{T}$ by the prime number theorem,
so $\tau_{\mathrm{env}}(T)$ grows doubly exponentially in $T$. The multiplier-free part alone gives
the threshold $2\pi e^{\Aconst}$, equal to $16.7$ and $59.5$ at $T=\frac12\log3$ and $\log2$ (the last
sign change of $\sT$ is at $10.53$ and $46.13$); for a certificate of total variation $\|\nu\|_{\TV}$
the threshold is $2\pi e^{\Aconst+\|\nu\|_{\TV}}$, about $10^{40}$ for the certificate of
\S\ref{sec:cert}, which is why its tail is handled by the periodic Lemma~\ref{lem:F1} and not by this
bound.
\end{proposition}
\begin{proof}
(i) Lemma~\ref{lem:mbounds} and $|\nuh|\le\|\nu\|_{\TV}$. (ii),(iii) as stated.
\end{proof}

Thus a pointwise argument can stop at $\tau_{\mathrm{env}}(T)$ and at no earlier point, and
$2\pi e^{\Aconst}$, doubly exponential in $T$, is the shortest interval such an argument can
reach; this is the dual counterpart of \cite[Thm.~1.4]{Zhu2026}. We do not know whether every
certificate must do work on all of $[0,2\pi e^{\Aconst}]$; the periodic tail lemmas show that a
lattice certificate is checked on $[0,\tau_1+NP]$, with $\tau_1+NP\approx442$ at $T=\frac12\log3$.

\subsection{The two-prime windows}\label{subsec:twoprime}
For $\frac12\log3<T\le\log2$ the contact set is $\{2,3\}$, $\Aconst=2.24883$ and $B_*=7.62101$.
No certificate of positive value is known in this range: the methods of this paper certify
nothing at any $T>\frac12\log3$, and numerically the lattice certificates of \S\ref{sec:cert}
remain negative there (Appendix~\ref{app:diag}). Positivity in this range rests on \cite{Paper42}
for the primitive form and on \cite[Thm.~1.2]{Zhu2026} for the complete form. By
Theorem~\ref{thm:Vp}, no mass-based rate theorem predicts the size of certificate class that would
be needed.


\section{The dual variable is the explicit formula}\label{sec:Z}

So far every certificate has been a list of numbers produced by an optimizer. This section
identifies these numbers. Up to an arbitrarily small uniform correction, every certificate of the
complete Weil form is the Weil explicit formula evaluated on a family of modulated test functions,
and its value at $\tau$ is a smoothing of the zero measure at $\tau$.

\subsection{Flat-top weights}
\begin{definition}\label{def:W}
$\Wcl$ is the set of even measurable $w:\R\to\R$ with $w\equiv1$ on $[-2T,2T]$ and
$|w(u)|\le Ce^{-(\frac12+\delta)|u|}$ for $|u|\ge2T$, for some $C,\delta>0$. $\Wcl^{\mathrm{reg}}\subset\Wcl$
consists of the $w\in C^\infty(\R)$ with compact support. For $w\in\Wcl$ define the even measure
\begin{equation}\label{eq:nuw}
d\nu_w=\Big[\tfrac12\big(1-w(u)\big)K(|u|)+2w(u)\cosh\tfrac u2\Big]\mathbf 1_{|u|\ge2T}\,du
-\sum_{n\notin\NT}\tfrac{v_n}2\,w(\log n)\,\big(\delta_{\log n}+\delta_{-\log n}\big),
\end{equation}
and $f_w:=\sT+\nuh_w+2\gh0$.
\end{definition}
The measure $\nu_w$ is finite ($K$ is integrable on $[2T,\infty)$, $|w|\cosh\frac u2$ is
integrable by the decay condition, and $\sum_nv_n|w(\log n)|\le2C\sum\Lambda(n)n^{-1-\delta}$) and
supported in $\{|u|\ge2T\}$. Hence $(\nu_w,2,0)$ is a dual feasible point, and by
Propositions~\ref{prop:weak} and \ref{prop:P}, $W_T\ge\inf f_w$ on $L^2(\IT)$.

\begin{theorem}\label{thm:Z}
\emph{(i)} For every $w\in\Wcl$ and $\tau\in\R$,
\begin{equation}\label{eq:fw}
\begin{split}
f_w(\tau)=C_0+\int_0^\infty\big(1-w(u)\cos\tau u\big)K(u)\,du&-\sum_{n\ge2}v_n\,w(\log n)\cos(\tau\log n)\\
&+2\int_\R w(u)\cosh\tfrac u2\cos(\tau u)\,du,
\end{split}
\end{equation}
which is the right-hand side of the Weil explicit formula for the even test function
$h_\tau(u)=w(u)\cos(\tau u)$.
\emph{(ii)} For $w\in\Wcl^{\mathrm{reg}}$, with $\wh(t)=\int w(u)e^{itu}du$ (an even entire function
of exponential type),
\begin{equation}\label{eq:Z}
f_w(\tau)=\sum_\rho\wh(\tau-\gamma_\rho),\qquad\rho=\tfrac12+i\gamma_\rho,
\end{equation}
the sum running over the non-trivial zeros of $\zeta$ with multiplicity, and converging absolutely.
Under RH all $\gamma_\rho$ are real, and $f_w(\tau)=\sum_\gamma\wh(\tau-\gamma)$ is the smoothing of
the zero-counting measure by the kernel $\wh$.
\end{theorem}
\begin{proof}
(i) By \eqref{eq:m}, $m(\tau)+\int_{2T}^\infty(1-w)K\cos\tau u\,du=C_0+\int_0^{2T}(1-\cos\tau u)K
+\int_{2T}^\infty K-\int_{2T}^\infty wK\cos\tau u=C_0+\int_0^\infty(1-w\cos\tau u)K$, using $w=1$ on
$[0,2T]$. The prime terms inside and outside the window combine into $-\sum_{n\ge2}v_nw(\log n)\cos(\tau\log n)$,
again because $w(\log n)=1$ for $n\in\NT$, and $2\gh0+2\int_{|u|\ge2T}w\cosh\frac u2\cos\tau u=2\int_\R w\cosh\frac u2\cos\tau u$.
(ii) For even $h\in C_c^\infty(\R)$ the Weil explicit formula \cite{Weil52}, \cite[Thm.~5.12]{IK}, in the
normalization of this paper, reads
\[
\sum_\rho\widehat h(\gamma_\rho)=2\int_\R h(u)\cosh\tfrac u2\,du-\sum_{n\ge2}v_n\,h(\log n)+\frac1{2\pi}\int_\R m(t)\,\widehat h(t)\,dt .
\]
Apply it to $h_\tau=w\cos(\tau\,\cdot)$, whose transform is $\frac12[\wh(t-\tau)+\wh(t+\tau)]$. The
zero sum is $\sum_\rho\wh(\tau-\gamma_\rho)$, since $\wh$ is even and the multiset $\{\gamma_\rho\}$ is
invariant under $\gamma\mapsto-\gamma$ ($\rho\mapsto1-\rho$). The polar and prime terms are the
last two terms of \eqref{eq:fw}. For the archimedean term insert \eqref{eq:m}: by Fourier
inversion $\frac1{2\pi}\int\wh(t\mp\tau)\,dt=w(0)=1$ and
$\frac1{2\pi}\int\wh(t\mp\tau)\cos(tu)\,dt=w(u)\cos(\tau u)$, and Fubini is justified by
$\int|\wh(t-\tau)|\int_0^\infty\min(2,t^2u^2)K(u)\,du\,dt<\infty$; the result is
$C_0+\int_0^\infty(1-w(u)\cos\tau u)K(u)\,du$. Absolute convergence of the zero sum holds because $\wh$ is
Schwartz on horizontal lines of the strip $|\Iml z|\le\frac12$, uniformly, and $N(t)\ll t\log t$.
\end{proof}

\begin{remark}[kernels on the Berkovich line]\label{rem:normline}
Let $\mathbb A^{1,\mathrm{an}}_\Z$ be the Berkovich analytic line over $\Z$ \cite{Poineau2010}, with
coordinate $X$, and put $\ell(x)=-\log|X|_x$. For an even distribution $H$ on $\R$ and measures $\mu_f$
with $\ell_*\mu_f=f\,du$, $f\in C_c^\infty(\IT)$, the kernel $G(x,y)=H(\ell(x)-\ell(y))$ has energy
$\iint G\,d\mu_f\,d\overline{\mu_f}=\frac1{2\pi}\int\widehat H|\widehat f|^2$. This equals $W_T(f)$ for
all such $f$ exactly when $H-k_T$ is supported in $\{|u|\ge2T\}$, where $\widehat{k_T}=\sT+2\gh0$
(Proposition~\ref{prop:P}). Thus positive definite kernels of this form realizing $W_T$ on the window
are the certificates of problem (D) at $\beta=(2,0)$, and by Theorem~\ref{thm:Zexh}(i) they exist for
every value $c<\lamW$ when $T\ge\frac12\log2$.
\end{remark}

\begin{theorem}[exhaustion by flat-top weights]\label{thm:Zexh}
Let $T\ge\frac12\log2$. \emph{(i)} $\sup_{w\in\Wcl^{\mathrm{reg}}}\inf_{\tau\in\R}f_w(\tau)=\lamW$.
\emph{(ii)} $\sup\{\inf_\tau[f_w+(\beta_0-2)\gh0+\beta_1\gh1]:\ w\in\Wcl^{\mathrm{reg}},\beta\in\R^2\}=\lamP$.
\end{theorem}
Thus every certificate of the complete Weil form is a flat-top explicit-formula certificate up to
an arbitrarily small uniform correction, and every certificate of the primitive form is one plus a
correction in the two primitivity multipliers.
\begin{proof}
``$\le$'' is weak duality in both cases. For ``$\ge$'' in (i): the proof of Theorem~\ref{thm:D1}
with the two $\widehat g$-constraints removed gives $\lamW=\sup\inf[\sT+2\gh0+\nuh]$ over finite
atomic $\nu$ outside the window, and the argument of Theorem~\ref{thm:E} (with $\beta=(2,0)$
fixed) replaces the atoms by a bounded even density $h$ with compact support in $\{|u|\ge2T\}$,
losing at most $\varepsilon/2$. It remains to realize $h\,du$ as $\nu_w$ up to total variation
$\varepsilon/2$, with $w\in\Wcl^{\mathrm{reg}}$; the transform then moves by at most $\varepsilon/2$
uniformly. On $u\ge2T$ the denominator $d_*(u)=2\cosh\frac u2-\frac12K(u)$ is positive and increasing
($D(0.3)>0$ and $D'>0$), because $2T\ge\log2>0.3$; put
$w_h(u)=\big(h(u)-\tfrac12K(u)\big)/d_*(u)$ for $u\ge2T$, extended evenly and by $1$ on the window.
Then $|w_h|\ll e^{-|u|}$ and the absolutely continuous part of $\nu_{w_h}$ is exactly $h\,du$, while
its atomic part is $-\sum_{n\notin\NT}\frac{v_n}2w_h(\log n)(\delta_{\pm\log n})$, of total variation
$\ll\sum_n\Lambda(n)n^{-3/2}<\infty$. Now modify, in this order: (a) replace $h$ on $[2T,2T+\eta]$
by a smooth function matching $d_*(u)+\frac12K(u)$ to infinite order at $2T$ (so that $w_h\equiv1$
smoothly across the window edge); (b) mollify $h$ and cut $w_h$ off smoothly beyond a large $U_1$
(where $w_h\ll e^{-u}$; the a.c.\ part of $\nu_w$ becomes $\frac12K$ there, of total variation
$\ll e^{-U_1/2}$, and all atoms with $\log n>U_1$ disappear because $w=0$ there); (c) for the finitely
many $n\notin\NT$ with $\log n\le U_1$, replace $h$ by $\frac12K$ on disjoint intervals of small total
length around $\log n$ (so that $w_h$ vanishes there and the remaining atoms disappear). Each step
changes $\nu_{w_h}$ by an arbitrarily small total variation, and the result is in
$\Wcl^{\mathrm{reg}}$. (ii) is the same argument applied to the certificates of Theorem~\ref{thm:E}
with their own $\beta$.
\end{proof}

\begin{corollary}[the Riemann Hypothesis in flat-kernel form]\label{cor:RH}
The following are equivalent: \emph{(a)} RH; \emph{(b)} for every $T\ge\frac12\log2$ there is
$w\in\Wcl^{\mathrm{reg}}$ with $\sum_\rho\wh(\tau-\gamma_\rho)\ge0$ for all $\tau\in\R$; \emph{(c)} for
every $T\ge\frac12\log2$ there is $w\in\Wcl^{\mathrm{reg}}$ such that the elementary function
\eqref{eq:fw} is nonnegative on $\R$.
\end{corollary}
\begin{proof}
(b)$\iff$(c) is Theorem~\ref{thm:Z}. (c)$\Rightarrow$(a): $f_w\ge0$ gives $W_T\ge0$ on $L^2(\IT)$
for every $T$, hence $W\ge0$ on all compactly supported test functions, which is Weil's
criterion \cite{Weil52,Bombieri2000}. (a)$\Rightarrow$(c): assume RH. For real $F$ in the form
domain let $F_\varepsilon=F*\eta_\varepsilon$ be the mollification in the proof of
Lemma~\ref{lem:roughEF}. By \cite[Thm.~8.4]{Paper45} and RH,
$W_{T+\varepsilon}(F_\varepsilon)=\sum_\gamma|\Fh(\gamma)|^2|\widehat{\eta_\varepsilon}(\gamma)|^2$ (zeros with
multiplicity), and the archimedean, polar and prime steps of that proof, which use only hypothesis (a),
give $W_{T+\varepsilon}(F_\varepsilon)\to W_T(F)$; since $\widehat{\eta_\varepsilon}\to1$ pointwise, Fatou's
lemma gives $\sum_\gamma|\Fh(\gamma)|^2\le W_T(F)$ (so hypothesis (b) holds and, by
Lemma~\ref{lem:roughEF}, with equality). The form $W_T$ is real, $W_T(u+iv)=W_T(u)+W_T(v)$ for real
$u,v$, so its infimum on the unit sphere, which is attained (\cite[Thm.~3]{Bombieri2000}, or the
compactness in the proof of Theorem~\ref{thm:O}), is attained at a real $F_\star$. If $W_T(F_\star)=0$
then $\Fh_\star(\gamma)=0$ at every ordinate $\gamma$. Under RH, $S(t)=O(\log t/\log\log t)$
\cite[Thm.~14.13]{Titchmarsh}, and a multiplicity is a jump of $S$, so a zero at height $t$ has
multiplicity $O(\log t/\log\log t)$; hence there are $\gg t\log\log t$ distinct ordinates in $(0,t]$,
while the nonzero entire function $\Fh_\star$ of exponential type $T$ has $O(t)$ zeros in $|z|\le t$ by
Jensen's formula. Hence $\lamW=W_T(F_\star)>0$, and
Theorem~\ref{thm:Zexh}(i) gives $w$ with $\inf f_w>0$.
\end{proof}

\begin{corollary}[certificates are statements about zeros]\label{cor:zeros}
If $w\in\Wcl^{\mathrm{reg}}$ and the elementary function \eqref{eq:fw} satisfies $f_w\ge c$ on $\R$,
then unconditionally $\sum_\rho\wh(\tau-\gamma_\rho)\ge c$ for every real $\tau$, the sum including
any zeros off the critical line through $\wh$ at complex arguments. Under RH, the
$\wh$-smoothing of the zero measure never falls below $c$. By Theorem~\ref{thm:finZ}, this
inequality cannot be verified on any finite list of zeros: every truncation of the sum to finitely
many ordinates takes negative values.
\end{corollary}
\begin{proof}
Immediate from Theorem~\ref{thm:Z}(ii), the zero sum being absolutely convergent.
\end{proof}


\subsection{Rigidity, smoothness, and finite zero sets}

Corollary~\ref{cor:zeros} raises the question whether $f_w\ge0$ could ever follow termwise from
$\wh\ge0$. For no flat-top weight can it.

\begin{lemma}[rigidity of flat-top weights]\label{lem:Rrig}
Let $w\in L^1(\R)$ be even and continuous with $w\equiv1$ on $[-2T,2T]$ for some $T>0$. Then
$\wh\ge0$ on $\R$ is impossible. Equivalently, if an even continuous $w\in L^1$ with $\wh\ge0$
satisfies $w\equiv1$ on an interval around the origin, then $w\equiv1$ on $\R$, which contradicts
integrability.
\end{lemma}
\begin{proof}
Suppose $w\in L^1$ is continuous with $\wh\ge0$. Then $\wh\in L^1$ and the inversion formula
$w(u)=\frac1{2\pi}\int\wh(t)e^{-itu}dt$ holds pointwise: this is the classical summability argument
(Gauss--Weierstrass or Fej\'er means converge to $w$ uniformly on compacta, and monotone convergence
applied to the nonnegative $\wh$ against the approximate identity gives $\int\wh=2\pi w(0)<\infty$;
see \cite[Ch.~VI]{Katznelson}). Hence $\wh\,dt/2\pi$ is a finite positive measure whose Fourier
transform is $w$, so by Bochner's theorem $w$ is positive definite, with $w(0)=1$. For a positive
definite $\phi$ the nonnegativity of the $3\times3$ Gram matrix at the points $0,-u,a$ gives the
standard inequality
\[
|\phi(u+a)-\phi(u)|^2\ \le\ 2\phi(0)\big(\phi(0)-\Rel\phi(a)\big)\qquad(u,a\in\R).
\]
Take any $a$ with $0<|a|\le2T$: then $w(a)=1=w(0)$, the right side vanishes, and $w(u+a)=w(u)$ for
every $u$. So $w$ is periodic with every period in $(0,2T]$; being continuous, it is constant, equal
to $w(0)=1$. A constant function is not in $L^1(\R)$.
\end{proof}

\begin{corollary}[positivity is never termwise]\label{cor:nontermwise}
For every even continuous $w\in\Wcl$ the kernel $\wh$ takes negative values. Consequently the
inequality $f_w=\sum_\rho\wh(\tau-\gamma_\rho)\ge0$ of Corollary~\ref{cor:RH} always expresses
cancellation across the zeros, never the termwise positivity of a kernel: there is no flat-top $w$
with $\wh\ge0$. \textup{(}The Li--Keiper kernels avoid this obstruction because they are not
flat-top; see Appendix~\ref{app:li}.\textup{)}
\end{corollary}
\begin{proof}
Lemma~\ref{lem:Rrig} applied to $w$, which lies in $L^1$ by the decay hypothesis of
Definition~\ref{def:W}, is even and equals $1$ on the window. A merely measurable $w$ with $\wh\ge0$
agrees almost everywhere with a continuous function, so nothing is lost.
\end{proof}

\begin{lemma}[finite zero sets and positive definiteness]\label{lem:fdict}
Let $T>0$, let $Z\subset\R\setminus\{0\}$ be a finite symmetric multiset, and let $w:\R\to\R$ be
even, continuous and integrable with $w\equiv1$ on $[-2T,2T]$. Put
$g_Z(u)=\sum_{\gamma\in Z}e^{i\gamma u}=\sum_{\gamma\in Z}\cos(\gamma u)$ and $\phi:=w\,g_Z$. Then
$\phi$ is continuous and integrable, $\widehat\phi(\tau)=\sum_{\gamma\in Z}\wh(\tau-\gamma)$,
$\phi(0)=|Z|$, $\phi\equiv g_Z$ on $[-2T,2T]$, and
\[
\sum_{\gamma\in Z}\wh(\cdot-\gamma)\ \ge\ 0\ \text{ on }\R
\qquad\Longleftrightarrow\qquad
\phi\ \text{ is positive definite.}
\]
\end{lemma}
\begin{proof}
Continuity, $|\phi|\le|Z|\,|w|\in L^1$ and the values on the window are immediate from $w\equiv1$
there; $\widehat\phi=\wh*\mu_Z$ with $\mu_Z=\sum_{\gamma\in Z}\delta_\gamma$ finite is the
convolution theorem. ($\Leftarrow$) If $\phi$ is positive definite, its spectral measure is
nonnegative, and $\phi\in L^1\cap C$ forces that measure to be $\widehat\phi\,d\tau/2\pi$, so
$\widehat\phi\ge0$. ($\Rightarrow$) If $\widehat\phi\ge0$, the summability argument of
Lemma~\ref{lem:Rrig} (Gauss--Weierstrass means, then monotone convergence against the nonnegative
$\widehat\phi$) gives $\int\widehat\phi=2\pi\phi(0)<\infty$, so $\widehat\phi\,d\tau/2\pi$ is a
finite positive measure with Fourier transform $\phi$, and $\phi$ is positive definite by Bochner.
\end{proof}

\begin{theorem}[finite zero sets are determinate]\label{thm:finDet}
Let $T>0$, let $Z\subset\R\setminus\{0\}$ be a finite symmetric multiset with $|Z|=\ell\ge2$
(multiplicities allowed), and put $g_Z(u)=\sum_{\gamma\in Z}\cos(\gamma u)$. Let $\nu$ be a
finite positive symmetric measure on $\R$ with $\nuh=g_Z$ on $(-2T,2T)$. Then
$\nu=\mu_Z=\sum_{\gamma\in Z}\delta_\gamma$.
\end{theorem}
No decay hypothesis on $\nu$ and no product structure is assumed: the extension problem for a
finite zero set is determinate.
\begin{proof}
Since $\nu$ is finite, $\phi:=\nuh$ is continuous, bounded, real, even and positive definite.

\emph{Step 1 (a lattice inside the window).} Let $w_1,\dots,w_r$ be the distinct elements of
$Z$, with multiplicities $m_1,\dots,m_r\ge1$, so $\sum_jm_j=\ell$ and $r\le\ell$. Put $z_j=e^{iw_ja}$ for
a step $a>0$, so that $g_Z(ka)=\sum_{j\le r}m_jz_j^{\,k}$ for all $k\in\Z$. Since $Z$ is symmetric
with multiplicity, the list $(z_j,m_j)$ is closed under conjugation. Each condition $z_j=z_{j'}$
($j\ne j'$) excludes a discrete set of $a$, so let $E\subset(0,2T/\ell)$ be that exceptional discrete
set and fix $a\in(0,2T/\ell)\setminus E$, so that $z_1,\dots,z_r$ are pairwise distinct. As $r\le\ell$ we
have $ra\le\ell a<2T$, hence $\phi(ka)=g_Z(ka)$ for $|k|\le r$.

\emph{Step 2 (a rank-$r$ Toeplitz block).} Put $c_k:=\phi(ka)\in\R$ ($\phi$ is real and even) and
$A:=[c_{i-j}]_{0\le i,j\le r}$. By Step 1, $A=V^*DV$ with $V_{j,i}=z_j^{\,i}$ and
$D=\operatorname{diag}(m_1,\dots,m_r)$, $m_j>0$; $V$ has rank $r$ because the $z_j$ are distinct, so
$\operatorname{rank}A=r$ and
\[
\ker A=\Big\{v\in\C^{r+1}:\ \textstyle\sum_{i=0}^{r}v_iz_j^{\,i}=0\ (j\le r)\Big\}=\C\cdot v,
\]
$v$ being the coefficient vector of $P(z)=\prod_{j\le r}(z-z_j)$. The root set is closed under
conjugation, so $v\in\R^{r+1}$, $v_r=1$ and $|v_0|=\prod_j|z_j|=1\ne0$.

\emph{Step 3 (the recursion forced by positive definiteness).} Fix $m>r$. The Gram matrix of $\phi$
at $\{0,a,\dots,ra,ma\}$ is positive semidefinite; in block form it is
$M=\begin{psmallmatrix}A&b\\b^*&c_0\end{psmallmatrix}$ with $b_i=\phi(ia-ma)=c_{m-i}$. For
$v\in\ker A$ and $x=(v,t)$ one has $x^*Mx=2\Rel(\bar t\,b^*v)+|t|^2c_0\ge0$ for every $t\in\C$;
choosing the phase of $t$ adversely and letting $|t|\to0$ gives $b^*v=0$, i.e.
\begin{equation}\label{eq:frec}
\sum_{i=0}^{r}v_i\,c_{m-i}=0 .
\end{equation}
The sequence $k\mapsto g_Z(ka)$ satisfies the same relation: $\sum_iv_iz^{m-i}=z^{m-r}z^rP(1/z)$
vanishes at $z=z_j$ because $1/z_j=\bar z_j$ is again a root of $P$; multiplying by $m_j$ and
summing over $j$ gives $\sum_iv_ig_Z((m-i)a)=0$.

\emph{Step 4 (induction).} By Step 1, $c_k=g_Z(ka)$ for $0\le k\le r$. If $c_k=g_Z(ka)$ for all
$0\le k\le m-1$ with $m>r$, then in \eqref{eq:frec} every term with $i\ge1$ equals the
corresponding term for $g_Z$; subtracting the two relations and using $v_0\ne0$ gives
$c_m=g_Z(ma)$. Hence $\phi(ma)=g_Z(ma)$ for every integer $m\ge0$ and every
$a\in(0,2T/\ell)\setminus E$.

\emph{Step 5 (a variable step).} Fix an integer $m\ge1$. By Step 4 the identity
$\phi(ma)=g_Z(ma)$ holds for every $a\in(0,2T/\ell)$ outside the discrete set $E$; both sides are
continuous in $a$, so it holds for every $a\in(0,2T/\ell)$. Since
$\bigcup_{m\ge1}m\cdot(0,2T/\ell)=(0,\infty)$ and both functions are even, $\phi=g_Z$ on $\R$. The
spectral measure of a continuous positive definite function is unique, so $\nu=\mu_Z$
(Proposition~\ref{prop:finuniv}(ii)).
\end{proof}

\begin{theorem}[no finite zero set carries a flat-top certificate]\label{thm:finZ}
Let $T>0$, let $Z\subset\R\setminus\{0\}$ be a finite symmetric multiset with $|Z|=\ell\ge2$
(multiplicities allowed), and let $w:\R\to\R$ be even, continuous and integrable with $w\equiv1$ on
$[-2T,2T]$ and $w(u)\to0$ as $|u|\to\infty$. Then
\[
f_{w,Z}(\tau):=\sum_{\gamma\in Z}\wh(\tau-\gamma)
\]
takes negative values on $\R$. In particular, for every $\Gamma>0$ and every $w\in\Wcl$ the
truncation $\sum_{|\gamma_\rho|\le\Gamma}\wh(\tau-\gamma_\rho)$ of a flat-top certificate is
not nonnegative.
\end{theorem}
\begin{proof}
Suppose $f_{w,Z}\ge0$. By Lemma~\ref{lem:fdict} the function $\phi=w\,g_Z$ is continuous and
positive definite with $\phi(0)=\ell$, and $\widehat\phi=f_{w,Z}\in L^1$ with $\int f_{w,Z}=2\pi\ell$, so
$\nu_w:=f_{w,Z}\,d\tau/2\pi$ is a finite positive symmetric measure with $\nuh_w=\phi$, which equals
$g_Z$ on $(-2T,2T)$ because $w\equiv1$ there. By Theorem~\ref{thm:finDet}, $\nu_w=\mu_Z$, i.e.
$w\,g_Z=g_Z$ on $\R$. But along a lattice $\{ma\}$ as in Step~1 of that proof, the $z_j$ are
distinct points of the unit circle, so the Ces\`aro averages of $(z_j\bar z_{j'})^m$ vanish for
$j\ne j'$ (a geometric sum) and
\[
\lim_{M\to\infty}\frac1M\sum_{m=1}^{M}|g_Z(ma)|^2=\sum_{j\le r}m_j^2\ \ge\ r\ \ge\ 1 ,
\]
whence $\limsup_m|g_Z(ma)|\ge1$, while $|w(ma)g_Z(ma)|\le\ell|w(ma)|\to0$, a contradiction. Hence
$f_{w,Z}\ge0$ is impossible; being continuous with $\int_\R f_{w,Z}=2\pi\ell>0$, it takes negative
values.
\end{proof}

The proof of Theorem~\ref{thm:finZ} does not locate the negative values: it contradicts positive
definiteness along a lattice in the variable $u$ and says nothing about $\tau$. Under additional
smoothness hypotheses, integration by parts locates them.

\begin{remark}[the location of the negativity]\label{rem:loc}
Let $w\in\Wcl$ be supported in $[-U,U]$ with $w^{(k)}$ continuous on $\R$, $w^{(j)}(u)\to0$ as
$|u|\to\infty$ for $0\le j\le k$, and $w^{(k+1)}$ of bounded variation with a jump of size $J\ne0$
at $u=U$ $(k\ge0)$; let $Z$ be a finite symmetric multiset of real ordinates with
$\sum_{\gamma\in Z}\cos(U\gamma)\ne0$. Integrating by parts $k+2$ times as in
Lemma~\ref{lem:smooth}(i), the boundary terms vanishing by the decay hypothesis, gives
\[
\wh(t)=\frac{2J}{|t|^{k+2}}\cos\big(Ut+\varphi_k\big)+o\big(|t|^{-k-2}\big),\qquad t\to+\infty,
\]
where $\varphi_k$ is the phase of $(-1)^ki^{-(k+2)}$ (the factor $(-1)^{k+2}=(-1)^k$ produced by
the $k+2$ integrations shifts the phase by $\pi$ for odd $k$). Summing over $Z$ and using
$\sum_\gamma\cos(U(\tau-\gamma)+\varphi_k)=\Rel[e^{i(U\tau+\varphi_k)}\sum_\gamma e^{-iU\gamma}]$
with $\sum_\gamma e^{-iU\gamma}=\sum_\gamma\cos(U\gamma)\ne0$, the bracket is a nonzero multiple of
$\cos(U\tau+\varphi)$, so $f_{w,Z}(\tau)<0$ for arbitrarily large $\tau$, at the rate $\tau^{-k-2}$
and the frequency $U$. Theorem~\ref{thm:finZ} gives the negativity under minimal hypotheses; this
remark gives its location and size under stronger ones. Neither statement contains the other.
\end{remark}

\begin{corollary}[finite sets of zeros do not suffice]\label{cor:infinite}
Let $T>0$. Every flat-top certificate of Corollary~\ref{cor:RH} is a statement about the complete,
infinite zero multiset of $\zeta$: no truncation of $\sum_\rho\wh(\tau-\gamma_\rho)$ to finitely many
zeros is nonnegative, for any window, any admissible weight and any truncation height. A finite
table of zeros can evaluate $f_w$ to any accuracy on a bounded range (Lemma~\ref{lem:smooth}(iii)),
but cannot establish $f_w\ge0$. Quantitatively, if $f_w\ge c>0$ for some $w\in\Wcl\cap C_c^k$
supported in $[-U,U]$ with $V=V(w^{(k+1)})$, and $\Gamma$ is large enough that
$C_UV\Gamma^{-k-1}\log\Gamma\le c/2$, then the truncation $f_{w,\Gamma}$ is $\ge c/2$ on
$|\tau|\le\Gamma/2$ while being negative somewhere; hence all of its negativity lies at
$|\tau|\gtrsim\Gamma/2$.
\end{corollary}
\begin{proof}
The first assertion is Theorem~\ref{thm:finZ} applied to the multiset
$Z=\{\gamma_\rho:|\gamma_\rho|\le\Gamma\}$, which is finite, symmetric and free of $0$. The
evaluation statement and the quantitative one follow from Lemma~\ref{lem:smooth}(iii), which bounds
$|f_w-f_{w,\Gamma}|$ by $C_UV\Gamma^{-k-1}\log\Gamma$ uniformly for $|\tau|\le\Gamma/2$.
\end{proof}

In the opposite direction, a finite zero set does admit a certificate, but only a trivial one.

\begin{proposition}[finite zero sets]\label{prop:finuniv}
Let $Z\subset\R\setminus\{0\}$ be a finite symmetric multiset and
$W_Z(h):=\sum_{\gamma\in Z}\widehat h(\gamma)$ the associated truncated Weil functional on even test
functions. Then, with $h=F\star\widetilde F$ and $\widetilde F(u)=\overline{F(-u)}$:
\emph{(i)} $W_Z(F\star\widetilde F)=\sum_{\gamma\in Z}|\Fh(\gamma)|^2\ge0$ whenever $Z\subset\R$;
\emph{(ii)} the measure $\mu_Z=\sum_{\gamma\in Z}\delta_\gamma$ is a dual feasible point of value
$0$ for the analogue of problem \textbf{(D)} of Definition~\ref{def:PD}, and it is the unique symmetric
positive tempered measure representing $W_Z$ (only the even part of a representing measure is
determined by even test functions);
\emph{(iii)} $W_Z\ge0$ on the cone of such $h$ $\iff$ $\widehat{\mu_Z}=g_Z$ is positive definite
$\iff$ $\mu_Z\ge0$, which holds precisely when $Z\subset\R$ (Bochner).
\end{proposition}
\begin{proof}
(i) $\widehat h=|\Fh|^2\ge0$ pointwise. (ii) $W_Z(h)=\int\widehat h\,d\mu_Z$ by definition, so
$\mu_Z$ is feasible of value $0$; if $\mu\ge0$ is symmetric tempered with
$\int\widehat h\,d\mu=\int\widehat h\,d\mu_Z$ for all even Schwartz $h$, then $\widehat\mu=
\widehat{\mu_Z}$ in $\mathcal S'$ (even Schwartz transforms are dense in the even Schwartz
functions), hence $\mu=\mu_Z$. (iii) $g_Z=\widehat{\mu_Z}$ with $\mu_Z\ge0$ finite is continuous
positive definite by Bochner; conversely the spectral measure of $g_Z$ is $\mu_Z$.
\end{proof}

For a finite zero set the dual certificate is therefore the spectrum itself: positivity is
immediate, the certificate is unique among symmetric measures, and it contains no arithmetic
information. By Theorem~\ref{thm:finZ}, the one-sided smoothing
$\sum_\rho\wh(\tau-\gamma_\rho)\ge0$ of Corollary~\ref{cor:RH} is never available for a finite zero
set, at any window and for any admissible weight. The mechanism is the following. A finite $Z$
makes the Toeplitz forms $[g_Z((i-j)a)]$ of finite rank; beyond that rank, positive semidefiniteness
of any extension propagates the linear recursion of $P$ indefinitely, so the extension is unique and
does not decay (Theorem~\ref{thm:finDet}), whereas a flat-top weight must decay. For the zero set of
$\zeta$ no such $P$ exists and the rank is never finite.

The tail estimate used in Corollary~\ref{cor:infinite} is the following elementary lemma. It also
extends Theorem~\ref{thm:Z}(ii), stated for $w\in C_c^\infty$, to the weights $w\in C^k_c$ that
arise in finite computations.

\begin{lemma}[smoothness and far zeros]\label{lem:smooth}
Let $w\in\Wcl$ be supported in $[-U,U]$, and let $k\ge0$ be an integer such that $w^{(k)}$ is
continuous on $\R$ (in particular across the window edge) and $w^{(k+1)}$ is of bounded variation,
with total variation $V=V(w^{(k+1)})$. Then:
\emph{(i)} for every $z\in\C$, $|\wh(z)|\le V\,e^{U|\Iml z|}|z|^{-k-2}$; in particular
$\wh(t)=O(|t|^{-k-2})$ on $\R$.
\emph{(ii)} The zero sum $\sum_\rho\wh(\tau-\gamma_\rho)$ converges absolutely and
Theorem~\ref{thm:Z}(ii) holds verbatim for such $w$; thus \eqref{eq:Z} is valid on
$\Wcl\cap C_c^k$ for every $k\ge0$, not only on $\Wcl^{\mathrm{reg}}$.
\emph{(iii)} Uniformly for $|\tau|\le\Gamma/2$,
\[
\sum_{|\gamma_\rho|>\Gamma}\big|\wh(\tau-\gamma_\rho)\big|\ \ll_{U}\ V\,\Gamma^{-k-1}\log\Gamma .
\]
\emph{(iv)} A zero $\rho=\frac12+i\gamma_\rho$ off the critical line, necessarily with
$|\Iml\gamma_\rho|\le\frac12$, contributes to the sum at most $e^{U/2}$ times the bound that
\emph{(i)} gives for a zero on the line with the same real ordinate.
\end{lemma}
\begin{proof}
(i) Integrating by parts $k+2$ times, the boundary terms vanishing because $w^{(j)}$ is continuous
and compactly supported for $j\le k$,
$\wh(z)=\int w(u)e^{izu}du=(iz)^{-(k+2)}\int e^{izu}\,dw^{(k+1)}(u)$, and $|e^{izu}|\le e^{U|\Iml z|}$
on $[-U,U]$. (ii) By (i) with $k\ge0$ we have $\wh(t)\ll V(1+|t|)^{-2}$ on $\R$ and
$|\wh(t+iy)|\ll Ve^{U/2}(1+|t|)^{-2}$ for $|y|\le\frac12$, while the number of zeros with
$|\gamma_\rho|\in[t,t+1]$ is $O(\log(2+t))$ by $N(t)\ll t\log t$; hence the sum converges absolutely,
and every step of the proof of Theorem~\ref{thm:Z}(ii) (the explicit formula for even $h$ of bounded
variation with compact support, and the Fubini justification, which used only
$\int|\wh(t-\tau)|\int_0^\infty\min(2,t^2u^2)K(u)du\,dt<\infty$) applies unchanged.
(iii) With $|\tau|\le\Gamma/2$ and $|\gamma_\rho|>\Gamma$ one has $|\tau-\gamma_\rho|\ge|\gamma_\rho|/2$,
so by (i) and partial summation against $N(t)=\frac t{2\pi}\log\frac t{2\pi e}+O(\log t)$,
$\sum_{|\gamma_\rho|>\Gamma}|\wh(\tau-\gamma_\rho)|\ll Ve^{U/2}\int_\Gamma^\infty t^{-k-2}\log t\,dt
\ll_UV\Gamma^{-k-1}\log\Gamma$. (iv) is the factor $e^{U|\Iml z|}$ of (i) at $|\Iml z|\le\frac12$.
\end{proof}

For $k=-1$, that is, a weight with a jump at the edge such as the box, the same computation gives
only $\wh(t)=O(|t|^{-1})$; the zero sum then converges conditionally, and (iii) becomes a statement
about oscillation rather than an absolute bound. For this reason the flat-top weights used in
practice are continuous across $u=\pm2T$.

\begin{remark}[the multiplier $\beta_0$, and the formal limit $w\uparrow1$]\label{rem:RHcert}
Theorem~\ref{thm:Z} explains why numerically computed certificates have $\beta_0$ close to $2$
(Appendix~\ref{app:diag}): the value $2$ is the polar term of the explicit formula, and by
Theorem~\ref{thm:Zexh}(ii) certificates of the primitive problem are flat-top certificates plus a
correction in $\beta$. Heuristically, as $w\uparrow1$ pointwise, the atomic part of $\nu_w$ tends to
$-\sum_{n\notin\NT}\frac{v_n}2(\delta_{\pm\log n})$, the primes outside the window, while its
absolutely continuous part tends to the infinite measure $2\cosh\frac u2\mathbf 1_{|u|\ge2T}du$, so
the limit leaves the class $\Wcl$ (compare Theorem~\ref{thm:O}). Since $\int\wh=2\pi w(0)=2\pi$,
formally $f_w\to2\pi\sum_\rho\delta_{\gamma_\rho}$, the zero measure, whose positivity is RH.
\end{remark}


\subsection{The decay tail of a flat-top certificate}

Proposition~\ref{prop:env} bounds the tail of any certificate at $2\pi e^{\Aconst+\|\nu\|_{\TV}}$,
and Lemma~\ref{lem:T} at a threshold exponential in the atomic mass. Both ignore the decay of
$\nuh_w$. For a $C^2$ flat-top weight this decay is available. The following theorem extends
Lemma~\ref{lem:T}: the single term $2\sum_j|a_j|$ is split into an atomic part, which does not decay,
and an absolutely continuous part, which decays like $\tau^{-2}$. It is also the smooth analogue of
Lemma~\ref{lem:F1}, with one essential difference recorded in Remark~\ref{rem:Dgrade}.

\begin{theorem}[decay tail bound for flat-top weights]\label{thm:Dtail}
Let $T>0$ and let $w\in\Wcl\cap C^2(\R)$ (so $w$ is even, $w\equiv1$ on $[-2T,2T]$, and $w'$, $w''$
match across $u=\pm2T$; in particular $w'(2T)=w''(2T)=0$), with $w''$ of bounded variation, and
assume either \emph{(a)} $\supp w\subset[-U,U]$, or \emph{(b)} $|w^{(j)}(u)|\le Ce^{-s|u|}$ for
$j\le2$ with $s>\frac12$. Put $W(u)=w(|u|)\cosh\frac u2$ and $R=(1-w)K$ on $[2T,\infty)$, and
\[
D_c(w)=4\int_0^\infty\Big|\big(w(u)\cosh\tfrac u2\big)''\Big|\,du,\qquad
D_K(w)=\int_{2T}^\infty\big|\big((1-w)K\big)''(u)\big|\,du,
\]
\[
D(w)=D_c(w)+D_K(w),\qquad E_{\mathrm{at}}(w)=\sum_{\log n>2T}v_n\,|w(\log n)| ,
\]
all four finite. Then, with $\Psi_w(\tau)=\int_{2T}^\infty(1-w)K\cos\tau u\,du$,
\begin{equation}\label{eq:Dsplit}
\nuh_w(\tau)+2\gh0(\tau)=2\widehat W(\tau)+\Psi_w(\tau)-\sum_{\log n>2T}v_n\,w(\log n)\cos(\tau\log n),
\end{equation}
and
\[
|2\widehat W(\tau)|\le D_c(w)\,\tau^{-2},\qquad |\Psi_w(\tau)|\le D_K(w)\,\tau^{-2}.
\]
Consequently, for $f=f_w+(\beta_0-2)\gh0+\beta_1\gh1$ and every $\tau>0$,
\[
f(\tau)\ \ge\ L_w(\tau):=\log\frac\tau{2\pi}-\frac1\tau-\frac1{2\tau^2}-\Aconst-E_{\mathrm{at}}(w)
-|\beta_0-2|k_0(\tau)-|\beta_1|k_1(\tau)-D(w)\,\tau^{-2},
\]
$L_w$ is nondecreasing on $(0,\infty)$, and positivity of $f$ on $\R$ reduces to the compact
interval $[0,\tau_{\mathrm{stop}}(w)]$, $\tau_{\mathrm{stop}}(w)=\min\{\tau:L_w(\tau)\ge0\}$.
\end{theorem}
\begin{proof}
\emph{The identity \eqref{eq:Dsplit}, and why the $\cosh$ part must be regrouped with $2\gh0$.}
By \eqref{eq:nuw} the absolutely continuous part of $\nu_w$ contributes
$\Psi_w(\tau)+4\int_{2T}^\infty w\cosh\frac u2\cos\tau u\,du$ and its atomic part contributes the
last sum of \eqref{eq:Dsplit}. Since $\gh0(\tau)=2\int_0^{2T}\cosh\frac u2\cos\tau u\,du$ and
$w\equiv1$ on the window, $4\int_{2T}^\infty w\cosh\frac u2\cos\tau u\,du+2\gh0(\tau)
=4\int_0^\infty w\cosh\frac u2\cos\tau u\,du=2\widehat W(\tau)$. The regrouping is necessary: the
density of $\nu_w$ at $u=2T^+$ equals $2\cosh T\ne0$, so $\nu_w$ has a jump at the window edge and
one integration by parts leaves the irreducible term $-4\cosh T\sin(2T\tau)/\tau$. Thus $\nuh_w$
alone is only $O(\tau^{-1})$, and no bound $|\nuh_w|\le D\tau^{-2}$ can hold. By
Lemma~\ref{lem:gbounds}, $2\gh0(\tau)=4\cosh T\sin(2T\tau)/\tau+O(\tau^{-2})$ cancels that term
exactly, with the same coefficient and the opposite sign; this is visible at once in
$2\widehat W$, since $W$ has no edge.

\emph{The two bounds.} $W$ is even and, by the hypotheses, globally $C^2$ on $\R$ with $W''$ of
bounded variation: on the window $W=\cosh\frac u2$, and $w,w',w''$ match across $\pm2T$. In case
(a) $W$ is compactly supported; in case (b) $|W^{(j)}|\le C'e^{-(s-\frac12)|u|}$ for $j\le2$ is
integrable, which is exactly where $s>\frac12$ is needed. Two integrations by parts, all boundary
terms vanishing, give $\widehat W(\tau)=-\tau^{-2}\int_\R e^{i\tau u}\,dW'(u)$, so
$|2\widehat W|\le2V(W')\tau^{-2}=4\int_0^\infty|W''|\,du\cdot\tau^{-2}=D_c(w)\tau^{-2}$ by evenness.
For $R$, extend by $0$ to $(-2T,2T)$ and evenly: $R(2T)=(1-w(2T))K(2T)=0$ and
$R'(2T)=-w'(2T)K(2T)+(1-w(2T))K'(2T)=0$ because $w\equiv1$ on the window and $w\in C^1$, so the
extension is $C^1$ on $\R$ with $R'$ of bounded variation, $V(R')=2\int_{2T}^\infty|R''|$ (in case
(a), $R=K$ beyond $U$, and $K,K',K''$ are integrable at infinity; in case (b),
$|R''|\le(|w''|+|w'|+|1-w|)\sup_{u\ge2T}(K+|K'|+|K''|)e^{-u/2}$). As $\Psi_w=\frac12\widehat R$,
two integrations by parts give $|\Psi_w|\le\frac12V(R')\tau^{-2}=D_K(w)\tau^{-2}$.

\emph{The atoms carry no decay}, and are bounded only by their mass:
$|\sum_{\log n>2T}v_nw(\log n)\cos(\tau\log n)|\le E_{\mathrm{at}}(w)$, the sum being finite in case
(a) and convergent in case (b) because $\sum_n\Lambda(n)n^{-\frac12-s}<\infty$ for $s>\frac12$,
the same threshold as above.

\emph{The lower bound.} $\sT(\tau)\ge\log\frac\tau{2\pi}-\frac1\tau-\frac1{2\tau^2}-\Aconst$ by
Lemma~\ref{lem:mbounds}(ii),(iii), the two multiplier terms are bounded by Lemma~\ref{lem:gbounds},
and the three pieces of \eqref{eq:Dsplit} by the above. Each term of $L_w$ is nondecreasing, hence so
is $L_w$, and $L_w(\tau_{\mathrm{stop}})\ge0$ gives $f\ge0$ on $[\tau_{\mathrm{stop}},\infty)$.
\end{proof}

\begin{remark}\label{rem:Dgrade}
The exponent $\tau^{-2}$ is the one Lemma~\ref{lem:smooth}(i) gives at $k=0$ after the two extra
derivatives, and it cannot be improved without $w\in C^3$: by Remark~\ref{rem:loc} the support
endpoint carries a jump of $w'''$ contributing a genuine $\tau^{-2}$ to $\widehat W$. The bound on
the atoms is sharp in the worst case by Lemma~\ref{lem:mbounds}(iii). This is where the method
differs from Lemma~\ref{lem:F1}: a lattice $\nuh$ is $P$-periodic, so the comb can be evaluated
exactly on $N$ periods and replaced by its infimum beyond, whereas a flat-top $\nuh_w$ is not
periodic and its comb can only be bounded by its mass. Theorem~\ref{thm:Dtail} handles the absolutely
continuous part by decay. We do not know how to combine the two techniques; the remaining
exponential dependence on the atomic mass comes from this gap.
\end{remark}

\begin{corollary}[comb-vanishing weights]\label{cor:comb}
Assume in addition $\supp w\subset[-U,U]$ and $w(\log n)=0$ for every prime power $n$ with
$2T<\log n\le U$ (a finite set of linear conditions). Then $E_{\mathrm{at}}(w)=0$, and for
$D(w)\le D_0$,
\[
\tau_{\mathrm{stop}}(w)\ \le\ \max\Big\{2\pi e^{\Aconst+1+|\beta_0-2|k_0(1)+|\beta_1|k_1(1)},\ \sqrt{D_0}+2\Big\},
\]
since for $\tau$ beyond the first entry one has $\log\frac\tau{2\pi}-\Aconst-|\beta_0-2|k_0-|\beta_1|k_1\ge1$
and then $L_w(\tau)\ge0$ as soon as $\tau^2\ge D_0+O(\tau)$. In particular
$\tau_{\mathrm{stop}}=O(D(w)^{1/2})$ as $D(w)\to\infty$: the verification interval is polynomial,
not exponential, in the size of the weight, compared with $2\pi e^{\Aconst+E_{\mathrm{at}}(w)}$ in
general and $2\pi e^{\Aconst+\|\nu\|_{\TV}}$ in Proposition~\ref{prop:env}.
\end{corollary}

\begin{corollary}[bounded profiles]\label{cor:bprof}
If $0\le w\le1$ with $\supp w\subset[-U,U]$ then
$D_c(w)\le4\cosh\frac U2\big(\|w''\|_1+\|w'\|_1+\frac14\|w\|_1\big)$ and
$D_K(w)\le\kappa_0(T)\|w''\|_1+2\kappa_1(T)\|w'\|_1+2|K'(2T)|$ with
$\kappa_j(T)=\sup_{u\ge2T}|K^{(j)}(u)|$, so $D(w)=O(e^{U/2})$. Boundedness, however, controls only
$D$: it bounds $E_{\mathrm{at}}$ by nothing better than $\sum_{2T<\log n\le U}v_n$, which grows like
$e^{U/2}$. Only the comb conditions of Corollary~\ref{cor:comb} control the atoms.
\end{corollary}


\section{The Kre\u{\i}n extension problem}\label{subsec:KEP}

The dual variable of Section~\ref{sec:Z} is a positive measure prescribed only on the window. This
is an extension problem for a positive definite datum given on an interval, in the sense of
Kre\u{\i}n \cite{Krein1940}, \cite[Ch.~2]{Sasvari1994}. In this section we make the identification
and restate Corollary~\ref{cor:RH} in its terms. The classical theory concerns continuous positive
definite functions, whereas the datum below is a distribution with atoms; we therefore import
nothing from it, and every statement is proved from Theorems~\ref{thm:Z} and \ref{thm:Zexh} and
elementary Fourier analysis.

\begin{definition}[the window datum, and $\mathrm{KEP}(T)$]\label{def:KEP}
Let $g$ be the even distribution defined on even $h\in C_c^\infty(\R)$ by
\begin{equation}\label{eq:gdatum}
\langle g,h\rangle=C_0\,h(0)+\int_0^\infty\big(h(0)-h(u)\big)K(u)\,du-\sum_{n\ge2}v_n\,h(\log n)
+2\int_\R h(u)\cosh\tfrac u2\,du,
\end{equation}
that is, the right-hand side of the explicit formula in the normalization of the proof of
Theorem~\ref{thm:Z}(ii): the L\'evy pair $C_0\delta_0-\frac12\mathrm{Pf}\,K(|u|)$, whose Fourier
transform is the archimedean multiplier $m$ of \eqref{eq:m}; the polar term $2\cosh\frac u2$; and
the prime atoms $-\frac{v_n}2\big(\delta_{\log n}+\delta_{-\log n}\big)$, $n\ge2$. For $T>0$ put
\[
\mathfrak S(T)=\Big\{\nu\ \text{a positive symmetric tempered measure on }\R:\
\langle\nuh-g,h\rangle=0\ \text{for every even }h\in C_c^\infty(-2T,2T)\Big\}.
\]
The \emph{Kre\u{\i}n extension problem} $\mathrm{KEP}(T)$ is \emph{solvable} if
$\mathfrak S(T)\ne\emptyset$, \emph{determinate} if $\#\mathfrak S(T)=1$, and \emph{indeterminate}
if $\#\mathfrak S(T)\ge2$.
\end{definition}
Symmetry is imposed for the reason of Proposition~\ref{prop:finuniv}(ii): even test functions
determine only the even part of a representing measure, so without it every problem would be
indeterminate for a trivial reason. We define ``the datum is positive definite on the window'' to
mean solvability; by the Schwartz--Bochner theorem this is equivalent to saying that the restriction
of $g$ to $(-2T,2T)$ extends to a positive definite tempered distribution on $\R$.

\begin{lemma}[one negative square in the polar term]\label{lem:polar}
At any $n\ge2$ distinct points $t_1,\dots,t_n$ the Gram matrix of the kernel
$(s,t)\mapsto2\cosh\frac{s-t}2$ is $ab^{\mathsf T}+ba^{\mathsf T}$ with $a_i=e^{t_i/2}$ and
$b_i=e^{-t_i/2}$; it has rank $2$ and signature $(1,1)$. Hence $2\cosh\frac u2$ is positive definite
on no interval, and the solvability of $\mathrm{KEP}(T)$ is automatic at no window, however small.
\end{lemma}
\begin{proof}
$2\cosh\frac{t_i-t_j}2=e^{t_i/2}e^{-t_j/2}+e^{-t_i/2}e^{t_j/2}$, which is the stated rank-two matrix.
It vanishes on $\{a,b\}^\perp$ and on $\mathrm{span}\{a,b\}$ has eigenvalues $a\cdot b\pm|a||b|$;
$a\cdot b=n$ while $|a||b|>n$ strictly by Cauchy--Schwarz, the vectors $a,b$ not being parallel when
the $t_i$ are distinct and $n\ge2$. So one eigenvalue is positive and one negative.
\end{proof}
The datum \eqref{eq:gdatum} is therefore the sum of a positive definite part ($m\ge0$), a part with
exactly one negative square (the pole), and a negative part (the prime atoms). This is the reason a
polar multiplier $\beta_0$ must be carried at all (Proposition~\ref{prop:P}).

\begin{theorem}[a certificate is an absolutely continuous solution]\label{thm:KEP1}
Let $T\ge\frac12\log2$ and let $w\in\Wcl^{\mathrm{reg}}$ satisfy $f_w\ge c$ on $\R$.
\emph{(i)} If $c\ge0$ then $\nu_w:=f_w(\tau)\,d\tau/2\pi$ belongs to $\mathfrak S(T)$; in particular
$\mathrm{KEP}(T)$ is solvable, by an absolutely continuous solution.
\emph{(ii)} If $c>0$ then $\mathfrak S(T)$ contains a continuum of pairwise distinct absolutely
continuous measures; in particular $\mathrm{KEP}(T)$ is indeterminate.
\emph{(iii)} If moreover RH holds then the zero measure $\mu_\zeta=\sum_\rho\delta_{\gamma_\rho}$
lies in $\mathfrak S(T)$ as well, and is purely atomic, hence distinct from every $\nu_w$.
\end{theorem}
\begin{proof}
(i) For $w\in\Wcl^{\mathrm{reg}}$ the product $w\cdot g$ is a tempered distribution and
Theorem~\ref{thm:Z} says $\widehat{w\cdot g}=f_w$, a continuous function with
$f_w(\tau)=\log|\tau|+O(1)$ (Lemma~\ref{lem:mbounds}, Theorem~\ref{thm:Dtail}), hence tempered;
Fourier inversion gives $\nuh_w=w\cdot g$, which agrees with $g$ on $(-2T,2T)$ because $w\equiv1$
there. Positivity of $\nu_w$ is $f_w\ge0$, and symmetry is the evenness of $w$ and $g$.
(ii) Fix an even $\chi\in C_c^\infty(\R)$, $\chi\ge0$, $\chi\not\equiv0$, with
$\supp\chi\subset\{|u|\ge\max(2T,\log2)\}$. For $\varepsilon\in\R$ put $w_\varepsilon=w+\varepsilon\chi
\in\Wcl^{\mathrm{reg}}$; then $f_{w_\varepsilon}=f_w+\varepsilon r_\chi$ with
\[
r_\chi(\tau)=2\int\chi(u)\cosh\tfrac u2\cos\tau u\,du-\sum_{n\ge2}v_n\chi(\log n)\cos(\tau\log n)
-\tfrac12\int_\R\chi(u)K(u)\cos\tau u\,du ,
\]
a bounded continuous function, say $|r_\chi|\le R$ ($\chi$ is compactly supported away from the
singularity of $K$, and $\sum_nv_n|\chi(\log n)|<\infty$). For $|\varepsilon|<c/R$ one has
$f_{w_\varepsilon}>0$, so $\nu_{w_\varepsilon}\in\mathfrak S(T)$ by (i). They are pairwise distinct:
$\nu_{w_\varepsilon}-\nu_{w_{\varepsilon'}}=\frac{\varepsilon-\varepsilon'}{2\pi}r_\chi\,d\tau$, and
$r_\chi\equiv0$ would force $\chi\cdot g=0$, whereas on $\supp\chi\subset\{|u|\ge\log2\}$ the
absolutely continuous part of $g$ is $\chi(u)d_*(u)$ with $d_*(u)=2\cosh\frac u2-\frac12K(u)>0$
(positivity and monotonicity of $d_*$ on $[0.3,\infty)$ are established in the proof of
Theorem~\ref{thm:Zexh}), which is nonnegative and not identically zero and which the singular prime
atoms cannot cancel.
(iii) Under RH, $\mu_\zeta$ is a positive symmetric tempered measure ($N(t)\ll t\log t$) and
$\widehat{\mu_\zeta}=g$ on all of $\R$ by Theorem~\ref{thm:Z}(ii).
\end{proof}

\begin{theorem}[solvability implies Weil positivity at that window]\label{thm:KEP2}
If $\mathfrak S(T)\ne\emptyset$ then $\lamW\ge0$.
\end{theorem}
\begin{proof}
Let $\nu\in\mathfrak S(T)$ and $F\in C_c^\infty(\IT)$, and put $h=F\star\widetilde F$ with
$\widetilde F(u)=\overline{F(-u)}$. Then $h$ is even, $\supp h\subset(-2T,2T)$ and
$\widehat h=|\Fh|^2\ge0$ is rapidly decreasing, so the pairing against the positive tempered measure
$\nu$ is legitimate. Since $h$ vanishes off the window, only the contacts $n\in\NT$ appear in
\eqref{eq:gdatum}, and comparison with \eqref{eq:Aform} and \eqref{eq:W} gives
$\langle g,h\rangle=\mathfrak a_T(F)+2\Rel\big(A(F)\overline{B(F)}\big)=W_T(F)$. By
Definition~\ref{def:KEP}, $W_T(F)=\int\widehat h\,d\nu=\int|\Fh|^2d\nu\ge0$. Finally
$C_c^\infty(\IT)$ is a form core: for $F$ in the form domain, the dilates $F_\lambda(t)=F(t/\lambda)$,
$\lambda\uparrow1$, are supported in $\lambda\IT$ and converge to $F$ in $L^2$ with
$\Fh_\lambda(\tau)=\lambda\Fh(\lambda\tau)$, so that
$\int\log(2+|\tau|)|\Fh_\lambda-\Fh|^2\to0$ (the weight is doubling), and mollifying $F_\lambda$ at a
scale below $(1-\lambda)T$ keeps the support inside $\IT$ and converges in the same norm. Hence
$W_T\ge0$ on the whole form domain and $\lamW\ge0$.
\end{proof}

\begin{theorem}[the dichotomy, and RH as indeterminacy]\label{thm:KEP3}
\emph{(a)} For every $T\ge\frac12\log2$,
\[
\lamW>0\ \Longrightarrow\ \mathrm{KEP}(T)\ \text{is indeterminate}\ \Longrightarrow\ \lamW\ge0 ,
\]
so certificates and indeterminacy coincide except on the set $\{T:\lamW=0\}$.
\emph{(b)} The following are equivalent: \emph{(1)} RH; \emph{(2)} $\mathrm{KEP}(T)$ is solvable for
every $T>0$; \emph{(3)} $\mathrm{KEP}(T)$ is indeterminate for every $T>0$; \emph{(4)} for every
$T>0$, $\mathfrak S(T)$ contains a continuum of absolutely continuous measures.
\end{theorem}
\begin{proof}
(a) If $\lamW>0$, Theorem~\ref{thm:Zexh}(i) gives $w\in\Wcl^{\mathrm{reg}}$ with $\inf f_w>0$, and
Theorem~\ref{thm:KEP1}(ii) applies; the converse implication is Theorem~\ref{thm:KEP2}.
(b) (1)$\Rightarrow$(4): under RH the proof of Corollary~\ref{cor:RH} gives $\lamW>0$, not merely
$\lamW\ge0$, for every $T\ge\frac12\log2$, by the Jensen/exponential-type argument, whence (4) at
those $T$ by (a). For $0<T<\frac12\log2$ use $\mathcal W_{T'}\subseteq\Wcl$ for $T'>T$
(\S\ref{sec:path}): a certificate at $T'=\frac12\log2$ is one at $T$, and the perturbation of
Theorem~\ref{thm:KEP1}(ii) applies, its $\chi$ being supported in $\{|u|\ge\log2\}$, which is where
$d_*>0$ is available. (4)$\Rightarrow$(3)$\Rightarrow$(2) are trivial. (2)$\Rightarrow$(1): by
Theorem~\ref{thm:KEP2}, $W_T\ge0$ for every $T$, which is Weil's criterion
\cite{Weil52,Bombieri2000}.
\end{proof}

\begin{remark}[the exceptional set]\label{rem:nonzero}
By \cite[Thms.~1, 3, 5]{Bombieri2000} and \cite[Thm.~11.3]{Paper45}, RH holds if and only if
$\lamW\ne0$ for every $T>0$. If RH fails, the exceptional set $\{T:\lamW=0\}$ of
Theorem~\ref{thm:KEP3}(a) is a nonempty compact interval, possibly a point. Indeed, under RH,
$\lamW>0$ by the proof of Corollary~\ref{cor:RH} and Lemma~\ref{lem:monoT}. If RH fails,
$\lambda^{\mathrm W}_{T_0}<0$ for some $T_0$ by Corollary~\ref{cor:margin2}(iii), while $\lamW>0$ for
$T\le0.8$ by \cite[Thm.~1.2]{Zhu2026}; Lemmas~\ref{lem:monoT} and \ref{lem:contT} and the
intermediate value theorem then give the claim. The statement concerns the complete form $\lamW$.
\end{remark}

\begin{corollary}[unconditional indeterminacy for $T\le0.8$]\label{cor:KEP08}
$\mathrm{KEP}(T)$ is indeterminate for every $T\le0.8$; in particular at $T=\frac12\log3$ and at
$T=\log2$.
\end{corollary}
\begin{proof}
$\lamW$ is nonincreasing in $T$, since $W_T$ is the restriction of $W_{T'}$ to
$L^2(\IT)\subset L^2(I_{T'})$ for $T'>T$. By \cite[Thm.~1.2]{Zhu2026},
$W_{0.8}\ge8.9\cdot10^{-18}\|\cdot\|^2$ for the complete form, so
$\lamW\ge\lambda^{\mathrm W}_{0.8}>0$ for $T\le0.8$, and Theorem~\ref{thm:KEP3}(a) applies for
$T\ge\frac12\log2$, which is the range on which that theorem is stated. For $0<T<\frac12\log2$ run
instead the argument of Theorem~\ref{thm:KEP3}(b), implication (1)$\Rightarrow$(4): a certificate at
$T'=\frac12\log2$ is one at $T$, and the perturbation of Theorem~\ref{thm:KEP1}(ii) applies, its
$\chi$ being supported in $\{|u|\ge\log2\}$ and therefore outside the window.
\end{proof}

\begin{remark}
Theorem~\ref{thm:C} alone does not give Corollary~\ref{cor:KEP08} at $T=\frac12\log3$. Its
multipliers are $\beta=(2.0025617,-0.0096544)\ne(2,0)$, and $g_0,g_1$ of \eqref{eq:g} are supported
inside the window. The certificate therefore solves the extension problem for the perturbed
datum $g+(\beta_0-2)g_0+\beta_1g_1$, and the defect $|\beta_0-2|\,\gh0(0)=5.9\cdot10^{-3}$ exceeds
$c=8.86\cdot10^{-5}$ by two orders of magnitude, so it cannot be absorbed into the margin.
Theorem~\ref{thm:C} certifies the primitive form; the unconditional statement above rests on
\cite{Zhu2026}, which certifies the complete one.
\end{remark}

By Theorem~\ref{thm:KEP3}(a), off the set $\{T:\lamW=0\}$ the solutions of the extension problem
and the certificates determine each other. A Nevanlinna-type parametrization of $\mathfrak S(T)$ in
the indeterminate case would therefore describe the set of certificates, but it presupposes one
known solution. In particular, the extension problem gives an equivalent form of RH, not an
independent proof of Theorem~\ref{thm:D1} for the complete datum.


\section{Lattice certificates and their ceilings}\label{sec:class}

This section bounds from above what a given class of dual objects can certify. The bounds are
used in \S\ref{sec:path}.

\subsection{A class-dependent bound}
\begin{problem}[class-dependent value bound]\label{prob:A}
For the lattice class of step $d=2T/M$ and range $U$, with the tail rows of Lemma~\ref{lem:F1},
find an explicit $D(T,d,U)\to0$ as $d\to0$, $U\to\infty$, such that the best certificate value
$c(T,d,U)$ satisfies $c\ge\lamP-D(T,d,U)$. Equivalently: a quantitative one-sided approximation
theorem for the function $-\sT-\beta\cdot\widehat g$ on $[0,\tau_{\mathrm{env}}]$ by $P$-periodic
trigonometric polynomials whose low modes $|j|<M$ vanish, in the norm $L^1(|\Fh_\star|^2d\tau/2\pi)$
of the primal minimizer.
\end{problem}
A solution would decide each window by a computation of predictable size and would turn
Theorem~\ref{thm:C} into a construction. By the identity $\lamP-c=\int(f-c)\,d\mu_\star$
(Proposition~\ref{prop:weak}) the deficit is the $\mu_\star$-average of the slack, so the
approximation is needed only where $|\Fh_\star|^2$ lives. Theorem~\ref{thm:Vp} shows that $D$ must
depend on the class, not on a radius alone. If Zhu's conjectural law $-\log\lamP\asymp e^{2T}$ holds,
the effort in Problem~\ref{prob:A} grows at least like $e^{2T}$: finite for each $T$, never uniform.

The converse of Problem~\ref{prob:A}, an upper bound on the value of a class, is a theorem.

\begin{proposition}[test-kernel ceiling for lattice classes]\label{prop:TK}
Let $d=2T/M$, $P=2\pi/d$, and let $\nu$ be any finite even measure supported on the lattice
$\{\pm jd:j\ge M\}$ (any range, any number of atoms). Put $S_\beta=\sT+\beta_0\gh0+\beta_1\gh1$ and,
for $\tau_0\in\R$, let $C_M(\tau_0,\beta)$ be the $M\times M$ Hermitian Toeplitz matrix with entries
$c_{j-k}$, $c_\ell=P^{-1}\int_{\tau_0}^{\tau_0+P}S_\beta(\tau)e^{i\ell d\tau}\,d\tau$. If
$\sT+\nuh+\beta\cdot\widehat g\ge c$ on $\R$ then
\[
c\ \le\ \lambda_{\min}\big(C_M(\tau_0,\beta)\big)\ \text{ for every }\tau_0,\qquad\text{hence}\qquad
c\ \le\ \operatorname{Ceil}(M,T):=\sup_{\beta\in\R^2}\inf_{\tau_0}\lambda_{\min}\big(C_M(\tau_0,\beta)\big),
\]
and $\beta\mapsto\inf_{\tau_0}\lambda_{\min}(C_M(\tau_0,\beta))$ is concave. The ceiling depends on
the class only through $M$.
\end{proposition}
\begin{proof}
$\nuh(\tau)=\sum_{j\ge M}2a_j\cos(jd\tau)$ is $P$-periodic with Fourier modes $|j|\ge M$ only; in
particular it has mean zero over any period. Let $q(z)=\sum_{k=0}^{M-1}q_kz^k$ and
$\varphi(\tau)=|q(e^{id\tau})|^2\ge0$, a trigonometric polynomial in $d\tau$ of degree $\le M-1$. Then
$\int_{\tau_0}^{\tau_0+P}\nuh\varphi\,d\tau=0$, so integrating $f\ge c$ against $\varphi$ gives
$c\int\varphi\le\int S_\beta\varphi$. Now $\int_{\tau_0}^{\tau_0+P}\varphi=P\sum|q_k|^2$ and
$\int_{\tau_0}^{\tau_0+P}S_\beta\varphi=P\sum_{j,k}\bar q_kq_jc_{j-k}$, so
$c\le\langle C_Mq,q\rangle/\|q\|^2$ for every $q\ne0$, i.e.\ $c\le\lambda_{\min}(C_M)$; by the
Fejér--Riesz theorem every nonnegative trigonometric polynomial of degree $\le M-1$ is of the form
$|q|^2$, so nothing is lost. For fixed $q$ and $\tau_0$ the Rayleigh quotient is affine in $\beta$
($S_\beta$ is), so $\inf_{q,\tau_0}$ is concave in $\beta$.
\end{proof}

The ceiling is the smallest eigenvalue of a small matrix and needs no optimizer; like
Theorem~\ref{thm:Vp}, it depends on the class and not on a radius. Numerically, at
$T=\frac12\log3$ it is negative for $M=8$ and positive for $M=12,24$, and the certificate of
Theorem~\ref{thm:C} realizes a fraction of about $15\%$ of its own class ceiling
(Appendix~\ref{app:diag}). The negative sign at $M\le8$ survives an interval enclosure; this is
Theorem~\ref{thm:N} below. Since the lattices $2T/2^kM_0$ are nested with dense union,
Theorem~\ref{thm:D1} gives $\sup_M\operatorname{Ceil}(M,T)\ge\lamP$; that the ceilings in fact
converge to $\lamP$, so that the test-kernel family is not lossy in the limit, is
Theorem~\ref{thm:R} and Corollary~\ref{cor:R3} below. We do not know whether the family is lossy at
fixed $M$, nor at what rate the ceilings converge; this is the tractable half of
Problem~\ref{prob:A}.


\subsection{Lattices of small order at $T=\tfrac12\log3$}\label{subsec:N}
The ceiling of Proposition~\ref{prop:TK} is a supremum over $\beta\in\R^2$, and turning a numerical
ceiling into a theorem requires control of $\beta$. The following a priori bound provides it: the
two primitivity multipliers of a certificate cannot be large, with a constant determined by the
geometry of the moment image of Lemma~\ref{lem:W}.

\begin{theorem}[a priori bound on the multipliers]\label{thm:Bbeta}
Fix $T$ and nonnegative weights $v$. Suppose there are $r_0>0$, $V^\star<\infty$ and finitely many
unit vectors $F_1,\dots,F_N\in H^1_0(\IT)$ such that, writing $w_i=(Q_0(F_i),Q_1(F_i))$ in the
notation of Lemma~\ref{lem:W} and $V_i=\mathfrak a_T(F_i)$,
\[
\{w\in\R^2:|w|\le r_0\}\subseteq\operatorname{conv}\{w_1,\dots,w_N\},\qquad V_i\le V^\star\ \ (1\le i\le N).
\]
Then every certificate $(\nu,\beta)$ of value $c$ satisfies $\|\beta\|_2\le(V^\star-c)/r_0$.
\end{theorem}
\begin{proof}
For a unit $F\in H^1_0(\IT)$ the measure $\mu_F=|\Fh|^2d\tau/2\pi$ is an even probability measure with
$\check\mu_F=\cF$ supported in $[-2T,2T]$ and $\cF(\pm2T)=0$, so $\int\nuh\,d\mu_F=\int\cF\,d\nu=0$ for
every measure $\nu$ supported in $\{|u|\ge2T\}$, exactly as in the proof of
Proposition~\ref{prop:weak}; moreover $\int\gh j\,d\mu_F=\int g_j\cF=Q_j(F)$ and
$\int\sT d\mu_F=\mathfrak a_T(F)$. Integrating $f\ge c$ against $\mu_F$ gives
$\mathfrak a_T(F)+\beta\cdot w(F)\ge c$. The relaxed feasible set (even probability measures whose
cosine transform vanishes off the window, with the two moment constraints dropped) is
convex, and both $\mu\mapsto\int\sT d\mu$ and $\mu\mapsto\mathfrak m(\mu)$ are affine; hence for
$w=\sum\lambda_iw_i$ in the convex hull the mixture $\mu=\sum\lambda_i\mu_{F_i}$ is admissible, has
moment $w$ and value at most $V^\star$, so $\beta\cdot w\ge c-V^\star$ for every $|w|\le r_0$.
Applying this to $w$ and $-w$ gives $|\beta\cdot w|\le V^\star-c$ on the disc of radius $r_0$, and
taking the supremum there, $r_0\|\beta\|_2\le V^\star-c$.
\end{proof}

\begin{lemma}[value of a witness]\label{lem:val}
For a unit $F\in H^1_0(\IT)$ with $L:=\|F'\|_2$ and every $u_0>0$,
\[
\mathfrak a_T(F)\ \le\ C_0+\Aconst+\frac{L^2}2\int_0^{u_0}u^2K(u)\,du+2\int_{u_0}^\infty K(u)\,du,
\]
the right side being minimized at $u_0=2/L$; the second integral is
$\log\coth\frac{u_0}2+\pi-2\arctan e^{u_0}$ by Theorem~\ref{thm:S}(iii).
\end{lemma}
\begin{proof}
By Theorem~\ref{thm:A}, $\mathfrak a_T(F)=C_0+\Lev(\cF)-\sum v_n\cF(\log n)\le C_0+\Aconst+
\int_0^\infty(1-\Rel\cF(u))K(u)du$, since $|\cF|\le\|F\|^2=1$. Now
$1-\Rel\cF(u)=\frac1{2\pi}\int(1-\cos\tau u)|\Fh|^2d\tau\le\min\{2,\frac{u^2}2\|F'\|^2\}$, the first
bound from $|\cF|\le1$ and the second from $1-\cos x\le x^2/2$ and Plancherel; split the integral at
$u_0$. The derivative in $u_0$ is $K(u_0)(\frac{L^2}2u_0^2-2)$, which vanishes and changes sign at
$u_0=2/L$.
\end{proof}

At $T=\frac12\log3$ we take as witnesses the unit vectors of the four-dimensional span of
$\psi_k(t)=(T^2-t^2)t^k$, $k=0,1,2,3$, inside $H^1_0(\IT)$: all the quantities entering
$Q_0=AB$, $Q_1=A'B-AB'$ and Lemma~\ref{lem:val} are elementary integrals of a polynomial against
$e^{\pm t/2}$. Taking, for each of $48$ directions, the top generalized eigenvector of the
corresponding combination of the two moment forms relative to the Gram matrix, computing the convex
hull of the $48$ moment vectors and the distance from the origin to its boundary gives
\[
r_0=8.4515\cdot10^{-5},\qquad V^\star=3.817495,\qquad
\|\beta\|_\infty\le\frac{V^\star}{r_0}=45169.4
\]
\[
\text{for every certificate of value }c\ge0 .
\]
The size of this constant is structural rather than a defect of the witnesses: on the window
$g_0(u)-\frac14g_1(u)=\cosh\frac u2-\frac u4\sinh\frac u2=1-\frac{u^4}{384}+O(u^6)$, so
$Q_0(F)-\frac14Q_1(F)=|\int F|^2+O(T^4)$ and the moment image is a thin sliver lying essentially on
one side of the line $Q_0=\frac14Q_1$; hence $r_0=O(T^4)$ for any witness family, and the bound is
$O(T^{-4})$. It is far too large to restrict a search, and good enough for the qualitative
statements below, because the coefficients it multiplies are of size $10^{-9}$.

\begin{theorem}[class impossibility at $T=\frac12\log3$]\label{thm:N}
Let $T=\frac12\log3$, $\NT=\{2\}$, $M\in\{1,\dots,8\}$,
$d=2T/M$, and let $\nu$ be any finite even signed measure supported on the lattice
$\{\pm jd:j\ge M\}$ (any range, any number of atoms), and $\beta\in\R^2$ arbitrary. Then
\[
\inf_{\tau\in\R}\big[\sT(\tau)+\nuh(\tau)+\beta_0\gh0(\tau)+\beta_1\gh1(\tau)\big]\ <\ 0 .
\]
Consequently no dual certificate in the sense of Definition~\ref{def:PD} supported on a lattice of
step $2T/M$ with $M\le8$ has nonnegative value, and none proves $\lamP>0$ at this window, for any
$\beta$, any range, any atom count and whichever of the tail lemmas of \S\ref{sec:tails} is used to
verify the inequality. The proof includes a computer verification in interval arithmetic
(Appendix~\ref{app:encl}).
\end{theorem}
\begin{proof}
Suppose $f:=\sT+\nuh+\beta\cdot\widehat g\ge c$ on $\R$ with $c\ge0$.

\emph{Step 1: $\beta$ is bounded.} By Theorem~\ref{thm:Bbeta} and the constants above,
$\|\beta\|_\infty\le45169.4$. All that is used below is $\|\beta\|_\infty\le B$ for some
$B\le9.7\cdot10^5$; the computed bound is about $4.5\cdot10^4$, so the hypothesis holds with a
factor $21$ to spare, and no sharp value of $V^\star/r_0$ is needed.

\emph{Step 2: one test kernel.} Fix $M$ and a pair $(\tau_0,q)$ with $\tau_0\in\Q$ and
$q\in\C^M$, and put $\varphi(\tau)=|\sum_{k<M}q_ke^{ikd\tau}|^2\ge0$, $W=[\tau_0,\tau_0+P]$,
$P=2\pi/d=\pi M/T$. As in the proof of Proposition~\ref{prop:TK}, $\nuh$ is $P$-periodic with Fourier
modes in $\{|j|\ge M\}$ while $\varphi$ has modes in $\{|j|\le M-1\}$, so $\int_W\nuh\varphi=0$;
integrating $f\ge c$ against $\varphi\ge0$ over $W$ and dividing by $\int_W\varphi=P\|q\|^2$ (exact,
by orthogonality over a full period) gives
\[
c\ \le\ a+b_0\beta_0+b_1\beta_1,\qquad a=\frac{\int_W\sT\varphi}{P\|q\|^2},\quad
b_j=\frac{\int_W\gh j\varphi}{P\|q\|^2},
\]
valid for every $\beta$. This is Proposition~\ref{prop:TK} at a single $(\tau_0,\varphi)$: no
infimum over $\tau_0$, no eigenvalue and no concave maximization are needed, which is what makes the
enclosure cheap. Any pair $(\tau_0,q)$ gives a valid bound, so the way the pair was found (in
practice, as a double-precision eigenvector) plays no role in the proof.

\emph{Step 3: the enclosure at $M=8$.} With $\tau_0=-27.223$ and the vector $q$ contained in the
certificate data accompanying this paper, interval arithmetic (Appendix~\ref{app:encl}) gives
\[
a\in[-6.009422\cdot10^{-3},\,-6.006218\cdot10^{-3}],\qquad |b_0|+|b_1|\le6.18\cdot10^{-9},
\]
so $c\le-6.006218\cdot10^{-3}+6.18\cdot10^{-9}B\le-5.70\cdot10^{-3}<0$ for every $B\le9.7\cdot10^5$,
contradicting $c\ge0$.

\emph{Step 4: the remaining $M$.} The same four enclosures with their own pairs $(\tau_0,q)$ give
$c\le-9.038987\cdot10^{-2}$ $(M=3)$, $-6.234677\cdot10^{-2}$ $(M=5)$, $-2.175677\cdot10^{-2}$ $(M=6)$
and $-4.242754\cdot10^{-2}$ $(M=7)$, all negative, the allowance $(|b_0|+|b_1|)B$ already included.
For $M\in\{1,2,4\}$ no computation is needed: $M\mid8$, so every multiple of $2T/M$ that is $\ge2T$ is
a multiple of $2T/8$ that is $\ge2T$, i.e.\ the class at step $2T/M$ is contained in the class at step
$2T/8$, which Step~3 already excludes.
\end{proof}

\begin{remark}
Theorem~\ref{thm:N} is not a statement about $\lamP(\frac12\log3)$, which is positive by
Theorem~\ref{thm:C}, a certificate at $M=24$; it says that the class $M\le8$ cannot detect this
positivity. It is the converse of the linear-program values of \S\ref{sec:cert}; the numerical
$M=8$ value lies below the certified ceiling $-5.70\cdot10^{-3}$, as it must
(Appendix~\ref{app:diag}).
\end{remark}

\subsection{The same impossibility at the two-prime windows}\label{subsec:Np}

The enclosure of Theorem~\ref{thm:N} is limited by the size of $|a|$, which at $\frac12\log3$ is of
order $10^{-3}$. At the two-prime windows the ceilings are four orders of magnitude smaller, and
reaching them needs a quadrature rule of high order and a test kernel designed against the
$\beta$-box rather than merely against $\sT$. Both are elementary, and they give the same kind of
statement at windows where no certificate is known.

\begin{theorem}[class impossibility at the two-prime windows]\label{thm:Np}
Let $T\in\{0.68,\log2\}$, so $\NT=\{2,3\}$ and
$\Aconst=2.24883$. Let $M\in\{1,\dots,24\}$, $d=2T/M$, let $\nu$ be any finite even signed measure
supported on the lattice $\{\pm jd:j\ge M\}$ (any range, any number of atoms, any signs), and
$\beta\in\R^2$ arbitrary. Then
\[
\inf_{\tau\in\R}\big[\sT(\tau)+\nuh(\tau)+\beta_0\gh0(\tau)+\beta_1\gh1(\tau)\big]\ <\ 0 .
\]
Consequently no dual certificate in the sense of Definition~\ref{def:PD} supported on a lattice of
step $2T/M$ with $M\le24$ has nonnegative value at either window, and none proves $\lamP>0$ there,
for any $\beta$, any range, any atom count and whichever tail lemma of \S\ref{sec:tails} is used to
verify the inequality. In particular
\[
\operatorname{Ceil}(24,0.68)\le-7.493129\cdot10^{-6},\qquad
\operatorname{Ceil}(24,\log2)\le-6.584323\cdot10^{-7}.
\]
The proof includes a computer verification in interval arithmetic (Appendix~\ref{app:encl}).
\end{theorem}
\begin{proof}
The three steps are those of Theorem~\ref{thm:N}. \emph{Step 1.} Theorem~\ref{thm:Bbeta}, with its
witnesses and moment polygon recomputed at these two windows, gives $r_0=2.4309\cdot10^{-4}$,
$V^\star=4.595697$ at $T=0.68$ and $r_0=2.6720\cdot10^{-4}$, $V^\star=4.550671$ at $T=\log2$, hence
$\|\beta\|_\infty\le B$ for every certificate of value $c\ge0$ with $B=1.89\cdot10^4$, resp.
$1.70\cdot10^4$. \emph{Step 2.} For a pair $(\tau_0,q)$ put $\varphi(\tau)=|\sum_{k<M}q_ke^{ikd\tau}|^2$,
$W=[\tau_0,\tau_0+P]$, $P=\pi M/T$; as in Proposition~\ref{prop:TK}, $\int_W\nuh\varphi=0$ and
$c\le a+b_0\beta_0+b_1\beta_1$ with $a=(P\|q\|^2)^{-1}\int_W\sT\varphi$ and
$b_j=(P\|q\|^2)^{-1}\int_W\gh j\varphi$, so $c\le a+(|b_0|+|b_1|)B$. \emph{Step 3.} Sixteen enclosures
per window, at $M\in\{6,7,9,11,13,14,\dots,24\}$, are all negative; the two displayed values are the
upper endpoints at $M=24$, the whole $\beta$-allowance included. For the remaining
$M\in\{1,2,3,4,5,8,10,12\}$ no computation is needed: the class at step $2T/M$ is contained in the
class at step $2T/(kM)$ whenever $k\in\N$, and every such $M$ divides $20$ or $24$, both of which are
enclosed.
\end{proof}

\begin{remark}
Theorem~\ref{thm:Np} is an upper bound on what a class can certify; the sign of $\lamP(0.68)$ and of
$\lamP(\log2)$ is not decided by the methods of this paper. It shows that the class $M=24$, which
carries the certificate of Theorem~\ref{thm:C} at $\frac12\log3$, contains no certificate of
nonnegative value at either of these windows, and neither does any lattice class of order $M\le24$.
\end{remark}

\subsection{Partial prime combs}
The corollary that produced the crude upper bound of \S\ref{sec:levy} applies more generally.
Corollary~\ref{cor:mono} uses only $v_n\ge0$ and $|\cF|\le\|F\|^2$, so it holds for the form
$\mathfrak a^{(v)}_T$ built from any nonnegative weights on any set of contacts.

\begin{theorem}[sub-comb negativity]\label{thm:SC}
Let $\mathcal N\subseteq\{n\ge2:\log n<2T\}$ be any set of contacts and $(v_n)_{n\in\mathcal N}$ any
nonnegative weights of total mass $A=\sum_{n\in\mathcal N}v_n$. Then the bottom of the corresponding
form on $\PT$ is at most $C_0+A+\varphi((3\pi/2T)^2)$. Since $\varphi$ is increasing with
$\varphi(0^+)=0$ and $T\mapsto\varphi((3\pi/2T)^2)$ decreases to $0$, if
\[
A<-C_0=5.3721834\ldots
\]
then the form is negative on $\PT$ for every $T>T_0(A)$, and nonpositive at $T=T_0(A)$, where $T_0(A)=\frac{3\pi}{2\sqrt{x_A}}$ and
$x_A$ is the unique solution of $\varphi(x_A)=-C_0-A$. For instance $T_0(0)=0.74921$,
$T_0(\frac12)=1.23295$, $T_0(1)=2.01848$, $T_0(2)=4.63567$, $T_0(4)=13.19461$, $T_0(5)=29.60905$.
\end{theorem}
\begin{proof}
The displayed bound is Corollary~\ref{cor:mono} with the weights $(v_n)_{n\in\mathcal N}$: the L\'evy
form $\Lev$ does not depend on the weights, its bottom on the codimension-two subspace $\PT$ is at
most $\varphi((3\pi/2T)^2)$, the comb contributes at most $A\|F\|^2$ and the constant is $C_0\|F\|^2$.
Monotonicity and the limits of $\varphi$ are immediate from the series of Theorem~\ref{thm:S}.
\end{proof}
\begin{remark}
The quantitative ancestors of Theorem~\ref{thm:SC} are the observations that the prime terms are
absent for $2T<\log2$ and that beyond that width they enter with a definite cost: Yoshida's prime-free
window \cite{Yoshida1992} and Bombieri's Theorem~12 \cite{Bombieri2000}. Theorem~\ref{thm:SC} adds
that no sub-collection of primes of total mass below $-C_0$ can carry positivity for large $T$.
\end{remark}


\begin{corollary}\label{cor:SC1}
Let the contacts be the ten primes $12049$, $124909$, $149333$, $322573$, $324011$, $373963$,
$417491$, $630229$, $782209$, $908861$, with their natural weights $v_p=2\log p/\sqrt p$ and all other
weights $0$. Then $A=0.563667\ldots$ and $T_0(A)=1.313456\ldots$. At
$T=\frac12\log908861+0.01=6.869974\ldots$, the smallest window in this family containing all ten
contacts, Theorem~\ref{thm:SC} gives
\[
\lambda_{\min}\big(A_T(v)\big|_{\PT}\big)\ \le\ -5.3721834+0.563667+2.685686=-2.122831 .
\]
\end{corollary}

Ten primes of size $10^5$--$10^6$, in a window that contains them, give a form that is negative by
$2.12$, about two fifths of the archimedean constant $|C_0|$. The complete comb at that window has
$\Aconst\sim4e^T\approx3.8\cdot10^3$, and its form is conjecturally nonnegative. Positivity, if it
holds, is thus a property of the complete initial segment $\{n<e^{2T}\}$: an argument that
establishes the inequality for a proper sub-comb and then adds contacts one at a time cannot work,
because every sub-comb of mass $<-C_0$ is negative in a large enough window. The mechanism is
explicit: each contact adds at most $v_n$ to the bound, the archimedean deficit $-C_0$ is a
constant, and the eigenvalue $\varphi((3\pi/2T)^2)$ of the window decays to $0$; only a comb whose
mass grows with $T$ keeps up. This concerns the operators $A_T(v)$, not $\zeta$. For comparison, the
window whose contact set is exactly the prime powers below $31$ (the sixteen numbers
$2,3,4,5,7,8,9,11,13,16,17,19,23,25,27,29$, among them the first ten primes) is
$T=\frac12\log31=1.71673$, where $\Aconst=16.640459$ exceeds $-C_0$ and Theorem~\ref{thm:SC} is
vacuous; the multiplier-free verification threshold there is
$2\pi e^{\Aconst}=1.06\cdot10^8$, against $16.7$ at $\frac12\log3$ and $59.5$ at $\log2$.


\subsection{The ceilings converge}\label{subsec:R}
Proposition~\ref{prop:TK} bounds a lattice class from above by one eigenvalue. The question left
open after it is whether this bound is asymptotically sharp, i.e.\ whether the test kernels see all
of the primal problem. They do, in the limit; the mechanism is Poisson summation used in the
direction opposite to the one that makes a lattice $\nuh$ periodic.

\begin{lemma}[the periodization of a primal measure is a test kernel]\label{lem:R1}
Let $\mu$ be feasible for \textup{(P)}, $M\in\N$, $d=2T/M$, $P=2\pi/d=\pi M/T$, and let $\rho$ be the
periodization of $\mu$ on $\R/P\Z$, $\rho(E)=\sum_{k\in\Z}\mu(E+kP)$. Then $\rho$ is a probability
measure with Fourier coefficients $\widehat\rho(j)=\check\mu(jd)$, which vanish for $|j|\ge M$. Hence
$\rho$ is absolutely continuous with a nonnegative trigonometric polynomial density of degree
$\le M-1$ in $\theta=d\tau$, and by the Fej\'er--Riesz theorem
$d\rho=P^{-1}|q(e^{id\tau})|^2d\tau$ for some polynomial $q$ of degree $\le M-1$: $\rho$ is a test
kernel of order $M$ in the sense of Proposition~\ref{prop:TK}.
\end{lemma}
\begin{proof}
$\rho$ is a finite positive measure of mass $\mu(\R)=1$, and
$\widehat\rho(j)=\int_{[\tau_0,\tau_0+P)}e^{-ijd\tau}d\rho=\int_\R e^{-ijd\tau}d\mu=\check\mu(jd)$,
because $e^{-ijd(\tau+kP)}=e^{-ijd\tau}$ ($dP=2\pi$) and $\mu$ is even. Since $|jd|\ge2T$ exactly when
$|j|\ge M$, and feasibility is the closed condition $\check\mu=0$ on $\{|u|\ge2T\}$ (so the boundary
index $j=M$ is included), only finitely many coefficients are nonzero; a measure with finitely many
nonzero Fourier coefficients is a trigonometric polynomial, here nonnegative, and Fej\'er--Riesz
factorizes it.
\end{proof}

\begin{theorem}[the ceilings converge to the primal value]\label{thm:R}
For $\beta\in\R^2$ write $\operatorname{Ceil}(M,T;\beta)=\inf_{\tau_0}\lambda_{\min}(C_M(\tau_0,\beta))$
and $S_\beta=\sT+\beta\cdot\widehat g$. Then for every $\mu$ feasible for \textup{(P)},
\[
\operatorname{Ceil}(M,T;\beta)\ \le\ \int\sT\,d\mu+\Delta(M,\beta)\,\mu\big(\{|\tau|>\pi M/2T\}\big),
\]
\[
\Delta(M,\beta)=\big[m(P/2)-C_0\big]+2\Aconst+2|\beta_0|\gh0(0)+2|\beta_1|\gh1(0),
\]
with $P/2=\pi M/2T$. Consequently $\limsup_{M\to\infty}\operatorname{Ceil}(M,T;\beta)\le\lamP$ for
every fixed $\beta$, and uniformly for $\beta$ in a bounded box.
\end{theorem}
\begin{proof}
Take $\tau_0=-P/2$, $W=[-P/2,P/2]$, and let $\pi_W(\tau)=\tau-kP\in W$ be the folding map. By
Lemma~\ref{lem:R1} and the Fej\'er--Riesz description of the matrix $C_M$ in the proof of
Proposition~\ref{prop:TK}, the infimum over test kernels is at most its value at $P\rho$, so
$\operatorname{Ceil}(M,T;\beta)\le\int_WS_\beta\,d\rho=\int_\R S_\beta(\pi_W\tau)\,d\mu(\tau)$. Now
$\int S_\beta d\mu=\int\sT d\mu$ because $\mu$ has both moments zero, and $\pi_W\tau=\tau$ on
$|\tau|\le P/2$, so
\[
\int S_\beta\circ\pi_W\,d\mu-\int\sT\,d\mu=\int_{|\tau|>P/2}\big(S_\beta\circ\pi_W-S_\beta\big)d\mu
\le\big(\sup_WS_\beta-\inf_\R S_\beta\big)\,\mu(|\tau|>P/2),
\]
and $\sup_WS_\beta\le m(P/2)+\Aconst+|\beta_0|\gh0(0)+|\beta_1|\gh1(0)$ ($m$ increasing,
$|\gh j|\le\gh j(0)$) while $\inf_\R S_\beta\ge C_0-\Aconst-|\beta_0|\gh0(0)-|\beta_1|\gh1(0)$, which
is $\Delta$. For the limit, fix $\delta>0$ and $\mu$ with $\int\sT d\mu\le\lamP+\delta$. Then
$\int m\,d\mu<\infty$ (since $\sT\ge m-\Aconst$), so by dominated convergence
$(m(R)-C_0)\mu(|\tau|>R)\le\int_{|\tau|>R}(m-C_0)d\mu\to0$ as $R\to\infty$; as
$\Delta(M,\beta)=(m(P/2)-C_0)+O_\beta(1)$ and $P/2\to\infty$, the product tends to $0$, uniformly for
$\beta$ in a bounded box.
\end{proof}

\begin{corollary}[the test-kernel family is not lossy in the limit]\label{cor:R3}
Suppose $\lamP>0$, and let $B=V^\star/r_0$ be the a priori bound of Theorem~\ref{thm:Bbeta}, so that
every certificate of value $c\ge0$ has $\|\beta\|_\infty\le B$. Put
$\operatorname{Ceil}_B(M,T)=\sup_{\|\beta\|_\infty\le B}\operatorname{Ceil}(M,T;\beta)$. Then
$\operatorname{Ceil}_B(M,T)\to\lamP$ along the nested lattices $2T/2^kM_0$ with growing range. In
particular the subfamily of test kernels $\{|q|^2\}$ used in Proposition~\ref{prop:TK} is not lossy in
the limit: its optimal value is the same $\lamP$ as that of the full relaxed-feasible family.
\end{corollary}
\begin{proof}
$\limsup_M\operatorname{Ceil}_B\le\lamP$ is Theorem~\ref{thm:R}, $\Delta$ being bounded uniformly on
the box. Conversely, by Theorem~\ref{thm:D1} there are, for every $\varepsilon>0$, lattice
certificates of value $c>\lamP-\varepsilon$ on the nested lattices (whose union is dense in
$[2T,\infty)$); eventually $c\ge0$, so by Theorem~\ref{thm:Bbeta} their $\beta$ lies in the box, and
Proposition~\ref{prop:TK} gives $c\le\operatorname{Ceil}_B(M,T)$.
\end{proof}
We do not know whether the family is lossy at fixed $M$, nor at what rate the ceilings converge. The
next corollary shows that no rate holds uniformly over the feasible set.

\begin{corollary}[no uniform rate; a rate on a sub-family]\label{cor:CP}
Let $\Delta(M,\beta)$ be as in Theorem~\ref{thm:R}, $P=\pi M/T$, and let
$\mu\in\mathcal F_\varepsilon$, so that Theorem~\ref{thm:R} reads
$\operatorname{Ceil}(M,T;\beta)\le\lamP+\varepsilon+\Delta(M,\beta)\,\mu(|\tau|>P/2)$.
\emph{(i)} Unconditionally, Theorem~\ref{thm:P3}(i) gives
$\operatorname{Ceil}(M,T;\beta)\le\lamP+\varepsilon+G(\varepsilon)[1+K_\beta/(m(P/2)-C_0)]$ with
$K_\beta=2\Aconst+2|\beta_0|\gh0(0)+2|\beta_1|\gh1(0)$, which is $O(1)$ and not $o(1)$: the
convergence of Theorem~\ref{thm:R} cannot be made quantitative uniformly over the feasible set.
\emph{(ii)} On the sub-family of feasible measures satisfying $\int(m-C_0)^{1+\delta}d\mu\le K$ for
some $\delta>0$ (non-empty in $\mathcal F_\varepsilon$ for every $\varepsilon>0$, since the $F_0$
of Theorem~\ref{thm:Vp}(ii) have Schwartz transforms), H\"older gives
$\mu(|\tau|>P/2)\le K(m(P/2)-C_0)^{-1-\delta}$ and hence
\[
\operatorname{Ceil}(M,T;\beta)\ \le\ \lamP+\varepsilon+O\big(K\,(\log M)^{-\delta}\big),
\]
since $m(P/2)-C_0=\log M+O_T(1)$. The constant $K=K(\varepsilon,\delta,T)$ is not explicit and, by
Remark~\ref{rem:P3exp}, cannot be made uniform over $\mathcal F_\varepsilon$.
\end{corollary}
\begin{proof}
Immediate from Theorem~\ref{thm:R} together with Theorem~\ref{thm:P3}(i) in case (i) and
H\"older's inequality in case (ii); the sub-family is non-empty because $\mu_{F_0}$ has a rapidly
decaying density.
\end{proof}

\subsection{Sampling a flat-top certificate}\label{subsec:lat}
Problem~\ref{prob:A} asks for the loss incurred by restricting the dual to a lattice class. For the
flat-top certificates of \S\ref{sec:cert} the loss of the sampled certificate can be computed
exactly. It is an edge effect of order $1/M$, not a smoothness effect, and it is carried entirely by
the polar block: the density of $\nu_w$ in \eqref{eq:nuw} is $\frac12(1-w)K+2w\cosh\frac u2$, whose
second summand does not vanish at $u=2T^+$ (it equals $2\cosh T$ there, the jump already met in
the proof of Theorem~\ref{thm:Dtail}), so no smoothness of $w$ can make the sampling error
$O(M^{-r})$.

\begin{theorem}[lattice sampling of a flat-top certificate]\label{thm:AL}
Let $r\ge2$ and let $w\in\Wcl$ be even with $w^{(r-1)}$ locally absolutely continuous and
$w^{(r)}\in \mathrm{BV}\cap L^1$, $\supp w\subseteq[-U,U]$ with $C^{r-1}$ contact with $0$ at $\pm U$,
and $w(\log n)=0$ for every $n$ with $\log n>2T$ \textup{(}comb-vanishing, finitely many linear
conditions\textup{)}. No sign condition and no bound $0\le w\le1$ is assumed. The boundary case
$\log n=2T$ is not excluded; it is treated by carrying the atom it produces \textup{(}see
$\nu_{\mathrm{lat}}$ below\textup{)}, which is legitimate because Definition~\ref{def:W} forces
$w(2T)=1$ there, so that atom is a fixed multiple of $\delta_{\pm2T}$ and can be reproduced exactly.
Put
$W(u)=w(|u|)\cosh\frac u2$, $R=\frac12(1-w)K\mathbf 1_{|u|\ge2T}$ and
$G=R+2W\mathbf 1_{|u|\ge2T}$, which by comb-vanishing is the whole absolutely continuous density
of $\nu_w$, the only surviving atom being the boundary one at $\pm2T$ when $2T\in\{\log n\}$; put
$V_0=\|G\|_{L^1}$, $V_r=\|R^{(r)}\|_{L^1}+2\|W^{(r)}\|_{L^1}$ and $c_w=\inf_\R f_w$. For $M\in\N$ set
$d=2T/M$, $P=2\pi/d=\pi M/T$ and
\[
\nu_{\mathrm{lat}}=\sum_{|j|\ge M}d\,\kappa_j\,G(jd)\,\delta_{jd}
-\sum_{n:\,\log n=2T}\tfrac{v_n}2\big(\delta_{2T}+\delta_{-2T}\big),\qquad
\kappa_{\pm M}=\tfrac12,\quad\kappa_j=1\ (|j|>M),
\]
an even signed measure with countably many atoms and finite total variation
$\|\nu_{\mathrm{lat}}\|_{\TV}=V_0+O(d)$. Then $(\nu_{\mathrm{lat}},\beta=(2,0))$ is dual feasible and
\[
c_{\mathrm{lat}}(M)\ \ge\ \min\Big\{c_w-E_{\mathrm{edge}}(M)-\varepsilon_r(M),\
m\big(\tfrac{\pi M}{2T}\big)-\Aconst-\|\nu_{\mathrm{lat}}\|_{\TV}-2\gh0(0)\Big\},
\]
\[
E_{\mathrm{edge}}(M)=\frac{8\cosh T}{P}+\frac{2\pi^2\sinh T}{P^2}+\frac{16.9\cosh T}{P^3}
=\frac4\pi\cosh T\cdot d+O(d^2),\qquad
\varepsilon_r(M)=2\big(1-2^{-r}\big)\zeta(r)V_r\Big(\frac{\pi M}{2T}\Big)^{-r}.
\]
Moreover $\sup_\tau|f_{\mathrm{lat}}-f_w|=\frac4\pi\cosh T\cdot d\,(1+O(d))$: the rate $M^{-1}$ and
its constant are attained, and no smoothness of $w$ improves them.
\end{theorem}
\begin{proof}
Feasibility: the atoms lie at $|u|=|j|d\ge Md=2T$, $(D)$ imposes no moment equality on $\nu$
(Proposition~\ref{prop:weak}), Definition~\ref{def:PD} allows $\nu$ to be signed, the boundary atom
at $\pm2T$ is harmless because $\cF(\pm2T)=0$, and $\beta$ is free. The atoms with $|j|<M$ are absent
because $G$ vanishes on the window. Finiteness of the total variation is
$\|\nu_{\mathrm{lat}}\|_{\TV}=d\sum_{|j|\ge M}\kappa_j|G(jd)|<\infty$, since $G\in L^1\cap
\mathrm{BV}$ on $[2T,\infty)$; the number of atoms is countably infinite, as $G=\frac12K>0$ beyond
$U$.

Value: since $w(\log n)=0$ for $\log n>2T$, $\nu_w=G(u)\,du-\sum_{\log n=2T}\frac{v_n}2(\delta_{2T}+\delta_{-2T})$,
and the second sum, present exactly when $2T\in\{\log n:n\ge2\}$ (as at the window $T=\frac12\log3$,
where it is $-\frac{v_3}2(\delta_{2T}+\delta_{-2T})$), has been put verbatim into
$\nu_{\mathrm{lat}}$, so it cancels in $f_{\mathrm{lat}}-f_w$ and may be ignored below.
Theorem~\ref{thm:Dtail}
gives $\nuh_w+2\gh0=2\widehat W+\Psi_w$ with $\Psi_w=\widehat R$. Because $2T=Md$ is a lattice point,
the samples of $G$ on $|j|\ge M$ split into those of $R$ and those of $2W$, so
\[
\nuh_{\mathrm{lat}}(\tau)=\Big[d\sum_{|j|\ge M}\kappa_jR(jd)e^{ijd\tau}\Big]
+\Big[d\sum_{j\in\Z}2W(jd)e^{ijd\tau}-T_d(\tau)\Big],
\]
where $T_d$ is the end-point-corrected trapezoidal sum, on the nodes of the window, of
$h=2\cosh\frac u2\mathbf 1_{[-2T,2T]}$, whose transform is $\widehat h=2\gh0$. Both $R$ and $2W$ are
$C^{r-1}$ with $r$-th derivative in $L^1$, so Poisson summation applies absolutely to each and the
two bracketed defects are bounded on $|\tau|\le P/2$ by
$V_r\sum_{k\ne0}|\tau+kP|^{-r}\le\varepsilon_r(M)$. Hence
$f_{\mathrm{lat}}-f_w=[2\gh0-T_d]+O(\varepsilon_r)$: the whole non-smooth cost is the quadrature
error of the polar block. Poisson summation for the compactly supported function $h$ of bounded
variation, with its symmetric values at the two jumps, gives $T_d=\sum_k\widehat h(\tau+kP)$, and
since $2TP=2\pi M$ every alias carries the same phase:
\[
2\gh0-T_d=-4\Big[\tfrac12\sinh T\cos2T\tau\sum_{k\ne0}\frac1{\xi_k^2+\frac14}
+\cosh T\sin2T\tau\sum_{k\ne0}\frac{\xi_k}{\xi_k^2+\frac14}\Big],\qquad\xi_k=\tau+kP .
\]
The first sum is at most $\pi^2/P^2$. For the second write
$\xi/(\xi^2+\frac14)=\xi^{-1}-\frac14\xi^{-1}(\xi^2+\frac14)^{-1}$ and use
$\sum_{k\ne0}\xi_k^{-1}=P^{-1}\big[\pi\cot\pi z-z^{-1}\big]$ with $z=\tau/P$, a function
regular at $z=0$ (the subtracted $z^{-1}$ is exactly the omitted $k=0$ term), odd and of
modulus at most $2/P$ on $|z|\le\frac12$; the correction is at most $4.21/P^3$. This is
$E_{\mathrm{edge}}$, uniformly in $\tau$, and it is attained where $\sin2T\tau=\pm1$. For
$|\tau|>P/2$ use $\sT\ge m(P/2)-\Aconst$ (Lemma~\ref{lem:mbounds}),
$|\nuh_{\mathrm{lat}}|\le\|\nu_{\mathrm{lat}}\|_{\TV}$ and $|2\gh0|\le2\gh0(0)$.
\end{proof}

\begin{corollary}\label{cor:ALdef}
At fixed $w$ and for $M\ge M_0(w)$ (large enough that the second branch of the minimum in
Theorem~\ref{thm:AL} does not bind, which happens for all large $M$ since $m(\pi M/2T)\to\infty$),
the lattice sampling costs
$c_w-c_{\mathrm{lat}}(M)\le\frac8\pi T\cosh T\cdot M^{-1}+O(M^{-2})+O_r(V_rM^{-r})$, and the
$M^{-1}$ is not removable: it is the alias of the jump of $2\cosh\frac u2\mathbf 1_{|u|<2T}$ at the
window edge, an object independent of $w$. Without the end-point correction
\textup{(}$\kappa_{\pm M}=1$\textup{)} the sharp constant is
$\cosh T\sqrt{4+16/\pi^2}\,d=2.3708\cosh T\cdot d$, so the boundary atom of weight $d\cosh T$ at
$\pm2T$ improves the constant by the factor $1.862$ and nothing improves the order.
\end{corollary}

Two remarks on Corollary~\ref{cor:ALdef}. First, for sampled certificates the whole lattice cost
sits at $u=\pm2T$, the same place as the non-attainment of Theorem~\ref{thm:O}; the smooth blocks $R$
and $W$ contribute only $O_r(M^{-r})$. We do not know whether the lattice class, with weights
optimized rather than sampled, also loses $\Theta(M^{-1})$; this is the remaining half of
Problem~\ref{prob:A}. Second, up to $O_r(M^{-r})$ one has $f_{\mathrm{lat}}-f_w=2\gh0-T_d$, the
discrepancy between the polar multiplier and its Riemann sum over the window; $T_d$ is a
$P$-periodic trigonometric sum of degree $M$ and $2\gh0$ is not a trigonometric polynomial, so the
difference is not one either. Heuristically, this is why Proposition~\ref{prop:TK}, whose ceiling
rests on orthogonality to every mode $|j|<M$, need not bind flat-top certificates; neither
statement implies this.

\section{The archimedean form}\label{sec:archstr}

With the primes deleted, the Weil form is an explicit convolution operator. This section studies
its thresholds, its negative index, and the number of even modes that the prime form can compete
with.

\subsection{Thresholds and nodal structure}\label{subsec:arch}

Theorem~\ref{thm:SC} with an empty contact set is a statement about the archimedean block alone. It
locates the window beyond which the primes are not a perturbation of a positive form but the source
of its positivity. Put, for $F\in L^2(\IT)$,
\[
\Ainf(F)=2A(F)B(F)+\frac1{2\pi}\int_\R m(\tau)|\Fh(\tau)|^2\,d\tau
=2A(F)B(F)+C_0\|F\|^2+\Lev(\cF),
\]
the complete Weil form with the prime comb deleted (the second equality is Theorem~\ref{thm:A} with
$\NT=\emptyset$), and write $\PiT(F)=-\sum_{n\in\NT}v_n\Rel\cF(\log n)$ for the prime form (the real part is redundant for real $F$), so that
$W_T=\Ainf+\PiT$. Let $\laminf(T)$ be the bottom of $\Ainf$ on the unit sphere of $L^2(\IT)$ and
$\laminfP(T)$ its bottom on the unit sphere of $\PT$, where the pole term vanishes identically and
$\Ainf=C_0\|F\|^2+\Lev(\cF)$.

\begin{proposition}[two archimedean thresholds]\label{prop:arch}
\emph{(i)} $\Ainf$ does not depend on $T$ and the feasible sets increase with $T$, so both
$\laminf$ and $\laminfP$ are nonincreasing (for $\laminfP$ strictly, Lemma~\ref{lem:strict}).
\emph{(ii)} $\laminf(T)\le\laminfP(T)\le C_0+\varphi\big((3\pi/2T)^2\big)$; the first inequality is
inclusion of feasible sets and the second is Theorem~\ref{thm:SC} with $\mathcal N=\emptyset$,
$A=0$. Hence both are nonpositive at $T=T_0(0)=0.74921$ and negative for $T>T_0(0)$ (at $T_0(0)$
itself, as for every $T\ge0.60$, $\laminf\le\laminfP<0$ by Corollary~\ref{cor:Tarch}), and $\laminfP(T)\downarrow C_0$ as
$T\to\infty$, since $\Lev\ge0$ and $\varphi((3\pi/2T)^2)\downarrow0$.
\emph{(iii)} For $2T<\log2$ the contact set is empty, so $W_T=\Ainf$ and $\laminf(T)\ge0$ by the
known positivity of the complete form at autocorrelation support $\le\log2$
\textup{\cite{Yoshida1992,Bombieri2000}} (and, certified, \cite[Thm.~1.2]{Zhu2026}).
\emph{(iv)} Consequently the thresholds
$T^*_\infty=\sup\{T:\laminf(T)>0\}$ and $T^*_{\infty,\mathrm P}=\sup\{T:\laminfP(T)>0\}$ exist, and
\[
T^*_\infty\ \le\ T^*_{\infty,\mathrm P}\ \le\ 0.74921 ;
\]
for $T^*_{\infty,\mathrm P}$ this is sharpened to $0.58\le T^*_{\infty,\mathrm P}\le0.60$ by
Corollary~\ref{cor:Tarch}, and for $T^*_\infty$ to $T^*_\infty\le0.4608<0.58\le T^*_{\infty,\mathrm P}$ by
Proposition~\ref{prop:archneg}.
\end{proposition}
\begin{proof}
(i) $K$, $C_0$ and the two Tate moments are $T$-free, and an infimum over a growing set is
nonincreasing. (ii) $\PT\subset L^2(\IT)$, so the restricted infimum is the larger one; the proof of
Corollary~\ref{cor:mono} bounds the bottom of $\Lev$ on the codimension-two subspace $\PT$ by
$\lambda_3\le\varphi((3\pi/2T)^2)$ and never mentions the primes, and the comb contributes $0$.
(iii) and (iv) are immediate; a nonincreasing function changes sign once, and the thresholds are
suprema. No continuity is used. For $\laminf$ it holds by the proof of Lemma~\ref{lem:contT}, and for
$\laminfP$ it is not asserted.
\end{proof}
\begin{remark}
In (iii) the strict inequality $\laminf>0$ is not re-proved here, and that is the only reason
(iv) is stated for $T^*_\infty$ as an upper bound: strict positivity at $T=\frac12\log2$,
available in \cite{Yoshida1992}, would add the lower bound $\frac12\log2\le T^*_\infty$. For
$T^*_{\infty,\mathrm P}$ a two-sided bound is proved below (Corollary~\ref{cor:Tarch}).
\end{remark}
\begin{remark}
The archimedean form by itself was studied by Connes and Consani \cite{ConnesConsani2021}, who prove
its positivity on a window by a Toeplitz/prolate-spheroidal compression and relate the archimedean
place to the trace formula. Proposition~\ref{prop:arch} looks at the same form from the other side:
it locates the widths beyond which the archimedean form alone stops being positive, first on
$L^2(I_T)$ and then on $\PT$, so that past $T^*_{\infty,\mathrm P}$ the prime form $\PiT$ has to
supply positivity rather than merely be dominated.
\end{remark}

The bottom eigenfunction of the archimedean form on $\PT$ is never of one sign, and this is forced
by the two Tate conditions alone.

\begin{proposition}[nodal structure of $\PT$]\label{prop:nodes}
Write $\Gamma_T(F)=C_0\|F\|^2+\Lev(\cF)$ for the archimedean form without the pole term on
$L^2(\IT)$, so that $\Ainf=\Gamma_T$ on $\PT$. Then:
\emph{(i)} the bottom of $\Gamma_T$ on all of $L^2(\IT)$ is a simple eigenvalue whose
eigenfunction can be taken a.e.\ strictly positive;
\emph{(ii)} say that $F\in L^2(\IT)$ has \emph{at most one sign change} if for some
$t_0\in[-T,T]$ one has $F\ge0$ a.e.\ on $(-T,t_0)$ and $F\le0$ a.e.\ on $(t_0,T)$, or the same
with the two inequalities exchanged (the endpoints $t_0=\pm T$ being the case of one sign).
Then every nonzero $F\in\PT$ has at least two sign changes in this sense; in particular no
nonzero element of $\PT$ is of one sign, and an even element has at least one nodal pair in
$(0,T)$;
\emph{(iii)} consequently the bottom eigenfunction of $\Gamma_T$ on $\PT$ is not of one sign,
and the Sturm index of $\Ainf|\PT$ is shifted by one nodal pair relative to $(i)$.
\end{proposition}
\begin{proof}
(i) For $F\in L^2(\IT)$ extended by zero, Theorem~\ref{thm:A} and Fubini give the pure-jump
Dirichlet form with killing
\[
\Lev(\cF)=\frac14\iint_{\IT\times\IT}K(|x-y|)\,|F(x)-F(y)|^2\,dx\,dy
+\frac12\int_{\IT}\kappa(x)|F(x)|^2\,dx,\qquad
\kappa(x)=\int_{|y|>T}K(|x-y|)\,dy<\infty,
\]
obtained by splitting $2\Lev(\cF)=\int_\R K(|u|)\big(\|F\|^2-\Rel\cF(u)\big)du
=\frac12\iint_{\R^2}K(|x-y|)|\widetilde F(x)-\widetilde F(y)|^2$ into the square $\IT\times\IT$
and the two strips on which the zero extension $\widetilde F$ vanishes; the divergence of $K$ at
the origin is absorbed by $|F(x)-F(y)|^2$ on the form domain, and no $\delta$-kernel exists,
$\int_0^{2T}K=+\infty$. The form is therefore a regular Dirichlet form of pure-jump type with
nonnegative jump density $K$ and killing $\kappa$, up to the additive constant $C_0\|F\|^2$,
which only shifts the spectrum. By the Beurling--Deny criteria the semigroup $e^{-t\Gamma_T}$ is
positivity preserving, and since $K>0$ everywhere it is irreducible, hence positivity improving;
Perron--Frobenius then gives a simple, a.e.\ strictly positive ground state, provided the bottom
is an eigenvalue. It is: $m$ is increasing (Lemma~\ref{lem:mbounds}(i)) and $m(\tau)\to\infty$, so
$\int_{|\tau|>R}|\Fh|^2\le2\pi\,(m(R)-C_0)^{-1}\Lev(\cF)$ uniformly on form-bounded sets, while all
competitors are supported in the fixed compact $\IT$; by the Riesz--Kolmogorov criterion the form
domain embeds compactly in $L^2(\IT)$ and the resolvent is compact.
(ii) The pair $\{e^{t/2},e^{-t/2}\}$ is a Chebyshev system of order $2$ on $(-T,T)$: a nonzero
combination $c_1e^{t/2}+c_2e^{-t/2}$ vanishes only where $e^{t}=-c_2/c_1$, i.e.\ at most once. Let
$0\ne F\in\PT$ and suppose $F$ had at most one sign change, at $t_0$. Choose
$G=c_1e^{t/2}+c_2e^{-t/2}$ vanishing exactly at $t_0$ (any nonzero $G$ of the family if
$t_0=\pm T$), with the sign of $G$ matched to that of $F$ on each side. Then $FG\ge0$ a.e.\ and
$FG$ is not a.e.\ zero, so $\int FG\ne0$; but $F\in\PT$ means $\langle F,e^{t/2}\rangle=\langle
F,e^{-t/2}\rangle=0$, whence $\int FG=0$, a contradiction. For even $F$ the two sign changes are
symmetric, so at least one lies in $(0,T)$.
(iii) Immediate from (ii).
\end{proof}

Numerically, at $T=\log2$ the bottom eigenvector of $\Ainf$ on $\PT$ has exactly one nodal pair
(Remark~\ref{rem:nodal}, Appendix~\ref{app:diag}), the minimum that Proposition~\ref{prop:nodes}(ii)
forces.

\emph{Parity sectors.} Let $(RF)(t)=F(-t)$. For even $F$ the two Tate moments coincide,
$A(F)=B(F)=\int F(t)\cosh\frac t2\,dt$, and for odd $F$ one has $A(F)=-B(F)=\int F(t)\sinh\frac t2\,dt$.
Write $\PT^{\mathrm{ev}}$ and $\PT^{\mathrm{odd}}$ for the even and the odd elements of $\PT$
($\PT^{\mathrm{ev}}$ is the space of Theorem~\ref{thm:arena} below), $\lambda_{\mathrm{ev}}(T)$ and
$\lambda_{\mathrm{odd}}(T)$ for the bottoms of $\Ainf=\Gamma_T$ on their unit spheres (so
$\lambda_{\mathrm{ev}}=\lambda^{\mathrm{ev}}_1$ in the notation of Theorem~\ref{thm:arena}),
$n_{\mathrm{ev}}(T)$ and $n_{\mathrm{odd}}(T)$ for the numbers of negative eigenvalues of $\Ainf$ on them,
and $T^*_{\mathrm{ev}}=\sup\{T:\lambda_{\mathrm{ev}}(T)>0\}$, $T^*_{\mathrm{odd}}=\sup\{T:\lambda_{\mathrm{odd}}(T)>0\}$.
By Lemma~\ref{lem:parity} of \S\ref{subsec:negindex}, $\laminfP=\min(\lambda_{\mathrm{ev}},\lambda_{\mathrm{odd}})$
and $n_-(\Ainf|\PT)=n_{\mathrm{ev}}+n_{\mathrm{odd}}$. For a parity label $\mathrm{par}\in\{\mathrm{ev},\mathrm{odd}\}$
we write $\PT^{\mathrm{par}}$, $\lambda_{\mathrm{par}}$, $n_{\mathrm{par}}$, $T^*_{\mathrm{par}}$.

\begin{lemma}[strict monotonicity]\label{lem:strict}
For every $T>0$ the bottoms $\lambda_{\mathrm{ev}}(T)$, $\lambda_{\mathrm{odd}}(T)$ and
$\laminfP(T)$ are attained, and all three are strictly decreasing on $(0,\infty)$. Consequently, for each
sector, $\lambda_{\mathrm{par}}(T)>0$ for $T<T^*_{\mathrm{par}}$ and $\lambda_{\mathrm{par}}(T)<0$ for
$T>T^*_{\mathrm{par}}$; $\lambda_{\mathrm{par}}(T)<0$ if and only if $n_{\mathrm{par}}(T)\ge1$;
$T^*_{\infty,\mathrm P}=\min(T^*_{\mathrm{ev}},T^*_{\mathrm{odd}})$; and $\laminfP$ is positive on
$(0,T^*_{\infty,\mathrm P})$ and negative on $(T^*_{\infty,\mathrm P},\infty)$.
\end{lemma}
\begin{proof}
\emph{Attainment.} $\Gamma_T-C_0\|\cdot\|^2=\Lev(c_\cdot)\ge0$ is a closed form whose domain embeds
compactly in $L^2(\IT)$ (proof of Proposition~\ref{prop:nodes}(i)); $\PT^{\mathrm{par}}$ and $\PT$ are
closed, $A$ and $B$ being continuous on $L^2(\IT)$. A minimizing sequence on the unit sphere has an
$L^2$-convergent subsequence whose limit has norm $1$ and lies in the subspace, and lower semicontinuity
of the closed form gives attainment. Since the symbol is real and even, $\Gamma_T(u+iv)=\Gamma_T(u)+\Gamma_T(v)$
for real $u,v$, and $u+iv\in\PT^{\mathrm{par}}$ if and only if $u,v\in\PT^{\mathrm{par}}$ ($A,B$ are real on
real functions), so the minimizer may be taken real.

\emph{Strictness.} Monotonicity is Proposition~\ref{prop:arch}(i), sector by sector: extension by zero
maps $\mathcal P^{\mathrm{par}}_{T'}$ isometrically into $\PT^{\mathrm{par}}$ for $T'<T$ with the same
values of $\Gamma$ (Lemma~\ref{lem:zeroext}). Suppose $T'<T$ and
$\lambda_{\mathrm{par}}(T)=\lambda_{\mathrm{par}}(T')=:\lambda$, and let $F^*$ be a real minimizer on
$\mathcal P^{\mathrm{par}}_{T'}$. Extended by zero it is a minimizer on $\PT^{\mathrm{par}}$, so
$\Gamma_T(F^*,G)=\lambda\langle F^*,G\rangle$ for every $G\in\PT^{\mathrm{par}}$ in the form domain. Take
real $G_0\in C_c^\infty(T',T)$ with $\int G_0e^{y/2}dy=\int G_0e^{-y/2}dy=0$ and $G=G_0+\epsilon RG_0$,
$\epsilon=\pm1$ the parity sign. Then $G\in\PT^{\mathrm{par}}$, $\langle F^*,G\rangle=0$ by disjoint
supports, and $\Gamma_T(F^*,G)=2\Gamma_T(F^*,G_0)$, since $\Gamma_T$ is $R$-invariant and
$RF^*=\epsilon F^*$. Polarizing the Dirichlet form of Proposition~\ref{prop:nodes}(i),
$\Lev(\cF)=\frac14\iint_{\R^2}K(|x-y|)|\widetilde F(x)-\widetilde F(y)|^2$, for real $F,G$ with
supports at positive distance gives
$\Gamma_T(F,G)=-\frac12\iint K(|x-y|)F(x)G(y)\,dx\,dy$, absolutely convergent. Hence
$\Phi(y):=\int_{I_{T'}}K(|x-y|)F^*(x)\,dx$ satisfies $\int\Phi G_0=0$ for every such $G_0$, so
$\Phi=\alpha e^{y/2}+\beta e^{-y/2}$ on $(T',T)$. For $y>T'\ge x$, Lemma~\ref{lem:K} and Tonelli give
$\Phi(y)=2\sum_{k\ge0}c_ke^{-(2k+\frac12)y}$ with $c_k=\int F^*(x)e^{(2k+\frac12)x}dx$,
$|c_k|\le\|F^*\|_1e^{(2k+\frac12)T'}$; so $e^{y/2}\Phi(y)=2\sum_kc_kw^k$, $w=e^{-2y}$, is a convergent
power series for $0<w<e^{-2T'}$, and the identity extends to all $y>T'$. As $y\to\infty$, $\Phi\to0$,
so $\alpha=0$; then $2\sum_kc_kw^k\equiv\beta$, so $c_k=0$ for $k\ge1$ (the value $c_0=\beta/2$ is not
needed). With $u=e^{2x}$, $c_k=\int f(u)u^k\,du$ for
$f(u)=\frac12F^*(\frac12\log u)u^{-3/4}\in L^1[e^{-2T'},e^{2T'}]$, so all moments of $uf(u)$ vanish and
$uf=0$ by Weierstrass, whence $F^*=0$, contradicting $\|F^*\|=1$. For $\laminfP$ use Lemma~\ref{lem:parity}.

\emph{Consequences.} If $T>T^*_{\mathrm{par}}$, pick $T'\in(T^*_{\mathrm{par}},T)$: then
$\lambda_{\mathrm{par}}(T')\le0$ and $\lambda_{\mathrm{par}}(T)<\lambda_{\mathrm{par}}(T')$. If
$T<T^*_{\mathrm{par}}$, some $T''>T$ has $\lambda_{\mathrm{par}}(T'')>0$, and monotonicity gives
$\lambda_{\mathrm{par}}(T)>0$. A negative vector forces $\lambda_{\mathrm{par}}<0$, and conversely an
attained negative bottom spans a negative line. The last two statements follow from
$\laminfP=\min(\lambda_{\mathrm{ev}},\lambda_{\mathrm{odd}})$, the sets $\{\lambda_{\mathrm{par}}>0\}$ being
intervals.
\end{proof}
The same argument, with the pole term of $\Ainf$ in place of the two multipliers, shows that $\laminf$
and the bottom of $\Gamma_T$ on $L^2(\IT)$ are strictly decreasing as well (attainment as above, the pole
term being a bounded form of rank two); this is not used below.

\begin{corollary}[the archimedean threshold]\label{cor:Tarch}
\[
0.58\ \le\ T^*_{\infty,\mathrm P}\ \le\ 0.60 .
\]
The bottom $\laminfP$ is strictly decreasing, positive on $(0,T^*_{\infty,\mathrm P})$ and negative on
$(T^*_{\infty,\mathrm P},\infty)$, and $\laminfP(T)<0$ for every $T\ge0.60$. In particular
$\laminfP(\frac12\log3)>0$: $\Ainf\ge\laminfP(\frac12\log3)\,\|F\|^2$ with a positive constant, i.e.\
the archimedean form is positive definite on $\mathcal P_{\frac12\log3}$ by itself. The proof
includes a computer verification in interval arithmetic (Appendix~\ref{app:encl}).
\end{corollary}
\begin{proof}
Theorem~\ref{thm:index}(iii) of \S\ref{subsec:negindex} gives
$T^*_{\infty,\mathrm P}=T^*_{\mathrm{ev}}\in[0.58,0.60]$, and Theorem~\ref{thm:index}(ii) gives
$n_{\mathrm{ev}}(T)\ge1$, hence $\laminfP(T)\le\lambda_{\mathrm{ev}}(T)<0$, for $T\ge0.60$
(Lemma~\ref{lem:strict}); the rest is Lemma~\ref{lem:strict}, with $\frac12\log3=0.5493<0.58$.
\end{proof}
Corollary~\ref{cor:Tarch} proves a sign, not a value; a numerical value of
$\laminfP(\frac12\log3)$ is in Appendix~\ref{app:diag}, and a lower bound on $\laminfP(0.58)$ would
require rerunning the upper certificate of Appendix~\ref{app:encl} with the minorant lowered by the
target. Its lower half is a proved lower bound on an infimum, which a Galerkin value cannot give. It
is obtained not from a Temple or Kato inclusion but from an explicit minorant of the symbol, an exact
finite section and a Schur complement over the tail controlled by the trace of Theorem~\ref{thm:P},
an intermediate problem in the classical sense. It improves the bound $0.74921$ of
Proposition~\ref{prop:arch}(iv).

On all of $L^2(\IT)$ the pole term $2A(F)B(F)$ is present, and one explicit odd test function makes
the archimedean form negative well before $T^*_{\infty,\mathrm P}$, improving the bound
$T^*_\infty\le T^*_{\infty,\mathrm P}\le0.60$ of Proposition~\ref{prop:arch}(iv) and
Corollary~\ref{cor:Tarch}.

\begin{proposition}[the archimedean complete form is negative from $T=0.4608$]\label{prop:archneg}
For $T>0$ let $F^{\mathrm{odd}}_T(t)=t$ on $\IT$ and put
\[
U(T):=\frac{2T^3}3\Big(C_0+\log\coth\frac T2+2\arctan e^{-T}\Big)+\frac{2T^3(3T+7)}9
-2\Big(4T\cosh\frac T2-8\sinh\frac T2\Big)^2 .
\]
Then $\Ainf(F^{\mathrm{odd}}_T)\le U(T)$, and $U(0.4608)\in[-7.0159575,-7.0159574]\cdot10^{-6}$. Hence
$\laminf(T)<0$ for every $T\ge0.4608$, and
\[
T^*_\infty\ \le\ 0.4608\ <\ 0.58\ \le\ T^*_{\infty,\mathrm P} .
\]
The single numerical evaluation in the proof is done in interval arithmetic.
\end{proposition}
\begin{proof}
Write $F=F^{\mathrm{odd}}_T$. It is real, odd and bounded, and $|\Fh(\tau)|=O(1/|\tau|)$, so it lies in the
form domain. By the parity paragraph after Proposition~\ref{prop:nodes}, $A(F)=-B(F)=s:=\int F\sinh\frac t2$,
so the pole term is $2A(F)B(F)=-2s^2$, and $s=2\int_0^Tt\sinh\frac t2\,dt=4T\cosh\frac T2-8\sinh\frac T2$.
For $0\le u\le2T$, $\cF(u)=\int_{u-T}^Tt(t-u)\,dt=(4T^3-6T^2u+u^3)/6$, so $\|F\|^2=\cF(0)=2T^3/3$ and
$\cF(0)-\cF(u)=u(6T^2-u^2)/6\ge0$; for $u\ge2T$, $\cF(u)=0$. On $u>0$ one has $K(u)\le\frac1u+\frac12$,
since $\sinh x\ge x$ and $\cosh x\ge1$. Against the nonnegative function $\cF(0)-\cF(u)$ this gives
\begin{align*}
\Lev(\cF)&\le\int_0^{2T}\frac{u(6T^2-u^2)}6\Big(\frac1u+\frac12\Big)du+\|F\|^2\int_{2T}^\infty K\\
&=\frac{2T^3(3T+7)}9+\frac{2T^3}3\Big(\log\coth\frac T2+2\arctan e^{-T}\Big),
\end{align*}
by Theorem~\ref{thm:S}(iii) for the tail, with $\pi-2\arctan e^T=2\arctan e^{-T}$.
Adding $C_0\|F\|^2$ and $-2s^2$ gives $\Ainf(F)\le U(T)$. The value of $U(0.4608)$ is obtained by an
outward-rounded interval evaluation at $40$ digits, with $C_0=-\gamma-\frac\pi2-3\log2-\log\pi$ and
$2\arctan e^{-T}$ enclosed by its alternating series with the remainder term; the expressions for
$\cF$, the integral over the window and $s$ are also checked symbolically. So
$\laminf(0.4608)\le\Ainf(F)/\|F\|^2<0$. The rest is Proposition~\ref{prop:arch}(i), $\laminf$ being
nonincreasing, together with Corollary~\ref{cor:Tarch}.
\end{proof}
The same evaluation gives $U(0.461)\in[-3.59383,-3.59382]\cdot10^{-5}$, $U(0.47)=-1.4048\cdot10^{-3}$
and $U(\frac12\log3)=-0.0200526$. Proposition~\ref{prop:archneg} separates the two thresholds
strictly. It is the empty-comb case of Corollary~\ref{cor:convfail}(i) below: on $[0.4608,0.58]$ the
archimedean form is positive definite on $\PT$ but not on $L^2(\IT)$. Galerkin computations give
numerical upper bounds on both thresholds, and suggest that at $T=\frac12\log3$ the primes lower the
bottom of the form on $\PT$ by two orders of magnitude while at $T=\log2$ they raise it on the
negative direction (Appendix~\ref{app:diag}); only the proved bounds are used in this paper.

The proved statements have the following consequence. At $T=\frac12\log3$ the archimedean form is
positive definite on $\PT$ by itself (Corollary~\ref{cor:Tarch}). For every $T\ge0.60$, in particular
at $T=\log2$, one has $\laminfP<0$ (Corollary~\ref{cor:Tarch}; for $T>0.74921$ already by
Proposition~\ref{prop:arch}(ii)), so every proof of $\lamA\ge0$ has to exhibit
$\PiT(F)\ge-\Ainf(F)>0$ on the vectors on which the archimedean block is negative; at $\log2$, by
Theorem~\ref{thm:index}, this set contains a line but no two-dimensional subspace. The threshold says
nothing about the size of $\PiT$ on that direction. A bound that treats the two blocks separately
cannot see this: $\PiT$ is bounded with
\[
\|\PiT\|\ \le\ \Aconst
\]
(each term is $-v_n$ times a compressed translation, and translations are unitary), so the triangle
bound $\lamA\ge\laminf-\|\PiT\|$ is far from the proved lower bound $8.86\cdot10^{-5}$ at
$T=\frac12\log3$ and vacuous at $\log2$, where the bottom of the form is not known even numerically
(Remark~\ref{rem:galerkin}; numerical values in Appendix~\ref{app:diag}).

\subsection{The negative index}\label{subsec:negindex}

The following bound on how much of the archimedean form can be negative constrains
Problem~\ref{prob:A} and any argument that works mode by mode.

\begin{proposition}[negative index of the archimedean form]\label{prop:NC}
Let $n_-(\Ainf|\PT)$ be the number of negative eigenvalues of $\Ainf$ on $\PT$, and for $\tau_1>0$ let
$S_{T,\tau_1}$ be the time--band limiting operator on $L^2(\IT)$, $\langle
S_{T,\tau_1}F,F\rangle=\frac1{2\pi}\int_{|\tau|\le\tau_1}|\Fh|^2d\tau$ for $\|F\|=1$. Let
$\tau_0=6.28983599\ldots$ be the zero of $m$ and, for $\tau_1>\tau_0$,
$\theta(\tau_1)=m(\tau_1)/(m(\tau_1)-C_0)\in(0,1)$. Then
\[
n_-(\Ainf|\PT)\ \le\ \#\{n:\lambda_n(S_{T,\tau_1})>\theta(\tau_1)\}\ \le\
\Big\lfloor\frac{2T\tau_1}{\pi\,\theta(\tau_1)}\Big\rfloor
\qquad\text{for every }\tau_1>\tau_0,
\]
and the right-hand side is minimized at $\tau_1=14.8815$, where it gives $n_-\le68.51\,T$: at most
$37$ at $T=\frac12\log3$, $47$ at $T=\log2$ and $54$ at $T=0.80$. The first inequality is far
stronger than the second; Corollary~\ref{cor:P3} below evaluates it.
\end{proposition}
\begin{proof}
For $\|F\|=1$ put $b=\langle S_{T,\tau_1}F,F\rangle$. Splitting the archimedean integral at $\tau_1$
and using $m$ increasing, $\Ainf(F)\ge C_0b+m(\tau_1)(1-b)$ on $\PT$, where the pole term vanishes.
Hence $\Ainf(F)<0$ forces $b>\theta(\tau_1)$. If $\Ainf$ were negative on an $n$-dimensional subspace
$V$, then every unit $F\in V$ would have $\langle S_{T,\tau_1}F,F\rangle>\theta$, so by min--max
$\lambda_{n-1}(S_{T,\tau_1})>\theta$, giving the first inequality. The second is Markov's inequality
with $\operatorname{tr}S_{T,\tau_1}=2T\tau_1/\pi$. Restricting from $L^2(\IT)$ to $\PT$ only decreases
the count.
\end{proof}
Proposition~\ref{prop:NC} is explicit in $T$. For $T\le0.90$ the count itself can be computed, sector
by sector, and it is $0$, $1$ or $2$ (Theorem~\ref{thm:index}).

\begin{lemma}[parity splitting]\label{lem:parity}
$\PT=\PT^{\mathrm{ev}}\oplus\PT^{\mathrm{odd}}$, the sum being orthogonal both in
$L^2(\IT)$ and for $\Gamma_T$, which equals $\Ainf$ on $\PT$. Hence
\[
n_-(\Ainf|\PT)=n_{\mathrm{ev}}(T)+n_{\mathrm{odd}}(T),\qquad
\laminfP(T)=\min\big(\lambda_{\mathrm{ev}}(T),\lambda_{\mathrm{odd}}(T)\big).
\]
Here, as in Proposition~\ref{prop:NC}, $n_-$ counts negative eigenvalues with multiplicity; for the closed,
lower-bounded form $\Gamma_T$, which has compact resolvent on each of these closed subspaces (proof of
Proposition~\ref{prop:nodes}(i)), this is the largest dimension of a subspace of the form domain on
which the form is negative definite.
\end{lemma}
\begin{proof}
$A(RF)=B(F)$ and $B(RF)=A(F)$, so for $F\in\PT$ both $F_{\mathrm e}=\frac12(F+RF)$ and
$F_{\mathrm o}=\frac12(F-RF)$ have vanishing Tate moments: $\PT$ is $R$-invariant and splits. The
parts are $L^2$-orthogonal, and
$\Gamma_T(F_{\mathrm e},F_{\mathrm o})=\frac1{2\pi}\int m\,\widehat{F_{\mathrm e}}\,\overline{\widehat{F_{\mathrm o}}}=0$,
because $m$ is even, $\widehat{F_{\mathrm e}}$ is even and $\widehat{F_{\mathrm o}}$ is odd. Negative
indices add over sums that are orthogonal for the norm and for the form, and bottoms take the minimum.
\end{proof}
The definition of $\Ainf$ writes the pole term as $2A(F)B(F)$, i.e.\ on real $F$, while $L^2(\IT)$ is
complex. Nothing depends on the choice: for $F=u+iv$ with $u,v$ real, $\Gamma_T(F)=\Gamma_T(u)+\Gamma_T(v)$
and $F\in\PT$ if and only if $u,v\in\PT$ (as in the proof of Lemma~\ref{lem:strict}), so the counts and
the bottoms are the same over $\R$ and over $\C$.

\begin{theorem}[the negative index of the archimedean form, $0<T\le0.90$]\label{thm:index}
In the notation of \S\ref{subsec:arch}:
\emph{(i)} $n_{\mathrm{ev}}(T)\le1$ and $n_{\mathrm{odd}}(T)\le1$ for $0<T\le0.90$;
$n_{\mathrm{ev}}(T)=n_{\mathrm{odd}}(T)=0$ and $\laminfP(T)>0$ for $T\le0.58$; $n_{\mathrm{odd}}(T)=0$
and $\lambda_{\mathrm{odd}}(T)>0$ for $T\le\log2$.
\emph{(ii)} $n_{\mathrm{ev}}(T)\ge1$ for $T\ge0.60$, and $n_{\mathrm{odd}}(T)\ge1$ for $T\ge0.81$.
\emph{(iii)} $0.58\le T^*_{\mathrm{ev}}\le0.60$, $\log2\le T^*_{\mathrm{odd}}\le0.81$,
$T^*_{\infty,\mathrm P}=T^*_{\mathrm{ev}}$, and
\[
n_-(\Ainf|\PT)=\begin{cases}0,&0<T<T^*_{\mathrm{ev}},\\ 1,&T^*_{\mathrm{ev}}<T<T^*_{\mathrm{odd}},\\
2,&T^*_{\mathrm{odd}}<T\le0.90.\end{cases}
\]
\emph{(iv)} At the windows: $n_-(\Ainf|\PT)=0$ at $T=\frac12\log3$ and $T=0.58$; $n_-(\Ainf|\PT)=1$ at
$T=0.60$ and $T=\log2$, the negative direction being even; $n_-(\Ainf|\PT)\in\{1,2\}$ at $T=0.80$,
with $n_{\mathrm{ev}}(0.80)=1$; and $n_-(\Ainf|\PT)=2$ at $T=0.81$, $0.82$ and $0.90$.

\noindent The proof includes a computer verification in interval arithmetic (Appendix~\ref{app:encl}).
\end{theorem}
\begin{proof}
(i) The upper certificates of Proposition~\ref{prop:haynsworth} (Appendix~\ref{app:encl}) give
$n_{\mathrm{par}}(T)\le r$ with $r=0$ at $T=0.58$ in both sectors and at $T=\log2$ in the odd sector,
and $r=1$ at $T=0.90$ in both sectors. Lemma~\ref{lem:zeroext} carries each bound to every smaller $T$,
and Lemma~\ref{lem:rzero} turns $r=0$ into $\lambda_{\mathrm{par}}>0$, at the certified window and
therefore, by monotonicity, below it. (ii) The explicit negative planes of
Proposition~\ref{prop:negplane} at $T=0.60$ (even) and $T=0.81$ (odd) give $n_{\mathrm{par}}\ge1$
there, and Lemma~\ref{lem:zeroext} carries each to every larger $T$, with no range condition at the
larger window. (iii) By (i), (ii) and Lemma~\ref{lem:strict}: $\lambda_{\mathrm{ev}}(0.58)>0$ gives
$T^*_{\mathrm{ev}}\ge0.58$, $\lambda_{\mathrm{ev}}(0.60)<0$ gives $T^*_{\mathrm{ev}}\le0.60$,
$\lambda_{\mathrm{odd}}(\log2)>0$ gives $T^*_{\mathrm{odd}}\ge\log2$, $\lambda_{\mathrm{odd}}(0.81)<0$
gives $T^*_{\mathrm{odd}}\le0.81$, and $T^*_{\infty,\mathrm P}=\min(T^*_{\mathrm{ev}},T^*_{\mathrm{odd}})
=T^*_{\mathrm{ev}}$ because $0.60<\log2$. By Lemma~\ref{lem:strict}, off the threshold points
$n_{\mathrm{par}}(T)\ge1$ exactly when $T>T^*_{\mathrm{par}}$, and $n_{\mathrm{par}}\le1$ on $(0,0.90]$
by (i); add over the sectors (Lemma~\ref{lem:parity}). (iv) is (i)--(iii) at the windows:
$\frac12\log3=0.5493<0.58$; at $0.60$ and $\log2$ the even count is $1$ (at least $1$ from $0.60$, at
most $1$ from $0.90$) and the odd count is $0$ (from $\log2$); at $0.80$ the even count is $1$ and the
odd count is at most $1$; at $0.81$, $0.82$, $0.90$ both counts are $1$.
\end{proof}
The certificate data are described in \S\ref{sec:repro}. The values $n_{\mathrm{par}}(T^*_{\mathrm{par}})$
at the two thresholds are not decided, and nothing depends on them; nor is it decided which of the two
values $n_-$ takes at $T=0.80$ (Appendix~\ref{app:diag}). Theorem~\ref{thm:index} concerns the
archimedean form with the prime comb deleted; it says nothing about the sign of $\mathfrak a_T$ or
$W_T$.

We return to bounds that are explicit for every $T$. The trace bound is crude. The same operator has
an explicit diagonal in the Legendre basis, and Ky Fan's maximum principle turns that diagonal into
an eigenvalue ceiling. The following is elementary and self-contained.

\begin{theorem}[a Ky Fan ceiling for the time--band limiting operator]\label{thm:P}
Let $S_c=Q_TB_sQ_T$ on $L^2(\IT)$, where $Q_T$ is restriction to $\IT$ and $B_s$ the band projection
to $|\tau|<s$, and put $c=Ts$. Let $j_n$ be the spherical Bessel functions and
$G_k(x)=\sum_{j<k}(2j+1)j_j(x)^2\in[0,1]$. Then, in the orthonormal Legendre basis
$e_j(t)=\sqrt{(2j+1)/2T}\,P_j(t/T)$,
\[
\langle S_ce_j,e_j\rangle=\frac{2(2j+1)}\pi\int_0^cj_j(x)^2dx,\qquad
\sum_{j\ge0}\langle S_ce_j,e_j\rangle=\frac{2c}\pi=\operatorname{tr}S_c ,
\]
and for every $k\ge0$
\[
\lambda_k(S_c)\ \le\ \Lambda_k(c):=\frac2\pi\int_0^c\big(1-G_k(x)\big)\,dx .
\]
In particular $\Lambda_0(c)=2c/\pi$ is the trace and $\Lambda_k(c)=O(c^{2k+1})$ as $c\to0$.
\end{theorem}
\begin{proof}
From $\int_{-1}^1P_j(u)e^{-ixu}du=2(-i)^jj_j(x)$ one gets
$\widehat{e_j}(\tau)=\sqrt{(2j+1)/2T}\,\cdot2T(-i)^jj_j(\tau T)$, so
$|\widehat{e_j}(\tau)|^2=2T(2j+1)j_j(\tau T)^2$ and Plancherel gives the displayed diagonal entry.
The addition theorem $\sum_{j\ge0}(2j+1)j_j(x)^2=1$ gives the trace identity, so the basis loses no
trace. By the Ky Fan maximum principle $\sum_{j<k}\lambda_j(S_c)\ge\sum_{j<k}\langle
S_ce_j,e_j\rangle$; since $S_c=(B_sQ_T)^*(B_sQ_T)\ge0$ its eigenvalues are nonnegative and
non-increasing, whence
$\lambda_k\le\sum_{n\ge k}\lambda_n=\operatorname{tr}S_c-\sum_{j<k}\lambda_j\le\sum_{j\ge k}\langle
S_ce_j,e_j\rangle=\Lambda_k(c)$.
\end{proof}

\begin{corollary}[quadrature-free form]\label{cor:P1}
$|j_n(x)|\le x^n/(2n+1)!!$ for $x\ge0$, hence
$\lambda_k(S_c)\le\Lambda^*_k(c):=\frac2\pi\sum_{j\ge k}c^{2j+1}/\big((2j+1)!!\big)^2$, a convergent
series of rationals in $c$: the bound is evaluated by arithmetic alone, with no quadrature.
\end{corollary}

\begin{corollary}[explicit negative-index counts]\label{cor:P3}
For every $\tau_1>\tau_0$, $n_-(\Ainf|\PT)\le n_-(\Gamma_T\ \text{on}\ L^2(\IT))\le\min\{k:\
\Lambda_k(T\tau_1)\le\theta(\tau_1)\}$, and the same with $\Lambda^*_k$ in place of $\Lambda_k$.
Minimizing over $\tau_1$ with the quadrature-free $\Lambda^*$ of Corollary~\ref{cor:P1} gives
\[
n_-(\Ainf|\PT)\ \le\ 5\ \ \big(T=\tfrac12\log3\big),\qquad 7\ \ (T=\log2),\qquad 8\ \ (T=0.80),
\]
against $37$, $47$ and $54$ from the trace bound of Proposition~\ref{prop:NC}.
\end{corollary}
\begin{proof}
If $\Lambda_k(T\tau_1)\le\theta(\tau_1)$ then $\lambda_k(S_{T,\tau_1})\le\theta(\tau_1)$ by
Theorem~\ref{thm:P}, so $\#\{n:\lambda_n>\theta\}\le k$, and the first inequality of
Proposition~\ref{prop:NC} applies for that $\tau_1$; each $\tau_1>\tau_0$ gives a valid bound, so the
minimum over $\tau_1$ does. The three values come from $\tau_1=6.6159$, $6.3576$, $6.3576$, where the
partial sums of $\Lambda^*$ fall below $\theta$; the omitted tails are below $10^{-170}$.
\end{proof}

\begin{remark}[the limit of this method]\label{rem:P4}
Numerically, at $T=\log2$ and the optimal $\tau_1$ the exact count
$\#\{n:\lambda_n(S_{T,\tau_1})>\theta(\tau_1)\}$ is already $5$, and a scan over $\tau_1$ does not
lower it (Appendix~\ref{app:diag}). So no improvement of the prolate estimate can bring the bound of
Proposition~\ref{prop:NC} below $5$ there: the distance to the true value $n_-(\Ainf|\PT)=1$
(Theorem~\ref{thm:index}) is due to the two-value relaxation $m\ge C_0\mathbf 1_{|\tau|<\tau_1}+
m(\tau_1)\mathbf 1_{|\tau|\ge\tau_1}$ and to $\operatorname{codim}\PT=2$, not to prolate theory. A
multi-band relaxation $m\ge\sum_ic_i\mathbf 1_{B_i}$ produces a different operator
$\sum_i(c_i-c_{i+1})S_{T,\tau_i}$ and a different counting problem, and in the limit of infinitely many
bands it is exact. It returns the exact count $1$ at $\log2$: with $28$ bands, the multi-band minorant
handled exactly on a $22$-mode Legendre section of each parity sector and Theorem~\ref{thm:P} used only
through the trace of the far tail, where it is sharp, rather than by counting eigenvalues of that operator
against a threshold, the certificates behind Theorem~\ref{thm:index} give $n_-(\Ainf|\mathcal P_{\log2})=1$
(Appendix~\ref{app:encl}).
\end{remark}

In the other direction the count admits a lower bound, from the same two ingredients as
Corollary~\ref{cor:mono}: Dirichlet modes and the concavity of $\varphi$.

\begin{theorem}[lower bound on the negative index]\label{thm:J}
Let $\Gamma_T$ be the form of Proposition~\ref{prop:nodes}, and let $\tau_0=6.28983599\ldots$ be the zero of $m$. Then
\[
n_-(\Gamma_T)\ \ge\ \#\{k\ge1:\ k\pi/2T<\tau_0\}=\Big\lceil\frac{2T\tau_0}{\pi}\Big\rceil-1
\quad\text{on }L^2(\IT),\qquad
n_-(\Ainf|\PT)\ \ge\ \Big\lceil\frac{2T\tau_0}{\pi}\Big\rceil-3 .
\]
\end{theorem}
\begin{proof}
Let $D_k\subset H^1_0(\IT)$ be the span of the first $k$ Dirichlet modes $\sin(j\pi(t+T)/2T)$,
$j\le k$. For $F\in D_k$, extended by zero, Plancherel gives
$\frac1{2\pi}\int\tau^2|\Fh|^2=\|F'\|^2\le(k\pi/2T)^2\|F\|^2$. By Theorem~\ref{thm:S}(i),
$m(\tau)-C_0=\varphi(\tau^2)$ with $\varphi$ complete Bernstein, hence concave and strictly
increasing; applying Jensen in the variable $\tau^2$ to the probability measure
$d\mu=|\Fh|^2d\tau/(2\pi\|F\|^2)$,
\[
\frac{\Gamma_T(F)}{\|F\|^2}=C_0+\int\varphi(\tau^2)\,d\mu\ \le\ C_0+\varphi\Big(\int\tau^2d\mu\Big)
=C_0+\varphi\big(\|F'\|^2/\|F\|^2\big)\ \le\ C_0+\varphi\big((k\pi/2T)^2\big).
\]
Since $m(\tau_0)=0$ means exactly $\varphi(\tau_0^2)=-C_0$, and $\varphi$ is strictly increasing,
$k\pi/2T<\tau_0$ forces $\Gamma_T<0$ on $D_k\setminus\{0\}$, i.e.\ $n_-(\Gamma_T)\ge k$; the
largest such $k$ is $\lceil2T\tau_0/\pi\rceil-1$. On $\PT$ the pole term vanishes and
$\Ainf=\Gamma_T$, and $\PT$ has codimension $2$, which costs at most $2$ in the count.
\end{proof}
The bound on $\PT$ is vacuous at the two windows $T=\frac12\log3$ and $T=\log2$: $2T\tau_0/\pi$
equals $2.1996$ and $2.7755$ there, so it reads $n_-(\Ainf|\PT)\ge0$ at both; on $L^2(\IT)$ it reads
$n_-(\Gamma_T)\ge2$ at both. Together with the min--max upper bound (the first inequality of
Proposition~\ref{prop:NC}, applied with $\tau_1$ decreasing to $\tau_0$ and combined with the
Landau--Widom asymptotics of the prolate spectrum of $S_{T,\tau_1}$), the two sides match:
\[
n_-(\Ainf|\PT)=\frac{2\tau_0}{\pi}\,T+O(\log T)=4.004\,T+O(\log T).
\]
Theorem~\ref{thm:J} and Corollary~\ref{cor:P3} together give two-sided brackets on $L^2(\IT)$:
$n_-(\Gamma_T)\in[2,5]$ at $T=\frac12\log3$, $[2,7]$ at $T=\log2$ and $[3,8]$ at $T=0.80$. On $\PT$
only the upper half is informative, $0\le n_-(\Ainf|\PT)\le5,7,8$ respectively. Theorem~\ref{thm:index}
replaces these three brackets on $\PT$ by the exact values $0$ at $\frac12\log3$ and $1$ at $\log2$,
and by $\{1,2\}$ at $0.80$, where Theorem~\ref{thm:J} supplies the lower value $1$ independently.

For $T_2>T^*_{\infty,\mathrm P}$ write $\mathrm H(T_2)$ for the statement that
$n_-(\Ainf|\PT)=1$ for every $T\in(T^*_{\infty,\mathrm P},T_2]$.

\begin{corollary}[one negative direction up to $T^*_{\mathrm{odd}}$]\label{cor:H1}
$\mathrm H(T_2)$ holds for every
$T_2\in(T^*_{\infty,\mathrm P},T^*_{\mathrm{odd}})$ and for every $T_2\in(T^*_{\infty,\mathrm P},\log2]$, and
fails for every $T_2>T^*_{\mathrm{odd}}$; the case $T_2=T^*_{\mathrm{odd}}>\log2$ is not decided. Hence the supremum of the admissible $T_2$ is
$T^*_{\mathrm{odd}}\in[\log2,0.81]$.
\end{corollary}
\begin{proof}
By Theorem~\ref{thm:index}(iii), $T^*_{\infty,\mathrm P}=T^*_{\mathrm{ev}}$, $n_-(\Ainf|\PT)=1$ on
$(T^*_{\mathrm{ev}},T^*_{\mathrm{odd}})$ and $n_-(\Ainf|\PT)=2$ on $(T^*_{\mathrm{odd}},0.90]$, which is
nonempty since $T^*_{\mathrm{odd}}\le0.81$. For $T_2<T^*_{\mathrm{odd}}$ the interval
$(T^*_{\mathrm{ev}},T_2]$ lies in the first range; for $T_2>T^*_{\mathrm{odd}}$ the interval
$(T^*_{\mathrm{ev}},T_2]$ meets the second. For $T_2\le\log2$ use instead Theorem~\ref{thm:index}(i):
$n_{\mathrm{odd}}=0$ on $(0,\log2]$, and $n_{\mathrm{ev}}=1$ on $(T^*_{\mathrm{ev}},0.90]$ by
Lemma~\ref{lem:strict} and (i).
\end{proof}

\begin{remark}
Theorem~\ref{thm:J} alone limits how far $\mathrm H(T_2)$ can reach:
$\lceil2T\tau_0/\pi\rceil-3\ge2$ as soon as $2T\tau_0/\pi>4$, so any admissible $T_2$ satisfies
$T_2\le2\pi/\tau_0=0.9989426\ldots$ Corollary~\ref{cor:H1} replaces that bound by
$\sup T_2=T^*_{\mathrm{odd}}\le0.81$: the second negative direction is odd and appears by $T=0.81$, where
Theorem~\ref{thm:J} still gives only $n_-(\Ainf|\PT)\ge1$. The window $\log2$ lies inside the admissible
range, and $\frac12\log3$ lies below $T^*_{\infty,\mathrm P}$, where $n_-=0$.
\end{remark}

\begin{remark}
Theorem~\ref{thm:J} and Corollary~\ref{cor:mono} rest on one fact: $\varphi$ is
concave, so Jensen applied to $|\Fh|^2$ always produces an upper bound for the archimedean form
and never a lower bound. A lower bound would require a lower bound on $\int\min(\tau^2,X)d\mu$, that
is, control of the spread of $|\Fh|^2$: a quantitative statement that a function supported in $\IT$
cannot concentrate its spectrum in $|\tau|<s$, which is an explicit non-asymptotic prolate
concentration estimate. This is why the moment arguments of this section bound $n_-$ from above
by tens and not by one.
\end{remark}

A Birman--Schwinger count is not available either. Writing $m_\pm=\max(\pm m,0)$ and
$M_\pm=R_Tm_\pm(D)R_T\ge0$ on $L^2(\IT)$, $R_T$ the restriction to $\IT$ of the extension by zero,
so that $\Ainf=M_+-M_-$ on $\PT$, the negative part is trace class with
\[
\operatorname{tr}M_-=\frac{2T}{\pi}\int_0^{\tau_0}\big(-m(\tau)\big)\,d\tau=4.495774\,T
\]
(since $\operatorname{tr}R_Tf(D)R_T=\frac{2T}{2\pi}\int_\R f$, and $m_-$ is even and supported in
$|\tau|<\tau_0$), while $M_+>0$ strictly (if $\langle M_+F,F\rangle=0$ then $\Fh$ vanishes off
$|\tau|<\tau_0$, and $\Fh$ is entire of exponential type $T$, so $F=0$) but without a spectral gap,
since $m_+$ vanishes identically on $|\tau|<\tau_0$ and near-band-limited quasi-modes make
$\langle M_+F,F\rangle$ arbitrarily small. A Birman--Schwinger count
$n_-(M_+-M_-)\le\operatorname{tr}(M_+^{-1}M_-)$ is therefore unavailable, and so is a
Cwikel--Lieb--Rozenblum bound, which fails in dimension one and in any case requires an elliptic
free part.

\subsection{Competing even modes}\label{subsec:arena}

No fixed finite number of archimedean modes captures the obstruction, and this is a theorem with an
explicit count. The archimedean gap along the even tower grows only logarithmically, while the prime
form is bounded by $\Aconst$ and does not decay; so the number $N^*(T)$ of even archimedean modes
whose level the prime form can reach, the competing even modes with which any truncation has to
work, is large, and grows doubly exponentially in the window.

\begin{theorem}[log-slow archimedean gaps]\label{thm:arena}
Let $\PT^{\mathrm{ev}}$ be the even part of $\PT$ and let
$\lambda^{\mathrm{ev}}_1\le\lambda^{\mathrm{ev}}_2\le\cdots$ be the eigenvalues of $\Ainf$ on it.
Then, for every $k\ge0$,
\[
\lambda^{\mathrm{ev}}_{k+1}\ \le\ C_0+\varphi\Big(\big(\tfrac{(2k+3)\pi}{2T}\big)^2\Big)
=m\Big(\frac{(2k+3)\pi}{2T}\Big)=\log k+O(1).
\]
\end{theorem}
\begin{proof}
$\Ainf$ commutes with parity, and on even functions the two Tate conditions coincide: both read
$\int F(t)\cosh\frac t2\,dt=0$. Hence $\PT^{\mathrm{ev}}$ has codimension $1$ in
$L^2_{\mathrm{ev}}(\IT)$, and the pole term vanishes on it, so $\Ainf=\Gamma_T$ there
(Proposition~\ref{prop:nodes}). Take the $k+2$ even Dirichlet modes $\cos((2j+1)\pi t/2T)$,
$j=0,\dots,k+1$, which vanish at $\pm T$ and satisfy
$\|F'\|^2\le((2k+3)\pi/2T)^2\|F\|^2$ on their span; intersecting that span with the single linear
constraint leaves a subspace of $\PT^{\mathrm{ev}}$ of dimension at least $k+1$. On it, exactly as
in the proof of Theorem~\ref{thm:J}, Jensen applied in the variable $\tau^2$ to the concave
$\varphi$ (Theorem~\ref{thm:S}(i)) gives
$\Gamma_T(F)/\|F\|^2\le C_0+\varphi(\|F'\|^2/\|F\|^2)\le C_0+\varphi(((2k+3)\pi/2T)^2)$, and
min--max bounds $\lambda^{\mathrm{ev}}_{k+1}$ by that number. The last equality is
$C_0+\varphi(\tau^2)=m(\tau)$ together with Lemma~\ref{lem:mbounds}(ii).
\end{proof}

\begin{corollary}[the number of competing even modes]\label{cor:arena}
Put $\tau_k=(2k+3)\pi/2T$ and let
$N^*(T)=\#\{j\ge1:\lambda^{\mathrm{ev}}_j\le\Aconst\}$ be the number of even archimedean modes
whose level the prime form, of norm at most $\Aconst$, can reach. Then
\[
N^*(T)\ \ge\ \#\{k\ge0:m(\tau_k)\le\Aconst\}\ =\ \Big\lfloor\frac{T\tau_*}{\pi}-\frac12\Big\rfloor,
\qquad \tau_*=m^{-1}(\Aconst).
\]
At $T=\log2$, where $\Aconst=2.24883$, Lemma~\ref{lem:mbounds}(ii) at the exact frequencies gives
$m(\tau_{11})\le2.2168\le\Aconst$ and $m(\tau_{12})\ge2.2596>\Aconst$, so $N^*\ge12$ and the
bound yields no more. At $T=\frac12\log3$, where $\Aconst=0.98026$, it gives
$m(\tau_1)\le0.8928\le\Aconst$ and $m(\tau_2)\ge1.1075>\Aconst$, so $N^*\ge2$. In general
$\Aconst=\sum_{n<e^{2T}}2\Lambda(n)n^{-1/2}=4e^{T}(1+o(1))$ by the prime number theorem and
partial summation, and $\tau_*=2\pi e^{\Aconst}(1+o(1))$ by Lemma~\ref{lem:mbounds}(ii), so
\[
N^*(T)\ \ge\ 2T\,e^{\Aconst}\,(1+o(1))\ =\ \exp\big(4e^{T}(1+o(1))\big):
\]
the number of competing even modes grows doubly exponentially in $T$.
\end{corollary}
\begin{proof}
By Theorem~\ref{thm:arena}, $\lambda^{\mathrm{ev}}_{k+1}\le m(\tau_k)$, so every $k$ with
$m(\tau_k)\le\Aconst$ contributes a mode. Since $m$ is increasing, $m(\tau_k)\le\Aconst$ is
$\tau_k\le\tau_*$, i.e.\ $k\le T\tau_*/\pi-\frac32$, and the number of integers $k\ge0$ with that
property is $\lfloor T\tau_*/\pi-\frac32\rfloor+1=\lfloor T\tau_*/\pi-\frac12\rfloor$. The two
window evaluations are Lemma~\ref{lem:mbounds}(ii) at $\tau_{11},\tau_{12}$ and at
$\tau_1,\tau_2$ respectively, with $\tau_k=(2k+3)\pi/2T$. The asymptotics are
$\log\Aconst$-free: $m(\tau)=\log(\tau/2\pi)+O(1/\tau)$ inverts to
$\tau_*=2\pi e^{\Aconst}(1+o(1))$, and $T\tau_*/\pi=2Te^{\Aconst}(1+o(1))$.
\end{proof}

Consequently no fixed finite truncation in the archimedean basis can decide the sign of the even
Schur complement: any argument that controls the prime form by the archimedean gap must carry at
least $N^*(T)$ even modes, because up to that index the gap, growing like $\log k$, is smaller than
the norm of a perturbation that neither decays nor is compact (Proposition~\ref{prop:shifts}). The
certified $E_7$ of Theorem~\ref{thm:nest} lies strictly inside a set of at least $12$ competing even
modes at $T=\log2$; this explains Remark~\ref{rem:galerkin} and the location of the near-zero
Galerkin bottom at the last even mode of every truncation. Numerically the bound $N^*(\log2)\ge12$ is
attained (Appendix~\ref{app:diag}).

\begin{remark}[no total positivity]\label{rem:PF2}
The property of $K$ that produces every upper bound in this paper also rules out the classical
route to a lower bound. By Lemma~\ref{lem:K}, $K$ is completely monotone, hence log-convex on
$(0,\infty)$ (Cauchy--Schwarz applied to the representing measure), and $K(u)=u^{-1}+O(1)$ at the
origin; so $u\mapsto K(|u|)$ is log-convex on each half-line and log-concave on no interval. By
Schoenberg's characterization of the P\'olya frequency functions of order $2$ (a nonnegative
integrable $f$ is $\mathrm{PF}_2$ if and only if it is log-concave \cite{Schoenberg1951}), the L\'evy
density $K(|\cdot|)$ is not $\mathrm{PF}_2$, so the translation kernel $K(x-y)$ is not totally
positive of order $2$ and Karlin's variation-diminishing and oscillation theory
\cite{Karlin1968} does not apply to it. This does not prove that no such argument exists: the
operator is $C_0+\Lev$, whose natural totally positive object is the transition density of the
associated semigroup rather than $K$ itself, and we do not know whether that density is
$\mathrm{PF}_2$ for large time. Complete monotonicity thus plays two roles: it makes $\varphi$ a
complete Bernstein function, whence by Jensen the upper bounds of Corollary~\ref{cor:mono},
Proposition~\ref{prop:TK}, Theorem~\ref{thm:J} and Theorem~\ref{thm:arena}, and it removes the
total-positivity hypothesis under which nodal counting would have ordered the spectrum.
\end{remark}


\section{Decidability at a fixed window and the decay of the margin}\label{sec:path}

This section and the next three describe what separates the results of this paper from a proof of
positivity for every $T$.

\subsection{The reduction chain}
Combining \cite[Thm.~11.3]{Paper45}, Theorems~\ref{thm:D1}, \ref{thm:Zexh} and Corollary~\ref{cor:RH},
\begin{multline*}
\mathrm{RH}\iff\forall T:\ \lamW\ge0\iff\forall T:\ \lamP\ge0\iff\forall T\ \exists\ \text{finite certificate}\\
\iff\forall T\ \exists\,w_T\in\Wcl^{\mathrm{reg}}:\ \sum_\rho\wh_T(\tau-\gamma_\rho)\ge0\ \text{on }\R .
\end{multline*}
These are equivalences between statements quantified over all $T$. At a single window one only has
$\lamW\le\lamP\le\lamA$, together with the exact relation between $\lamW$ and $\lamA$ of
Theorem~\ref{thm:parityschur} in \S\ref{subsec:alignineq}; the per-window converse
$\mathrm{WP}_T\Rightarrow\lamW\ge0$ fails in general for comb weights (Corollary~\ref{cor:convfail}). The chain
has the following consequence.

\begin{proposition}[the sign at each window is decidable]\label{prop:decide}
Fix a computable $T$ with $\lamP\ne0$. Then the sign of $\lamP$ is decided by a procedure that
terminates. In the positive case it produces a finite certificate, hence $\lamA>0$; in the negative
case it produces $\Phi\in\Kp$ with $C_0+\Lev(\Phi)-\sum v_n\Phi(\log n)<0$, which decides $\lamP<0$
but, since $\lamP\le\lamA$, not the sign of $\lamA$. By Theorem~\ref{thm:O} the procedure in the
positive case never terminates \emph{at} the value $\lamP$.
\end{proposition}
\begin{proof}
Run two searches in alternation. (a) Enumerate rational data $(u_j,a_j,\beta)$ with $u_j\ge2T$ on a
lattice $u_j=j\cdot2T/M$ and rational targets $c>0$; for each, attempt the finite verification of
$f\ge c$ by Lemma~\ref{lem:F1} and Proposition~\ref{prop:TAY} with interval enclosures of the finitely
many constants involved. If $\lamP>0$, Theorem~\ref{thm:D1} gives atomic data with
$\inf f>\lamP/2$; the verification of a strict inequality by a cover with rational data is
robust under small rational perturbation of the data (the functions are Lipschitz in the data),
so some enumerated item is verified in finite time. (b) Enumerate finite trial families
$\Phi=c_F$ with $F$ a step function with rational values on a rational partition of $\IT$, and
evaluate $C_0+\Lev(\Phi)-\sum v_n\Phi(\log n)$ by enclosing the convergent integral of the
explicit completely monotone density $K$ against the piecewise-polynomial $1-\Phi$ and the finitely
many values $\Phi(\log n)$; if $\lamP<0$ such a $\Phi$ with negative value exists by density of step
functions in the form domain, and a strict inequality is again found in finite time. One of (a),
(b) terminates.
\end{proof}

Thus RH is equivalent to an infinite family of finite, decidable problems, one per window, and a
certificate settling any one of them has a known meaning (Theorem~\ref{thm:Z}). What is missing is
uniformity in $T$. The following sections make this precise: at each window the statement concerns
finitely many numbers, no classical property of the primes decides it
(Proposition~\ref{prop:uniform}, Corollary~\ref{cor:void}), and the margin to be certified tends to $0$
(Corollary~\ref{cor:margin2}).

\subsection{The relevant range of ordinates}\label{subsec:where}
By Proposition~\ref{prop:env} a pointwise argument certifies every certificate beyond
$\tau_{\mathrm{env}}(T)=2\pi e^{\Aconst+\|\nu\|_{\TV}+o(1)}$, $\Aconst\sim4e^T$, and no pointwise
argument lowers $2\pi e^{\Aconst}$; the periodic lemmas stop earlier for lattice certificates, but
what they certify beyond their range is that $\sT$ has become positive and stays so. The certificate
therefore has to do work only on $[0,2\pi e^{\Aconst}]$. Heuristically, three scales appear there.

The resolution scale $2\pi e^{2T}$: for $w$ with flat top $[-2T,2T]$ the main lobe of $\wh$ has width
$\approx\pi/T$, and the mean spacing of the zeros at height $\tau$ is $2\pi/\log(\tau/2\pi)$; the two
agree at $\tau\approx2\pi e^{2T}$. Below this height $\sum_\rho\wh(\tau-\gamma_\rho)$ depends on individual
zeros; above it the sum sees the density, $\sum_\gamma\wh(\tau-\gamma)\approx\log(\tau/2\pi)=m(\tau)$, with
fluctuations expressed by the explicit formula through the primes. The same scale appears as
$T^*=2\pi e^{2L}$ in Zhu's conjectural Landau--Widom law \cite[Conj.~12.1]{Zhu2026}
(see also Appendix~\ref{app:diag}).

The alignment scale: between $2\pi e^{2T}$ and $\tau_{\mathrm{env}}$ one has $\sT(\tau)<0$ if and only if
$\sum_{n\in\NT}v_n\cos(\tau\log n)>m(\tau)$. Such $\tau$ are ordinates at which the phases $\tau\log p$,
$p<e^{2T}$, are simultaneously close to $0$ modulo $2\pi$; by Kronecker's theorem they exist at every
level of alignment, and by Theorem~\ref{thm:Z} they are exactly the ordinates at which the zeros must be
locally denser than average by the amount by which the comb exceeds $\log(\tau/2\pi)$. The function
$\nuh$ of a certificate must be large there; the contact comb of \S\ref{sec:cert} is its visible trace.
Beyond the last sign change of $\sT$, which lies below $\tau_{\mathrm{env}}$ but is of the same order,
the certificate only has to keep $\nuh$ from spoiling the positivity of $\sT$; this is the role of the
tail lemmas.

Dependence on $T$: with $\NT=\{n\ge2:\log n<2T\}$ the mass $\Aconst=\sum_{n\in\NT}v_n$ is a
left-continuous step function of $T$, constant between consecutive points $\frac12\log n$ with
$\Lambda(n)>0$, with jump $v_n=2\Lambda(n)/\sqrt n$ at $T=\frac12\log n$. Hence the multiplier-free
threshold $2\pi e^{\Aconst}$ of Proposition~\ref{prop:env}(iii) (the threshold for a certificate carries
the further factor $e^{\|\nu\|_{\TV}}$) is multiplied by $e^{v_n}$ when $n$ enters: by $2$ at $n=4$, $4.22$
at $5$, $4.35$ at $7$, $1.63$ at $8$, $2.08$ at $9$ and $4.25$ at $11$. The ratio
$2\pi e^{\Aconst}/2\pi e^{2T}=e^{\Aconst-2T}$ equals $0.89$ at $T=\frac12\log3$ and $456$ at $T=1.20$, and
tends to infinity doubly exponentially because $\Aconst\sim4e^T$. Finally $\mathcal W_{T'}\subseteq\Wcl$
for $T'>T$ (a weight flat on $[-2T',2T']$ is flat on $[-2T,2T]$, and the decay condition does not
involve $T$), so $\sup_{w\in\Wcl}\inf_\R f_w=\lamW$ is nonincreasing in $T$, which is
Corollary~\ref{cor:mono} on the dual side; it is continuous in $T$ by Lemma~\ref{lem:contT} and
\cite[Thm.~5]{Bombieri2000}. The lattice sub-classes used in numerical optimization are not nested in
$T$, and no monotonicity or continuity is asserted for them.

\subsection{The margin decays}\label{subsec:margin}

The trivial half of Zhu's conjectural law $-\log\lamW\asymp e^{2T}$ \cite[Conj.~12.1]{Zhu2026}, namely that
the margin tends to $0$ at an explicit rate, can be proved from three facts about the zeros of $\zeta$.
They are the only numerical facts about zeros used in a proof in this paper, and nothing outside this
subsection depends on them.
\begin{itemize}
\item[(X19)] $\rho_1=\frac12+i\gamma_1$, $\gamma_1=14.134725141734\ldots$, is a simple zero on the critical line,
and $\zeta$ has no non-trivial zero $\rho\ne\rho_1,\bar\rho_1$ with $|\Iml\rho|\le21$
\cite{RosserYoheSchoenfeld1968,PlattTrudgian2021}. Only $\gamma_1\in[14.13,14.14]$ and the gap $(\gamma_1,21]$
are used.
\item[(X19$'$)] Every non-trivial zero $\rho$ with $|\Iml\rho|\le H_{\mathrm{PT}}:=3\cdot10^{12}$ lies on the
critical line \cite{PlattTrudgian2021}.
\item[(X-N)] For $t\ge2$, $\big|N(t)-\frac t{2\pi}\log\frac t{2\pi e}-\frac78\big|\le0.137\log t+0.443\log\log t+1.588$,
where $N(t)$ counts the zeros with $0<\Iml\rho\le t$ with multiplicity \cite[Thm.~19]{Rosser1941}.
\end{itemize}
Besides these, we use that the non-trivial zeros lie in $0<\Rel s<1$, Lemma~\ref{lem:roughEF} and
\cite[Thm.~11.3]{Paper45}.
\begin{lemma}[a crude zero count]\label{lem:Ncount}
$N(t)\le\frac t{2\pi}\log t$ for $t\ge21$.
\end{lemma}
\begin{proof}
$\frac t{2\pi}\log\frac t{2\pi e}=\frac t{2\pi}\log t-\frac{1+\log2\pi}{2\pi}t$ and
$\frac{1+\log2\pi}{2\pi}=0.45166\ldots$, so by (X-N)
$\frac t{2\pi}\log t-N(t)\ge0.45166\,t-\frac78-1.588-0.137\log t-0.443\log\log t$, which is $6.11\ldots$ at
$t=21$ and has derivative $0.45166-0.137/t-0.443/(t\log t)>0$ for $t\ge21$.
\end{proof}

\begin{lemma}[zero sums by counting]\label{lem:zerosum}
For $p\ge2$ and $H\ge21$, with the sum over the ordinates $\gamma=\Iml\rho>0$ of the non-trivial zeros
counted with multiplicity,
\[
\Sigma_p(H):=\sum_{\gamma>H}\gamma^{-p}\ \le\ \frac p{2\pi}\,H^{1-p}\Big[\frac{\log H}{p-1}+\frac1{(p-1)^2}\Big].
\]
\end{lemma}
\begin{proof}
$\sum_{\gamma>H}\gamma^{-p}=\int_{H^+}^\infty t^{-p}dN(t)=\big[t^{-p}N(t)\big]_H^\infty+p\int_H^\infty t^{-p-1}N(t)\,dt$;
the term at $\infty$ vanishes and the one at $H$ is $-H^{-p}N(H)\le0$ by Lemma~\ref{lem:Ncount}, which
also bounds the integral by $\frac p{2\pi}\int_H^\infty t^{-p}\log t\,dt=\frac p{2\pi}H^{1-p}\big[\frac{\log H}{p-1}+\frac1{(p-1)^2}\big]$.
\end{proof}

\begin{lemma}[the test family]\label{lem:testfam}
Fix $T>0$ and an integer $k\ge5$; let $b_k=\frac k2\mathbf 1_{[-1/k,1/k]}$ and $\omega_k=b_k^{*k}$ (the
$k$-fold convolution).
\emph{(i)} $\omega_k\ge0$, $\int\omega_k=1$, $\supp\omega_k\subseteq[-1,1]$, $\omega_k\in C^{k-2}(\R)$ with
$\omega_k^{(j)}(\pm1)=0$ for $j\le k-2$, and $\|\omega_k\|_2^2\ge\frac12$.
\emph{(ii)} $\widehat\omega_k(x)=\big(\sin(x/k)/(x/k)\big)^k$; $|\widehat\omega_k(x)|\le\min\{1,(k/|x|)^k\}$ for real
$x$, and $|\widehat\omega_k(z)|\le(k/|\Rel z|)^ke^{|\Iml z|}$ for complex $z$.
\emph{(iii)} Put $G(t)=\omega_k(t/T)$, $\varpi(c)=(c^2-\frac14)(c^2+\gamma_1^2)$, $\partial=d/dt$, and
\[
F:=\varpi(\partial)G=G''''+(\gamma_1^2-\tfrac14)G''-\tfrac{\gamma_1^2}4G .
\]
Then $F$ is real, supported in $[-T,T]$, bounded, and $F\in\PT$: four integrations by parts, the
boundary terms carrying $G^{(j)}(\pm T)=0$ for $j\le3$, give $\int Fe^{ct}dt=\varpi(c)\int Ge^{ct}dt$
for every $c\in\C$, and $\varpi(\pm\frac12)=0$.
\emph{(iv)} $\Fh(\tau)=(\tau^2+\frac14)(\tau^2-\gamma_1^2)\,T\,\widehat\omega_k(T\tau)$,
$\Phi_F(s)=\varpi(s-\frac12)\Phi_G(s)$ and $\Phi_G(\frac12+\beta+i\gamma)=T\,\widehat\omega_k\big(T(\gamma-i\beta)\big)$;
in particular $\Phi_F(\rho_1)=\Phi_F(\bar\rho_1)=0$.
\emph{(v)} With $Y:=T\gamma_1/\sqrt2$ and $\varepsilon_k(Y):=\frac1\pi k^{2k}Y^{1-2k}/(2k-1)$,
$\|F\|^2\ge\frac{\gamma_1^4}{64}\,T\,\big(\frac12-\varepsilon_k(Y)\big)$ whenever $\varepsilon_k(Y)<\frac12$.
\end{lemma}
\begin{proof}
(i) $\omega_k$ is a spline of degree $k-1$ supported in $[-1,1]$ and vanishing there to order $k-1$;
$1=(\int_{-1}^1\omega_k)^2\le2\int\omega_k^2$. (ii) $\widehat b_k(x)=\sin(x/k)/(x/k)$; $|\sin y/y|\le\min(1,1/|y|)$
on $\R$, and $|\sin(a+ib)|\le\cosh b\le e^{|b|}$. (iii) $G\in C^3$ with $G'''$ absolutely continuous
and $G''''\in L^\infty$ ($k=5$), and $G^{(j)}(\pm T)=0$ for $j\le3$; $\int G''e^{ct}=c^2\int Ge^{ct}$.
(iv) is (iii) with $c=i\tau$, resp.\ $c=s-\frac12$, and $\varpi(\pm i\gamma_1)=0$. (v) On
$|\tau|\le\gamma_1/\sqrt2$ one has $(\tau^2+\frac14)^2\ge\frac1{16}$ and $(\tau^2-\gamma_1^2)^2\ge\gamma_1^4/4$,
so $\|F\|^2=\frac1{2\pi}\int|\Fh|^2\ge\frac{\gamma_1^4}{64}\,T\,\frac1{2\pi}\int_{|x|\le Y}|\widehat\omega_k|^2$, and
$\frac1{2\pi}\int_{|x|\le Y}|\widehat\omega_k|^2\ge\|\omega_k\|_2^2-\frac1\pi\int_Y^\infty(k/x)^{2k}dx\ge\frac12-\varepsilon_k(Y)$.
\end{proof}

\begin{theorem}[the margin decays]\label{thm:margin}
Let $T>0$, $k\ge5$ an integer, $p:=2k-8$, $Y:=T\gamma_1/\sqrt2$, and suppose $\varepsilon_k(Y)<\frac12$. Then
\begin{equation}\label{eq:Bmargin}
\lamW\ \le\ \lamP\ \le\ \lamA\ \le\ \mathcal B(T,k):=
\frac{2\cdot1.0012\cdot T^2(k/T)^{2k}\big[\Sigma_p(21)+e^T\,\Sigma_p(H_{\mathrm{PT}})\big]}
{(14.13^4\,T/64)\,\big(\frac12-\varepsilon_k(Y)\big)},
\end{equation}
with $\Sigma_p(H)$ bounded by Lemma~\ref{lem:zerosum}. Consequently $\limsup_{T\to\infty}\lamA\le0$, and the
same for $\lamP$ and $\lamW$.
\end{theorem}
\begin{proof}
Take $F$ of Lemma~\ref{lem:testfam}. (1) $F\in\PT$, so $\lamA\le\mathfrak a_T(F)/\|F\|^2$, and
$\lamW\le\lamP\le\lamA$ by Theorem~\ref{thm:A} and Remark~\ref{rem:Kprim}.
(2) For $|\tau|\ge21$, $(\tau^2+\frac14)^2\le1.0012\,\tau^4$ and $(\tau^2-\gamma_1^2)^2\le\tau^4$, so
$|\Fh(\tau)|^2\le1.0012\,\tau^8T^2(k/(T|\tau|))^{2k}=O(|\tau|^{8-2k})$ with $8-2k\le-2$: $F$ is in the
form domain, which is hypothesis (a) of Lemma~\ref{lem:roughEF}; hypothesis (b) is verified in step (6).
Hence $\mathfrak a_T(F)=\sum_\rho\Phi_F(\rho)\Phi_F(1-\rho)$.
(3) $\rho_1,\bar\rho_1$ contribute $0$ (Lemma~\ref{lem:testfam}(iv)); by (X19) no other zero has $|\Iml\rho|\le21$.
(4) A zero with $21<|\gamma|\le H_{\mathrm{PT}}$ lies on the line by (X19$'$) and contributes
$|\Fh(\gamma)|^2\le1.0012\,T^2(k/T)^{2k}|\gamma|^{8-2k}$.
(5) A zero $\rho=\frac12+\beta+i\gamma$ with $|\gamma|>H_{\mathrm{PT}}$ has $|\beta|<\frac12$;
$|\varpi(\beta+i\gamma)|\le(\gamma^2+\beta^2+\frac14)(\gamma^2+\beta^2+\gamma_1^2)\le(1+10^{-22})\gamma^4$, $\varpi$ is even,
and $|\Phi_G(\rho)|,|\Phi_G(1-\rho)|\le T(k/(T|\gamma|))^ke^{T|\beta|}$ by Lemma~\ref{lem:testfam}(ii),(iv);
so $|\Phi_F(\rho)\Phi_F(1-\rho)|\le1.0012\,T^2(k/T)^{2k}|\gamma|^{8-2k}e^T$.
(6) Summing (4)--(5) over $\pm\gamma$ with Lemma~\ref{lem:zerosum} ($p=2k-8\ge2$) gives
$\sum_\rho|\Phi_F(\rho)\Phi_F(1-\rho)|\le2\cdot1.0012\,T^2(k/T)^{2k}\big[\Sigma_p(21)+e^T\Sigma_p(H_{\mathrm{PT}})\big]<\infty$,
which is hypothesis (b) of Lemma~\ref{lem:roughEF} and bounds $\mathfrak a_T(F)$.
(7) Divide by Lemma~\ref{lem:testfam}(v) with $\gamma_1^4\ge14.13^4$. The limit statement is
Corollary~\ref{cor:margin1}.
\end{proof}

\begin{corollary}[closed form]\label{cor:margin1}
For every $T\ge1$,
\begin{equation}\label{eq:closedmargin}
\lamW\ \le\ \lamP\ \le\ \lamA\ \le\ 4\cdot10^9\,T\,e^{-15.449\,T},\qquad 15.449<\tfrac{42}e=15.4494\ldots
\end{equation}
\end{corollary}
\begin{proof}
Take $k:=\lceil21T/e\rceil\ge8$ and write $k=21T/e+\theta$, $\theta\in[0,1)$. Then
$(k/21T)^{2k}=\exp\big(2k[\log(1+\theta e/21T)-1]\big)\le\exp(-42T/e+2\theta^2e/21T)\le e^{-15.4494\,T+0.259}$.
Next $(k/T)^{2k}\Sigma_p(21)\le21^9(k/21T)^{2k}\cdot\frac1{2\pi}\frac{2k-8}{2k-9}\big[\log21+\frac1{2k-9}\big]\le21^9\cdot0.5797\,(k/21T)^{2k}$
for $k\ge8$. The bound for the $H_{\mathrm{PT}}$-term is at most $10\,e^T(21/H_{\mathrm{PT}})^{2k-9}$ times that for the first, with
$21/H_{\mathrm{PT}}=7\cdot10^{-12}$ and $2k-9\ge15.449\,T-9$, hence less than $10^{-70}$ of it for $T\ge1$.
Finally $\varepsilon_k(Y)\le0.01566$ for all $T\ge1$: for fixed $k$, $\varepsilon_k$ decreases in $T$, so it
suffices to evaluate it at the left end of each $k$-range, namely $T=1$ for $k=8$ ($\varepsilon_8=0.0060$)
and $T=(k-1)e/21$ for $k\ge9$, where the largest value is $\varepsilon_9=0.01566$ at $T=8e/21$, the values
decreasing thereafter. So $\frac12-\varepsilon_k\ge0.484$, and
$\mathcal B(T,k)\le\frac{2\cdot1.0012\cdot0.5797\cdot21^9e^{0.259}}{(14.13^4/64)\cdot0.484}\,T\,e^{-15.4494\,T}(1+10^{-70})
\le3.97\cdot10^9\,T\,e^{-15.4494\,T}$.
\end{proof}

\begin{corollary}[the structure of the margin]\label{cor:margin2}
\emph{(i)} There is no $c>0$ with $\lamA\ge c$ for all $T$, nor for $\lamP$ or $\lamW$; for $T\ge1$ every
lower bound on any of the three lies below $4\cdot10^9Te^{-15.449T}$.
\emph{(ii)} The limits $L^{\mathrm A}:=\lim_{T\to\infty}\lamA$, $L^{\mathrm{prim}}:=\lim_{T\to\infty}\lamP$ and
$L^{\mathrm W}:=\lim_{T\to\infty}\lamW$ exist in $[-\infty,0]$, and $L^{\mathrm W}\le L^{\mathrm{prim}}\le L^{\mathrm A}$.
\emph{(iii)} $\mathrm{RH}\iff L^{\mathrm W}=0\iff L^{\mathrm{prim}}=0\iff L^{\mathrm A}=0\iff L^{\mathrm A}\ge0$.
Under RH, therefore, $\lamW\to0$, at least at the rate \eqref{eq:closedmargin}. We do not know whether $L^{\mathrm A}=-\infty$
when RH fails.
\end{corollary}
\begin{proof}
(i) is Corollary~\ref{cor:margin1}. (ii) Monotone limits (Lemma~\ref{lem:monoT}) bounded above by
Theorem~\ref{thm:margin}. (iii) Under RH, $W_T\ge0$ on $C_c^\infty$ tests by Weil's criterion, and for
$F$ in the form domain $W_T(F)=\lim_{\varepsilon\to0}W_{T+\varepsilon}(F*\eta_\varepsilon)\ge0$ by the
mollification in the proof of Lemma~\ref{lem:roughEF} (dominated convergence in the archimedean term,
finitely many comb terms, $A,B$ continuous); so $0\le\lamW\le\lamP\le\lamA\le\mathcal B(T,k)\to0$ and all
three limits are $0$. If RH fails, \cite[Thm.~11.3]{Paper45} gives a real primitive $F\in C_c^\infty$,
$\supp F\subseteq[-T_0,T_0]$, with $\mathfrak a_{T_0}(F)=W_{T_0}(F)<0$; then $\lambda_{\min}(A_{T_0}|\mathcal P_{T_0})<0$,
and $L^{\mathrm A}\le\lambda_{\min}(A_{T_0}|\mathcal P_{T_0})<0$ by Lemma~\ref{lem:monoT}. The remaining
equivalences follow from $L^{\mathrm W}\le L^{\mathrm{prim}}\le L^{\mathrm A}\le0$.
\end{proof}


\begin{remark}[comparison with the conjectural law]\label{rem:margin}
Theorem~\ref{thm:margin} gives, whenever $\lamA>0$, $-\log\lamA\ge15.449\,T-\log(4\cdot10^9T)$, which is linear in
$T$, whereas Zhu's law asserts $-\log\lamW\asymp e^{2T}$; the bound is exponentially weaker and is not evidence
for the law. It uses the positions of the zeros up to $21$, their lying on the line up to $3\cdot10^{12}$, and
the counting bound beyond. It is an asymptotic statement, vacuous at the windows where the sign is known
(Appendix~\ref{app:diag}). The loss of the proof on its own test function is quantified in
Remark~\ref{rem:bandtable}, and the bound is expressed in terms of the weights of the comb in
Remark~\ref{rem:resweights}.
\end{remark}

\section{A fixed window is a finite condition}\label{subsec:finite}

Fix $T$. For a nonnegative arithmetic function $a$ on $n\ge2$ write $v_n(a)=2a(n)/\sqrt n$ and
$\lambda_{\min}(A_T(a)|\PT)$ for the bottom on $\PT$ of the form \eqref{eq:form} built with these weights
(as in Corollary~\ref{cor:SC1}). We say that \emph{Weil positivity holds at $T$ for $a$}, written
$\mathrm{WP}_T(a)$, if $\lambda_{\min}(A_T(a)|\PT)\ge0$, and we write $\mathrm{WP}_T:=\mathrm{WP}_T(\Lambda)$, which is
$\lamA\ge0$. ($\mathrm{WP}_T$ is not the dual value $D_T$ of Definition~\ref{def:PD}.) For a real vector
$v=(v_n)_{n\in\NT}$, signed or not, $\mathfrak a^{(v)}_T$ denotes the form \eqref{eq:form} with the comb
$\sum_{n\in\NT}v_n\cos(\tau\log n)$. With $v(a):=(v_n(a))_{n\in\NT}$, the \emph{comb-level complete form} is
\[
W_T(a)(F):=\mathfrak a^{(v(a))}_T(F)+2\Rel\big(A(F)\overline{B(F)}\big),
\]
the form \eqref{eq:W} built with the weights of $a$, and $\lambda^{\mathrm W}_T(a)$ is its bottom on the unit sphere of
$L^2(\IT)$. So $W_T(\Lambda)=W_T$, and on $\PT$ the form $W_T(a)$ agrees with $\mathfrak a^{(v(a))}_T$. For
$a\ne\Lambda$, $W_T(a)$ is not the zero sum of $\zeta_a$: for the modifications of Theorem~\ref{thm:finmod}
below, $\zeta_a=\zeta E$ has exactly the zeros of $\zeta$, whose zero sum is $W_T$ (Lemma~\ref{lem:roughEF}),
while $W_T(a)(F)-W_T(F)=\sum_{n\in\NT}(v_n(\Lambda)-v_n(a))\cF(\log n)$, the contribution of $E'/E$, is in
general nonzero. Statements about $\lambda^{\mathrm W}_T(a)$ and $\mathrm{WP}_T(a)$ for $a\ne\Lambda$ are therefore
statements about forms, not about the zeros of $\zeta$ or of any $\zeta_a$.

\begin{proposition}[finite dependence]\label{prop:finite}
Put $v_T(a):=(v_n(a))_{n\in\NT}\in\R^{\NT}$, $|\NT|=\lceil e^{2T}\rceil-2$. Then:
\emph{(i)} $\mathrm{WP}_T(a)$ depends on $a$ only through $v_T(a)$;
\emph{(ii)} the set $\mathcal V_T:=\{v\in\R^{\NT}:\ \mathfrak a^{(v)}_T\ge0$ on $\PT\cap$ form domain$\}$ is closed
and convex, namely
\[
\mathcal V_T=\bigcap_{F}\Big\{v:\ \sum_{n\in\NT}v_n\,\cF(\log n)\le C_0+\Lev(\cF)\Big\},
\]
the intersection over the real unit vectors $F\in\PT$ in the form domain;
\emph{(iii)} $0\notin\mathcal V_T$ for $T>T_0(0)=0.74921$, indeed for $T\ge0.60$, and a vector $v\ge0$ with
$\sum_nv_n<-C_0-\varphi((3\pi/2T)^2)$ is not in $\mathcal V_T$ (for signed $v$ the same holds with
$\sum_n|v_n|$);
\emph{(iv)} for $a=\Lambda$ the numbers entering $\mathrm{WP}_T$ are exactly $\Lambda(n)$, $2\le n<e^{2T}$,
together with the one explicit function $m$ (equivalently $C_0$ and $K$), and nothing else.
\end{proposition}
\begin{proof}
(i) and (iv) are the definition of $\sT$. (ii) Real tests suffice: for $F=F_1+iF_2$ with $F_j$ real,
$|\Fh|^2=|\Fh_1|^2+|\Fh_2|^2+2\Iml(\Fh_1\overline{\Fh_2})$, and the last term is odd and integrates to $0$
against the even symbol, so the form splits as the sum of its values at $F_1,F_2$; and $A(F)=A(F_1)+iA(F_2)$
with $A(F_j)$ real, likewise for $B$, so $F\in\PT$ if and only if $F_1,F_2\in\PT$. For real unit $F$ in the
form domain, \eqref{eq:Aform} reads $\mathfrak a^{(v)}_T(F)=C_0+\Lev(\cF)-\sum v_n\cF(\log n)$ with
$\Lev(\cF)<\infty$, a closed half-space condition on $v$. (iii) The first statement is
Proposition~\ref{prop:arch}(ii) ($\laminfP<0$ for $T>T_0(0)$), and Corollary~\ref{cor:Tarch} for $T\ge0.60$; the second Theorem~\ref{thm:SC} with these
weights; for signed $v$ the comb term of that proof is bounded by $\sum|v_n|\,\|F\|^2$.
\end{proof}

\begin{remark}[regimes, and what is decidable]\label{rem:regimes}
If $n_1<n_2<\cdots$ enumerate the prime powers, then on each regime $T\in(\frac12\log n_k,\frac12\log n_{k+1}]$
the vector $v_T(\Lambda)$, the mass $\Aconst$ and the symbol $\sT$ are constant, $\mathrm{WP}_T$ varies only
through $\IT$, and $\lamA$ is nonincreasing across all regimes (Lemma~\ref{lem:monoT}). At a regime
boundary one atom $v_{n_{k+1}}$ enters, and the sign of its effect is not determined by anything here:
$\cF(\log n)$ takes both signs on $\PT$ (Lemma~\ref{lem:cpos}). Per window, $\mathrm{WP}_T$ is a closed
condition on a known vector, semi-decidable from the negative side by exhibiting $F$;
Proposition~\ref{prop:decide} decides the sign of $\lamP$ when $\lamP\ne0$, not that of $\lamA$.
\end{remark}

\begin{lemma}[a contact of either sign]\label{lem:cpos}
Let $0<u_0<2T$. There are real $F_\pm\in\PT\cap C_c^\infty(\IT)$ with $\pm c_{F_\pm}(u_0)>0$.
\end{lemma}
\begin{proof}
Choose $0<\ell<\min(u_0,2T-u_0)$ and $J=(-T,-T+\ell)$: then $J$, $J+u_0$, $J+2u_0$ are pairwise disjoint
and $J\cup(J+u_0)\subset\IT$. For real $\phi\in C_c^\infty(J)$ put $F^\pm_\phi=\phi\pm\phi(\cdot-u_0)$. Then
$c_{F^\pm_\phi}(u_0)=\langle\phi\pm\phi_{u_0},\phi_{u_0}\pm\phi_{2u_0}\rangle=\pm\|\phi\|^2$, the three cross terms
vanishing by disjointness. The map $\phi\mapsto F^\pm_\phi$ is injective, so its image is
infinite-dimensional, and its intersection with $\PT$ has codimension at most $2$ in it and contains some
$F^\pm_\phi$ with $\phi\ne0$.
\end{proof}

\begin{corollary}[each local prime form is indefinite on $\PT$]\label{cor:Qp}
For a prime $p$ with $\log p<2T$ let $Q_p(F)=-\sum_{k\ge1,\,k\log p<2T}v_{p^k}\Rel\cF(k\log p)$ be the
$p$-part of the prime form $\PiT$ of \S\ref{subsec:arch}. With $u_0=\log p$, $J$ and $F^\pm_\phi$ as in
the proof of Lemma~\ref{lem:cpos}, one has $Q_p(F^\pm_\phi)=\mp v_p\|\phi\|^2$ for every real
$\phi\in C_c^\infty(J)$. Hence, for every $T>\frac12\log p$, $Q_p$ is negative definite on the
infinite-dimensional real subspace $\{F^+_\phi:\phi\ \text{real}\}\cap\PT$ of $\PT\cap C_c^\infty(\IT)$,
and positive definite on $\{F^-_\phi:\phi\ \text{real}\}\cap\PT$.
\end{corollary}
\begin{proof}
$\ell<u_0$ makes the translates $J+ku_0$, $k\ge0$, pairwise disjoint. So for $k\ge2$ the supports
$J\cup(J+u_0)$ of $F^\pm_\phi$ and $(J+ku_0)\cup(J+(k+1)u_0)$ of $F^\pm_\phi(\cdot-ku_0)$ are disjoint,
and $c_{F^\pm_\phi}(k\log p)=0$. The case $k=1$ is Lemma~\ref{lem:cpos}. The subspaces are
infinite-dimensional by the last sentence of that proof.
\end{proof}


Thus no single $Q_p$ is $\ge0$ on $\PT$, nor is the archimedean form for $T\ge0.60$
(Corollary~\ref{cor:Tarch}); for finite sets of places see Theorem~\ref{thm:SC} and
Proposition~\ref{prop:finsupp}. Indefinite summands do not forbid a positive sum (on $\PT$ one has
$\mathfrak a_T=\Ainf+\PiT$ with both summands indefinite for $T\ge0.60$); the corollary only rules out
inferring the sign from a single local form.

\begin{theorem}[finite modifications]\label{thm:finmod}
Let $\psi_a(x):=\sum_{n\le x}a(n)$ (so $\psi_\Lambda$ is Chebyshev's function) and, for $\Rel s>1$,
$\zeta_a(s):=\exp\big(\sum_{n\ge2}a(n)n^{-s}/\log n\big)$.
\emph{(i) Deletion.} Let $N\ge5$ and $a:=\Lambda\mathbf 1_{n\ge N}$. For $T\le\frac12\log N$ the comb of $a$ is
empty, so $\mathrm{WP}_T(a)\iff\laminfP(T)\ge0$; hence $\mathrm{WP}_T(a)$ holds for every $T\le0.58$
and fails for every $T$ with $0.60\le T\le\frac12\log N$ (Theorem~\ref{thm:index}(i), Corollary~\ref{cor:Tarch}; for $T<\frac12\log2$
and for $T>T_0(0)$ already by Proposition~\ref{prop:arch}(iii) and (ii)); on $(0.58,0.60)$ it holds below
$T^*_{\infty,\mathrm P}$ and fails above it (Lemma~\ref{lem:strict}), the threshold not being located more
precisely here.
\emph{(ii) Amplification.} Let $n_0\ge2$, $c>0$ and $a:=\Lambda+c\,\mathbf 1_{n=n_0}$. For every $T$ with
$2T>\log n_0$ there is $c_*(T,n_0)<\infty$ such that $\mathrm{WP}_T(a)$ fails for all $c>c_*$. For $n_0=2$
this holds at every $T>\frac12\log2$, including $(\frac12\log2,\log2]$, where $\mathrm{WP}_T(\Lambda)$ holds
\cite{Paper42}.
\emph{(iii) What is preserved.} Both modifications keep $a\ge0$ and change $\psi_a$ by a constant for
$x\ge N$ (resp.\ $x\ge n_0$), so they preserve every partial-sum statement $\psi_a(x)=x+O(\varepsilon(x))$
with $\inf\varepsilon>0$; and $\zeta_a=\zeta\cdot E$ with
$E(s)=\exp\big(-\sum_{n<N}\Lambda(n)n^{-s}/\log n\big)$ (resp.\ $\exp(c\,n_0^{-s}/\log n_0)$) entire and zero-free,
so $\zeta_a$ continues meromorphically with \emph{exactly the zeros and the pole of $\zeta$}. Deletion also
preserves $a\le\Lambda$. Neither preserves the functional equation, $E(s)\ne E(1-s)$.
\end{theorem}
\begin{proof}
(i) $n\in\NT$ means $n<e^{2T}\le N$, so $a=0$ on $\NT$ and the form is the archimedean form on $\PT$, whose
bottom is $\laminfP(T)$ (\S\ref{subsec:arch}); Proposition~\ref{prop:arch}(iii) gives
$\laminfP\ge\laminf\ge0$ for $T<\frac12\log2$, and Proposition~\ref{prop:arch}(ii) gives
$\laminfP(T)\le C_0+\varphi((3\pi/2T)^2)<0$ for $T>T_0(0)$; Theorem~\ref{thm:index}(i) and Corollary~\ref{cor:Tarch} give $\laminfP>0$ on
$(0,0.58]$ and $\laminfP<0$ on $[0.60,\infty)$; $N\ge5$ makes $\frac12\log N\ge0.8047>T_0(0)$.
(ii) Take $F=F_+$ of Lemma~\ref{lem:cpos} with $u_0=\log n_0$. Being smooth, it lies in the form domain,
and $\mathfrak a^{(a)}_T(F)=\mathfrak a_T(F)-(2c/\sqrt{n_0})\,\cF(\log n_0)$, which is negative for
$c>c_*:=\sqrt{n_0}\max(\mathfrak a_T(F),0)/(2\cF(\log n_0))$. (iii) The exponent of $\zeta_a$ differs from that
of $\zeta$ by a finite Dirichlet polynomial, an entire function.
\end{proof}

\begin{remark}[resolution in the weights]\label{rem:resweights}
This remark quantifies Theorem~\ref{thm:finmod}(ii) by Corollary~\ref{cor:margin1}, first for the
primitive form. Fix $T$ with $\lamA>0$ and $n_0\in\NT$. A real unit minimizer $F_\star$ of
$\mathfrak a_T$ on $\PT$ exists: the sublevel sets are compact (proof of Theorem~\ref{thm:O}), $\PT$ is
closed, and the splitting in the proof of Proposition~\ref{prop:finite}(ii) gives a real one. Here
$c_{F_\star}$ is its autocorrelation, not the transform $\Phi_{F_\star}$ of \eqref{eq:W}. Suppose
$c_{F_\star}(\log n_0)\ne0$. Put $s=\operatorname{sgn}c_{F_\star}(\log n_0)$ and, for $\varepsilon>0$,
$a=\Lambda+s\varepsilon\mathbf 1_{n=n_0}$. If $s=-1$, this requires $\varepsilon\le\Lambda(n_0)$, so
that $a\ge0$. In particular, a contact with $\Lambda(n_0)=0$ admits only $s=+1$, and the class
$\mathcal C_{\mathrm{BT}}$ admits only $s=-1$. Apply the computation in the proof of
Theorem~\ref{thm:finmod}(ii) to $F_\star$ in place of $F_+$:
\[
\lambda_{\min}(A_T(a)|\PT)\ \le\ \mathfrak a^{(v(a))}_T(F_\star)\ =\ \lamA-\frac{2\varepsilon}{\sqrt{n_0}}\,|c_{F_\star}(\log n_0)| .
\]
So $\mathrm{WP}_T(a)$ fails for every admissible $\varepsilon>\varepsilon_*(T,n_0):=\sqrt{n_0}\,\lamA/(2|c_{F_\star}(\log n_0)|)$.
For $T\ge1$, Corollary~\ref{cor:margin1} and $\sqrt{n_0}<e^T$ give
\[
\varepsilon_*(T,n_0)\ \le\ \frac{2\cdot10^9\,T\,e^{-14.449\,T}}{|c_{F_\star}(\log n_0)|},\qquad
\varepsilon_*(T,2)\ \le\ \frac{2\sqrt2\cdot10^9\,T\,e^{-15.449\,T}}{|c_{F_\star}(\log2)|}.
\]
For the complete form, take a real unit minimizer $F^{\mathrm W}$ of $W_T$ on $L^2(\IT)$. It is attained,
as in the proof of Corollary~\ref{cor:RH}. Take $s$ to be the sign of $c_{F^{\mathrm W}}(\log n_0)$. The same
computation gives $\lambda^{\mathrm W}_T(a)\le\lamW-(2\varepsilon/\sqrt{n_0})|c_{F^{\mathrm W}}(\log n_0)|$, and
the same bounds hold with $c_{F^{\mathrm W}}$ in place of $c_{F_\star}$, since $\lamW\le\lamA$. That is a
different number, and it needs the minimizer of $W_T$, not that of $\mathfrak a_T$.

These are upper bounds on the threshold, not its size, and the true threshold may be far smaller.
Since $|c_{F_\star}|\le\|F_\star\|^2=1$, the factor $1/|c_{F_\star}(\log n_0)|$ only enlarges the bound, and
no lower bound on $\max_{n\in\NT}|c_{F_\star}(\log n)|$ is known. Even when $|c_{F_\star}(\log n_0)|=1$ the
first bound exceeds $1$ for $T\le1.51$, so the remark is vacuous there, as in Remark~\ref{rem:margin}. For
$\mathcal C_{\mathrm{PNT}}(\varepsilon)$ and $\mathcal C_{\mathrm{zeros}}$, Corollary~\ref{cor:void}
already applies with no size restriction. The remark adds content only for classes closed under a
change of one weight by $\varepsilon_*(T,n_0)$, for instance inputs known only to within a given width.
We do not know whether $c_{F_\star}(\log n_0)\ne0$ for some minimizer and some contact. If this fails,
every minimizer vanishes on $\log\NT$, and Lemma~\ref{lem:degen} describes that case.
\end{remark}

\begin{lemma}[the degenerate case]\label{lem:degen}
Fix $T$. For $v\in\R^{\NT}$ let
$V_T(v)=\lambda_{\min}(A^{(v)}_T|\PT):=\inf\{\mathfrak a^{(v)}_T(F):\ F\in\PT\ \text{in the form domain},\ \|F\|=1\}$, so that
$V_T(v_T(\Lambda))=\lamA$. Then $V_T$ is finite and concave on $\R^{\NT}$. If $F_\star$ is a real unit
minimizer at $v^0=v_T(\Lambda)$, then
$V_T(v)\le V_T(v^0)-\sum_{n\in\NT}(v_n-v^0_n)\,c_{F_\star}(\log n)$ for all $v$. Hence, if some real unit
minimizer has $c_{F_\star}(\log n)=0$ for every $n\in\NT$, then $v_T(\Lambda)$ is a global maximizer of
$V_T$ on $\R^{\NT}$. In particular, $\lamA\ge V_T(0)=\laminfP(T)$.
\end{lemma}
\begin{proof}
By \eqref{eq:Aform} and the splitting in the proof of Proposition~\ref{prop:finite}(ii), $V_T$ is the
infimum, over real unit $F\in\PT$ in the form domain, of the affine functions
$v\mapsto C_0+\Lev(\cF)-\sum_nv_n\cF(\log n)$. So $V_T$ is concave. It is bounded below by
$C_0-\sum_n|v_n|$, because $\Lev\ge0$ and $|\cF|\le1$, and bounded above by its value at any fixed $F$.
The displayed inequality is $V_T(v)\le\mathfrak a^{(v)}_T(F_\star)$, written out with \eqref{eq:Aform}.
\end{proof}

We do not know whether the degenerate case occurs at any $T$. If it did, the comb of $\zeta$ would be
the optimal comb for the primitive bottom at that window among all real weight vectors.

\begin{proposition}[finitely supported combs]\label{prop:finsupp}
Let $a\ge0$ have finite support $S$ and put $A_S=\sum_{n\in S}v_n(a)$. Then
$\limsup_{T\to\infty}\lambda_{\min}(A_T(a)|\PT)\le C_0-A_S<0$; in particular $\mathrm{WP}_T(a)$ fails for all
large $T$. This applies to every truncation $\Lambda\mathbf 1_{n<N}$, and for the empty comb it recovers
a bound consistent with $\laminfP\downarrow C_0$ (Proposition~\ref{prop:arch}(ii)).
\end{proposition}
\begin{proof}
For real $g\in C_c^\infty(-1,1)$, $g\ne0$, put $G_T(t)=g(t/T)$ and $F_T=(-\partial^2+\frac14)G_T\in C_c^\infty(\IT)$.
Two integrations by parts give $A(F_T)=\int G_T(-\partial^2+\frac14)e^{t/2}dt=0$ and likewise $B(F_T)=0$, so
$F_T\in\PT$, and $\Fh_T(\tau)=(\tau^2+\frac14)T\widehat g(T\tau)$. Let $2T>\max_S\log n$, so that
$\sigma:=m-\sum_Sv_n(a)\cos(\tau\log n)$ does not depend on $T$. Substituting $s=T\tau$,
\[
\frac{\mathfrak a^{(a)}_T(F_T)}{\|F_T\|^2}=\frac{\int\sigma(s/T)(s^2/T^2+\frac14)^2|\widehat g(s)|^2ds}{\int(s^2/T^2+\frac14)^2|\widehat g(s)|^2ds}
\ \longrightarrow\ \sigma(0)=C_0-A_S\qquad(T\to\infty)
\]
by dominated convergence: for $T\ge1$, $|\sigma(s/T)|\le C(1+\log(2+|s|))$ (Lemma~\ref{lem:mbounds} and
$|\text{comb}|\le A_S$), $(s^2/T^2+\frac14)^2\le(s^2+\frac14)^2$, $\widehat g$ is Schwartz, and $\sigma$ is
continuous at $0$.
\end{proof}

\begin{corollary}[classical prime-side hypotheses decide no window]\label{cor:void}
Let $\mathcal C$ be a class of nonnegative arithmetic functions containing $\Lambda$. If $\mathcal C$ is
closed under deletion of finitely many terms, then for every $T>T_0(0)$ some $a\in\mathcal C$ fails
$\mathrm{WP}_T(a)$, namely $a=\Lambda\mathbf 1_{n\ge N}$ with $N\ge\max(5,e^{2T})$. If $\mathcal C$ is closed under
amplification of one term, then for every $T>\frac12\log2$ some $a\in\mathcal C$ fails $\mathrm{WP}_T(a)$, namely
$a=\Lambda+c\,\mathbf 1_{n=2}$ with $c>c_*(T,2)$. The following classes contain $\Lambda$ and have these closure
properties:
\begin{itemize}
\item the class $\mathcal C_{\mathrm{PNT}}(\varepsilon)=\{a\ge0:\psi_a(x)=x+O(\varepsilon(x))\}$, for any error term
$\varepsilon$ with $\inf\varepsilon>0$ that is known for $\Lambda$ (e.g.\ de la Vall\'ee Poussin's), is closed under
both operations;
\item the class $\mathcal C_{\mathrm{zeros}}=\{a\ge0:\ \zeta_a=\zeta E,\ E$ entire and zero-free$\}$ of all $a$ with
the zeros and the pole of $\zeta$, and hence with every zero-free region and every zero-density estimate
of $\zeta$, is closed under both operations;
\item the class $\mathcal C_{\mathrm{BT}}=\{a:0\le a\le\Lambda\}$, which satisfies every upper-bound sieve
inequality satisfied by $\Lambda$, is closed under deletion.
\end{itemize}
Consequently no argument that uses about the primes only membership of $\Lambda$ in such a class (for
instance a prime number theorem with error term, a zero-free region, a zero-density estimate, or an
upper-bound sieve inequality) can prove $\mathrm{WP}_T$ at a window $T>T_0(0)$, since it would equally prove
$\mathrm{WP}_T(a)$ for the member $a$ above, which is false; for $\mathcal C_{\mathrm{PNT}}(\varepsilon)$ and
$\mathcal C_{\mathrm{zeros}}$ the same holds at every window $T>\frac12\log2$.
\end{corollary}
\begin{proof}
Theorem~\ref{thm:finmod}(i),(ii) give the failing members; the closure properties are
Theorem~\ref{thm:finmod}(iii).
\end{proof}

This concerns classes of weights, not $\zeta$ itself: if RH is true, every true statement implies it.

\begin{remark}[the functional equation]\label{rem:twohalves}
A similar obstruction comes from the other side of the explicit formula. The functional equation and
the pole are shared by every Dedekind zeta function, and Epstein zeta functions of class number $>1$,
which share the functional equation, reality, the polar set and nonnegative coefficients, have zeros off
the critical line \cite{DavenportHeilbronn1936}, so their Weil forms are negative on some window. Hence a
proof of $\mathrm{WP}_T$ for $\zeta$ must use a property of $\zeta$ that (a) is not preserved by the modified
weights of Theorem~\ref{thm:finmod}, which fail $\mathrm{WP}_T$ at the windows listed in
Corollary~\ref{cor:void}, so that it depends on the actual values $\Lambda(n)$, $n\in\NT$, and not only on
their asymptotics; and (b) is not shared by all Epstein zeta functions, so that it does not follow from
the functional equation, the pole, reality and nonnegative coefficients alone. The known proofs at
individual windows are computations with the actual vector: \cite{Paper42} for $T\le\log2$ (the
primitive form), \cite{Zhu2026} at support $1.6$ (the complete form, which implies $\mathrm{WP}_T$), and
Theorem~\ref{thm:C} at $\frac12\log3$ (the primitive form).
\end{remark}

\begin{proposition}[the arithmetic content is uniformity in $T$]\label{prop:uniform}
\[
\mathrm{RH}\iff\forall T:\ \mathrm{WP}_T\iff\forall T:\ v_T(\Lambda)\in\mathcal V_T ,
\]
where each $\mathcal V_T$ is a closed convex subset of a finite-dimensional space defined by $\Gamma$-data
alone (Proposition~\ref{prop:finite}). Moreover: \emph{(a)} for fixed $T$, $\mathrm{WP}_T$ is a statement
about the finite, exactly known list $\{\Lambda(n):n<e^{2T}\}$ and the function $m$; \emph{(b)} no finite
initial segment of $\Lambda$ decides $\forall T\,\mathrm{WP}_T$: for every $N$ there are $a\ge0$ with $a=\Lambda$
on $n<N$ for which $\mathrm{WP}_T(a)$ fails for all large $T$ (namely $\Lambda\mathbf 1_{n<N}$), and, for each
$T>\frac12\log N$, others for which it fails at that $T$ (namely $\Lambda+c\mathbf 1_{n=N}$ with $c>c_*(T,N)$). So
the only arithmetic assertion in positivity for every $T$ is that the sequence $(\Lambda(n))_{n\ge2}$ meets
every $\mathcal V_T$; at a single window there is none beyond the finite vector.
\end{proposition}
\begin{proof}
The first equivalence is \cite[Thm.~11.3]{Paper45} with Weil's criterion, as in the introduction; the
second is Proposition~\ref{prop:finite}(ii). (a) is Proposition~\ref{prop:finite}(iv). (b) is
Proposition~\ref{prop:finsupp} and Theorem~\ref{thm:finmod}(ii).
\end{proof}

\subsection{The negative set of the symbol}\label{subsec:band}

Fix a regime, so that $\sigma=\sT$ and $A=\Aconst$ are fixed, and put
$m_{\mathrm{lo}}(\tau)=\log\frac\tau{2\pi}-\frac1\tau-\frac1{2\tau^2}$, the lower envelope of
Lemma~\ref{lem:mbounds}(ii), which increases from $-\infty$ to $+\infty$ on $(0,\infty)$.

\begin{proposition}[the negative set is a finite band]\label{prop:band}
Let $\tau^{\mathrm{lo}}(A)$ be the unique positive root of $m_{\mathrm{lo}}(\tau)=A$. Then \emph{(i)} $\sigma>0$ on
$(\tau^{\mathrm{lo}}(A),\infty)$; \emph{(ii)} $\sigma(0)=C_0-A<0$; \emph{(iii)} $\sigma$ is real-analytic and
nonconstant on $\R$, so its zero set in $[0,\infty)$ is finite; \emph{(iv)} $\{\sigma<0\}$ is a finite union of
bounded open intervals, the first containing $0$, and the last sign change
$\tau_{\mathrm{last}}:=\sup\{\tau:\sigma(\tau)<0\}$ satisfies $\tau_{\mathrm{last}}\le\tau^{\mathrm{lo}}(A)$;
\emph{(v)} $\tau^{\mathrm{lo}}(A)=2\pi e^A\,e^{1/\tau^{\mathrm{lo}}+1/(2(\tau^{\mathrm{lo}})^2)}>2\pi e^A$.
\end{proposition}
\begin{proof}
(i) $m(\tau)\ge m_{\mathrm{lo}}(\tau)>A\ge\sum_{n\in\NT}v_n\cos(\tau\log n)$ for $\tau>\tau^{\mathrm{lo}}$
(Lemma~\ref{lem:mbounds}(ii),(iii)). (ii) $m(0)=C_0$. (iii) $\Rel\psi(\frac14+\frac{i\tau}2)=\frac12[\psi(\frac14+\frac{i\tau}2)+\psi(\frac14-\frac{i\tau}2)]$
is real-analytic in $\tau\in\R$, the comb is a trigonometric polynomial, and $\sigma\to+\infty$; by (i) the
zeros lie in the compact $[0,\tau^{\mathrm{lo}}]$. (iv) follows from (i)--(iii); (v) rearranges the equation.
\end{proof}

For $\NT=\{2,3\}$ ($A=2.248826$) one has $\tau^{\mathrm{lo}}=60.5432$, and for $\NT=\{2,3,4\}$
($A=2.941974$, $v_4=\log2$) $\tau^{\mathrm{lo}}=120.0865$ (roots of the strictly increasing explicit
function $m_{\mathrm{lo}}$, computed in double precision), against $2\pi e^A=59.543$ and $119.087$. So the
``threshold $2\pi e^{\Aconst}$'' of Proposition~\ref{prop:env}(iii) is the multiplier-free envelope zero up
to the factor $e^{1/\tau^{\mathrm{lo}}+\dots}$, which the $o(1)$ there covers.

\begin{proposition}[the alignment structure at two and three contacts]\label{prop:align}
Let $\NT=\{2,3\}$, i.e.\ $T\in(\frac12\log3,\log2]$, so $v_2=0.980258$, $v_3=1.268568$, and let $\tau_a=23.3409$,
$\tau_b=17.7446$ be the roots of $m_{\mathrm{lo}}=v_3$ and $m_{\mathrm{lo}}=v_2$. Then
$\{\sigma<0\}=\{\tau:v_2\cos(\tau\log2)+v_3\cos(\tau\log3)>m(\tau)\}$, and every $\tau>\tau_a$ with
$\sigma(\tau)<0$ satisfies
\[
\cos(\tau\log2)>\frac{m_{\mathrm{lo}}(\tau)-v_3}{v_2}>0,\qquad\cos(\tau\log3)>\frac{m_{\mathrm{lo}}(\tau)-v_2}{v_3}>0,
\]
so that there are unique integers $k,l\ge1$ with $|\tau\log2-2\pi k|<\delta_2(\tau)$ and
$|\tau\log3-2\pi l|<\delta_3(\tau)$, where $\delta_2=\arccos\frac{m_{\mathrm{lo}}-v_3}{v_2}$ and
$\delta_3=\arccos\frac{m_{\mathrm{lo}}-v_2}{v_3}$ are $<\frac\pi2$ and decreasing; consequently
\[
|k\log3-l\log2|<\frac{\delta_2\log3+\delta_3\log2}{2\pi}<\frac{\log6}4=0.4479,\qquad\text{i.e.}\qquad
0.639<2^l/3^k<1.565 .
\]
Every component of $\{\sigma<0\}\cap(\tau_a,\tau^{\mathrm{lo}}]$ lies in one arc $\{|\tau\log2-2\pi k|<\delta_2\}$ of
length $<\pi/\log2$, and at most $6$ such arcs meet $(\tau_a,\tau^{\mathrm{lo}}]$. For $\NT=\{2,3,4\}$,
$T\in(\log2,\frac12\log5]$, the same holds with $|\tau\log2-2\pi k|<\arccos y_0=1.1185$ ($y_0=0.43702$ the
positive root of $2v_4y^2+v_2y-v_4$) for $\tau>\tau_a$, with the bracket constant
$(1.1185\log3+\frac\pi2\log2)/2\pi=0.369$ on $(\tau_b',\tau^{\mathrm{lo}}]$, $\tau_b'=34.4909$ the root of
$m_{\mathrm{lo}}=v_2+v_4$, and with at most $12$ arcs.
\end{proposition}
\begin{proof}
If $\sigma(\tau)<0$ then $v_2\cos(\tau\log2)>m(\tau)-v_3\cos(\tau\log3)\ge m_{\mathrm{lo}}(\tau)-v_3$, which is
positive for $\tau>\tau_a$; symmetrically for $3$, since $\tau_b<\tau_a$. The arcs of $\cos x>c>0$ are
$|x-2\pi k|<\arccos c<\frac\pi2$, and $k,l\ge1$ because $\tau\log2>\frac\pi2$ there. Dividing the two arc
conditions by $\log2$ and $\log3$ and subtracting gives $2\pi|k/\log2-l/\log3|<\delta_2/\log2+\delta_3/\log3$;
multiply by $\log2\log3/2\pi$, and use $\delta_j<\frac\pi2$. Consecutive arcs are $2\pi/\log2$ apart, whence
the count $\lfloor(\tau^{\mathrm{lo}}-\tau_a)\log2/2\pi\rfloor+2=6$. For $\{2,3,4\}$ the comb is
$g(\tau\log2)+v_3\cos(\tau\log3)$ with $g(x)=v_2\cos x+v_4\cos2x=2v_4y^2+v_2y-v_4$, $y=\cos x$, increasing for
$y>-v_2/4v_4$ and positive exactly for $y>y_0$ on $[-1,1]$; for $\tau>\tau_a$ a negative $\tau$ has $g>0$, and
for $\tau>\tau_b'$ also $v_3\cos(\tau\log3)>m_{\mathrm{lo}}-v_2-v_4>0$; the rest is as before, with the count
$\lfloor10.67\rfloor+2=12$.
\end{proof}

The last sign change is located with a proof. By the upper bound
$m\le m_{\mathrm{up}}:=\log\frac\tau{2\pi}+\frac1\tau+\frac12\log(1+\frac1{4\tau^2})$ of
Lemma~\ref{lem:mbounds}(ii): at $\tau=45.64206$ the comb for $\{2,3\}$ equals $2.21553>m_{\mathrm{up}}=2.00492$,
and at $\tau=108.72265$ the comb for $\{2,3,4\}$ equals $2.93679>m_{\mathrm{up}}=2.86013$; so $\sigma<0$ at these
points and, with Proposition~\ref{prop:band}, $\tau_{\mathrm{last}}\in[45.642,60.544]$ for $\{2,3\}$ and
$\tau_{\mathrm{last}}\in[108.723,120.087]$ for $\{2,3,4\}$ (up to double-precision evaluation, with margins
$0.21$ and $0.077$). The measured components of $\{\sigma<0\}$ are in Remark~\ref{rem:bandtable}.

\begin{proposition}[the negative part is trace class, and the band does not decide]\label{prop:negtrace}
Write $\sigma=\sigma_+-\sigma_-$ with $\sigma_\pm=\max(\pm\sigma,0)$, so that $\sigma_-$ is bounded and supported in
$[-\tau_{\mathrm{last}},\tau_{\mathrm{last}}]$, and let $\mathrm{Pos}_T=R_T\sigma_+(D)R_T\ge0$ (a form on the form
domain) and $\mathrm{Neg}_T=R_T\sigma_-(D)R_T\ge0$ on $L^2(\IT)$, in the notation of \S\ref{subsec:negindex}. Then
\emph{(i)} $\mathrm{Neg}_T$ is trace class, with $\operatorname{tr}\mathrm{Neg}_T=\frac{2T}{2\pi}\int_\R\sigma_-(\tau)\,d\tau$;
\emph{(ii)} $\mathfrak a_T(F)\ge-\frac1{2\pi}\int_{|\tau|\le\tau_{\mathrm{last}}}\sigma_-|\Fh|^2$ for every $F$ in the
form domain: all negativity of the form is produced on the band;
\emph{(iii)} $\lamA\ge-\|\mathrm{Neg}_T\|\ge-\operatorname{tr}\mathrm{Neg}_T$;
\emph{(iv)} $\mathrm{WP}_T\iff\langle\mathrm{Neg}_TF,F\rangle\le\mathrm{Pos}_T(F)$ for all $F\in\PT$ in the form domain;
\emph{(v)} the sign of the bottom on $\PT$ is not a function of $\sigma|_{[-\tau_{\mathrm{last}},\tau_{\mathrm{last}}]}$
alone: with $\lambda_0:=\lambda_0(S_{T,\tau_{\mathrm{last}}})<1$ the top eigenvalue of the time--band limiting
operator of Proposition~\ref{prop:NC}, $s_+:=\max_{|\tau|\le\tau_{\mathrm{last}}}\sigma_+$ and $\mathbf 1_{\mathrm{off}}$ the
indicator of $|\tau|>\tau_{\mathrm{last}}$, the symbol $\sigma\mathbf 1_{|\tau|\le\tau_{\mathrm{last}}}-c\,\mathbf 1_{\mathrm{off}}$
gives a form negative on all of $L^2(\IT)\setminus\{0\}$ for $c>s_+/(1-\lambda_0)$, and the symbol
$\sigma+c\,\mathbf 1_{\mathrm{off}}$ gives a form $\ge(c(1-\lambda_0)-\|\mathrm{Neg}_T\|)\|F\|^2>0$ on the form domain for
$c>\|\mathrm{Neg}_T\|/(1-\lambda_0)$, while both agree with $\sigma$ on the band.
\end{proposition}
\begin{proof}
(i) $\mathrm{Neg}_T=(R_T\sigma_-^{1/2}(D))(\sigma_-^{1/2}(D)R_T)$ is a product of two Hilbert--Schmidt operators,
since $\sigma_-$ is bounded and compactly supported; its kernel $k(t-s)$, $k=\frac1{2\pi}\int\sigma_-e^{i\tau\cdot}$, is
continuous, and the trace is $\int_{\IT}k(0)\,dt$, as for $\operatorname{tr}M_-$ in \S\ref{subsec:negindex}.
(ii)--(iv) are the split $\mathfrak a_T=\mathrm{Pos}_T-\mathrm{Neg}_T$ and $\mathrm{Pos}_T\ge0$.
(v) $S=S_{T,\tau_{\mathrm{last}}}$ is compact, $0\le S\le1$, and $\lambda_0=\|S\|$ is attained at an eigenvector
$F_0$; $\lambda_0=1$ would give $\int_{|\tau|>\tau_{\mathrm{last}}}|\Fh_0|^2=0$, impossible for the entire
function $\Fh_0\ne0$. Hence $\frac1{2\pi}\int_{\mathrm{off}}|\Fh|^2=\|F\|^2-\langle SF,F\rangle\ge(1-\lambda_0)\|F\|^2$. For the
first symbol the form is $\le s_+\langle SF,F\rangle-c(1-\lambda_0)\|F\|^2\le(s_+-c(1-\lambda_0))\|F\|^2$; for the second
it is $\mathrm{Pos}_T(F)-\langle\mathrm{Neg}_TF,F\rangle+\frac c{2\pi}\int_{\mathrm{off}}|\Fh|^2\ge(c(1-\lambda_0)-\|\mathrm{Neg}_T\|)\|F\|^2$.
Both signs hold on all of $L^2(\IT)$, hence on $\PT$. The statement is about symbols, not about $\zeta$.
\end{proof}

So the band $[0,\tau^{\mathrm{lo}}]$ carries all the negativity of $\mathfrak a_T$, it is finite and of
explicit length, and beyond $\tau_a\approx23.3$ its negative part sits at the near-solutions of
$2^l\approx3^k$ in the two computed regimes; this makes the heuristic alignment scale of
\S\ref{subsec:where} precise at those regimes. Positivity, however, is a global balance of $\mathrm{Pos}_T$
against $\mathrm{Neg}_T$, which the band alone does not decide.

\section{The complete form at a single window}\label{subsec:alignineq}

\begin{definition}[the alignment inequality]\label{def:Al}
For $w\in\Wcl$ let $\mathrm{Al}_T(w)$ be the statement: for every $\tau\in\R$,
\[
-2\Rel\sum_{n\ge2}\frac{\Lambda(n)w(\log n)}{n^{\frac12+i\tau}}\ \ge\
-\Big[C_0+\int_0^\infty\big(1-w(u)\cos\tau u\big)K(u)\,du\Big]-2\int_\R w(u)\cosh\tfrac u2\cos(\tau u)\,du .
\]
The left side is the smoothed prime sum on the critical line, $-\sum_nv_nw(\log n)\cos(\tau\log n)$; by
\eqref{eq:fw}, $\mathrm{Al}_T(w)$ is exactly $f_w\ge0$ on $\R$.
\end{definition}

\begin{lemma}[the alignment inequality lives on a band]\label{lem:Alband}
For $w\in\Wcl$ let $\tau_{\mathrm{env}}(w)$ be the unique positive root of
$m_{\mathrm{lo}}(\tau)-\Aconst-\|\nu_w\|_{\TV}-2k_0(\tau)=0$ (a definite number, in place of the
$2\pi e^{\Aconst+\|\nu\|_{\TV}+o(1)}$ of Proposition~\ref{prop:env}). Then $f_w\ge0$ on $[\tau_{\mathrm{env}}(w),\infty)$,
so $\mathrm{Al}_T(w)$ is equivalent to its restriction to $|\tau|\le\tau_{\mathrm{env}}(w)$.
\end{lemma}
\begin{proof}
$m_{\mathrm{lo}}$ increases and $k_0$ decreases (Lemmas~\ref{lem:mbounds}, \ref{lem:gbounds}), so the root is
unique; Proposition~\ref{prop:env}(i) for the feasible triple $(\nu_w,2,0)$ gives
$f_w\ge m_{\mathrm{lo}}-\Aconst-\|\nu_w\|_{\TV}-2k_0$, and $f_w$ is even.
\end{proof}

\begin{lemma}[non-attainment for the complete form]\label{lem:Onat}
Let $T\ge\frac12\log2$ and $w\in\Wcl$. Then $\inf_\R f_w<\lamW$.
\end{lemma}
\begin{proof}
Weak duality gives $c:=\inf f_w\le\lamW$ (the paragraph after Definition~\ref{def:W}). Suppose $c=\lamW$. The
bottom of $W_T$ on the unit sphere of $L^2(\IT)$ is attained at some $F_\star$ (\cite[Thm.~3]{Bombieri2000},
or the compactness argument in the proof of Theorem~\ref{thm:O}, which uses only that the symbol tends to
$+\infty$). Then $\varrho:=\sT+\nuh_w+2\gh0-\lamW\ge0$ is continuous, and
$\frac1{2\pi}\int\varrho|\Fh_\star|^2=W_T(F_\star)-\lamW=0$ by Proposition~\ref{prop:P}, since
$\frac1{2\pi}\int\nuh_w|\Fh_\star|^2=\int c_{F_\star}d\nu_w=0$. As $\Fh_\star$ is entire and nonzero,
$\varrho\equiv0$ on $\R$. Taking inverse Fourier transforms on $0<|u|<2T$, where $\nu_w=0$ and $g_0=\cosh\frac u2$,
and where the constant contributes nothing: $2\cosh\frac u2=\frac12K(|u|)+\sum_{n\in\NT}\frac{v_n}2(\delta_{\log n}+\delta_{-\log n})$
there, impossible since $K(u)\sim1/u$ as $u\downarrow0$ while $\cosh$ is bounded.
\end{proof}

\begin{theorem}[the alignment inequality is $\lamW>0$]\label{thm:Al}
\emph{(a)} If $\mathrm{Al}_T(w)$ holds for some $w\in\Wcl$, regular or not, then $W_T\ge0$ on the form domain of
$L^2(\IT)$; hence $\lamW\ge0$ and $\mathrm{WP}_T$.
\emph{(b)} Let $T\ge\frac12\log2$. Then $[\exists\,w\in\Wcl:\ \mathrm{Al}_T(w)]\iff\lamW>0$, and $w$ can then be
taken in $\Wcl^{\mathrm{reg}}$ with $\inf f_w>0$. In particular, if $\lamW=0$ no $w\in\Wcl$ satisfies
$\mathrm{Al}_T(w)$. Likewise, if $\lamP>0$ there are $w\in\Wcl^{\mathrm{reg}}$ and $\beta\in\R^2$ with
$\inf_\tau[f_w+(\beta_0-2)\gh0+\beta_1\gh1]>0$.
\emph{(c)} $[\forall T\ge\frac12\log2\ \exists\,w\in\Wcl^{\mathrm{reg}}:\ \mathrm{Al}_T(w)]\iff\mathrm{RH}\iff[\forall T:\ \mathrm{WP}_T]$,
through the chain
\[
\forall T\ \mathrm{WP}_T\ \Longrightarrow\ \mathrm{RH}\ \Longrightarrow\ \forall T\ge\tfrac12\log2:\ \lamW>0
\ \Longrightarrow\ \forall T\ge\tfrac12\log2\ \exists\,w:\ \mathrm{Al}_T(w)\ \Longrightarrow\ \forall T\ \mathrm{WP}_T .
\]
\end{theorem}
\begin{proof}
(a) $(\nu_w,2,0)$ is dual feasible and $W_T\ge\inf f_w$ on $L^2(\IT)$ (the paragraph after
Definition~\ref{def:W}), and $W_T=\mathfrak a_T$ on $\PT$. (b) ``$\Leftarrow$'' and the regular choice:
Theorem~\ref{thm:Zexh}(i). ``$\Rightarrow$'': $\lamW\ge\inf f_w\ge0$ by weak duality, and $\lamW>\inf f_w$ by
Lemma~\ref{lem:Onat}. The last sentence is Theorem~\ref{thm:Zexh}(ii). (c) The first arrow is
\cite[Thm.~11.3]{Paper45} with Weil's criterion; the second is the Jensen/exponential-type step in the proof
of Corollary~\ref{cor:RH}; the third is (b); the fourth is (a) for $T\ge\frac12\log2$ and
Lemma~\ref{lem:monoT} below it.
\end{proof}

\emph{The per-window converse.} At a single window, the converse of Theorem~\ref{thm:Al}(a) is the
implication $\mathrm{WP}_T\Rightarrow\lamW\ge0$. It reduces to the signs of two numbers. Fix $T$. We work
with real functions. This suffices for bottoms and negative indices, as in the remark after
Lemma~\ref{lem:parity}: for $F=u+iv$ with $u,v$ real, $W_T(F)=W_T(u)+W_T(v)$ and
$\mathfrak a_T(F)=\mathfrak a_T(u)+\mathfrak a_T(v)$. Let $\mathcal D_T$ be the form domain of \eqref{eq:form}
and $\mathcal D^{\mathrm{par}}_T$ its elements of parity $\mathrm{par}\in\{\mathrm{ev},\mathrm{odd}\}$. Put
$\varphi_{\mathrm{ev}}=\cosh\frac t2\,\mathbf 1_{\IT}$, $\varphi_{\mathrm{odd}}=\sinh\frac t2\,\mathbf 1_{\IT}$,
$\ell_{\mathrm{par}}(F)=\langle F,\varphi_{\mathrm{par}}\rangle$, $\varepsilon_{\mathrm{ev}}=+1$ and
$\varepsilon_{\mathrm{odd}}=-1$. Then $A=\ell_{\mathrm{ev}}+\ell_{\mathrm{odd}}$ and
$B=\ell_{\mathrm{ev}}-\ell_{\mathrm{odd}}$ on real $F$ (the parity paragraph after
Proposition~\ref{prop:nodes}). Let $\mathcal Q^{\mathrm{par}}_T:=\PT^{\mathrm{par}}\cap\mathcal D_T
=\ker(\ell_{\mathrm{par}}|\mathcal D^{\mathrm{par}}_T)$, and write $\lambda_{\min}(A_T|\PT^{\mathrm{par}})$ for the
bottom of $\mathfrak a_T$ on its unit sphere. This is the bottom with the primes, not the prime-free
$\lambda_{\mathrm{par}}$ of \S\ref{subsec:arch}. Write $\lambda^{\mathrm W}_{T,\mathrm{par}}$ for the bottom of
$W_T$ on the unit sphere of $\mathcal D^{\mathrm{par}}_T$. As in Lemma~\ref{lem:parity}, $n_-$ is the largest
dimension of a subspace of the form domain on which the form is negative definite. The letters $h$,
$\mathfrak g$, $\mathfrak G$, $J$, $S^{\mathrm{Sch}}$, $\mathcal N$ are local to this section.

\begin{theorem}[the parity Schur complement of the complete form]\label{thm:parityschur}
Fix $T>0$.
\emph{(o) Parity splitting.} On real $F=F_{\mathrm e}+F_{\mathrm o}\in\mathcal D_T$,
\[
W_T(F)=\big[\mathfrak a_T(F_{\mathrm e})+2\ell_{\mathrm{ev}}(F_{\mathrm e})^2\big]
+\big[\mathfrak a_T(F_{\mathrm o})-2\ell_{\mathrm{odd}}(F_{\mathrm o})^2\big],
\]
the sum being orthogonal for $L^2$, for $\mathfrak a_T$ and for $W_T$. Hence
$\lamA=\min_{\mathrm{par}}\lambda_{\min}(A_T|\PT^{\mathrm{par}})$ and $\lamW=\min_{\mathrm{par}}\lambda^{\mathrm W}_{T,\mathrm{par}}$.

\emph{(a) Non-degenerate sector.} Suppose $0$ is not in the spectrum of the self-adjoint operator of the
closed form $\mathfrak a_T|\mathcal Q^{\mathrm{par}}_T$. This holds in particular if
$\lambda_{\min}(A_T|\PT^{\mathrm{par}})>0$. Then there is a unique $h_{\mathrm{par}}\in\mathcal D^{\mathrm{par}}_T$
with $\ell_{\mathrm{par}}(h_{\mathrm{par}})=1$ and $\mathfrak a_T(h_{\mathrm{par}},G)=0$ for all
$G\in\mathcal Q^{\mathrm{par}}_T$. Put $\mathfrak g_{\mathrm{par}}(T):=W_T(h_{\mathrm{par}})=\mathfrak a_T(h_{\mathrm{par}})+2\varepsilon_{\mathrm{par}}$.
Then:
\begin{itemize}
\item[(a1)] $\mathcal D^{\mathrm{par}}_T=\mathcal Q^{\mathrm{par}}_T\oplus\R h_{\mathrm{par}}$, orthogonally for
$\mathfrak a_T$ and for $W_T$, and $W_T(G+xh_{\mathrm{par}})=\mathfrak a_T(G)+\mathfrak g_{\mathrm{par}}x^2$ for
$G\in\mathcal Q^{\mathrm{par}}_T$, $x\in\R$;
\item[(a2)] $n_-(W_T|\mathcal D^{\mathrm{par}}_T)=n_-(\mathfrak a_T|\mathcal Q^{\mathrm{par}}_T)+[\mathfrak g_{\mathrm{par}}<0]$.
The radical of $W_T$ on $\mathcal D^{\mathrm{par}}_T$ is $\R h_{\mathrm{par}}$ if $\mathfrak g_{\mathrm{par}}=0$ and $0$
otherwise;
\item[(a3)] if $\lambda_{\min}(A_T|\PT^{\mathrm{par}})>0$, then
$\mathfrak a_T(h_{\mathrm{par}})=\min\{\mathfrak a_T(F):F\in\mathcal D^{\mathrm{par}}_T,\ \ell_{\mathrm{par}}(F)=1\}$ and
$\mathfrak g_{\mathrm{par}}=\min\{W_T(F):F\in\mathcal D^{\mathrm{par}}_T,\ \ell_{\mathrm{par}}(F)=1\}$;
\item[(a4)] $\lambda^{\mathrm W}_{T,\mathrm{par}}\le\mathfrak g_{\mathrm{par}}/\|h_{\mathrm{par}}\|^2$ and
$\|h_{\mathrm{par}}\|\ge1/\|\varphi_{\mathrm{par}}\|$. If $\lambda_{\min}(A_T|\PT^{\mathrm{par}})>0$ \emph{and}
$\mathfrak g_{\mathrm{par}}>0$, then
$\lambda^{\mathrm W}_{T,\mathrm{par}}\ge\frac12\min\big(\lambda_{\min}(A_T|\PT^{\mathrm{par}}),\mathfrak g_{\mathrm{par}}/\|h_{\mathrm{par}}\|^2\big)$.
If in addition $2\lambda^{\mathrm W}_{T,\mathrm{par}}<\lambda_{\min}(A_T|\PT^{\mathrm{par}})$, then
$\lambda^{\mathrm W}_{T,\mathrm{par}}\le\mathfrak g_{\mathrm{par}}/\|h_{\mathrm{par}}\|^2\le2\lambda^{\mathrm W}_{T,\mathrm{par}}$.
\end{itemize}

\emph{(b) Degenerate sector.} Suppose $\lambda_{\min}(A_T|\PT^{\mathrm{par}})=0$, and let
$\mathcal N=\ker(\mathfrak a_T|\mathcal Q^{\mathrm{par}}_T)$. It is nonzero and finite-dimensional. For each
$F_0\in\mathcal N$ there is a unique $\alpha(F_0)\in\R$ with $\mathfrak a_T(F_0,G)=\alpha(F_0)\ell_{\mathrm{par}}(G)$
for all $G\in\mathcal D^{\mathrm{par}}_T$. If $\alpha\not\equiv0$ on $\mathcal N$, then $W_T$ is unbounded below on
$\mathcal D^{\mathrm{par}}_T$, so $\lamW<0$. If $\alpha\equiv0$ on $\mathcal N$ (the \emph{range condition}), then
$\mathcal N$ lies in the radical of $W_T$ on $\mathcal D^{\mathrm{par}}_T$ and (a) applies on the $L^2$-orthogonal
complement of $\mathcal N$ in $\mathcal D^{\mathrm{par}}_T$. With the $h_{\mathrm{par}},\mathfrak g_{\mathrm{par}}$ so obtained,
$W_T\ge0$ on $\mathcal D^{\mathrm{par}}_T$ if and only if $\mathfrak g_{\mathrm{par}}\ge0$.

\emph{(c) Assembled.} Put $\mathfrak G_T:=\mathrm{diag}(\mathfrak g_{\mathrm{ev}}(T),\mathfrak g_{\mathrm{odd}}(T))$. It is
defined whenever each sector is as in (a), or as in (b) with the range condition. If $0$ is not in the
spectrum of the operator of $\mathfrak a_T|(\PT\cap\mathcal D_T)$, then
$n_-(W_T)=n_-(\mathfrak a_T|\PT)+n_-(\mathfrak G_T)$, and the radical of $W_T$ has dimension $n_0(\mathfrak G_T)$. At every window,
without further hypotheses,
\[
\lamW>0\iff\big[\lamA>0\ \text{and}\ \mathfrak G_T\succ0\big],
\]
\[
\lamW\ge0\iff\big[\mathrm{WP}_T,\ \text{the range condition in every sector with }
\lambda_{\min}(A_T|\PT^{\mathrm{par}})=0,\ \text{and}\ \mathfrak G_T\succeq0\big].
\]
In particular, if $\lamA>0$, then $\lamW\ge0$ if and only if $\mathfrak g_{\mathrm{ev}}(T)\ge0$ and
$\mathfrak g_{\mathrm{odd}}(T)\ge0$.
\end{theorem}
\begin{proof}
(o) The symbol $\sT$ is even. For real even $F_{\mathrm e}$ and real odd $F_{\mathrm o}$, $\widehat{F_{\mathrm e}}$ is even and
$\widehat{F_{\mathrm o}}$ is odd, so $\mathfrak a_T(F_{\mathrm e},F_{\mathrm o})=0$ as in the proof of
Lemma~\ref{lem:parity}. Also $\ell_{\mathrm{ev}}$ vanishes on odd functions and $\ell_{\mathrm{odd}}$ on even ones, so
$A(F)B(F)=(\ell_{\mathrm{ev}}(F_{\mathrm e})+\ell_{\mathrm{odd}}(F_{\mathrm o}))(\ell_{\mathrm{ev}}(F_{\mathrm e})-\ell_{\mathrm{odd}}(F_{\mathrm o}))
=\ell_{\mathrm{ev}}(F_{\mathrm e})^2-\ell_{\mathrm{odd}}(F_{\mathrm o})^2$, with no cross term. Polarizing gives the
orthogonality. $F\in\PT$ if and only if $\ell_{\mathrm{ev}}(F_{\mathrm e})=\ell_{\mathrm{odd}}(F_{\mathrm o})=0$, and bottoms over an
orthogonal sum take the minimum.

(a) Write $\mathcal Q=\mathcal Q^{\mathrm{par}}_T$, $\ell=\ell_{\mathrm{par}}$ and $\varepsilon=\varepsilon_{\mathrm{par}}$. On
$\mathcal D^{\mathrm{par}}_T$, $W_T(F,G)=\mathfrak a_T(F,G)+2\varepsilon\,\ell(F)\ell(G)$ by (o).

\emph{Solvability.} First, $\mathfrak a_T=\Gamma_T+\PiT$ with $\|\PiT\|\le\Aconst$ (\S\ref{subsec:arch}). The form
$\Gamma_T\ge C_0\|\cdot\|^2$ is closed, with domain $\mathcal D_T$ (Theorem~\ref{thm:A}), and this domain embeds
compactly in $L^2(\IT)$ (proof of Proposition~\ref{prop:nodes}(i)). So $\mathfrak a_T$ is closed and bounded
below by $-M$, where $M:=\Aconst-C_0$. Second, $\mathcal Q$ is closed in the form norm, because $\ell$ and the
parity projection are $L^2$-continuous. Third, $\mathcal Q$ is dense in
$H:=\ker\ell\cap L^2_{\mathrm{par}}(\IT)$. Indeed, let $x\in H$, let $d_n\in\mathcal D^{\mathrm{par}}_T$ with $d_n\to x$ in
$L^2$ (smooth functions of parity $\mathrm{par}$ are dense), and let $e\in\mathcal D^{\mathrm{par}}_T$ with $\ell(e)=1$.
Then $d_n-\ell(d_n)e\in\mathcal Q$ tends to $x$. Hence $\mathfrak a_T|\mathcal Q$ is a closed, semibounded, densely
defined form in $H$ with compact resolvent. Its spectrum is discrete, and its negative spectral subspace
$\mathcal Q_-$ is finite-dimensional and $\mathfrak a_T$-orthogonal to $\mathcal Q_+:=\mathcal Q\cap\mathcal Q_-^\perp$. If
$0$ is not in the spectrum, then $\mathfrak a_T\ge\delta\|\cdot\|^2$ on $\mathcal Q_+$ for some $\delta>0$. Hence
$\mathfrak a_T(x)\ge\frac{\delta}{\delta+M+1}\big(\mathfrak a_T(x)+(M+1)\|x\|^2\big)$ there, so $\mathfrak a_T$ is coercive
on $\mathcal Q_+$ in the form norm and the Lax--Milgram lemma applies on $\mathcal Q_+$. On $\mathcal Q_-$ the form is a
negative definite finite matrix. So for every form-bounded linear functional $f$ on $\mathcal Q$ there is a unique
$Y\in\mathcal Q$ with $\mathfrak a_T(Y,\cdot)=f$ on $\mathcal Q$. Uniqueness holds because the operator has trivial
kernel. Now fix $H_0\in\mathcal D^{\mathrm{par}}_T$ with $\ell(H_0)=1$, for instance a smooth bump of parity
$\mathrm{par}$. The functional $f=\mathfrak a_T(H_0,\cdot)$ is form-bounded, by Cauchy--Schwarz for the nonnegative
form $\mathfrak a_T+M\|\cdot\|^2$. Put $h:=H_0-Y$. Then $\ell(h)=1$ and $\mathfrak a_T(h,\cdot)=0$ on $\mathcal Q$. Two such
$h$ differ by an element of $\mathcal Q$ that is $\mathfrak a_T$-orthogonal to $\mathcal Q$, which is $0$.

(a1) Any $F\in\mathcal D^{\mathrm{par}}_T$ decomposes as $F=(F-\ell(F)h)+\ell(F)h$ with $F-\ell(F)h\in\mathcal Q$. For
$G\in\mathcal Q$, $W_T(h,G)=\mathfrak a_T(h,G)+2\varepsilon\ell(h)\ell(G)=0$. Expand, using $\ell(G)=0$ and $\ell(h)=1$.

(a2) By (a1), $\mathcal D^{\mathrm{par}}_T=(\mathcal Q_-\oplus[\R h\text{ if }\mathfrak g<0])\oplus(\mathcal Q_+\oplus[\R h\text{ if }\mathfrak g\ge0])$,
orthogonally for $W_T$. The first summand is negative definite and has the stated dimension. The second
summand carries $W_T\ge0$. A negative definite subspace meets the second summand only in $0$, so its
projection onto the first summand along the second is injective. For the radical: if
$W_T(G+xh,\cdot)=0$, then $\mathfrak a_T(G,\cdot)=0$ on $\mathcal Q$, so $G=0$, and then $\mathfrak g x=0$.

(a3) If the bottom is positive, then $\mathfrak a_T(G+h)=\mathfrak a_T(G)+\mathfrak a_T(h)\ge\mathfrak a_T(h)$ and
$W_T(G+h)=\mathfrak a_T(G)+\mathfrak g\ge\mathfrak g$.

(a4) The upper bound is the Rayleigh quotient at $h$, and $1=\ell(h)\le\|h\|\|\varphi_{\mathrm{par}}\|$. Put
$\lambda:=\lambda_{\min}(A_T|\PT^{\mathrm{par}})$ and assume $\lambda>0$ and $\mathfrak g>0$. With
$\mu:=\min(\lambda,\mathfrak g/\|h\|^2)$,
\[
W_T(G+xh)=\mathfrak a_T(G)+\mathfrak gx^2\ge\lambda\|G\|^2+\frac{\mathfrak g}{\|h\|^2}\|xh\|^2
\ge\mu\big(\|G\|^2+\|xh\|^2\big)\ge\tfrac12\mu\|G+xh\|^2 .
\]
If also $2\lambda^{\mathrm W}_{T,\mathrm{par}}<\lambda$, then $\frac12\mu\le\lambda^{\mathrm W}_{T,\mathrm{par}}<\frac12\lambda$ forces
$\mu=\mathfrak g/\|h\|^2$. The hypothesis $\mathfrak g>0$ is used in the second inequality and cannot be dropped.

(b) The bottom $0$ is attained (compact resolvent), so $\mathcal N\ne0$, and $\mathcal N$ is finite-dimensional. For
$F_0\in\mathcal N$ and $G\in\mathcal Q$, Cauchy--Schwarz for the nonnegative form $\mathfrak a_T|\mathcal Q$ gives
$\mathfrak a_T(F_0,G)=0$. So $\mathfrak a_T(F_0,\cdot)$ is a linear functional on $\mathcal D^{\mathrm{par}}_T$ vanishing on
$\ker\ell$, hence a multiple $\alpha(F_0)\ell$. Since $\ell(F_0)=0$, we get $W_T(F_0,H)=\alpha(F_0)\ell(H)$ and
$W_T(F_0)=0$. So $W_T(xF_0+H)=W_T(H)+2x\alpha(F_0)\ell(H)$, which is unbounded below in $x$ if
$\alpha(F_0)\ne0$ and $\ell(H)=1$. Now suppose $\alpha\equiv0$. Then $\mathcal N$ is $W_T$-orthogonal to
$\mathcal D^{\mathrm{par}}_T$. Put $\mathcal D':=\mathcal D^{\mathrm{par}}_T\cap\mathcal N^\perp$ and $\mathcal Q':=\mathcal Q\cap\mathcal N^\perp$.
Then $\mathcal Q'=\ker(\ell|\mathcal D')$, and $\ell\ne0$ on $\mathcal D'$: subtract from $H_0$ its $L^2$-projection onto
$\mathcal N\subset\mathcal Q$, which $\ell$ annihilates. On $\mathcal Q'$ the form is bounded below by the next eigenvalue,
which is $>0$. So (a) applies with $\mathcal D',\mathcal Q'$ in place of $\mathcal D^{\mathrm{par}}_T,\mathcal Q$. Moreover
$\mathcal D^{\mathrm{par}}_T=\mathcal N\oplus\mathcal D'$, $W_T$-orthogonally. Conclude by (a1).

(c) By (o), everything adds over the two sectors, and the inertia formula is (a2) in each sector. If
$\lamW>0$, then $\lamA\ge\lamW>0$ (recall $\lamW\le\lamA$), so both sectors are as in (a), and
$\mathfrak g_{\mathrm{par}}=W_T(h_{\mathrm{par}})\ge\lamW\|h_{\mathrm{par}}\|^2>0$. The converse is (a4) in each sector. If $\lamW\ge0$,
then $\lamA\ge\lamW\ge0$, which is $\mathrm{WP}_T$. In a sector with bottom $0$, (b) excludes $\alpha\not\equiv0$, and
in every sector $\mathfrak g_{\mathrm{par}}=W_T(h_{\mathrm{par}})\ge0$. Conversely, (a1) and $\mathfrak a_T\ge0$ on
$\mathcal Q^{\mathrm{par}}_T$ give $W_T\ge0$ on a sector with positive bottom when $\mathfrak g_{\mathrm{par}}\ge0$, and (b) does the
same on a sector with bottom $0$.
\end{proof}

The range condition says that the column $\rho(G):=\mathfrak a_T(H_0,G)$ of the bordered form, with
$\ell_{\mathrm{par}}(H_0)=1$, vanishes on $\ker(\mathfrak a_T|\mathcal Q^{\mathrm{par}}_T)$, since $\rho(F_0)=\alpha(F_0)$. This is the
range condition of a bordered Schur complement. Because $\mathfrak a_T|\mathcal Q^{\mathrm{par}}_T$ has a spectral gap above
$0$, it is equivalent to $\rho$ lying in the range of the square root of the operator. The theorem takes
the complement over $\PT$ and not over $L^2(\IT)$. So only $\mathfrak a_T|\PT$ has to be invertible, which is
the natural hypothesis for a converse whose premise is on the primitive side.

\begin{proposition}[$\mathfrak G_T$ is the polar capacity]\label{prop:polarcap}
Assume $\mathrm{WP}_T$, and the range condition of Theorem~\ref{thm:parityschur}(b) in every
sector with $\lambda_{\min}(A_T|\PT^{\mathrm{par}})=0$. Let
$S^{\mathrm{Sch}}_T:=\mathrm{diag}\big(\mathfrak a_T(h_{\mathrm{ev}}),\mathfrak a_T(h_{\mathrm{odd}})\big)$. This is the compression of
$\mathfrak a_T$ to $\mathrm{span}(h_{\mathrm{ev}},h_{\mathrm{odd}})$, the $\mathfrak a_T$-orthogonal complement of the primitive part,
i.e.\ the Schur complement of $\mathfrak a_T$ over $\PT$. Let $J:=\mathrm{diag}(2,-2)$. Then
\[
\mathfrak G_T=S^{\mathrm{Sch}}_T+J .
\]
For $(\beta,\alpha)\in\R^2$ put
$E_T(\beta,\alpha):=\inf\{W_T(F):\ F\in\mathcal D_T\ \text{real},\ \Phi_F(0)=B(F)=\beta,\ \Phi_F(1)=A(F)=\alpha\}$. Then
\[
E_T(\beta,\alpha)=(\alpha,\beta)\,\mathfrak G^{AB}_T\,(\alpha,\beta)^{\mathsf T},\qquad
\mathfrak G^{AB}_T=\frac14\begin{pmatrix}\mathfrak g_{\mathrm{ev}}+\mathfrak g_{\mathrm{odd}}&\mathfrak g_{\mathrm{ev}}-\mathfrak g_{\mathrm{odd}}\\
\mathfrak g_{\mathrm{ev}}-\mathfrak g_{\mathrm{odd}}&\mathfrak g_{\mathrm{ev}}+\mathfrak g_{\mathrm{odd}}\end{pmatrix},
\]
and the infimum is attained. In these coordinates $J$ becomes $\big(\begin{smallmatrix}0&1\\1&0\end{smallmatrix}\big)$,
the form $2\alpha\beta=2\Phi_F(0)\Phi_F(1)$, with the one negative square of Lemma~\ref{lem:polar}. Hence
$\lamW\ge0$ if and only if $E_T\ge0$ on $\R^2$. Under $\mathrm{WP}_T$, the complete form is nonnegative
exactly when the least complete Weil energy needed to prescribe the values of $\Phi_F$ at the pole
$s\in\{0,1\}$ is nonnegative.
\end{proposition}
\begin{proof}
$\mathfrak G_T=S^{\mathrm{Sch}}_T+J$ is the definition $\mathfrak g_{\mathrm{par}}=\mathfrak a_T(h_{\mathrm{par}})+2\varepsilon_{\mathrm{par}}$. In the
degenerate case with the range condition, $\mathfrak a_T(h_{\mathrm{par}})$ does not depend on the choice of $h_{\mathrm{par}}$
modulo $\mathcal N$. For real $F$, $\Phi_F(1)=A(F)$ and $\Phi_F(0)=B(F)$ (proof of Lemma~\ref{lem:roughEF}). The
constraints are $\ell_{\mathrm{ev}}(F_{\mathrm e})=c:=\frac{\alpha+\beta}2$ and $\ell_{\mathrm{odd}}(F_{\mathrm o})=s:=\frac{\alpha-\beta}2$,
and they split over the sectors, as does $W_T$ by Theorem~\ref{thm:parityschur}(o). In each sector,
Theorem~\ref{thm:parityschur}(a1) (or (b)) gives $W_T(G+xh_{\mathrm{par}})=\mathfrak a_T(G)+\mathfrak g_{\mathrm{par}}x^2$ with $x$ the
prescribed value of $\ell_{\mathrm{par}}$ and $\mathfrak a_T(G)\ge0$ by $\mathrm{WP}_T$. So the sector minimum is
$\mathfrak g_{\mathrm{par}}x^2$, attained at $G=0$. Hence $E_T=\mathfrak g_{\mathrm{ev}}c^2+\mathfrak g_{\mathrm{odd}}s^2$, which is the displayed
matrix. The form $2c^2-2s^2=2\alpha\beta$ gives $J$ in the new coordinates. The last statement is
Theorem~\ref{thm:parityschur}(c), since $E_T\ge0$ if and only if $\mathfrak G_T\succeq0$.
\end{proof}

The converse is therefore the pair of inequalities $\mathfrak a_T(h_{\mathrm{odd}})\ge2$ and
$\mathfrak a_T(h_{\mathrm{ev}})\ge-2$. In the odd direction the primitive Schur complement must beat the negative
polar square. In the even direction the positive polar square may absorb a negative Schur value down
to $-2$. Neither is automatic. The next corollary shows that the pair is not implied by $\mathrm{WP}_T$ at the level of combs. Its witnesses are odd, so it is the odd inequality that is refuted.
It uses the comb-level form $W_T(a)$ of \S\ref{subsec:finite}. Put $F^{\mathrm{odd}}_T(t)=t$ on $\IT$ as in
Proposition~\ref{prop:archneg}.

\begin{corollary}[the converse fails for modified weights]\label{cor:convfail}
The proof uses Theorem~\ref{thm:index}, Corollary~\ref{cor:Tarch} and Theorem~\ref{thm:C}, whose proofs include
computer verifications, and two further evaluations in interval arithmetic. The clause $\theta_*<1$ in (ii)
uses \cite[Thm.~1.2]{Zhu2026}.
\emph{(i) Deletion.} Let $a=\Lambda\mathbf 1_{n\ge5}$. For every $T\in[0.4608,0.58]$, $\mathrm{WP}_T(a)$ holds
strictly while $\lambda^{\mathrm W}_T(a)<0$:
\[
\lambda_{\min}(A_T(a)|\PT)=\laminfP(T)>0,\qquad\lambda^{\mathrm W}_T(a)=\laminf(T)<0 .
\]
On $[0.4608,\log2]$ the odd-sector implication fails. There the bottom of $\mathfrak a^{(v(a))}_T$ on $\PT^{\mathrm{odd}}$ is
$\lambda_{\mathrm{odd}}(T)>0$, while $W_T(a)$ is negative on an odd vector. On $(T^*_{\infty,\mathrm P},\log2]$,
$\mathrm{WP}_T(a)$ fails, so the full implication holds there vacuously and only its odd-sector form is refuted.
\emph{(ii) One contact.} Let $T=\frac12\log3$, so $\NT=\{2\}$, and for $\theta\in[0,1]$ let
$a_\theta:=\Lambda-(1-\theta)\Lambda(2)\mathbf 1_{n=2}$. Its comb on $\NT$ is $\theta v_2$. Then
$\lambda_{\min}(A_T(a_\theta)|\PT)>0$ for every $\theta\in[0,1]$, and $\lambda^{\mathrm W}_T(a_\theta)<0$ for every
$\theta\in[0,0.8148]$. The set of $\theta\in[0,1]$ with $\lambda^{\mathrm W}_T(a_\theta)\ge0$ is an interval
\emph{(iii) Consequences.} Let $\mathcal C$ be a class of nonnegative arithmetic functions containing $\Lambda$.
If $\mathcal C$ is closed under deletion of finitely many terms, it contains $\Lambda\mathbf 1_{n\ge5}$, for which
the implication $\mathrm{WP}_T(a)\Rightarrow\lambda^{\mathrm W}_T(a)\ge0$ fails at every $T\in[0.4608,0.58]$ by (i). If
$\mathcal C$ contains $a_\theta$ for some $\theta\le0.8148$ (for instance, if $\mathcal C$ is closed under lowering
one term within $[0,\Lambda(n)]$), the implication fails for that member at $T=\frac12\log3$ by (ii). The classes
$\mathcal C_{\mathrm{PNT}}(\varepsilon)$, $\mathcal C_{\mathrm{zeros}}$ and $\mathcal C_{\mathrm{BT}}$ of
Corollary~\ref{cor:void} are closed under deletion and under lowering one term. Consequently no argument
that uses about the primes only a property preserved by these modifications (for instance a prime number
theorem with error term, a zero-free region, a zero-density estimate, or an upper-bound sieve inequality)
can prove $\mathrm{WP}_T\Rightarrow\lamW\ge0$ at any $T\in[0.4608,0.58]$ or at $T=\frac12\log3$. At $T=\frac12\log3$
such an argument must distinguish $\Lambda(2)$ from every weight $\theta\Lambda(2)$ with $\theta\le0.8148$, that is,
from weights at least $18.5\%$ below it; nothing is claimed about weights above $\Lambda(2)$.
\end{corollary}
\begin{proof}
(i) For $T\le\log2<\frac12\log5$ one has $\NT\subseteq\{2,3,4\}$, where $a=0$. So the comb of $a$ is empty,
$W_T(a)=\Ainf$ and $\mathfrak a^{(v(a))}_T=\Gamma_T$ on $\PT$, as in the proof of Theorem~\ref{thm:finmod}(i). For
$T\le0.58$, $\laminfP(T)>0$ by Theorem~\ref{thm:index}(i), and for $T\ge0.4608$, $\laminf(T)<0$ by
Proposition~\ref{prop:archneg}. On $[0.4608,\log2]$, $\lambda_{\mathrm{odd}}(T)>0$ by Theorem~\ref{thm:index}(i). The
odd vector $F^{\mathrm{odd}}_{0.4608}$, extended by zero, lies in $L^2(\IT)$ with
$\Ainf(F^{\mathrm{odd}}_{0.4608})\le U(0.4608)<0$ ($\Ainf$ is $T$-free). $\mathrm{WP}_T(a)$ fails on
$(T^*_{\infty,\mathrm P},\log2]$ by Lemma~\ref{lem:strict}.

(ii) The comb: $\NT=\{n:\log n<\log3\}=\{2\}$ and $v_2(a_\theta)=\theta v_2$. So
$W_T(a_\theta)=\Ainf+\theta\PiT$, and $\mathfrak a^{(v(a_\theta))}_T=\Gamma_T+\theta\PiT$ on $\PT$.

\emph{The primitive side.} $\theta\mapsto\lambda_{\min}(A_T(a_\theta)|\PT)=\inf_{F\in\PT,\|F\|=1}[\Gamma_T(F)+\theta\PiT(F)]$
is an infimum of affine functions of $\theta$. It is finite, since $\|\PiT\|\le v_2$, and hence concave on
$\R$. At $\theta=0$ it is $\laminfP(\frac12\log3)>0$ (Corollary~\ref{cor:Tarch}). At $\theta=1$ it is
$\lamA\ge8.8577\cdot10^{-5}$ (Theorem~\ref{thm:C}, which is stated for the actual comb $v_2$ at this window). A
concave function on $[0,1]$ is at least the smaller of its endpoint values.

\emph{The complete side.} Let
$F_5(t)=t-\frac{2634299}{500000}t^3+\frac{1919331}{250000}t^5$ on $\IT$, with $T=\frac12\log3$. As in
Proposition~\ref{prop:archneg}, $\Ainf(F_5)\le U_0$. Here $U_0$ is obtained by bounding $K\le\frac1u+\frac12$
against $c_{F_5}(0)-c_{F_5}(u)\ge0$ (Cauchy--Schwarz) on $(0,2T)$, by Theorem~\ref{thm:S}(iii) on $[2T,\infty)$,
and by $2A(F_5)B(F_5)=-2s(F_5)^2$. The autocorrelation $c_{F_5}$ is an exact polynomial in $(u,T)$ with rational
coefficients. Evaluating in outward-rounded interval arithmetic at $50$ digits with $T=\frac12\log3$ gives
\begin{gather*}
U_0\in[-0.00613778443569004,\,-0.00613778443569003],\\
\PiT(F_5)=-v_2c_{F_5}(\log2)\in[0.00753283781432786,\,0.00753283781432787],
\end{gather*}
and $U_0+0.8148\,\PiT(F_5)\in[-2.81845757,-2.81845756]\cdot10^{-8}$. Since $\PiT(F_5)>0$, we get
$W_T(a_\theta)(F_5)=\Ainf(F_5)+\theta\PiT(F_5)\le U_0+0.8148\,\PiT(F_5)<0$ for every $\theta\le0.8148$. (The bare
witness $F^{\mathrm{odd}}_T$ gives the same for $\theta\le0.4741$.) Finally, $\theta\mapsto\lambda^{\mathrm W}_T(a_\theta)$
is concave and finite on $\R$ for the same reason, hence continuous. At $\theta=1$ it is
$\lamW\ge\lambda^{\mathrm W}_{0.8}>0$ by \cite[Thm.~1.2]{Zhu2026}, as in the proof of Corollary~\ref{cor:KEP08}. So
$\{\theta\in[0,1]:\lambda^{\mathrm W}_T(a_\theta)\ge0\}$ is a closed interval that contains $1$ and is disjoint from
$[0,0.8148]$. It has the form $[\theta_*,1]$, and $\theta_*<1$ by continuity.

(iii) By (i) and (ii), the implication is false for $\Lambda\mathbf 1_{n\ge5}$ (delete $n=2,3,4$) and for $a_\theta$.
For the closure properties: lowering one term keeps $0\le a\le\Lambda$; it changes $\psi_a$ by a constant for
$x\ge n$; and $\zeta_a=\zeta\exp\big(-(1-\theta)\Lambda(n)n^{-s}/\log n\big)$ with an entire zero-free factor, as in
Theorem~\ref{thm:finmod}(iii).
\end{proof}

As in Corollary~\ref{cor:void}, the form $W_T(a)$ is not the zero sum of $\zeta_a$ (\S\ref{subsec:finite}), so
$\lambda^{\mathrm W}_T(a)<0$ says nothing about the zeros of $\zeta_a$, which are those of $\zeta$; moreover $\zeta_a$
has no functional equation (Theorem~\ref{thm:finmod}(iii)). We have not examined a model with the
functional equation (Remark~\ref{rem:twohalves}). For the empty comb, (i) is exactly the strict separation
$T^*_\infty<T^*_{\infty,\mathrm P}$, proved by Proposition~\ref{prop:archneg} with Corollary~\ref{cor:Tarch}
(see also the numerical values of \S\ref{subsec:arch}).

\emph{The chain at one window.} For $T\ge\frac12\log2$, Theorem~\ref{thm:Al}(b),
Theorem~\ref{thm:parityschur}(c) and Theorem~\ref{thm:KEP3}(a) give
\begin{multline*}
\big[\exists\,w\in\Wcl:\ \mathrm{Al}_T(w)\big]\iff\lamW>0\iff\big[\lamA>0\ \text{and}\ \mathfrak G_T\succ0\big]\\
\Longrightarrow\ \mathrm{KEP}(T)\ \text{indeterminate}\ \Longrightarrow\ \lamW\ge0\\
\iff\big[\mathrm{WP}_T,\ \text{range condition},\ \mathfrak G_T\succeq0\big].
\end{multline*}
Indeterminacy of $\mathrm{KEP}(T)$ is not known to be equivalent to either $\lamW>0$ or $\lamW\ge0$. The two
ends of the chain can differ only on $\{T:\lamW=0\}$. So, at a single window, $\mathrm{WP}_T$ implies $\lamW\ge0$ if
and only if $\mathfrak G_T\succeq0$, together with the range condition in a degenerate sector; that is, the
converse of Theorem~\ref{thm:Al}(a) is the pair of scalar inequalities $\mathfrak g_{\mathrm{ev}}(T)\ge0$,
$\mathfrak g_{\mathrm{odd}}(T)\ge0$. By Corollary~\ref{cor:convfail} it fails for suitably modified weights at every
$T\in[0.4608,0.58]$ and at $\frac12\log3$. For $\zeta$ it holds at every $T\le0.8$, but only because its conclusion
$\lamW>0$ is known there by \cite[Thm.~1.2]{Zhu2026}, whence $\mathfrak G_T\succ0$ by
Theorem~\ref{thm:parityschur}(c).

We do not know whether, at a fixed window, $\mathrm{WP}_T$ implies $\lamP\ge0$; the inequalities
$\lamW\le\lamP\le\lamA$ are the only ones between the three bottoms proved at a fixed $T$. Nor do we know
whether the per-window converse holds for $\zeta$ at any $T>0.8$, i.e.\ whether
$\mathrm{WP}_T\Rightarrow\mathfrak G_T\succeq0$ (with the range condition) there; the only known route from
$\mathrm{WP}_T$ to $\lamW\ge0$ at such a window passes through all $T$ at once, via RH. Finally, we do not know an
effective window shift for $\zeta$, i.e.\ an implication $\mathrm{WP}_{T'}\Rightarrow\lamW\ge0$ with an explicit
$T'=T'(T)$. At the level of combs, a shift valid for every nonnegative comb needs $T'\ge T^*_{\infty,\mathrm P}$
whenever $T\ge0.4608$: for $a=\Lambda\mathbf 1_{n\ge N}$ with $N\ge\max(5,e^{2\max(T,T')})$ both combs are empty,
$\mathrm{WP}_{T'}(a)\iff\laminfP(T')\ge0$, and $\lambda^{\mathrm W}_T(a)=\laminf(T)<0$.

\begin{remark}[the per-window quantity]\label{rem:excess}
By Theorem~\ref{thm:Zexh}(i), $\sup_{w\in\Wcl^{\mathrm{reg}}}\inf_\tau f_w=\lamW$ for $T\ge\frac12\log2$: the
largest uniform margin with which the alignment inequality can hold is $\lamW$ itself, a number determined by $v_T(\Lambda)$ and $m$ (Proposition~\ref{prop:finite}), and RH is the
statement $\lim_{T\to\infty}\lamW\ge0$ (Corollary~\ref{cor:margin2}(iii)). The trivial envelope
$-\sum_nv_n|w(\log n)|$ of the prime sum certifies $\mathrm{Al}_T(w)$ wherever the right side of
Definition~\ref{def:Al} is at most it, and by Proposition~\ref{prop:env}(ii) the constant $\Aconst$ of
that envelope cannot be lowered; on the complement the actual phases $\tau\log p$ modulo $2\pi$ enter, and in
the two regimes of Proposition~\ref{prop:align} that complement is the set of near-solutions of
$2^l\approx3^k$. For $w\in\Wcl^{\mathrm{reg}}$, by Theorem~\ref{thm:Z}(ii), $\mathrm{Al}_T(w)$ at $\tau$ is the
inequality $\sum_\rho\wh(\tau-\gamma_\rho)\ge0$: the primes are not over-aligned there relative to the zero
density that the explicit formula assigns to them.
\end{remark}

\begin{remark}[the prime side of a primitive test is the Chebyshev remainder]\label{rem:Cheb}
By Theorem~\ref{thm:Z}, the prime-side inequality $f_w\ge0$ is a zero-side inequality: a proof of it from the
primes below $\tau_{\mathrm{env}}$ is a lower bound for the smoothed prime sum
$-2\Rel\sum_n\Lambda(n)n^{-\frac12-i\tau}w(\log n)$ on the critical line beyond what the envelope gives,
which is exactly the content that the explicit formula assigns to the zeros in that range. On $\PT$ this
has an exact form. Let $\psi_\Lambda(x)=\sum_{n\le x}\Lambda(n)$ be Chebyshev's function and
$r_\Lambda(u):=\psi_\Lambda(e^u)-e^u$. For real $F\in\PT$ in the form domain,
\begin{equation}\label{eq:Cheb}
\mathfrak a_T(F)=C_0\|F\|^2+\Lev(\cF)+2\int_0^{2T}\cF(u)e^{-u/2}\,du-2\int_{(0,2T)}e^{-u/2}\cF(u)\,dr_\Lambda(u).
\end{equation}
Indeed $\sum_{n\in\NT}v_n\cF(\log n)=2\int_{(0,2T)}e^{-u/2}\cF(u)\,d\psi_\Lambda(e^u)$, since
$u\mapsto\psi_\Lambda(e^u)$ jumps by $\Lambda(n)$ at $u=\log n$ and $\cF(2T)=0$; split
$d\psi_\Lambda(e^u)=e^udu+dr_\Lambda(u)$; the identity $\int_{-2T}^{2T}\cF(u)e^{u/2}du=A(F)B(F)=0$ on $\PT$
(substitute $s=t-u$) and the evenness of $\cF$ give $\int_0^{2T}\cF e^{u/2}du=-\int_0^{2T}\cF e^{-u/2}du$; insert
into \eqref{eq:Aform}. So on primitive tests the whole arithmetic of the comb is the last term of
\eqref{eq:Cheb}, and the unconditional control of $r_\Lambda$ is the zero-free-region theory. Heuristically,
a zero-free-region bound controls the last term only up to an error that grows with $T$,
whereas any $F\in\PT$ with $\mathfrak a_T(F)=o(\|F\|^2)$ (Corollary~\ref{cor:margin2}) needs
$2\int e^{-u/2}\cF\,dr_\Lambda$ to equal $C_0\|F\|^2+\Lev(\cF)+2\int_0^{2T}\cF e^{-u/2}du$ up to
$o(\|F\|^2)$; by Lemma~\ref{lem:roughEF} that precision is a statement about
$\sum_\rho\Phi_F(\rho)\Phi_F(1-\rho)$.
\end{remark}

\section{One direction at $T=\log2$}\label{subsec:schur}

The form $W_T$ on $\PT$ splits along any single direction by the classical Schur--Feshbach
congruence. We use the direction singled out by the archimedean block: at $T=\log2$ the archimedean
form has exactly one negative direction on $\PT$, and it is even (Theorem~\ref{thm:index}(iv); by
Corollary~\ref{cor:H1} the same count holds on all of $(T^*_{\infty,\mathrm P},\log2]$). The explicit pivot
$F^*$ below is a numerical approximation of the bottom vector of the even sector; nothing depends on the
quality of the approximation. The technique (Schur complement and parity splitting, with interval
enclosures of the entries) is taken from earlier unpublished work of the author, where it is applied to
a spectral cut-off; the pivot and the enclosure of its value are new.

\begin{proposition}[Schur reduction]\label{prop:schur}
Let $w$ be a real quadratic form on a form domain $Q\subseteq\PT$, $e\in Q$ a unit vector with
$w_{11}:=w[e]\ne0$, $V=e^\perp\cap\PT$, $Q_V=Q\cap V$, $\beta(G):=w(e,G)$, and let
$s[G]:=w[G]-\beta(G)^2/w_{11}$ be the Schur form on $Q_V$. Then, algebraically and for every
$F=xe+G$ with $x\in\R$, $G\in Q_V$,
\[
w[F]=w_{11}\Big(x+\frac{\beta(G)}{w_{11}}\Big)^2+s[G].
\]
Consequently \emph{(i)} if $w_{11}>0$ then $w\ge0$ on $Q$ if and only if $s\ge0$ on $Q_V$, and
$w\ge c\|\cdot\|^2$ implies $s\ge c\|\cdot\|^2$; the converse implication for the coercive form needs
$\beta$ bounded, $|\beta(G)|\le C\|G\|$, and then holds with a constant $c(c',w_{11},C)>0$.
\emph{(ii)} If $w\ge0$ on $Q$ then $w_{11}\ge0$ --- and no more: $w_{11}=0$ forces $\beta\equiv0$, and
then $w\ge0$ if and only if $w|_{Q_V}\ge0$.
\emph{(iii)} If $e\in D(W)$ for $W=\Ainf+\PiT$, then $b:=P_VWe\in V$, $\beta=\langle b,\cdot\rangle$
is bounded, $s$ is the form of $S=W|_V-(b\otimes b)/w_{11}$, and
$LWL^*=w_{11}\oplus S$ with $L=I-w_{11}^{-1}(b\otimes e)$ boundedly invertible: the reduction is a
congruence, so it preserves the sign of the form and not its eigenvalues.
Finally $W_T$ commutes with the reflection $(RF)(t)=F(-t)$ on $\PT$, so for an even pivot
$S=S_{\mathrm{even}}\oplus W_T|_{V_{\mathrm{odd}}}$ and the Schur correction lives in the even sector
alone.
\end{proposition}
\begin{proof}
The identity is the completion of the square, $x^2w_{11}+2x\beta(G)=w_{11}(x+\beta(G)/w_{11})^2-\beta(G)^2/w_{11}$;
it is pointwise on $Q$ and uses no topology. (i) $(\Leftarrow)$ both terms are $\ge0$;
$(\Rightarrow)$ take $x=-\beta(G)/w_{11}$, and in the coercive case
$s[G]=w[xe+G]\ge c\|xe+G\|^2\ge c\|G\|^2$. (ii) Take $G=0$, $x=1$; if $w_{11}=0$ and
$\beta(G_0)\ne0$ then $w[xe+G_0]=2x\beta(G_0)+w[G_0]\to-\infty$. (iii) Substitute
$\beta(G)=\langle We,G\rangle$ and read off the block matrix; $L^{-1}=I+w_{11}^{-1}(b\otimes e)$
because $\langle b,e\rangle=0$. Parity: $\cF$ is unchanged by $F\mapsto RF$, so $\Ainf$ and $\PiT$
are $R$-invariant, and $R$ swaps the constraints $A(F)=0$ and $B(F)=0$, so $\PT$ is $R$-invariant;
for an even $e$, $\beta$ vanishes on $V_{\mathrm{odd}}$.
\end{proof}

\begin{theorem}[a value at $T=\log2$]\label{thm:w11}
Let $T=\log2$ and let $F^*$ be the explicit even polynomial of degree $17$ in $\PT$ whose coefficients
in the orthonormal Legendre basis of $L^2(\IT)$ are given in the certificate data accompanying this
paper. Then
\[
W_{\log2}[F^*]\ \in\ [0.0342113430,\ 0.0342182500],
\]
in particular $W_{\log2}[F^*]>0.0342$. The proof includes a computer verification in interval arithmetic
(Appendix~\ref{app:encl}).
\end{theorem}
\begin{proof}
Everything is interval arithmetic on the L\'evy form of Theorem~\ref{thm:A},
$W_T[F]=2A(F)B(F)+C_0c(0)+\int_0^{2T}(c(0)-c(u))K(u)du+c(0)\int_{2T}^\infty K-\sum_{n\in\NT}v_nc(\log n)$,
$c=c_{F^*}$, $\NT=\{2,3\}$. The coefficient vector is read as exact decimal data; how it was
obtained is irrelevant to the statement. (a) $c$ is computed as an exact polynomial of degree
$\le35$ with interval coefficients, by rational arithmetic on the Legendre recurrence, the binomial
expansion of $p(x-s)$, the convolution with $p$ and term-by-term integration, with no quadrature and no
interpolation. (b) $g:=c(0)-c\ge0$ on $[0,2T]$ by Cauchy--Schwarz, which allows integrating an
enclosure of $K$ against $g$ term by term. (c) On the body $[\delta,2T]$, $K$ is squeezed between two
affine functions on each panel, since $1/(2\sinh(u/2))$ is convex and $1/(2\cosh(u/2))$ is concave for
$u/2<\operatorname{arccosh}\sqrt2=0.88137$, which covers the range since $T<0.88137$ (this sandwich is not
available at windows $T\ge\operatorname{arccosh}\sqrt2$); and
$\int g$, $\int ug$ are evaluated exactly from the interval antiderivatives of $c$: no quadrature
remainder enters. (d) On $[0,\delta]$, $\delta=10^{-6}$, use $K(u)\le u^{-1}+\frac12$ and
$g(u)\le G_1u$ with $G_1$ obtained from the exact Chebyshev expansion of $g$, for an upper
contribution $4.8\cdot10^{-6}$ and a lower one $0$. (e) The remaining terms are closed forms in
intervals: $C_0=-\gamma-3\log2-\frac\pi2-\log\pi$, $\int_{2T}^\infty K=\log\coth\frac T2+\pi-2\arctan e^{T}$,
and the two atoms by interval Horner. (f) The enclosure does not assume $F^*\in\PT$: it pays for the
defect, $|A|,|B|\le1.30\cdot10^{-14}$ by exact interval integration of the polynomial against
$e^{\pm t/2}$, so $|2AB|\le3.4\cdot10^{-28}$. The total width is $6.91\cdot10^{-6}$.
\end{proof}

This is a statement about one explicit direction at $\log2$, namely the easy half, $w_{11}>0$, of
Proposition~\ref{prop:schur} at $e=F^*$; it is not a positivity statement for the form.

\begin{remark}[what the numerics do not say]\label{rem:galerkin}
Galerkin (Rayleigh--Ritz) values on nested subspaces are upper bounds on the infimum of the form on
$\PT$, and at the two-prime windows they have not stabilized in the basis size: at fixed subspace
dimension the computed bottom moves by more than an order of magnitude between $N=18$ and $N=36$
Legendre modes, and keeps falling. No numerical value of $\inf W_{\log2}|_{\PT}$ is supported by
this paper, positive or negative. The lower bound of Theorem~\ref{thm:C} at $T=\frac12\log3$ is a
certified interval bound and is unaffected.
\end{remark}

The other half of Proposition~\ref{prop:schur} is not reachable by a norm inequality: numerically, even the
condition $\lambda_3-c^2/w_{11}>0$, in which the prime form is dropped, fails at $T=0.58,0.62,0.65,\log2$, so
positivity on $V$ would be a cancellation between the Schur correction and the prime form, and no
inequality that separates the archimedean block from the prime block can establish it
(Appendix~\ref{app:diag}).

\begin{theorem}[positive definiteness on nested even subspaces at $T=\log2$]\label{thm:nest}
Let $T=\log2$ and let $e_1=F^*,e_3,e_5,\dots,e_{13}$ be the seven explicit even polynomials of degree $17$
in $\PT$ given in the certificate data accompanying this paper (the even eigenvectors of the archimedean
form $\Ainf$ in a Galerkin basis of $\PT$; their Gram matrix differs from the identity by at most
$1.07\cdot10^{-15}$ per entry). Then $W_{\log2}$ is positive definite on
every $E_k=\mathrm{span}\{e_1,e_3,\dots,e_{2k-1}\}$, $k=1,\dots,7$, with certified smallest
eigenvalues
\[
3.4211\cdot10^{-2},\quad2.1136\cdot10^{-3},\quad4.0453\cdot10^{-4},\quad9.0763\cdot10^{-5},\quad
4.0951\cdot10^{-5},\quad2.2156\cdot10^{-5},\quad6.5687\cdot10^{-6};
\]
in particular $W_{\log2}[F]\ge6.56\cdot10^{-6}\|F\|^2$ for every $F\in E_7$, and the determinant of
the $2\times2$ block on $\mathrm{span}\{F^*,e_3\}$ is at least $2.10\cdot10^{-4}$. The same machinery
certifies $E_5$ at $T=0.62$ with $\lambda_{\min}\ge3.92\cdot10^{-4}$ and at $T=0.65$ with
$\lambda_{\min}\ge1.70\cdot10^{-4}$.
The proof includes a computer verification in interval arithmetic (Appendix~\ref{app:encl}).
\end{theorem}
\begin{proof}
Only quadratic values of $W_T$ are ever enclosed: the off-diagonal entries are obtained by
polarization, $W_T(e_i,e_j)=\frac14\big(W_T[e_i+e_j]-W_T[e_i-e_j]\big)$, which is exact because $W_T$
is a quadratic form (Theorem~\ref{thm:A}; $A,B$ are linear, so $2AB$ is quadratic) and $\PT$ is a
linear subspace. This matters: the affine sandwich of the proof of Theorem~\ref{thm:w11} is licensed
by $c(0)-c\ge0$, which is Cauchy--Schwarz for an autocorrelation and false for a cross-correlation.
Each quadratic value is enclosed exactly as in Theorem~\ref{thm:w11}: exact interval polynomial
autocorrelation, the L\'evy form of Theorem~\ref{thm:A}, $K$ squeezed between two affine functions on
each panel of $[\delta,2T]$, the one-sided bound on $[0,\delta]$, closed forms in intervals for the
remaining terms, and the pole defect paid for. Because the panelwise bounds do not depend on the
integrand's coefficients, the panel work is done once, as two fixed linear functionals on the
monomial coefficients, which is what makes $k(k+1)/2$ enclosures affordable. Interval inclusion
monotonicity then gives an interval matrix $\boldsymbol M$ containing the true matrix, and an interval
Cholesky factorization of $\boldsymbol M-\mu I$ that runs to completion with every pivot interval
strictly positive certifies $\lambda_{\min}>\mu$ for every real symmetric matrix in
$\boldsymbol M$; bisection on $\mu$ gives the stated bounds. The Rayleigh transfer to $\|F\|^2$ costs
nothing at this precision, by the Gram bound.
\end{proof}

This is positivity on an explicit finite-dimensional subspace, not a positivity statement for
$W_{\log2}$ on $\PT$; it extends Theorem~\ref{thm:w11} from a line to a $7$-plane.

\begin{proposition}[the Schur coupling is purely arithmetic]\label{prop:bprime}
In the notation of Proposition~\ref{prop:schur} with $e=F^*$ an eigenvector of $\Ainf$ on $\PT$,
\[
b\ =\ P_VW_TF^*\ =\ P_V\PiT F^* ,
\]
i.e.\ the coupling between the pivot and its orthogonal complement is carried by the prime form alone.
\end{proposition}
\begin{proof}
$\Ainf$ is diagonal in its own eigenbasis, so $\Ainf F^*$ is a multiple of $F^*$ and is annihilated by
$P_V$; $W_T=\Ainf+\PiT$.
\end{proof}

\begin{proposition}[the prime form is a sum of truncated shifts]\label{prop:shifts}
Let $R_T$ be the restriction to $\IT$ and $\tau_a$ the translation by $a$. Then
\[
\PiT\ =\ -\sum_{n\in\NT}v_n\,\Rel\big(R_T\tau_{\log n}R_T\big),
\]
a finite sum of compressed translations. Each $R_T\tau_aR_T$, $0<a<2T$, is isometric on
$L^2(-T,T-a)$ and vanishes on $L^2(T-a,T)$, so it is a partial isometry with
infinite-dimensional defect; $\PiT$ is bounded, $\|\PiT\|\le\Aconst$, and is not compact: on functions
supported in $(-T,-T+a_0)$, $a_0=\min_{n\in\NT}\log n$, every backward compressed shift vanishes, so
$\PiT$ acts there as a fixed nonzero combination of isometries on an infinite-dimensional subspace.
Consequently, for any orthonormal system $(e_i)$ of $\PT$, the entries
$\pi_{ij}=\PiT(e_i,e_j)$ satisfy $|\pi_{ij}|\le\Aconst$ and $\sum_j\pi_{ij}^2\le\Aconst^2$.
\end{proposition}
\begin{proof}
$\cF(\log n)=\langle F,\tau_{\log n}F\rangle$ symmetrized, whence the displayed formula; the
isometry/defect statement is immediate from $(\tau_aG)(t)=G(t-a)$ read on $\IT$. Boundedness is the
triangle inequality with $\sum v_n=\Aconst$; the row bounds are Cauchy--Schwarz and Bessel applied to
$\PiT e_i\in L^2(\IT)$, $\|\PiT e_i\|\le\Aconst$.
\end{proof}

The atoms of the prime matrix are one-sided overlaps, and their signs are nodal statements; the
cancellation between the two of them at $\log2$ is measured in Appendix~\ref{app:diag}.

\begin{lemma}[one-sided overlap]\label{lem:N}
Let $F,e\in L^2(\IT)$ be even and $0<a<2T$. Then
\[
\big\langle\Rel\big(R_T\tau_aR_T\big)F,\,e\big\rangle=\int_{a-T}^{T}F(z-a)\,e(z)\,dz=:I_{F,e}(a),
\]
and consequently, for even $e_i,e_j$,
\[
\PiT(e_i,e_j)=-\sum_{n\in\NT}v_n\int_{\log n-T}^{T}e_i(z-\log n)\,e_j(z)\,dz .
\]
\end{lemma}
\begin{proof}
With $\widetilde F$ the extension of $F$ by zero, $(R_T\tau_aR_TF)(z)=\widetilde F(z-a)$ for
$z\in\IT$, so $\langle R_T\tau_aR_TF,e\rangle=\int_{a-T}^{T}F(z-a)e(z)\,dz$ for $0<a<2T$. The
adjoint term is the same expression with $a$ replaced by $-a$, namely
$\int_{-T}^{T-a}F(z+a)e(z)\,dz$, which the substitution $z\mapsto-z$ and the evenness of $F$ and $e$
turn into $\int_{a-T}^{T}F(z-a)e(z)\,dz$ again. The second display is
Proposition~\ref{prop:shifts}.
\end{proof}

\begin{proposition}[sign forcing]\label{prop:G}
Let $F,e\in L^2(\IT)$ be even, real and continuous, and suppose their nodal sets are contained in
$(-z_F,z_F)$ and $(-z_e,z_e)$ respectively, with $e>0$ on $(z_e,T)$ and $F<0$ on $(z_F,T)$. Then for
every $a$ with $T+\max(z_F,z_e)<a<2T$ one has $I_{F,e}(a)<0$, so that the corresponding atom
$-v_nI_{F,e}(\log n)$ is positive. The sign of a contact is therefore forced once the contact lies
beyond $a^*:=T+\max(z_F,z_e)$, and is free only for contacts below it.
\end{proposition}
\begin{proof}
For $z\in(a-T,T)$ one has $z>a-T>z_e$, so $e(z)>0$; and $z-a\in(-T,T-a)$ with $T-a<-z_F$, so
$z-a$ lies in $(-T,-z_F)$, where $F<0$ by evenness. The integrand is therefore negative throughout
the interval of integration, which is nonempty since $a<2T$.
\end{proof}


\section{Positivity at the window \texorpdfstring{$T=\tfrac12\log5$}{T = (1/2) log 5}}\label{sec:log5}

Put $T_5:=\frac12\log5=0.804718\ldots$. Since $\NT=\{n\ge2:\log n<2T\}$, the contact set is
$\{2,3,4\}$ for $\frac12\log4\le T\le T_5$, and $5$ enters only for $T>T_5$; at $T_5$,
$\Aconst=v_2+v_3+v_4=2.94186\ldots$. The window $T_5$ is therefore the largest one with this contact set,
and it lies slightly beyond the range $T\le0.8$ of \cite[Thm.~1.2]{Zhu2026}. The notation
$\varphi_{\mathrm{par}},\ell_{\mathrm{par}},\varepsilon_{\mathrm{par}},\mathcal Q^{\mathrm{par}}_T,
\mathfrak g_{\mathrm{par}},\mathfrak G_T$ is that of \S\ref{subsec:alignineq}, and
$\lambda^{\mathrm W}_{T,\mathrm{par}}$ is the bottom of $W_T$ on the unit sphere of
$\mathcal D^{\mathrm{par}}_T$ (Theorem~\ref{thm:parityschur}(o)).

\begin{theorem}[the primitive form at $\frac12\log5$]\label{thm:log5A}
$\lambda_{\min}(A_{T_5}|\mathcal P^{\mathrm{ev}}_{T_5})\ge2.5\cdot10^{-12}$ and
$\lambda_{\min}(A_{T_5}|\mathcal P^{\mathrm{odd}}_{T_5})\ge3.0\cdot10^{-10}$. Hence
$\lamA\ge2.5\cdot10^{-12}$ for every $0<T\le\frac12\log5$.
\end{theorem}

\noindent The proof includes a computer verification in ball arithmetic (Appendix~\ref{subsec:encllog5}).

\begin{proof}
By Theorem~\ref{thm:parityschur}(o), $\lambda_{\min}(A_{T_5}|\mathcal P_{T_5})$ is the minimum of the two
parity bottoms, which are bounded below by Proposition~\ref{prop:lowercert} with $c=0$ and the
certificate data of Appendix~\ref{subsec:encllog5}. The bound for $T<T_5$ follows from
Lemma~\ref{lem:monoT}.
\end{proof}

\begin{theorem}[the complete form at $\frac12\log5$]\label{thm:log5W}
$W_{T_5}\ge10^{-18}\|\cdot\|^2$ on $\mathcal D^{\mathrm{ev}}_{T_5}$ and $W_{T_5}\ge10^{-16}\|\cdot\|^2$ on
$\mathcal D^{\mathrm{odd}}_{T_5}$. Hence $\lamW\ge10^{-18}$ for every $0<T\le\frac12\log5$, and
$\mathfrak G_T\succ0$ there.
\end{theorem}

\noindent The proof includes a computer verification in ball arithmetic (Appendix~\ref{subsec:encllog5}).

\begin{proof}
By Theorem~\ref{thm:parityschur}(o), $\lambda^{\mathrm W}_{T_5}=\min_{\mathrm{par}}\lambda^{\mathrm W}_{T_5,\mathrm{par}}$,
and each parity is bounded below by Proposition~\ref{prop:lowercert} with $c=-\frac12\varepsilon_{\mathrm{par}}$.
Lemma~\ref{lem:monoT} gives the bound for $T<T_5$. Since $\lamW>0$, Theorem~\ref{thm:parityschur}(c)
gives $\mathfrak G_T\succ0$.
\end{proof}

\begin{corollary}[indeterminacy for $T\le\frac12\log5$]\label{cor:KEPlog5}
$\mathrm{KEP}(T)$ is indeterminate for every $0<T\le\frac12\log5$.
\end{corollary}
\begin{proof}
The proof of Corollary~\ref{cor:KEP08}, with Theorem~\ref{thm:log5W} in place of
\cite[Thm.~1.2]{Zhu2026}: $\lamW>0$ for $T\le T_5$, so Theorem~\ref{thm:KEP3}(a) applies for
$\frac12\log2\le T\le T_5$, and the argument given there covers $0<T<\frac12\log2$.
\end{proof}

\begin{remark}[scope]\label{rem:log5scope}
The gain over \cite[Thm.~1.2]{Zhu2026} is small, from $0.8$ to $0.80472$, and it is obtained by the
window variable alone: on $[0.8,T_5]$ the contact set is still $\{2,3,4\}$. The bound
$10^{-18}$ is consistent with Zhu's $8.9\cdot10^{-18}$ at $T=0.8$ and with the monotonicity of
$\lamW$. The complete bottom is about six orders of magnitude below the primitive one, which is the
size of the polar capacity $\mathfrak g_{\mathrm{ev}}/\|h_{\mathrm{ev}}\|^2$ of
Theorem~\ref{thm:parityschur}(a4). No certified upper bound is given; the values at which the
verification fails (Appendix~\ref{subsec:encllog5}) are not claims of negativity. Both theorems are
certificates, not a mechanism. By Theorem~\ref{thm:margin} and the doubly exponential growth of the
verification threshold (Proposition~\ref{prop:env}), the same method becomes rapidly more expensive as $T$
grows; the first window with the new contact $5$ is not treated.
\end{remark}

\emph{The converse across windows.} Theorem~\ref{thm:parityschur}(c) is a statement at one window. The
following lemma lets a single window control all smaller ones.

\begin{lemma}[monotonicity of the polar values]\label{lem:gmono}
Let $0<T<T'$ and suppose $\lambda_{\min}(A_S|\mathcal P^{\mathrm{par}}_S)>0$ for $S=T$ and $S=T'$. Then
$\mathfrak g_{\mathrm{par}}(T')\le\mathfrak g_{\mathrm{par}}(T)$.
\end{lemma}
\begin{proof}
By Theorem~\ref{thm:parityschur}(a3),
$\mathfrak g_{\mathrm{par}}(S)=\min\{W_S(F):F\in\mathcal D^{\mathrm{par}}_S,\ \ell_{\mathrm{par}}(F)=1\}$. Extension by
zero maps $\mathcal D^{\mathrm{par}}_T$ into $\mathcal D^{\mathrm{par}}_{T'}$; it preserves the parity, the value of
$\ell_{\mathrm{par}}$ (the function $\varphi_{\mathrm{par}}$ on $I_{T'}$ restricts to $\varphi_{\mathrm{par}}$ on $\IT$),
and the value of $W$, by the proof of Lemma~\ref{lem:monoT}. So the feasible set at $T'$ contains that at
$T$, and the minimum does not increase.
\end{proof}

\begin{proposition}[one-window converse]\label{prop:onewindow}
Suppose $\lamA>0$ for $0<T\le T_1$. Then
\emph{(i)} $\lamW\ge0$ for all $0<T\le T_1$ if and only if $\mathfrak g_{\mathrm{ev}}(T_1)\ge0$ and
$\mathfrak g_{\mathrm{odd}}(T_1)\ge0$;
\emph{(ii)} $\lamW>0$ for all $0<T\le T_1$ if and only if $\mathfrak G_{T_1}\succ0$.
In either case $\mathfrak g_{\mathrm{par}}(T)\ge\mathfrak g_{\mathrm{par}}(T_1)$ for $0<T\le T_1$.
\end{proposition}
\begin{proof}
By Lemma~\ref{lem:monoT}, the conditions on $(0,T_1]$ are equivalent to the same conditions at $T_1$. By
Theorem~\ref{thm:parityschur}(c), since $\lamA>0$ at $T_1$, $\lambda^{\mathrm W}_{T_1}\ge0$ if and only if
$\mathfrak g_{\mathrm{ev}}(T_1),\mathfrak g_{\mathrm{odd}}(T_1)\ge0$, and $\lambda^{\mathrm W}_{T_1}>0$ if and only if
$\mathfrak G_{T_1}\succ0$. Since $\lamA>0$ implies that both parity bottoms are positive
(Theorem~\ref{thm:parityschur}(o)), Lemma~\ref{lem:gmono} gives the last statement.
\end{proof}

By Theorem~\ref{thm:log5A}, Proposition~\ref{prop:onewindow} applies with $T_1=T_5$, and
Theorem~\ref{thm:log5W} is the statement that its condition (ii) holds there.



\section{Inertia, remainders and barrier theorems}\label{sec:inertia}

This section records what the \emph{inertia} of the forms of this paper, that is, the number of
negative directions rather than the value of the bottom, can and cannot decide. All forms are
on $L^2(\IT)$ with the metric of $L^2(\IT)$; $n_-(\mathfrak q|V)$ is the largest dimension of a
subspace of $V\cap$(form domain) on which $\mathfrak q<0$, as in Lemma~\ref{lem:parity}. Recall from
\S\ref{subsec:arch} that $\mathfrak a_T=\Gamma_T+\PiT$ on $L^2(\IT)$, with
$\Gamma_T(F)=\frac1{2\pi}\int m|\Fh|^2$ and $\PiT(F)=-\sum_{n\in\NT}v_n\Rel\cF(\log n)$. Nothing in this
section or in the two that follow bears on the sign of $\mathfrak a_T$ at any window.

\subsection{The Sonin split of the truncated form}
Connes and Consani \cite[Thm.~4.7, (83)--(84)]{ConnesConsani2021} define the Sonin projection $\mathbf S$
and a real, even, continuous function $\varepsilon$ on $\R_+^*$ with
$\operatorname{Tr}(\vartheta(f)\mathbf S)=W_\infty(f)+\int f(\rho^{-1})\varepsilon(\rho)\,d^*\rho$. In the additive
variable ($f=g*g^*$, $F(t)=g(e^t)$ after the conjugation $\Delta^{1/2}$ of their Appendix~A) we write
\[
\mathrm{Son}(F):=\operatorname{Tr}\big(\vartheta(g)\mathbf S\vartheta(g)^*\big)\ge0,\qquad
\mathfrak e(F):=\int_{-2T}^{2T}\cF(x)\,\varepsilon(e^{|x|})\,dx=\langle C_\varepsilon F,F\rangle ,
\]
where $C_\varepsilon$ is the integral operator on $L^2(\IT)$ with the continuous kernel $\varepsilon(e^{|x-y|})$;
it is Hilbert--Schmidt, and we put $\kappa_T:=\|C_\varepsilon\|\le4T\max_{[1,e^{2T}]}|\varepsilon|$.

\begin{theorem}[the Sonin split]\label{thm:inertia-split}
For every $T>0$ and every $F\in C_c^\infty(\IT)$, with no moment condition,
\[
\mathfrak a_T(F)=\mathrm{Son}(F)+\mathfrak r_T(F),\qquad
\mathfrak r_T(F):=-\mathfrak e(F)+\PiT(F)
=-\int_{-2T}^{2T}\cF(x)\varepsilon(e^{|x|})\,dx-\sum_{n\in\NT}v_n\Rel\cF(\log n).
\]
\end{theorem}
\begin{proof}
In the dictionary of
\cite[(152)--(154)]{ConnesConsani2021}, $W_\infty(g*g^*)=\frac1{2\pi}\int h_+|\Fh|^2$ with
$h_+(\tau)=\Rel\psi(\frac14+\frac{i\tau}2)-\log\pi$, which is $m$ of \eqref{eq:symbol}; no polar term occurs, so
$W_\infty(g*g^*)=\Gamma_T(F)$, not $\Ainf(F)$. By (83) and cyclicity, $\Gamma_T=\mathrm{Son}-\mathfrak e$. Adding
$\PiT$ gives $\mathfrak a_T=\Gamma_T+\PiT$, which is \eqref{eq:form} since
$\frac1{2\pi}\int\cos(\tau\log n)|\Fh|^2=\Rel\cF(\log n)$ and $v_n=2\Lambda(n)/\sqrt n$; the semilocal prime term
of the trace formula, $-\sum_{p}2\log p\sum_{k\log p<2T}p^{-k/2}\Rel\cF(k\log p)$, is the same sum over $n=p^k\in\NT$.
\end{proof}

The prime form is bounded ($\|\PiT\|\le\Aconst$, \S\ref{subsec:arch}) but is not a compact perturbation.

\begin{proposition}[the prime form at the first contact]\label{thm:inertia-B2}
Let $\frac12\log2<T<\frac12\log3$, so $\NT=\{2\}$ and $\Aconst=v_2=\sqrt2\log2$. Then
$\operatorname{spec}(\PiT)=\{-\frac{\log2}{\sqrt2},0,\frac{\log2}{\sqrt2}\}$, each point of infinite multiplicity,
and $\|\PiT\|=\Aconst/2$. For every $T>\frac12\log2$, $\PiT$ has infinite positive and infinite negative index
on $\PT\cap\{\int F=0\}$.
\end{proposition}
\begin{proof}
Put $a=\log2$. For $r\in[0,a)$ the map $F\mapsto(F(-T+r+ja))_{j}$ is unitary from $L^2(\IT)$ onto
$\int^\oplus_{[0,a)}\C^{n(r)}dr$ with $n(r)=2$ for $r<2T-a$ and $n(r)=1$ otherwise. Since
$\cF(a)=\int F(t)\overline{F(t-a)}dt$ couples only $j$ and $j-1$ in a fibre, $\PiT$ acts on a $2$-dimensional fibre
as $-v_2\begin{psmallmatrix}0&1/2\\1/2&0\end{psmallmatrix}$ and is $0$ on a $1$-dimensional one; both sets of $r$
have positive measure. The last assertion is Theorem~\ref{thm:inertia-nogo} with $k=0$, and with the sign of one
bump in each pair reversed for the positive index.
\end{proof}

\subsection{Infinite index of every bounded remainder, and finite index with a fraction of Sonin}

\begin{theorem}[infinite negative index]\label{thm:inertia-nogo}
Let $T>\frac12\log2$ and let $k\in L^\infty(-2T,2T)$ be real and even, with $C_k$ the operator with kernel
$k(x-y)$ on $L^2(\IT)$. For every $N\ge1$ there is an $N$-dimensional subspace of
$C_c^\infty(\IT)\cap\PT$ on which $-\langle C_kF,F\rangle+\PiT(F)<0$ for $F\ne0$. Hence
$n_-\big(-C_k+\PiT\,\big|\,\PT\cap\{\int F=0\}\big)=\infty$. In particular (with $k=\varepsilon(e^{|\cdot|})$,
Theorem~\ref{thm:inertia-split}) $n_-(\mathfrak a_T-\mathrm{Son}\,|\,\PT)=\infty$, and no inequality
$\mathfrak a_T-\mathrm{Son}\ge-b|v\rangle\langle v|+c(I-|v\rangle\langle v|)$, $c>0$, nor any with finitely many
negative directions, holds on $\PT$ for any inner product.
\end{theorem}
\begin{proof}
Fix $c_0\in C_c^\infty(-1,1)$ real, $c_w=c_0(\cdot/w)$, $\beta=c_w''-\frac14c_w$. Two
integrations by parts give $\int\beta e^{\pm t/2}=0$, so translates of $\beta$ lie in $\PT$. Put
$h_j=\beta(\cdot-t_j)+\beta(\cdot-t_j-a)$, $a=\log2$, with $t_j$ chosen so that all mutual differences avoid
$\{\pm\log n\pm\{0,a\}:n\in\NT\}$ by more than $4w$. For $x=\sum c_jh_j$ the only nonzero value of $c_x(\log n)$,
$n\in\NT$, is the inner one at $\log2$, so $\PiT(x)=-v_2\|\beta\|_2^2\sum|c_j|^2=-\sqrt2\log2\,\|\beta\|_2^2\sum|c_j|^2$.
Also $|\langle C_kx,x\rangle|\le\|k\|_\infty\|x\|_1^2\le4N\|k\|_\infty\|\beta\|_1^2\sum|c_j|^2$ and
$\|\beta\|_1^2/\|\beta\|_2^2=O(w)$. For small $w$ the form is $\le-\frac12\sqrt2\log2\,\|\beta\|_2^2\sum|c_j|^2$.
Imposing $\int x=0$ costs one dimension. A bound with $k'$ negative directions fails on a $(k'+1)$-dimensional
such subspace by dimension count. That $\varepsilon(e^{|\cdot|})$ is bounded on $[0,2T]$ is its continuity
\cite[(84)]{ConnesConsani2021}.
\end{proof}
\begin{remark}
The theorem uses neither the sign of $\mathfrak a_T$ nor the identity of the prime: it says that the Sonin
trace, subtracted \emph{in full}, leaves a remainder that no finite-rank correction controls once a contact is
present. It does not contradict $\mathfrak a_T\ge0$: on the test pairs above,
$\mathrm{Son}$ grows like $w^{-3}\log(1/w)$, the prime term like $w^{-3}$. At $T\le\frac12\log2$ the comb is
empty and $\mathfrak r_T=-C_\varepsilon$ is compact, which is the regime of \cite[\S6]{ConnesConsani2021}.
The one-directional bound of \cite[\S5]{ConnesConsani2021} for $T\le\frac12\log2$ is stated in their
$\xi$-metric; it is not an $L^2(\IT)$ statement, and nothing in that metric enters here.
\end{remark}

\begin{theorem}[relative bound $0$ and finite index]\label{thm:inertia-finite}
Let $T>0$, $\eta\in(0,1]$, and let $k$ be as in Theorem~\ref{thm:inertia-nogo} with $\kappa:=\|C_k\|$. Put
$A':=\kappa+\Aconst$, let $s_*$ be the least $s\ge0$ with $\eta\,m(s)\ge2A'$, put $L:=A'-\eta C_0>0$ and
$c=Ts_*$. Then on the form domain
\[
\eta\,\Gamma_T-C_k+\PiT\ \ge\ A'\,I-(A'+L)\,S_c ,
\]
with $S_c$ the time--band operator of Theorem~\ref{thm:P}. Consequently
\[
n_-\big(\eta\Gamma_T-C_k+\PiT\,\big|\,L^2(\IT)\big)\le\#\{j:\lambda_j(S_c)\ge A'/(A'+L)\}
\le\frac{A'+L}{A'}\cdot\frac{2Ts_*}{\pi},
\]
and the Ky Fan bound $\lambda_j(S_c)\le\Lambda_j(c)$ of Theorem~\ref{thm:P} applies. With
$k=(1-\eta)\varepsilon(e^{|\cdot|})$, $\eta\Gamma_T-C_k+\PiT=\mathfrak a_T-(1-\eta)\mathrm{Son}$, so
$n_-(\mathfrak a_T-(1-\eta)\mathrm{Son})<\infty$ for every $\eta>0$, while it is $\infty$ for $\eta=0$
(Theorem~\ref{thm:inertia-nogo}). Since $m(s)=\log(|s|/2\pi)+O(s^{-2})$, $s_*=2\pi e^{2A'/\eta}(1+o(1))$ as $\eta\downarrow0$.
\end{theorem}
\begin{proof}
By $\|C_k\|=\kappa$ and
$\|\PiT\|\le\Aconst$, the form is $\ge\frac1{2\pi}\int(\eta m-A')|\Fh|^2$. By Theorem~\ref{thm:S}, $m-C_0=\varphi(\tau^2)$
is increasing in $|\tau|$, so $\eta m-A'\ge A'$ for $|\tau|\ge s_*$ and $\eta m-A'\ge\eta C_0-A'=-L$ for
$|\tau|<s_*$. Hence the form is $\ge A'\langle(I-S_c)F,F\rangle-L\langle S_cF,F\rangle$. On the span of the
eigenvectors of $S_c$ with eigenvalue $<A'/(A'+L)$ this is positive definite, whence the count; the trace of $S_c$
is $2c/\pi$ (Theorem~\ref{thm:P}).
\end{proof}
We do not assert any rate for $n_-(\mathfrak a_T-(1-\eta)\mathrm{Son})$ as $\eta\downarrow0$ beyond the
upper bound above; a matching lower bound would need a decay estimate for the inverse transform of $m$ that we
have not written out.

\subsection{Quartets and the zero-side inertia}
For a hypothetical zero $\rho=\frac12+z$, $z=\beta+i\gamma$, $0<\beta<\frac12$, $\gamma>0$, the four zeros
$\frac12\pm z,\frac12\pm\bar z$ form a \emph{quartet}. For real $F$ put $\mathsf A_z(F):=\Phi_F(\frac12+z)$.

\begin{lemma}[signature of a quartet]\label{thm:inertia-quartet}
For real $F\in L^2(\IT)$ with $F=F_{\mathrm e}+F_{\mathrm o}$, the contribution of a quartet to
$\sum_\rho\Phi_F(\rho)\Phi_F(1-\rho)$ in \eqref{eq:W} is
\[
q_z(F)=4\Rel\big(\mathsf A_z(F_{\mathrm e})^2\big)-4\Rel\big(\mathsf A_z(F_{\mathrm o})^2\big),
\]
with no cross term between parities. For every $T>0$, $q_z$ has signature exactly $(2,2)$ on the real space
$\PT$: one negative direction in $\PT^{\mathrm{ev}}$ and one in $\PT^{\mathrm{odd}}$.
\end{lemma}
\begin{proof}
Normalization: $\Phi_F(s)=\int F(t)e^{(s-\frac12)t}dt$ as
in \eqref{eq:W}; $\Phi_F(\frac12-z)=\Phi_{F(-\cdot)}(\frac12+z)$, so the four terms are
$2\mathsf A_z(F)\Phi_F(\frac12-z)+2\overline{(\cdot)}=4\Rel\big((\mathsf A_z(F_{\mathrm e})+\mathsf A_z(F_{\mathrm o}))(\mathsf A_z(F_{\mathrm e})-\mathsf A_z(F_{\mathrm o}))\big)$.
$\Rel(w^2)$ has signature $(1,1)$ in $(\Rel w,\Iml w)$. Independence of $\Rel\mathsf A_z,\Iml\mathsf A_z$ on
$\PT^{\mathrm{ev}}$: otherwise $a\cosh\beta t\cos\gamma t+b\sinh\beta t\sin\gamma t=c\cosh\frac t2$ on $(0,T)$, hence on
$\R$ by analyticity; growth ($\beta<\frac12$) forces $c=0$, then $t=0$ gives $a=0$ and $b=0$. Likewise on
$\PT^{\mathrm{odd}}$ with $\sinh\frac t2$.
\end{proof}

\begin{proposition}[finite modifications of the zero set]\label{thm:inertia-rank}
Let $Z'$ arise from the zero multiset of $\zeta$ by adding $m$ quartets and removing $m$ pairs
$\frac12\pm i\gamma^0_k$, and write $W^{Z'}_T(F)=\sum_{\rho\in Z'}\Phi_F(\rho)\Phi_F(1-\rho)$. On real $L^2(\IT)$,
$W^{Z'}_T-W_T$ is a form of rank $\le6m$, and so are its changes on every Schur complement taken with invertible
blocks; hence $|n_-(W^{Z'}_T|V)-n_-(W_T|V)|\le6m$ for every subspace $V$.
\end{proposition}
\begin{proof}
Each quartet contributes $q_z$ (rank $4$), each removed pair $-2|\Phi_F(\frac12+i\gamma^0)|^2$ (rank $2$); Weyl
interlacing.
\end{proof}

\subsection{What these theorems exclude}\label{subsec:inertia-excl}
\emph{Inertia alone.} By Proposition~\ref{thm:inertia-rank}, any argument that uses only the decomposition
$\mathfrak a_T=\Gamma_T-C_\varepsilon+\PiT$ with the operator properties of the pieces (the archimedean multiplier;
$\PiT$ with exactly the weights $v_n$; a compact remainder continuous in $T$), inertia arithmetic (Sylvester,
Haynsworth, interlacing, Lemma~\ref{lem:monoT}) and any property of $\PiT$ alone, applies verbatim to a model with
the same $\Gamma_T$ and $\PiT$ and a finite-rank change of the compact term coming from $m$ added quartets.
In such a model the quartet forms have $2m$ negative directions on $\PT$ for every $T$
(Lemma~\ref{thm:inertia-quartet}), and these directions carry weight $e^{2\beta T}$ against the $O(1)$ of the
on-line part; we therefore expect (this is argued, not proved here) $n_-(\mathfrak a^{Z'}_T|\PT)\ge2m$ for all large $T$.
Granting this, an argument of the class just described cannot prove ``$\sup_Tn_-(\mathfrak a_T|\PT)<\infty
\Rightarrow n_-\equiv0$'': applied to the model it would contradict the model's index. The exclusion concerns the class
\emph{as defined}: it does not cover arguments that use quantitative properties of the compact term itself (a
$T$-uniform bound on $\|C_\varepsilon\|$, or a sign), which the model does not share (its compact term has norm
$\asymp e^{2\beta T}$).

\emph{Sonin-norm bounds.} By Theorem~\ref{thm:inertia-nogo} no bound
$\mathfrak a_T\ge\mathrm{Son}-(\text{finite rank})$ holds past $\frac12\log2$; by
Theorem~\ref{thm:inertia-finite}, ceding any fraction $\eta>0$ of $\mathrm{Son}$ restores a finite index, but only
up to $N\asymp Te^{2A'/\eta}$ directions, and positivity is then the positivity of an $N\times N$ Schur complement,
the finite arena of \S\ref{subsec:arena} with a larger $N$. A bound of the form
$\|\mathfrak a_T-\mathrm{Pos}_T\|_{\PT}\le e^{o(T)}$ with $\mathrm{Pos}_T\ge0$, if computed from the local data
($\Gamma$, the weights $v_n$), is the same for $\zeta$ and for the models of Proposition~\ref{thm:inertia-rank};
if computed as the difference of the forms, its lower half contains $\lamA\ge-e^{o(T)}$, which is a statement about
$L^{\mathrm A}$ (Corollary~\ref{cor:margin2}(ii)) that no local computation decides.

\emph{The witness decorrelator.} The first-contact witness of \S\ref{sec:firstcontact} cannot be excluded by an
inequality on the weights that is linear in $\Lambda$ (one test) or ranges over a family of power size, and is
true for $\zeta$, unless it controls $\sum_n\Lambda(n)n^{-1/2+i\gamma}w(\log n)$ beyond the window with the exact
strength of the absence of the zero; in the normalization of \eqref{eq:fw} this sum is
$-\frac12$ times the prime term of $f_w$ at $\tau=-\gamma$. Together with Corollary~\ref{cor:void} (properties
stable under finite modification of the weights decide no window), such inequalities are circular.

\subsection{What remains}
Within the witness architecture, one case is not covered by any of the barrier statements above or in
Corollary~\ref{cor:void}: an inequality on the vector $(v_n)_{n\ge2}$, used beyond the window, that is
\emph{non-linear} in $\Lambda$, is not a family of power size, and is not invariant under finite modifications of
the weights. In the language of Problem~\ref{prob:B}: a proof of $f_{w_T}\ge0$ on
$[0,\tau_{\mathrm{env}}(T)]$ would have to use such a property of the primes outside the window. It is open
whether any such property exists. Outside that architecture, the finiteness questions on $n_-(\mathfrak a_T|\PT)$
as $T\to\infty$, and the value of $L^{\mathrm A}$ in Corollary~\ref{cor:margin2}, remain open.

\begin{remark}[conventions of the sources]\label{rem:inertia-conv}
The Fourier convention of \cite{ConnesConsani2021} ($e^{-2\pi ixy}$) differs from ours ($\Fh(\tau)=\int Fe^{i\tau t}$);
this is immaterial, since only $|\Fh|^2$ and even symbols enter. Their truncated form $E_S$ has no polar term and
equals $\mathfrak a_T$, not $W_T$; the two agree on $\PT$. Throughout, $\PiT$ is written with $\Rel\cF$, as in
\S\ref{subsec:arch}; for real $F$ the real part is redundant.
\end{remark}

\section{The semilocal Sonin decomposition at the first prime}\label{sec:semilocal}

Section~\ref{sec:inertia} showed that the archimedean Sonin trace, subtracted in full, leaves a
remainder of infinite negative index as soon as the contact $n=2$ is present
(Theorem~\ref{thm:inertia-nogo}). Here we replace the archimedean Sonin space by the semilocal one of
Connes and Consani for $S=\{\infty,2\}$ and show that the comb of $2$ is then absorbed \emph{exactly},
with a Hilbert--Schmidt remainder. The objects $F_S$, the quasi-inner function and the semilocal Sonin
space are those of \cite{CCquasi2020,CCM2023}; what is new is listed in Remark~\ref{rem:sl-new}. No
statement of this section bears on the sign of $\mathfrak a_T$.

\subsection{Two models and their dictionary}\label{subsec:sl-dict}
This paper works on $L^2(\IT)$ in the additive variable $t$, with $\Fh(\tau)=\int Fe^{i\tau t}dt$,
autocorrelation $\cF$, weights $v_n=2\Lambda(n)/\sqrt n$, and the forms
$\Gamma_T(F)=\frac1{2\pi}\int m|\Fh|^2$, $\PiT(F)=-\sum_{n\in\NT}v_n\Rel\cF(\log n)$ of
\S\ref{subsec:arch}. The operators below live in the \emph{u-model} $H=L^2(\R_+,du)$. The dictionary is:
\begin{enumerate}
\item[(D1)] $L^2(\R)^{\mathrm{ev}}\to H$, $\xi\mapsto\sqrt2\,\xi|_{\R_+}$, is unitary; it carries the
Fourier transform $\xi\mapsto\int\xi(x)e^{-2\pi ixy}dx$ of \cite{ConnesConsani2021} to $F_\infty$, the
operator with kernel $2\cos(2\pi ux)$, and the cut-off $\mathbf 1_{|x|\le\lambda}$ to
$P=\mathbf 1_{(0,\lambda]}(u)$.
\item[(D2)] $W:H\to L^2(\R,dt)$, $(W\xi)(t)=e^{t/2}\xi(e^t)$, is unitary; $P\mapsto\mathbf 1_{t\le a}$,
$a=\log\lambda$; dilations become translations, and $D:=-i\,d/dt$.
\item[(D3)] We use the scaling action of \cite[(61)]{ConnesConsani2021},
$(\vartheta(\rho)\xi)(v)=\rho^{-1/2}\xi(\rho^{-1}v)$, $\vartheta(f)=\int f(\rho)\vartheta(\rho)d^*\rho$. The
inverted convention $\int h(a)a^{1/2}\xi(a\,\cdot)\,d^*a=\vartheta(\tilde h)$, $\tilde h(\rho)=h(\rho^{-1})$, gives the
same identities: every functional below ($W_\infty$, the comb, $\int f(\rho^{-1})\varepsilon(\rho)d^*\rho$ with
$\varepsilon$ inversion-symmetric) is invariant under $f\mapsto\tilde f$.
\item[(D4)] For $F\in C_c^\infty(\IT)$ put $g_F(\rho):=F(\log\rho)$ and $f_F:=g_F*g_F^*$, so that
$f_F(e^x)=\cF(x)$ and $\int f_F(\rho^{-1})\varepsilon(\rho)d^*\rho=\int_{-2T}^{2T}\cF(x)\varepsilon(e^{|x|})dx
=\langle C_\varepsilon F,F\rangle$ (notation of Theorem~\ref{thm:inertia-split}); and
$W_\infty(f_F)=\Gamma_T(F)$ (proof of Theorem~\ref{thm:inertia-split}).
\item[(D5)] For a projection $\Pi$ on $H$ we write $\mathrm{Son}_\Pi(F):=\operatorname{Tr}(\vartheta(g_F)\Pi\vartheta(g_F)^*)
=\operatorname{Tr}(\vartheta(f_F)\Pi)\ge0$.
\end{enumerate}
\emph{Metric warning.} All operator statements of \S\S\ref{subsec:sl-dict}--\ref{subsec:sl-trace} are in $H$
(equivalently $L^2(\R,dt)$); the forms of \S\ref{subsec:sl-split} are on $L^2(\IT)$ with the metric of
$L^2(\IT)$, through (D4). No statement uses the $\xi$-metric of \cite[\S6]{ConnesConsani2021}.

\subsection{The semilocal Fourier transform}
\begin{definition}[\cite{CCM2023,CCquasi2020}]\label{def:FS}
Let $D_2\xi(u)=\xi(2u)$ and $E:=\sum_{k\ge0}D_2^k$ on $H$, $E^{-1}=1-D_2$. The semilocal Fourier transform is
$F_S:=EF_\infty E^{-1}$ (\cite[Prop.~4.1]{CCM2023}: $F_S\circ\eta_S=\eta_S\circ F_\infty$, with $\eta_S=E$ in the
u-model). In Mellin form ($\operatorname{Re}s=\frac12$) $F_S$ is $F_\infty$ composed with the multiplier
$(1-2^{-s})/(1-2^{-(1-s)})$; this is the factor $\rho_2$ of the quasi-inner function
$u=\rho_\infty\prod_p\rho_p$, $\rho_p(z)=\gamma_p(z)/\gamma_p(1-z)$, of \cite[Lemma~3.1]{CCquasi2020}.
Put $\widehat P^S:=F_SPF_S$, $Q_1:=1-P$, $Q_2:=1-\widehat P^S$, $A_S:=PF_SP$, and let $\Pi_S$ be the
orthogonal projection on $\operatorname{ran}Q_1\cap\operatorname{ran}Q_2$, the semilocal Sonin space
(\cite{CCquasi2020}, main theorem: Sonin $=\ker u_{22}$).
\end{definition}

\begin{lemma}[u-model form; scattering phase]\label{lem:sl-phase}
In $L^2(\R,dt)$, $F_\infty=I\,m_\infty(D)$ with $I$ the reflection $t\mapsto-t$ and $|m_\infty|=1$, and
$F_S=F_\infty U=I\,M(D)$, $M=m_\infty U$, $U=e^{i\varphi(D)}$ with
\[
\varphi'(t)=-\sum_{k\ge1}v_{2^k}\cos(kt\log2),\qquad\text{and for finite }S:\quad
\varphi_S'(t)=-\sum_{p\in S}\sum_{k\ge1}v_{p^k}\cos(kt\log p).
\]
The comb $\varphi_S'$ is untruncated; in pairings with $\cF$ only $n\in\NT$ survive, since $\supp\cF\subset[-2T,2T]$.
\end{lemma}
\begin{proof}
The identification of $U$ with $\rho_2$ is \cite[Lemma~3.1]{CCquasi2020}; we verify the normalization.
$WD_2W^{-1}=q\,T$, $q=2^{-1/2}$, $T$ translation by $\log2$, so $E^{-1}=1-qT$, $F_\infty EF_\infty=E^*$ and
$U=E^*E^{-1}=(1-qT)/(1-qT^{-1})$, a quotient of conjugates: $\varphi(t)=2\arg(1-qe^{it\log2})=-2\sum_kq^k\sin(kt\log2)/k$,
whence $\varphi'$ with $2q^k\log2=v_{2^k}$. $IF_\infty$ commutes with translations and is unitary, hence a
multiplier; its value is the Tate factor. For several primes $U_S=\prod_pU_p$.
\end{proof}

\begin{proposition}\label{prop:sl-AS}
\emph{(i)} $F_S$ is a unitary, self-adjoint involution and $F_S\vartheta(\rho)F_S=\vartheta(\rho^{-1})$; $F_S$ is real.
\emph{(ii)} $A_S$ is compact, self-adjoint, $\|A_S\|<1$ and $\pm1$ are not eigenvalues.
\emph{(iii)} For $S=\{\infty,2\}$: $A_S$ is not Hilbert--Schmidt, and $A_S\in\mathcal S_p$ for every $p>2$.
\end{proposition}
\begin{proof}
(i) $F_S$ is similar to the involution $F_\infty$, and unitary by Lemma~\ref{lem:sl-phase}; unitary plus involutive is
self-adjoint. $F_\infty\vartheta(\rho)F_\infty=\vartheta(\rho^{-1})$ and $U$ commutes with $\vartheta(\rho)$.
(ii) Compactness is \cite[Thms.~2.1, 4.4]{CCquasi2020}; directly: the grouped kernel
$\sum_{k\le K}(2c_k-c_{k-1})$, $c_k=\cos(2\pi2^kux)$, gives Hilbert--Schmidt operators converging in norm, with
$\|PD_2^kF_\infty E^{-1}P\|\le(1+q)q^k$. No eigenvalue $\pm1$: by Paley--Wiener applied to $E^{-1}\xi$.
(iii) If $A_S$ had a kernel $\kappa\in L^2$, test against $\chi c_m$, $\chi\in C_c^\infty((\delta,\lambda-\delta)^2)$:
non-stationary phase in the lacunary sum gives $\langle\kappa,\chi c_m\rangle\to\frac12\iint\chi\ne0$, while
Riemann--Lebesgue gives $\to0$. For $p>2$:
$\|Pc_kP\|_{\mathcal S_p}\le\|Pc_kP\|^{1-2/p}\|Pc_kP\|_{\mathrm{HS}}^{2/p}\le(\tfrac12q^k)^{1-2/p}\lambda^{2/p}$ is summable.
For $|S|\ge3$ the lacunary cross sums in (iii) have not been carried out; (i) and (ii) hold for every finite $S$.
\end{proof}

\begin{lemma}[trace formula for the linear part]\label{lem:sl-linear}
Let $h$ be such that $y\mapsto h(e^y)$ is Schwartz, and $S$ finite. Then $\vartheta(h)(\widehat P^S-\widehat P_\infty)$ is
trace class, and for every $\lambda>0$
\[
\operatorname{Tr}\vartheta(h)(\widehat P^S-\widehat P_\infty)=\sum_{p\in S}\sum_{k\ge1}\log p\;p^{-k/2}\big(h(p^k)+h(p^{-k})\big)
=\tfrac12\sum_{p\in S}\sum_{k\ge1}v_{p^k}\big(h(p^k)+h(p^{-k})\big).
\]
\end{lemma}
\begin{proof}
Write $U=\sum_ju_jT_{h_j}$ with $|U|=1$ and $\sum_j(1+|h_j|)|u_j|<\infty$ (for $S$ finite, $h_j$ ranges over
$\pm\log n$, $n$ in the multiplicative semigroup of $S$, and $u_j$ is a product of geometric sequences, so the
condition holds; the set of $h_j$ need not be a lattice). Then $UP_aU^*-P_a=\sum_{j,l}u_j\bar u_l(P_{a-h_j}-P_a)T_{h_j-h_l}$,
$P_{a-h}-P_a=-\operatorname{sgn}(h)\mathbf 1_{I_h}$ with $|I_h|=|h|$. By \cite[Thm.~4.5]{Simon2005},
$\|\widehat k(D)\mathbf 1_I\|_1\le C_k(1+|I|)$, so the series converges in trace norm, and
$\operatorname{Tr}=-\sum h_ju_j\bar u_lk(h_j-h_l)=-\frac1{2\pi}\int\widehat k\,\varphi'$. Conjugating by $F_\infty$
reverses $D$; Lemma~\ref{lem:sl-phase} gives the comb.
\end{proof}
This is the dual form of the classical phase-shift trace formula; it is the local
term at $p$ of Connes' semilocal trace formula, and is not new in substance. It concerns the \emph{linear} part
$1-\widehat P$ only, not the Sonin projection $\Pi_S$.

\subsection{Gap and trace class}\label{subsec:sl-trace}
Let $\{\xi_n\}$ be a real orthonormal eigenbasis of $A_S$ in $\operatorname{ran}P$, $A_S\xi_n=\lambda_n\xi_n$, ordered by
$|\lambda_n|$ decreasing, $\lambda_1:=\|A_S\|$; for $\lambda_n\ne0$ put $\zeta_n:=(1-P)F_S\xi_n/(1-\lambda_n^2)^{1/2}$ and
$\tau_n:=\lambda_n(1-\lambda_n^2)^{-1/2}$.

\begin{theorem}[operator gap]\label{thm:sl-gap}
In the Halmos decomposition of $(Q_1,Q_2)$ the angle operator satisfies $\cos^2\theta_n=\lambda_n^2$; the
$\zeta_n$ are orthonormal; $Q_1Q_2Q_1=\Pi_S+\sum_n\lambda_n^2|\zeta_n\rangle\langle\zeta_n|$;
$\operatorname{spec}(Q_1Q_2Q_1|_{\operatorname{ran}Q_1})\subset\{0\}\cup\{\lambda_n^2\}\cup\{1\}$ with
$\sup_n\lambda_n^2=\lambda_1^2<1$. Hence $1$ is isolated, $\Pi_S$ is the Riesz projection of $Q_1Q_2Q_1$ at $1$, and
$\|(Q_1Q_2Q_1)^n-\Pi_S\|=\lambda_1^{2n}$.
\end{theorem}
\begin{proof}
On the generic part $Q_1=\mathrm{diag}(1,0)$, $Q_2=\begin{psmallmatrix}c^2&cs\\cs&s^2\end{psmallmatrix}$, so
$Q_1Q_2Q_1=c^2\oplus0$ and $P\widehat P^SP=0\oplus c^2$; and $P\widehat P^SP=A_S^2$ by
Proposition~\ref{prop:sl-AS}(i). $\|(1-P)F_S\xi_n\|^2=1-\lambda_n^2$ and $Q_1\widehat P^SQ_1\zeta_n=(1-\lambda_n^2)\zeta_n$.
The gap is Proposition~\ref{prop:sl-AS}(ii).
\end{proof}
The gap is spectral (at $1$, for $Q_1Q_2Q_1$); it is not a gap in heights of the positivity
requirement discussed in Remark~\ref{rem:sl-not}. Since $\sum\lambda_n^2=\infty$, the operator
$Q_1Q_2Q_1-\Pi_S$ is not trace class; it is Hilbert--Schmidt by Proposition~\ref{prop:sl-AS}(iii).

\begin{lemma}[trace class]\label{lem:sl-TC}
For $f\in C_c^\infty(\R_+^*)$: $\vartheta(g)Q_1Q_2$ is Hilbert--Schmidt, $\vartheta(f)Q_1Q_2Q_1$, $\vartheta(f)\Pi_S$ and
$\vartheta(f)(\widehat P^S-\widehat P_\infty)$ are trace class, and $P\vartheta(f)P\,A_S^2\in\mathcal S_1$; in particular
$\sum_n\lambda_n^2\langle\xi_n|\vartheta(f)\xi_n\rangle$ converges absolutely although $\sum\lambda_n^2=\infty$.
\end{lemma}
\begin{proof}
By Lemma~\ref{lem:sl-phase} and (D2), $Q_1Q_2=I\,\mathbf 1_{t<-a}M(D)\mathbf 1_{t>a}M(D)^*I$. Split
$\widehat{\tilde g}(D)\mathbf 1_{t<-a}M(D)\mathbf 1_{t>a}$ into a commutator, of Hilbert--Schmidt norm
$\int|k|^2|x|dx$, and the Hankel block $\mathbf 1_{t<-a}(\widehat{\tilde g}M)(D)\mathbf 1_{t>a}$ of norm
$\int|K_g|^2(-2a-x)_+dx$, finite because $\widehat{\tilde g}M\in\mathcal S$ ($m_\infty\in\mathcal O_M$ by Stirling,
$U\in C_b^\infty$). Trace class of $\vartheta(f)Q_1Q_2Q_1$ via $f=g_1*g_2$, \cite[Lemma~D.1]{ConnesConsani2021} for the
commutators and the Dixmier--Malliavin factorization; $\vartheta(f)\Pi_S=\vartheta(f)Q_1Q_2Q_1\Pi_S$.
The last claim: $P\vartheta(f)XP=\Theta_{00}(A_S^2-A_\infty^2)+\Theta_{01}X_{10}$ with $X=\widehat P^S-\widehat P_\infty$,
$\Theta_{01}=P\vartheta(f)(1-P)\in\mathcal S_1$ \cite[Thm.~4.5]{Simon2005}, and $A_\infty\in\mathcal S_1$.
\end{proof}

\begin{theorem}[structure of the Halmos density]\label{thm:sl-L2}
Let $S=\{\infty,2\}$. The series $E_S(f):=\sum_n\tau_n\big(\langle\xi_n|\vartheta(f)\zeta_n\rangle+\langle\zeta_n|\vartheta(f)\xi_n\rangle\big)$
converges absolutely for every $f\in C_c^\infty(\R_+^*)$ and equals $\int f(\rho^{-1})d_S(\rho)d^*\rho$, where
$d_S(\rho^{-1})=d_S(\rho)$ and, for $\rho\ge1$,
$d_S(\rho)=\rho^{1/2}\int_{\lambda/\rho}^{\lambda}G(\rho x,x)\,dx$, $G=(1-P)F_SA_S(1-A_S^2)^{-1}P$. For every $R>1$,
\[
d_S(\rho)=\sum_{0\le j\le\log_2R}\beta_j\log^+\frac1{|\rho-2^j|}+O_R(1)\quad(1\le\rho\le R),\qquad
\beta_j=2^{j/2}(1-2^{-j})\frac{\lambda^2}{2\log2},
\]
and $d_S(\rho)\sim b\,(\rho-1)\log\frac1{\rho-1}$ as $\rho\downarrow1$, $b=\lambda^2/(2\log2)$. Hence
$d_S\in L^p_{\mathrm{loc}}(\R_+^*)$ for every $p<\infty$ and $d_S\notin L^\infty_{\mathrm{loc}}$; the pointwise series
diverges to $+\infty$ at every $\rho=2^j$, $j\ge1$.
\end{theorem}
\begin{proof}
Absolute convergence: the two cross series are traces of $\Theta_{01}G$ and $\Theta_{10}G^*$ in an orthonormal basis.
$G=G_0(1+R)$, $R=A_S^2(1-A_S^2)^{-1}$ Hilbert--Schmidt by Proposition~\ref{prop:sl-AS}(iii), so the $R$-part has
bounded density. The kernel of $G_0=(1-P)F_SPF_SP$ is
$\frac12\sum_{k,l\ge-1}a_ka_l[S(2^ku-2^lz)+S(2^ku+2^lz)]$, $S(w)=\sin(2\pi\lambda w)/(2\pi w)$, $a_{-1}=-1$, $a_k=1$;
the lacunary sum $\sum_{k\ge k_0}S(2^kt)=\lambda\,\#\{k\ge k_0:2^k|t|\le1\}+O(1)$, uniformly, gives
$G_0(\rho x,x)=\frac{\lambda}{2\log2}\sum_i\log^+\frac1{|\rho-2^i|}+O_R(1)$ uniformly in $x\in[\lambda/R,\lambda]$;
integrate over $[\lambda/\rho,\lambda]$. Two estimates are used here as stated and not written out in detail: the
uniform bound $\sup_u\|G_0(u,\cdot)\|_{L^2(0,\lambda)}<\infty$ in the region $x\to0$, and the constant in the lacunary
sum. Divergence at $2^j$: the dyadic quasi-covariance
$K_S(2^jx,y)=K_S(x,y)-B_j(x,y)$ ($B_j$ entire) reduces the pointwise term to $\|\mathbf 1_{I}A_S\|_{\mathrm{HS}}^2=\infty$
(Proposition~\ref{prop:sl-AS}(iii) on $I=[2^{-j}\lambda,\lambda]$).
\end{proof}

\begin{theorem}[the semilocal trace formula]\label{thm:sl-84}
Let $S=\{\infty,2\}$, $\lambda>0$, $f\in C_c^\infty(\R_+^*)$, and
$\mathrm{comb}_2(f):=\sum_{k\ge1}\log2\cdot2^{-k/2}(f(2^k)+f(2^{-k}))$.
\emph{(i) Relative form (unconditional, every $\lambda$).}
\[
\operatorname{Tr}\vartheta(f)\Pi_S-E_S(f)=\operatorname{Tr}\vartheta(f)\Pi_\infty-E_\infty(f)-\mathrm{comb}_2(f).
\]
\emph{(ii) Absolute form.} If the archimedean formula \cite[Thm.~4.7, (83)]{ConnesConsani2021},
$\operatorname{Tr}\vartheta(f)\Pi_\infty=W_\infty(f)+\int f(\rho^{-1})\varepsilon_\infty(\rho)d^*\rho$, holds at level $\lambda$, then
\[
\operatorname{Tr}\vartheta(f)\Pi_S=W_\infty(f)-\mathrm{comb}_2(f)+\int f(\rho^{-1})\,d_S(\rho)\,d^*\rho .
\]
At $\lambda=1$ this hypothesis is \cite[Thm.~4.7]{ConnesConsani2021} (Sonin space $S(1,1)$).
\end{theorem}
\begin{proof}
For $f$ real with $\tilde f=f$, $[F,\vartheta(f)]=0$ for $F\in\{F_\infty,F_S\}$, and expanding
$\langle\xi_n|F\vartheta(f)F\xi_n\rangle$ with $F\xi_n=\lambda_n\xi_n+(1-\lambda_n^2)^{1/2}\zeta_n$ gives, for each $n$,
$\langle\zeta_n|\vartheta(f)\zeta_n\rangle=\langle\xi_n|\vartheta(f)\xi_n\rangle-\tau_n(c_n+c_n')$ (\cite[(87)]{ConnesConsani2021},
exact, no summation). The four series (diagonal in $\xi$, two cross, diagonal in $\zeta$) converge absolutely by
Lemma~\ref{lem:sl-TC}, so $\operatorname{Tr}\vartheta(f)(Q_1Q_2Q_1-\Pi_F)$ and the cross traces combine term by term into
$\operatorname{Tr}\vartheta(f)Q_1Q_2Q_1-\operatorname{Tr}\vartheta(f)\Pi_F=\ldots-E_F(f)$; the diagonal $\lambda_n^2$ terms cancel
$n$ by $n$ (the semilocal form of \cite[(89)--(90)]{ConnesConsani2021}). Subtracting the two cases,
$Q_1Q_2^SQ_1-Q_1Q_2^\infty Q_1=-(1-P)X(1-P)$, whose $\vartheta(f)$-trace is computed blockwise and gives
$-\mathrm{comb}_2(f)$ by Lemma~\ref{lem:sl-linear}. Antisymmetric $f$: both sides vanish ($F\Pi_F=\Pi_FF$). (ii) Insert
(83).
\end{proof}

\begin{remark}[the level $\lambda$]\label{rem:sl-level}
\cite[Thm.~4.7]{ConnesConsani2021} is stated for Sonin's space $S(1,1)$ (cut-offs $|q|\le1$, $|p|\le1$), which under
(D1) is the level $\lambda=1$ in the u-model; $S(\lambda,\lambda)$ is not a dilate of $S(1,1)$ unless $\lambda=1$
(dilation maps $S(1,1)$ to $S(a,1/a)$). The absolute form (ii) is therefore asserted at $\lambda=1$ only; all
proofs of this section are $\lambda$-uniform, and earlier working notes that stated (ii) ``at $\lambda=\frac12$''
used the wrong level.
\end{remark}

\subsection{The split of the truncated form}\label{subsec:sl-split}
For $T>0$ let $\Pi^{(2)}_T(F):=-\sum_{k:\,2^k\in\NT}v_{2^k}\Rel\cF(k\log2)$ and
$\Pi^{(\ne2)}_T:=\PiT-\Pi^{(2)}_T$, so that $\mathfrak a_T=\Gamma_T+\Pi^{(2)}_T+\Pi^{(\ne2)}_T$, and
$\Pi^{(\ne2)}_T=0$ exactly when $T<\frac12\log3$. With (D4), $\mathrm{comb}_2(f_F)=-\Pi^{(2)}_T(F)$, since
$v_{2^k}=2\log2\cdot2^{-k/2}$ and $\cF(-x)=\overline{\cF(x)}$.

\begin{corollary}[semilocal Sonin split]\label{cor:sl-split}
Let $\lambda=1$, $S=\{\infty,2\}$, $\varepsilon_S:=d_S$. For every $T>0$ and $F\in C_c^\infty(\IT)$,
\[
\Gamma_T(F)+\Pi^{(2)}_T(F)=\mathrm{Son}_{\Pi_S}(F)-\langle C_{\varepsilon_S}F,F\rangle ,
\]
and $C_{\varepsilon_S}$, the operator on $L^2(\IT)$ with kernel $\varepsilon_S(e^{|x-y|})$, is Hilbert--Schmidt. In
particular, for $T<\frac12\log3$, $\mathfrak a_T=\mathrm{Son}_{\Pi_S}-C_{\varepsilon_S}$ with Hilbert--Schmidt remainder, on
$L^2(\IT)$ and on $\PT$; for every $\delta>0$ the number of eigenvalues of $-C_{\varepsilon_S}|_{\PT}$ below $-\delta$ is
finite, so $n_-\big(\mathfrak a_T-\mathrm{Son}_{\Pi_S}+\delta\|\cdot\|^2\,\big|\,\PT\big)<\infty$. The part
$-C_{\varepsilon_S}+C_{\varepsilon_\infty}$ is Hilbert--Schmidt for every $\lambda$, unconditionally.
\end{corollary}
\begin{proof}
Theorem~\ref{thm:sl-84}(ii) with $f=f_F$ and (D4), (D5). $\int_{I_T^2}|\varepsilon_S(e^{|x-y|})|^2\le2T\int_{-2T}^{2T}|\varepsilon_S(e^{|x|})|^2dx<\infty$
by Theorem~\ref{thm:sl-L2} ($\log^+\frac1{|e^x-2^j|}\asymp\log^+\frac1{|x-j\log2|}$) and continuity of $\varepsilon_\infty$
\cite[(84)]{ConnesConsani2021}. A compact self-adjoint operator has finitely many eigenvalues below $-\delta$.
The last sentence uses only Theorem~\ref{thm:sl-84}(i) and Theorem~\ref{thm:sl-L2}.
\end{proof}
We do not claim $n_-(-C_{\varepsilon_S}|_{\PT})<\infty$ itself: a compact self-adjoint operator may have infinitely
many negative eigenvalues, so compactness gives only the $\delta$-version above. (That $n_-(\mathfrak a_T|\PT)<\infty$
is trivial, Theorem~\ref{thm:inertia-finite} with $\eta=1$.) Likewise $\Gamma_T+\Pi^{(2)}_T$ equals $\mathfrak a_T$ only
for $T<\frac12\log3$; the identity of the corollary holds for every $T$, but it is a statement about
$\Gamma_T+\Pi^{(2)}_T$, not about $\mathfrak a_T$, beyond that window.

\begin{remark}[what this does not give]\label{rem:sl-not}
(1) \emph{No sign.} The corollary writes $\mathfrak a_T$ ($T<\frac12\log3$) as a positive trace form minus a compact
form; positivity of $\mathfrak a_T$ on $\PT$ is the statement that $\mathrm{Son}_{\Pi_S}$ dominates
$C_{\varepsilon_S}$ there, window by window. Nothing here decides it; the singular part
$\sum\beta_j\log^+$ of $\varepsilon_S$ is positive, so it lowers $\mathfrak a_T-\mathrm{Son}_{\Pi_S}$.
(2) \emph{Comparison with Theorem~\ref{thm:inertia-nogo}.} For $\frac12\log2<T<\frac12\log3$,
$\mathfrak a_T-\mathrm{Son}_{\Pi_\infty}$ has infinite negative index with eigenvalues bounded away from $0$, while
$\mathfrak a_T-\mathrm{Son}_{\Pi_S}$ is compact: the contact $2$ is absorbed by the semilocal Sonin space. For
$T>\frac12\log3$ the obstruction returns with $3$: $\mathfrak a_T-\mathrm{Son}_{\Pi_S}=-C_{\varepsilon_S}+\Pi^{(\ne2)}_T$,
and the pairs of Theorem~\ref{thm:inertia-nogo} at distance $\log3$ are weakly null, so a compact $C_{\varepsilon_S}$
contributes $o(\|x\|^2)$ while $\Pi^{(\ne2)}_T(x)=-v_3\|\beta\|_2^2\sum|c_j|^2$; the index is infinite. Each new contact
requires enlarging $S$; Lemma~\ref{lem:sl-linear} holds for every finite $S$, but
Proposition~\ref{prop:sl-AS}(iii) and Theorem~\ref{thm:sl-L2} are proved only for $S=\{\infty,2\}$
(for $\{\infty,2,3\}$ one expects $\log^2$ singularities on $\{2^a3^b\}$; this is not proved, and the case of a set
$S(T)$ growing with $T$ is open).
(3) \emph{The inertia of the remainder} must be absorbed by $\mathrm{Son}_{\Pi_S}$ window by window; there is no
uniform-in-$T$ control of $\|C_{\varepsilon_S}\|$ ($\beta_j$ grows like $2^{j/2}$).
(4) \emph{The gap requirement.} Theorem~\ref{thm:sl-gap} is a spectral gap at $1$; a positivity argument through
$\Pi_S$ would need a gap of the density of $\Pi_S$ in the height variable, which is not addressed here.
(5) \emph{Absolute form at $\lambda\ne1$} is conditional on an extension of \cite[Thm.~4.7]{ConnesConsani2021}
(Remark~\ref{rem:sl-level}).
\end{remark}

\begin{remark}[attribution]\label{rem:sl-new}
From \cite{CCquasi2020,CCM2023,CCM2024}: $F_S$, the quasi-inner function and the factor $\rho_2$
(Lemma~\ref{lem:sl-phase}), compactness of the Hankel block, the semilocal Sonin space, and the moment/Jacobi
structure for $S=\{p,\infty\}$ \cite{CCM2024}. From \cite{ConnesConsani2021}: the archimedean (83)--(90) whose
algebra Theorem~\ref{thm:sl-84} transports. New here: Proposition~\ref{prop:sl-AS}(iii) (non-Hilbert--Schmidt,
$\mathcal S_p$, $p>2$), Theorem~\ref{thm:sl-gap} as stated, Lemma~\ref{lem:sl-TC} (in particular the absolute
convergence of the diagonal series), Theorem~\ref{thm:sl-L2} with the constants $\beta_j$ (with the two estimates
noted in its proof), and Theorem~\ref{thm:sl-84}. Lemma~\ref{lem:sl-linear} is not new in substance.
Of the statements of this section, Theorem~\ref{thm:sl-84}(ii) and Corollary~\ref{cor:sl-split} are proved
given \cite[Thm.~4.7]{ConnesConsani2021} at $\lambda=1$; the rest is unconditional.
\end{remark}

\section{Archimedean coercivity on Tate-orthogonal blocks and the finite reduction}\label{sec:coer}

This section localizes the positivity question for $\mathfrak a_T$ on $\PT$ to an explicit
subspace of dimension $2k-2$, the \emph{piecewise Tate exponentials}, by proving that the
archimedean form is coercive, with explicit constants, on the complementary subspace of functions
that satisfy the two Tate conditions \emph{on each of $k$ consecutive blocks}. The coercivity is a
high-pass phenomenon ($m(\tau)\sim\log|\tau|$, Lemma~\ref{lem:mbounds}); the prime form is handled
by a norm bound. Nothing in this section bears on the sign of $\mathfrak a_T$: the finite statement
of Theorem~\ref{thm:coer-schur} inherits the margin of Theorem~\ref{thm:margin} in full
(Remark~\ref{rem:coer-honest}).

\emph{Notation.} All forms are on $L^2(\IT)$ with its metric. The form domain of \eqref{eq:form} is
$\mathcal D_T=\{F\in L^2(\IT):\int\log(2+|\tau|)|\Fh|^2<\infty\}$. By Theorem~\ref{thm:A},
\[
\mathfrak a_T=\Gamma+\PiT,\qquad
\Gamma(F):=\frac1{2\pi}\int_\R m(\tau)|\Fh(\tau)|^2d\tau=C_0\|F\|^2+\Lev(\cF),
\]
\[
\PiT(F)=-\sum_{n\in\NT}v_n\Rel\cF(\log n),
\]
where $\Gamma$ is the form $\Ainf$ of \S\ref{subsec:arch} without its polar term ($\Gamma=\Ainf$ on
$\PT$) and does not depend on $T$ (it is $\Gamma_T$ of \S\ref{sec:inertia}). For $F,G\in\mathcal D_T$ we write $\Gamma(F,G)$,
$\mathfrak a_T(F,G)$ for the polarized (sesquilinear) forms. For a window $\tau>0$ we write
$\lambda_{\mathrm A}(\tau):=\lambda_{\min}(A_\tau|\mathcal P_\tau)$ (so $\lambda_{\mathrm A}(T)=\lamA$);
it is the bottom of $\mathfrak a_\tau$ on the unit sphere of $\mathcal P_\tau\cap\mathcal D_\tau$, it is
attained (the operator $H_\tau$ of the paragraph preceding Lemma~\ref{lem:contA} has compact
resolvent), and it is continuous in $\tau$ (Lemma~\ref{lem:contA}). For $2\tau\le\log2$ the contact
set $\mathcal N_\tau$ is empty, so $\mathfrak a_\tau=\Gamma$ on $L^2(I_\tau)$ and
$\lambda_{\mathrm A}(\tau)$ is a purely archimedean quantity. We use throughout that
$|\cF(u)|\le\cF(0)=\|F\|^2$ and that $\cF(u)=0$ for $|u|\ge2T$.

\subsection{Blocks, the Tate-orthogonal subspace and the reduced kernel}\label{subsec:coer-blocks}

\begin{definition}\label{def:coer-V}
Let $k\ge1$ and let $B_1,\dots,B_k$ be the $k$ consecutive open intervals of length $2\tau$,
$\tau:=T/k$, that partition $\IT$ up to endpoints, $B_j=(-T+2(j-1)\tau,\,-T+2j\tau)$. Put
\[
\mathcal V_k:=\Big\{F\in L^2(\IT):\ \int_{B_j}F(t)e^{t/2}dt=\int_{B_j}F(t)e^{-t/2}dt=0\ \ (1\le j\le k)\Big\},
\qquad E_k:=\PT\ominus\mathcal V_k ,
\]
and $e_j^\pm:=e^{\pm t/2}\mathbf 1_{B_j}$. (Any partition into intervals of lengths $\le\log2$ could be
used in Lemmas~\ref{lem:coer-decomp}--\ref{lem:coer-sector}; equal blocks are needed only from
Lemma~\ref{lem:coer-lowpass} on.)
\end{definition}

\begin{lemma}[the decomposition $\PT=\mathcal V_k\oplus E_k$]\label{lem:coer-decomp}
\emph{(i)} $\mathcal V_k$ is a closed subspace of $\PT$. \emph{(ii)} The orthogonal complement of
$\mathcal V_k$ in $L^2(\IT)$ is $\operatorname{span}\{e_j^\pm\}$, of dimension $2k$.
\emph{(iii)} $E_k=\operatorname{span}\{e_j^\pm\}\cap\PT=\{\sum_jc_j^+e_j^++c_j^-e_j^-:\ A=B=0\}$ has
dimension $2k-2$, and $\PT=\mathcal V_k\oplus E_k$ (orthogonal in $L^2(\IT)$).
\emph{(iv)} $E_k\subset\mathcal D_T$, and $\PT\cap\mathcal D_T=(\mathcal V_k\cap\mathcal D_T)\oplus E_k$.
\end{lemma}
\begin{proof}
(i) The $2k$ conditions are continuous linear functionals; summing them over $j$ gives $A(F)=B(F)=0$.
(ii) $\mathcal V_k$ is the joint kernel of the functionals $F\mapsto\langle F,e_j^\pm\rangle$, and the $2k$
functions $e_j^\pm$ are linearly independent (disjoint supports; $e^{t/2},e^{-t/2}$ independent on an
interval). (iii) $E_k=\mathcal V_k^\perp\cap\PT$ since $\mathcal V_k\subset\PT$. The two functionals
$A,B$ restricted to $\operatorname{span}\{e_j^\pm\}$ are independent: on the two-dimensional space
$\operatorname{span}\{e_1^+,e_1^-\}$ their matrix is
$\begin{psmallmatrix}\int_{B_1}e^t&|B_1|\\|B_1|&\int_{B_1}e^{-t}\end{psmallmatrix}$, of determinant
$\int_{B_1}e^t\int_{B_1}e^{-t}-|B_1|^2>0$ (Cauchy--Schwarz, strict since $e^t$ is not constant). Hence
$\dim E_k=2k-2$. If $F\in\PT$ and $F=v+w$ with $v\in\mathcal V_k$, $w\in\mathcal V_k^\perp$, then
$w=F-v\in\PT$, so $w\in E_k$. (iv) $\widehat{e_j^\pm}(s)=\big[e^{(\pm\frac12+is)t}/(\pm\frac12+is)\big]_{\partial B_j}$,
so $|\widehat{e_j^\pm}(s)|\le2e^{T/2}(\frac14+s^2)^{-1/2}$ and $\int\log(2+|s|)|\widehat{e_j^\pm}|^2<\infty$.
Since $\mathcal D_T$ is a linear space containing $E_k$, the $\mathcal V_k$-component of an $F\in\PT\cap\mathcal D_T$
lies in $\mathcal D_T$.
\end{proof}

\begin{lemma}[the reduced kernel]\label{lem:coer-K1}
Put $K_1(u):=K(u)-2e^{-u/2}=\dfrac{2e^{-5u/2}}{1-e^{-2u}}=2\sum_{j\ge1}e^{-(2j+\frac12)u}$ ($u>0$), with
$K$ as in Lemma~\ref{lem:K}. Then $0<K_1(u)\le u^{-1}$ for all $u>0$, and $K_1(u)=u^{-1}-\frac32+O(u)$ as
$u\downarrow0$. Let $F,G\in L^2(\IT)$ be real, $\supp F\subset\overline B$, $\supp G\subset\overline{B'}$,
with $B,B'$ disjoint intervals and $B'$ to the right of $B$. Then
\begin{equation}\label{eq:coer-cross}
\Gamma(F+G)=\Gamma(F)+\Gamma(G)-\iint F(x)G(y)K(y-x)\,dx\,dy ,
\end{equation}
and if $\int F(x)e^{x/2}dx=0$ or $\int G(y)e^{-y/2}dy=0$, the kernel $K$ in \eqref{eq:coer-cross} may be
replaced by $K_1$.
\end{lemma}
\begin{proof}
The three expressions for $K_1$ agree by Lemma~\ref{lem:K} ($K=2e^{-u/2}\sum_{j\ge0}e^{-2ju}$).
$K_1\le u^{-1}$ is $e^{5u/2}-e^{u/2}\ge2u$: the function $g(u)=e^{5u/2}-e^{u/2}-2u$ has $g(0)=g'(0)=0$ and
$g''=\frac{25}4e^{5u/2}-\frac14e^{u/2}>0$. The expansion follows from
$\frac1{2\sinh(u/2)}=u^{-1}-\frac u{24}+O(u^3)$, $\frac1{2\cosh(u/2)}=\frac12+O(u^2)$ and
$2e^{-u/2}=2-u+O(u^2)$. For \eqref{eq:coer-cross}, apply \eqref{eq:Aform} with $\NT=\emptyset$ to
$F+G$: $c_{F+G}(u)=\cF(u)+c_G(u)+\int F(t)G(t-u)dt+\int G(t)F(t-u)dt$, and for $u>0$ the first cross
integral vanishes ($t\in B$, $t-u\in B'$ is impossible), while the second equals
$\int\!\!\int F(x)G(y)\mathbf 1_{y-x=u}$; integrating against $K(u)du$ over $u>0$ gives
\eqref{eq:coer-cross} (the term $\|F+G\|^2=\|F\|^2+\|G\|^2$ has no cross part). Finally
$2e^{-(y-x)/2}=2e^{-y/2}e^{x/2}$, so $\iint F(x)G(y)\,2e^{-(y-x)/2}dxdy=2\int Fe^{x/2}\int Ge^{-y/2}$,
which vanishes under either Tate condition.
\end{proof}
The cross term in \eqref{eq:coer-cross} carries the factor $1$ for the \emph{unordered} pair $\{B,B'\}$; the
polarized form is $\Gamma(F,G)=-\frac12\iint FGK$. This normalization is used in Remark~\ref{rem:coer-nogo}.

\begin{lemma}[a block sector is a short window]\label{lem:coer-sector}
Let $B\subset\IT$ be an interval of length $2\tau$ with $2\tau\le\log2$. Then
\[
\inf\Big\{\mathfrak a_T(F):\ F\in\mathcal D_T,\ \supp F\subset\overline B,\ \int_BFe^{\pm t/2}=0,\ \|F\|=1\Big\}
=\lambda_{\mathrm A}(\tau),
\]
independently of $T$ and of the position of $B$; the infimum is attained and depends continuously on
$\tau$. In particular, for $F\in\mathcal V_k$ with $\tau=T/k$, $2\tau\le\log2$, and $F_j:=F\mathbf 1_{B_j}$,
\begin{equation}\label{eq:coer-diag}
\sum_{j=1}^k\Gamma(F_j)\ \ge\ \lambda_{\mathrm A}(\tau)\,\|F\|^2 .
\end{equation}
\end{lemma}
\begin{proof}
Let $B=I_\tau+a$. For $F$ supported in $\overline B$ and $F_a(t):=F(t+a)$, supported in $\overline{I_\tau}$,
one has $c_{F_a}=\cF$, hence $\Gamma(F_a)=\Gamma(F)$ by \eqref{eq:Aform}, and
$\int F_ae^{\pm t/2}=e^{\mp a/2}\int Fe^{\pm t/2}$, so the two Tate conditions on $B$ are those of
$\mathcal P_\tau$ for $F_a$. Since $\cF(\log n)=0$ for $\log n\ge2\tau$ and $\log2\ge2\tau$, also
$\PiT(F)=0$, so $\mathfrak a_T(F)=\Gamma(F)=\mathfrak a_\tau(F_a)$ and the infimum is
$\lambda_{\mathrm A}(\tau)$. Attainment and continuity are the facts recalled in the notation paragraph.
\eqref{eq:coer-diag} is the sum over blocks.
\end{proof}

\subsection{A lower bound for the short window}\label{subsec:coer-short}

\begin{lemma}\label{lem:coer-K}
$K(u)\ge u^{-1}$ for $0<u\le2$.
\end{lemma}
\begin{proof}
By Lemma~\ref{lem:K}, $K(u)-u^{-1}=\frac{u-2\sinh(u/2)}{2u\sinh(u/2)}+\frac1{2\cosh(u/2)}$. From
$\tanh x\le x$, $2\sinh(u/2)\le u\cosh(u/2)$, so $u-2\sinh(u/2)\ge-u(\cosh(u/2)-1)=-2u\sinh^2(u/4)$ and
the first term is $\ge-\frac{\sinh(u/4)}{2\cosh(u/4)}\ge-\frac u8$. Hence
$K(u)-u^{-1}\ge\frac1{2\cosh(u/2)}-\frac u8\ge0$ whenever $u\cosh(u/2)\le4$, which holds for $u\le2$
($2\cosh1=3.09$).
\end{proof}

\begin{lemma}[the tail]\label{lem:coer-tail}
For $\tau>0$, $\int_{2\tau}^\infty K(u)\,du\ge\log\frac2\tau+\frac\pi2-\tau$.
\end{lemma}
\begin{proof}
By Theorem~\ref{thm:S}(iii) the integral is $\log\coth\frac\tau2+\pi-2\arctan e^\tau$; $\coth x\ge x^{-1}$
gives the first term, and $\frac{d}{d\tau}(\pi-2\arctan e^\tau)=-\frac1{\cosh\tau}\ge-1$ with value
$\frac\pi2$ at $\tau=0$ gives the second.
\end{proof}

\begin{lemma}[the truncated logarithm]\label{lem:coer-w}
Put $w(s):=\int_0^2\frac{1-\cos sv}v\,dv=\log(2|s|)+\gamma-\operatorname{Ci}(2|s|)$ for $s\ne0$, $w(0)=0$.
Then $w$ is even, increasing in $|s|$, and for every $0<R\le\pi$ and all $s\in\R$,
$w(s)\ge w(R)\min\{1,s^2/R^2\}$.
\end{lemma}
\begin{proof}
$w'(s)=\int_0^2\sin(sv)\,dv=\frac{1-\cos2s}s\ge0$ for $s>0$. Next $w(s)/s^2=\int_0^2v\,\frac{1-\cos sv}{(sv)^2}dv$,
and $x\mapsto\frac{1-\cos x}{x^2}=\frac12\big(\frac{\sin(x/2)}{x/2}\big)^2$ is decreasing on $(0,2\pi)$;
for $0<s\le R\le\pi$ and $v\le2$ one has $sv\le2\pi$, so $w(s)/s^2\ge w(R)/R^2$. For $s\ge R$ use monotonicity.
\end{proof}

\begin{lemma}[a double Poincar\'e inequality]\label{lem:coer-poinc}
Let $h\in L^2(\R)$ with $\supp h\subset[-1,1]$ and $h''\in L^2(\R)$ (distributional derivative). Then
$\|h''\|^2\ge\pi^2\|h'\|^2\ge\frac{\pi^4}4\|h\|^2$.
\end{lemma}
\begin{proof}
$h,h''\in L^2(\R)$ give $h\in H^2(\R)\subset C^1(\R)$, and the support condition gives $h(\pm1)=h'(\pm1)=0$.
The second inequality is the Dirichlet Poincar\'e inequality for $h\in H^1_0(-1,1)$ (first eigenvalue
$(\pi/2)^2$). For the first, $g:=h'\in H^1_0(-1,1)$ with $\int_{-1}^1g=h(1)-h(-1)=0$. Expand $g=\sum_{n\ge1}c_n\phi_n$
in the orthonormal Dirichlet eigenbasis $\phi_n(x)=\sin\frac{n\pi(x+1)}2$ of $(-1,1)$, so that
$\|g\|^2=\sum c_n^2$ and $\|g'\|^2=\sum(\frac{n\pi}2)^2c_n^2$. Since $\int\phi_n=\frac4{n\pi}$ for odd $n$ and
$0$ for even $n$, the constraint reads $c_1=-\sum_{n\ge3\ \mathrm{odd}}c_n/n$. The even-indexed terms already
satisfy $(\frac{n\pi}2)^2\ge\pi^2$, and for the odd-indexed ones
\[
\sum_{n\ \mathrm{odd}}\Big(\frac{n^2\pi^2}4-\pi^2\Big)c_n^2=\frac{\pi^2}4\Big[\sum_{n\ge3\ \mathrm{odd}}(n^2-4)c_n^2-3c_1^2\Big]\ \ge\ 0,
\]
because by Cauchy--Schwarz $c_1^2\le S\sum_{n\ge3\ \mathrm{odd}}(n^2-4)c_n^2$ with
$S=\sum_{n\ge3\ \mathrm{odd}}\frac1{n^2(n^2-4)}=\frac14\big(\frac43-\frac{\pi^2}8\big)=0.0249<\frac13$
(partial fractions and $\sum_{n\ \mathrm{odd}}n^{-2}=\pi^2/8$). Hence $\|g'\|^2\ge\pi^2\|g\|^2$. (The constant
$\pi^2$ is attained by $g=\phi_2$; the constraint touches only the even part of $g$.)
\end{proof}

\begin{theorem}[the short window]\label{thm:coer-A}
Let $0<\tau\le1$ with $2\tau\le\log2$, and put $a:=\tau^2/4$. For every $R\in(0,\pi]$,
\begin{equation}\label{eq:coer-A}
\lambda_{\mathrm A}(\tau)\ \ge\ C_0+\log\frac2\tau+\frac\pi2-\tau+w(R)\Big[1-\frac{16(R^2+a)^3}{27\pi^4R^2}\Big].
\end{equation}
In particular, with $R=2$ ($w(2)=\log4+\gamma-\operatorname{Ci}(4)=2.1045$),
\[
\lambda_{\mathrm A}(\tau)\ \ge\ \log\frac1\tau-1.209-\tau\qquad(0<\tau\le0.1),
\]
so $\lambda_{\mathrm A}(\tau)\ge\log\frac1\tau-1.31$ on $(0,0.1]$.
\end{theorem}
\begin{proof}
Let $F\in\mathcal P_\tau\cap\mathcal D_\tau$, $\|F\|=1$. Since $\mathcal N_\tau=\emptyset$ and $\cF(u)=0$ for
$u\ge2\tau$, \eqref{eq:Aform} gives
\[
\mathfrak a_\tau(F)=C_0+\int_{2\tau}^\infty K+\int_0^{2\tau}\big(1-\Rel\cF(u)\big)K(u)\,du .
\]
The tail is bounded by Lemma~\ref{lem:coer-tail}. Put $G(y):=\tau^{1/2}F(\tau y)$, so that $\supp G\subset[-1,1]$,
$\|G\|=1$, $\cF(u)=c_G(u/\tau)$, and
$\int_0^{2\tau}(1-\Rel\cF)K=\int_0^2(1-\Rel c_G(v))\,\tau K(\tau v)\,dv\ge\int_0^2\frac{1-\Rel c_G(v)}v\,dv$,
using $1-\Rel c_G\ge0$ and Lemma~\ref{lem:coer-K} ($\tau v\le2\tau\le2$). By Plancherel,
$1-\Rel c_G(v)=\frac1{2\pi}\int(1-\cos sv)|\widehat G(s)|^2ds$, so by Tonelli the last integral equals
$\frac1{2\pi}\int w|\widehat G|^2$, and by Lemma~\ref{lem:coer-w}, for $0<R\le\pi$,
\[
\frac1{2\pi}\int w|\widehat G|^2\ \ge\ w(R)\Big[1-\frac1{2\pi}\int\Big(1-\frac{s^2}{R^2}\Big)^+|\widehat G(s)|^2ds\Big].
\]
The Tate conditions $\int Fe^{\pm t/2}=0$ read $\int Ge^{\pm\tau y/2}=0$, i.e.\ $\widehat G(\mp i\tau/2)=0$.
$\widehat G$ is entire of exponential type $\le1$ and in $L^2(\R)$, so $H(s):=\widehat G(s)/(s^2+a)$ is entire of
exponential type $\le1$ and in $L^2(\R)$ ($s^2+a\ge a>0$ on $\R$); by the Paley--Wiener theorem $H=\widehat h$ with
$h\in L^2(\R)$, $\supp h\subset[-1,1]$, and $G=-h''+ah$, so $h''\in L^2$. Then
$1=\|G\|^2=\frac1{2\pi}\int(s^2+a)^2|H|^2=\|h''\|^2+2a\|h'\|^2+a^2\|h\|^2\ge\|h''\|^2\ge\frac{\pi^4}4\|h\|^2$
(Lemma~\ref{lem:coer-poinc}). On the other hand $(1-\frac{s^2}{R^2})^+(s^2+a)^2\le M_R:=\frac{4(R^2+a)^3}{27R^2}$
(the maximum of $(1-x/R^2)(x+a)^2$ on $0\le x\le R^2$ is at $x=(2R^2-a)/3$), so
$\frac1{2\pi}\int(1-s^2/R^2)^+|\widehat G|^2\le M_R\|h\|^2\le\frac{4M_R}{\pi^4}$. Collecting terms gives
\eqref{eq:coer-A}. For $R=2$ and $a\le\frac1{400}$ the bracket is $\ge0.9025$ and $w(2)\cdot0.9025\ge1.899$;
with $C_0+\log2+\frac\pi2=-3.1082$ this gives $\log\frac1\tau-1.209-\tau$.
\end{proof}

\begin{remark}\label{rem:coer-A}
(1) The only numerical input is the value $w(2)$ (a tabulated elementary quantity, $\operatorname{Ci}(4)=-0.1410$).
Optimizing $R$ in \eqref{eq:coer-A} gives $R=2.165$ and $\log\frac1\tau-1.181-\tau$ (numerical
supremum). (2) The true value is $\lambda_{\mathrm A}(\tau)=\log\frac1\tau-0.510\pm0.001$ for $\tau\in[0.02,0.1]$
(Galerkin; numerical, upper bounds, not proved). The loss of the proof is in
$\int_{|s|<R}|H|^2\le\|h\|^2$ and in the constant $\pi^4/4$ (the clamped-beam constant $31.285$ would
improve it; not used). (3) Theorem~\ref{thm:coer-A} is a complement to Lemma~\ref{lem:contA} at the
left end: $\lambda_{\mathrm A}(\tau)\to+\infty$ as $\tau\to0$ at the rate $\log\frac1\tau$, consistently with
Corollary~\ref{cor:mono} ($\lambda_{\mathrm A}(\tau)\le C_0+\varphi((3\pi/2\tau)^2)=\log\frac1\tau+O(1)$).
(4) The theorem is about $\lambda_{\mathrm A}(\tau)$, the bottom over $\mathcal P_\tau$ (two Tate conditions
on $F$), not about $\lambda^{\mathrm{prim}}_\tau$ of Definition~\ref{def:PD} (two conditions on the autocorrelation,
Remark~\ref{rem:Kprim}); no statement about $\lambda^{\mathrm{prim}}_\tau$ is made here.
\end{remark}

\begin{remark}[precedents]\label{rem:coer-prec}
Coercivity of the Weil form on short windows in finite codimension is classical in qualitative form:
Yoshida \cite[Lemma~3]{Yoshida1992} proves that for given $a_0,\mu>0$ there is $N$ with $\langle\phi,\phi\rangle_W\ge\mu\|\phi\|^2$
for $\phi$ supported in $[-a,a]$, $0<a\le a_0$, orthogonal to $N$ Fourier modes (no rate in $a$); Connes and Consani
\cite[Thm.~3.6, Cor.~3.8, Rem.~3.9]{ConnesConsani2021} write the archimedean part as $-2\cdot\mathrm{Id}$ plus a
Hilbert--Schmidt operator and obtain positivity for $|I|\le2\cdot0.1404$ after finitely many linear conditions.
The leading term $\log(1/\tau)$ of the short-window law has been observed numerically for Suzuki's operator (without
primes) by Kim et al.\ \cite{Kim2026}. Finite truncations of forms with a difference kernel on $[-L,L]$ are also studied
in \cite{CvS2025}. What is new in Theorems~\ref{thm:coer-A} and \ref{thm:coer-B} is the rate $\log(1/\tau)$ with an
explicit constant, under the two Tate conditions, and the block structure.
\end{remark}

\subsection{Coercivity of the archimedean form on $\mathcal V_k$}\label{subsec:coer-B}

For $R>0$ let $P_R$ be the orthogonal projection of $L^2(\R)$ onto the band $\{|s|\le R\}$,
$\widehat{P_RF}=\Fh\mathbf 1_{|s|\le R}$.

\begin{lemma}[low-pass leakage of $\mathcal V_k$]\label{lem:coer-lowpass}
Let $F\in\mathcal V_k$ with equal blocks of half-length $\tau$, and $R>0$. Then
\[
\|P_RF\|^2\le\delta(R\tau)\,\|F\|^2,\qquad \delta(\theta):=\frac{\cosh^2\tau}{6}\Big(\theta^2+\frac{\tau^2}4\Big)^2 .
\]
\end{lemma}
\begin{proof}
$G:=P_RF$ is entire. Fix a block $B$ with centre $c$ and put, on $B$,
$u(x):=G(c)\cosh\frac{x-c}2+2G'(c)\sinh\frac{x-c}2\in\operatorname{span}\{e^{x/2},e^{-x/2}\}$, so that
$u(c)=G(c)$, $u'(c)=G'(c)$, $u''=u/4$. Then $r:=G-u$ satisfies $r(c)=r'(c)=0$ and
$r''-\frac r4=G''-\frac G4=:H$, whence $r(x)=\int_c^x2\sinh\frac{x-t}2\,H(t)\,dt$. By Cauchy--Schwarz and
$2\sinh\frac z2\le z\cosh\frac z2\le z\cosh\tau$ ($0\le z\le\tau$),
$|r(x)|^2\le\cosh^2\tau\,\frac{|x-c|^3}3\,\|H\|_{L^2(B)}^2$, so $\int_Br^2\le\cosh^2\tau\,\frac{\tau^4}6\|H\|_{L^2(B)}^2$.
Since $F_B:=F\mathbf 1_B$ is orthogonal to $e^{\pm x/2}\mathbf 1_B$, $\langle F_B,G\rangle=\langle F_B,r\rangle\le\|F_B\|\,\|r\|_{L^2(B)}$.
Summing over the blocks and using Cauchy--Schwarz once more,
\begin{align*}
\|G\|^2=\langle F,G\rangle&\le\|F\|\Big(\sum_B\|r\|_{L^2(B)}^2\Big)^{1/2}\le\|F\|\cosh\tau\,\frac{\tau^2}{\sqrt6}\,\|G''-\tfrac14G\|_{L^2(\R)}\\
&\le\|F\|\cosh\tau\,\frac{\tau^2}{\sqrt6}\Big(R^2+\frac14\Big)\|G\|,
\end{align*}
the first equality because $P_R$ is an orthogonal projection, the last inequality because
$\widehat G$ is supported in $|s|\le R$. Squaring gives the claim, as $\tau^4(R^2+\frac14)^2=((R\tau)^2+\frac{\tau^2}4)^2$.
\end{proof}

\begin{theorem}[coercivity on $\mathcal V_k$]\label{thm:coer-B}
Let $k\ge1$, $\tau=T/k\le1$, $\kappa:=\cosh^2\tau$, $a:=\tau^2/4$, $\delta$ as in
Lemma~\ref{lem:coer-lowpass}, $\theta_*:=(\sqrt{6/\kappa}-a)^{1/2}$ (the root of $\delta(\theta_*)=1$), and
$\ell(\rho):=\log\frac\rho{2\pi}-\frac1\rho-\frac1{2\rho^2}$ (the lower bound of Lemma~\ref{lem:mbounds}(ii)).
Then for every $F\in\mathcal V_k\cap\mathcal D_T$,
\begin{equation}\label{eq:coer-B}
\Gamma(F)\ \ge\ \Big(\log\frac1\tau-C_B(\tau)\Big)\|F\|^2,\qquad
C_B(\tau):=\log\frac1\tau-C_0\,\delta(\tau)-\frac{2\kappa}3\big[\Phi_\tau(\theta_*)-\Phi_\tau(\tau)\big],
\end{equation}
where $\Phi_\tau$ is the elementary primitive
\[
\Phi_\tau(\theta):=\Big(\frac{\theta^4}4+\frac{a\theta^2}2\Big)\big(\log\theta-\log(2\pi\tau)\big)-\frac{\theta^4}{16}-\frac{a\theta^2}4
-\tau\Big(\frac{\theta^3}3+a\theta\Big)-\frac{\tau^2\theta^2}4-\frac{a\tau^2}2\log\theta
\]
of $(\theta^3+a\theta)\,\ell(\theta/\tau)$. Moreover $C_B(\tau)\to\frac14-\log\frac{6^{1/4}}{2\pi}=1.6399$ as
$\tau\to0$, and $C_B(0.02)=1.6573$, $C_B(0.0619)=1.6957$, $C_B(0.1)=1.7331$ (evaluations of the closed form).
\end{theorem}
\begin{proof}
Let $\|F\|=1$ and $\mu:=\frac1{2\pi}|\Fh|^2ds$, a probability measure, so that $\Gamma(F)=\int m\,d\mu$ by
\eqref{eq:form} with $\NT=\emptyset$. By Lemma~\ref{lem:mbounds}(i) and \eqref{eq:m}, $m$ is even, real-analytic,
increasing on $[0,\infty)$ with $m(0)=C_0$; hence $m(s)=C_0+\int_0^{|s|}m'(\rho)d\rho$ with $m'\ge0$ and, by Tonelli,
\[
\int m\,d\mu=C_0+\int_0^\infty m'(\rho)\,\mu(|s|\ge\rho)\,d\rho\ \ge\ C_0+\int_0^{\rho_*}m'(\rho)\big(1-\delta(\rho\tau)\big)d\rho,
\qquad\rho_*:=\theta_*/\tau,
\]
since $\mu(|s|\ge\rho)=1-\|P_{\rho}F\|^2+\mu(|s|=\rho)\ge1-\delta(\rho\tau)$ by Lemma~\ref{lem:coer-lowpass} and the
integrand is nonnegative on $[0,\rho_*]$, where $\delta(\rho\tau)\le1$. Integrating by parts on $[0,\rho_*]$, with
vanishing boundary terms ($1-\delta(\theta_*)=0$ and $m(0)-C_0=0$),
\[
\int_0^{\rho_*}m'(\rho)(1-\delta(\rho\tau))d\rho=\int_0^{\theta_*}\big(m(\theta/\tau)-C_0\big)\delta'(\theta)\,d\theta,
\]
so that
\[
\int m\,d\mu\ \ge\ C_0\,\delta(0)+\int_0^{\theta_*}m(\theta/\tau)\,\delta'(\theta)\,d\theta ,
\]
using $\int_0^{\theta_*}\delta'=1-\delta(0)$. Here $\delta'(\theta)=\frac{2\kappa}3\theta(\theta^2+a)\ge0$. By
Lemma~\ref{lem:mbounds}(i),(ii), $m\ge\max\{C_0,\ell\}$ on $(0,\infty)$; we use $m\ge C_0$ on $[0,\tau]$ (i.e.\ $\rho\le1$)
and $m\ge\ell$ on $[\tau,\theta_*]$ (note $\ell(1)=-3.34>C_0$, and $\ell$ is increasing, so no further loss occurs
there). Hence
\[
\int m\,d\mu\ \ge\ C_0\,\delta(0)+C_0\big(\delta(\tau)-\delta(0)\big)+\int_\tau^{\theta_*}\ell(\theta/\tau)\,\delta'(\theta)\,d\theta
=C_0\,\delta(\tau)+\frac{2\kappa}3\big[\Phi_\tau(\theta_*)-\Phi_\tau(\tau)\big],
\]
which is \eqref{eq:coer-B}; the primitive is checked by differentiation. As $\tau\to0$: $\kappa\to1$, $a\to0$,
$\theta_*\to6^{1/4}$, $\delta(\tau)\to0$, $\frac{2\kappa}3\Phi_\tau(\tau)\to0$, and
$\frac{2\kappa}3\Phi_\tau(\theta_*)=\log\theta_*-\log(2\pi\tau)-\frac14+o(1)$ (the $a$- and $\tau$-terms are $o(1)$),
whence $C_B(\tau)\to\log2\pi-\log6^{1/4}+\frac14$. The three values are evaluations of the closed form (double precision).
\end{proof}

\begin{remark}\label{rem:coer-B}
(1) Using $m\ge\max\{C_0,\ell\}$ with the exact crossover instead of the split at $\rho=1$ changes $C_B$ by less than
$5\cdot10^{-5}$ on $(0,0.1]$. (2) The uniform statement ``$C_B(\tau)\le1.74$ on $(0,0.1]$'' holds on a grid of
twelve points where $C_B$ is increasing (numerical, not proved); monotonicity of $C_B$ is not proved and is not used:
Corollary~\ref{cor:coer-L1} only needs $C_B$ at the specific $\tau=T/k$. (3) The bound loses about $1.0$ against
the Galerkin value of $\inf\Gamma|_{\mathcal V_k}$ (e.g.\ $1.087$ versus $2.093$ at $T=\frac12\log5$, $k=13$,
numerical); the loss is in Lemma~\ref{lem:coer-lowpass}, whose constant is about $25$ times larger than needed
on random elements of $\mathcal V_k$. (4) The theorem uses the whole multiplier $m$ and the Tate
conditions only through the geometry of $\mathcal V_k$ (Lemma~\ref{lem:coer-lowpass}); it does not separate the
block-diagonal part \eqref{eq:coer-diag} from the coupling \eqref{eq:coer-cross}. See Remark~\ref{rem:coer-nogo}
for why that separation is not viable.
\end{remark}

\begin{proposition}[the norm of the prime form]\label{prop:coer-prime}
\emph{(i)} $|\PiT(F)|\le\Aconst\|F\|^2$ for all $F\in L^2(\IT)$ and all $T$.
\emph{(ii)} For $\ell>0$ let $L(T,\ell):=\lceil2T/\ell\rceil$. Then $|\Rel\cF(\ell)|\le\cos\frac{\pi}{L(T,\ell)+1}\,\|F\|^2$,
hence
\[
\|\PiT\|\ \le\ \sum_{n\in\NT}v_n\cos\frac\pi{L(T,\log n)+1}
=\begin{cases}
\Aconst/2, & T\le\log2,\\[2pt]
\dfrac{v_2}{\sqrt2}+\dfrac{v_3+v_4}2=1.6740, & \log2<T\le\tfrac12\log5 .
\end{cases}
\]
Here $\frac{v_2}{\sqrt2}+\frac{v_3+v_4}2=\log2+\frac{\log3}{\sqrt3}+\frac{\log2}2$; in the first case the bound is
$\frac{v_2}2=0.4901$ for $T\le\frac12\log3$ and $\frac{v_2+v_3}2=1.1244$ for $\frac12\log3<T\le\log2$. The bounds hold on all of $L^2(\IT)$, in particular on $\mathcal V_k$.
\end{proposition}
\begin{proof}
(i) is $|\cF(\log n)|\le\|F\|^2$. (ii) Let $S_\ell$ be the shift $F\mapsto F(\cdot-\ell)\mathbf 1_{\IT}$ on
$L^2(\IT)$, so that $\Rel\cF(\ell)=\frac12\langle(S_\ell+S_\ell^*)F,F\rangle$. The map
$F\mapsto\big(F(-T+r+j\ell)\big)_{0\le j<n(r)}$, $r\in[0,\ell)$, $n(r):=\#\{j\ge0:-T+r+j\ell<T\}\le L(T,\ell)$, is unitary
from $L^2(\IT)$ onto $\int^\oplus_{[0,\ell)}\C^{n(r)}dr$, and it carries $S_\ell$ to the fibrewise shift $j\mapsto j+1$.
Hence $S_\ell+S_\ell^*$ acts on the fibre $\C^{n}$ as the adjacency matrix of the path with $n$ vertices, whose
eigenvalues are $2\cos\frac{j\pi}{n+1}$, $1\le j\le n$; its norm is $2\cos\frac\pi{n+1}\le2\cos\frac\pi{L+1}$. For
$T\le\log2$ one has $2T/\log n\le2$ for all $n\ge2$, so $L=2$ and $\cos\frac\pi3=\frac12$. For $\log2<T\le\frac12\log5$:
$\NT=\{2,3,4\}$, $L(T,\log2)=3$ since $2<2T/\log2\le\log5/\log2=2.32$, and $L(T,\log3)=L(T,\log4)=2$ since
$2T\le\log5<2\log3$; $\cos\frac\pi4=\frac1{\sqrt2}$. The values use $v_2=\sqrt2\log2$, $v_3=2\log3/\sqrt3$, $v_4=\log2$.
\end{proof}
For $\frac12\log2<T<\frac12\log3$ the bound $\Aconst/2$ is the exact norm (Proposition~\ref{thm:inertia-B2}, same
fibre decomposition). The bound $\|\PiT\|\le\Aconst/2$ does \emph{not} hold for all $T$: once $T>\log2$ the orbit of
$\log2$ has length $3$ and the path norm is $\sqrt2>1$; the proposition is the corrected statement.

\begin{corollary}[coercivity of the truncated form on $\mathcal V_k$]\label{cor:coer-L1}
Let $k\ge1$ with $\tau=T/k\le1$, and put $c(T,k):=\inf\{\mathfrak a_T(F):F\in\mathcal V_k\cap\mathcal D_T,\ \|F\|=1\}$. Then
\begin{equation}\label{eq:coer-L1}
c(T,k)\ \ge\ \log\frac kT-C_B\Big(\frac Tk\Big)-\|\PiT\| ,
\end{equation}
with $\|\PiT\|$ bounded by Proposition~\ref{prop:coer-prime}. In particular (evaluations of the closed form):
\begin{center}
\begin{tabular}{@{}lllll@{}}
\toprule
$T$ & $k$ & $\dim E_k$ & bound for $\|\PiT\|$ & lower bound for $c(T,k)$\\
\midrule
$\frac12\log3$ & $8$ & $14$ & $0.4901$ & $0.486$\\
$\log2$ & $12$ & $22$ & $1.1244$ & $0.035$\\
$0.69$ & $14$ & $26$ & $1.1244$ & $0.202$\\
$\frac12\log5$ & $25$ & $48$ & $1.6740$ & $0.094$\\
$\frac12\log5$ & $26$ & $50$ & $1.6740$ & $0.134$\\
$\frac12\log5$ & $30$ & $58$ & $1.6740$ & $0.281$\\
$\frac12\log5$ & $85$ & $168$ & $\Aconst=2.9420$ & $0.070$\\
\bottomrule
\end{tabular}
\end{center}
Moreover, for every $T>0$ there is $k_0(T)$ such that $c(T,k)>0$ for all $k\ge k_0(T)$; indeed the right-hand side
of \eqref{eq:coer-L1} tends to $+\infty$ with $k$, because $C_B(T/k)\to1.6399$ (Theorem~\ref{thm:coer-B}).
\end{corollary}
\begin{proof}
$\mathfrak a_T=\Gamma+\PiT$ on $\mathcal D_T$, Theorem~\ref{thm:coer-B} and Proposition~\ref{prop:coer-prime}.
\end{proof}

\begin{remark}\label{rem:coer-size}
The size of $k_0(T)$: by \eqref{eq:coer-L1} and Proposition~\ref{prop:coer-prime}(i), $c(T,k)>0$ as soon as
$\log\frac kT>C_B(T/k)+\Aconst$, so $k_0(T)\le Te^{\Aconst+C_B(T/k_0)}$, where $C_B(T/k_0)$ is within $0.1$ of
its limit once $T/k_0\le0.1$ (numerical; the limit itself is proved). Since $\Aconst\asymp e^T$, the dimension
$2k-2$ of $E_k$ grows doubly exponentially in $T$; on the declared range $T\le\frac12\log5$ it stays below $60$.
The Galerkin values of $c(T,k)$ are larger than the bounds by about $1.0$ (Remark~\ref{rem:coer-B}(3)):
e.g.\ $c(\frac12\log5,13)=1.241$ (numerical, converged in the Galerkin degree), where \eqref{eq:coer-L1} gives
$-0.59$.
\end{remark}

\subsection{The finite reduction}\label{subsec:coer-schur}

\begin{theorem}[reduction to the piecewise Tate exponentials]\label{thm:coer-schur}
Fix $T>0$ and $k\ge1$ with $c:=c(T,k)>0$ (Corollary~\ref{cor:coer-L1}: every $k\ge k_0(T)$). Put
$\mathcal V:=\mathcal V_k\cap\mathcal D_T$ and $M:=\Aconst-C_0$.
\emph{(i)} $\mathfrak a_T$ is an inner product on $\mathcal V$ for which $\mathcal V$ is complete, and for every
$e\in E_k$ the functional $v\mapsto\mathfrak a_T(e,v)$ is bounded on $(\mathcal V,\mathfrak a_T)$; hence there is a unique
$Ke\in\mathcal V$ with $\mathfrak a_T(Ke,v)=\mathfrak a_T(e,v)$ for all $v\in\mathcal V$. The map $K:E_k\to\mathcal V$ is linear,
and $\|K\|:=\sup\{\|Ke\|:\|e\|=1\}<\infty$.
\emph{(ii)} For $v\in\mathcal V$ and $e\in E_k$,
\begin{equation}\label{eq:coer-schur}
\begin{split}
\mathfrak a_T(v+e)&=\mathfrak a_T(v+Ke)+\mathfrak s_{T,k}(e),\\
\mathfrak s_{T,k}(e)&:=\mathfrak a_T(e)-\mathfrak a_T(Ke)=\mathfrak a_T(e)-\sup_{u\in\mathcal V}\big[2\Rel\mathfrak a_T(e,u)-\mathfrak a_T(u)\big].
\end{split}
\end{equation}
Let $S_{T,k}$ be the matrix of the real quadratic form $\mathfrak s_{T,k}$ in an orthonormal basis of $E_k$ (a
real symmetric $(2k-2)\times(2k-2)$ matrix) and $\lambda_S$ its least eigenvalue. Then
\begin{equation}\label{eq:coer-sandwich}
\frac{\min\{c,\lambda_S\}}{(1+\|K\|)^2}\ \le\ \lamA\ \le\ \min\{c,\lambda_S\}.
\end{equation}
\emph{(iii)} Consequently $\lamA>0$ if and only if $S_{T,k}\succ0$, and $\lamA\ge0$ if and only if $S_{T,k}\succeq0$.
\end{theorem}
\begin{proof}
(i) By Lemma~\ref{lem:mbounds}, $\sigma_T\ge C_0-\Aconst$, so $\mathfrak a_T+M\|\cdot\|^2\ge0$ on $\mathcal D_T$, and
by Lemma~\ref{lem:mbounds}(ii) and $m\ge C_0$ there are constants $c_1,c_2>0$ (depending on $T$) with
$c_1\int\log(2+|s|)|\Fh|^2\le\mathfrak a_T(F)+(M+1)\|F\|^2\le c_2\int\log(2+|s|)|\Fh|^2$ on $\mathcal D_T$. On $\mathcal V$,
$\mathfrak a_T\ge c\|\cdot\|^2$ gives $\mathfrak a_T(v)\ge\frac c{c+M+1}\big(\mathfrak a_T(v)+(M+1)\|v\|^2\big)$, so
$\mathfrak a_T$ is equivalent on $\mathcal V$ to the weighted norm $\int\log(2+|s|)|\widehat v|^2$, for which
$\mathcal D_T$ is complete and $\mathcal V_k\cap\mathcal D_T$ is closed ($\mathcal V_k$ is $L^2$-closed and the weighted norm
dominates the $L^2$ norm). For $e\in E_k\subset\mathcal D_T$ (Lemma~\ref{lem:coer-decomp}(iv)), the Cauchy--Schwarz
inequality for the nonnegative form $\mathfrak a_T+M\langle\cdot,\cdot\rangle$ gives
$|\mathfrak a_T(e,v)|\le(\mathfrak a_T(e)+M\|e\|^2)^{1/2}(\mathfrak a_T(v)+M\|v\|^2)^{1/2}+M\|e\|\|v\|\le C_e\,\mathfrak a_T(v)^{1/2}$,
so $Ke$ exists by the Riesz representation theorem; $K$ is linear, and $\|K\|<\infty$ since $\dim E_k<\infty$.
(ii) $\mathfrak a_T(v+Ke)=\mathfrak a_T(v)+2\Rel\mathfrak a_T(Ke,v)+\mathfrak a_T(Ke)=\mathfrak a_T(v)+2\Rel\mathfrak a_T(e,v)+\mathfrak a_T(Ke)$,
which rearranges to \eqref{eq:coer-schur}; the variational expression is $\mathfrak a_T(Ke)=\sup_u[2\Rel\mathfrak a_T(Ke,u)-\mathfrak a_T(u)]$
(maximum at $u=Ke$). For \eqref{eq:coer-sandwich}: every $F\in\PT\cap\mathcal D_T$ is $v+e$ (Lemma~\ref{lem:coer-decomp}(iv)),
$\|F\|^2=\|v\|^2+\|e\|^2$, and $v\mapsto v+Ke$ is a bijection of $\mathcal V$. Taking $e=0$ gives $\lamA\le c$; taking
$v=-Ke$ gives $\mathfrak a_T(F)=\mathfrak s_{T,k}(e)$ with $\|F\|^2\ge\|e\|^2$, so $\lamA\le\lambda_S$. Conversely, if
$\lambda_S>0$, then $\mathfrak a_T(v+e)\ge c\|v+Ke\|^2+\lambda_S\|e\|^2$, and
$\|v\|^2+\|e\|^2\le(\|v+Ke\|+\|K\|\|e\|)^2+\|e\|^2\le(1+\|K\|)^2(\|v+Ke\|^2+\|e\|^2)$, which gives the lower bound; if
$\lambda_S\le0$ the lower bound is implied by $\lamA\le\lambda_S$. (iii) follows from \eqref{eq:coer-sandwich}
and, for the nonstrict case, from $\mathfrak a_T(v+e)\ge\mathfrak s_{T,k}(e)$ together with $v=-Ke$.
\end{proof}

Finite reductions of the window problem exist in several sufficient (one-directional) forms, in Legendre
bases: \cite[Thms.~1.1, 1.4]{Zhu2026} and \cite{Groskin2026}. Theorem~\ref{thm:coer-schur} is an exact reduction
in an explicit basis of explicit dimension, with $k$ supplied by Corollary~\ref{cor:coer-L1}; it is not the first
finite reduction, and it does not improve any certified value.

\begin{remark}[what the reduction is and is not]\label{rem:coer-honest}
(1) Theorem~\ref{thm:coer-schur} is the Schur complement of $\mathfrak a_T$ with respect to $\PT=\mathcal V_k\oplus E_k$,
in the form-domain setting required because $\mathfrak a_T$ is unbounded on $L^2$. Its content for the
positivity question is a \emph{localization}: the sign of $\lamA$ is the sign of an explicit real symmetric matrix of
size $2k-2$, with $k=25$ at $T=\frac12\log5$, whose index set is the piecewise Tate exponentials
$e^{\pm t/2}\mathbf 1_{B_j}$. The Galerkin near-minimizers of $\mathfrak a_T$ on $\PT$ have $99\%$ or more of their mass in
$E_k$ for $k\ge5$ (numerical); this is approximation theory (piecewise Tate exponentials approximate
smooth functions to $O(\tau^2)$), not arithmetic structure.
(2) \emph{No gain in precision.} By \eqref{eq:coer-sandwich}, $\lamA\le\lambda_S\le(1+\|K\|)^2\lamA$ whenever
$\lambda_S\le c$; so $S_{T,k}$ inherits the decay of Theorem~\ref{thm:margin} and the smallness already present at
the declared range: at $T=\frac12\log5$, $\lambda_S\approx\lamA\approx3\cdot10^{-12}$ with entries of $S_{T,k}$ of size
$O(1)$ (numerical, Galerkin; the certified value is Theorem~\ref{thm:log5A}). The entries of $S_{T,k}$ are not in
closed form: by \eqref{eq:coer-schur} the correction $\mathfrak a_T(Ke)$ is a supremum over $\mathcal V$, so Galerkin in
$\mathcal V$ gives lower bounds for the correction, i.e.\ upper bounds for $S_{T,k}$; a certified lower bound for $S_{T,k}$
needs a near-maximizer $u_N$ together with the dual estimate $\|\mathfrak a_T(e,\cdot)-\mathfrak a_T(u_N,\cdot)\|^2_{\mathcal V^*}/c$,
which is legitimate precisely because $c>0$ is proved. A certification by interval arithmetic along these lines is
realistic for $T\le0.6$, where $\lamA\ge10^{-5}$ (numerical), and is not attempted here.
(3) The decomposition with a \emph{fixed} number of blocks (the orbit blocks of the first contact, $k=3$, $\dim E=4$)
is not stable in $T$: the Galerkin value of $c(T,3)$ falls to $0.013$ at $T=0.69$ (numerical); only
$k\asymp T/\tau_0$ equal blocks give a uniform $c$ on a bounded range, by Corollary~\ref{cor:coer-L1}.
\end{remark}

\begin{remark}[the route by operator norms is not viable]\label{rem:coer-nogo}
By Lemmas~\ref{lem:coer-sector} and \ref{lem:coer-K1}, for $F\in\mathcal V_k$,
$\Gamma(F)=\sum_j\Gamma(F_j)+\mathfrak q(F)$ with $\sum_j\Gamma(F_j)\ge\lambda_{\mathrm A}(\tau)\|F\|^2$ and the coupling form
$\mathfrak q(F)=-\sum_{i<j}\iint F_i(x)F_j(y)K_1(|x-y|)\,dxdy$, so that
$c(T,k)\ge\lambda_{\mathrm A}(\tau)-\sup_{\mathcal V_k}|\mathfrak q|/\|\cdot\|^2-\|\PiT\|$. This splitting is sharp in neither
term and cannot replace Theorem~\ref{thm:coer-B}:
\begin{itemize}
\item[(a)] For two adjacent blocks the bilinear form $(F_1,F_2)\mapsto\iint F_1F_2K_1(|x-y|)$ on Tate-orthogonal
functions has norm exactly $\pi$: $\le\pi$ by $0<K_1(u)\le u^{-1}$ and Hilbert's inequality
$\iint\frac{f(x)g(y)}{x+y}\le\pi\|f\|\|g\|$ on $(0,\infty)$ \cite[Thm.~316]{HLP}; $\ge\pi$ by the Weyl sequence
$f_M(x')=x'^{-1/2}\mathbf 1_{[2\tau e^{-M},2\tau]}$ in the distance $x'$ to the common endpoint, placed on both blocks and
projected onto the orthogonal complement of $\{e^{\pm x/2}\mathbf 1_B\}$: in logarithmic variables
$\iint f_Mf_M/(x'+y')=\int_{[-M,0]^2}\frac{du\,dv}{2\cosh\frac{u-v}2}=\pi M-O(1)$, $\|f_M\|^2=M$, the bounded part
$K_1-u^{-1}=-\frac32+O(u)$ contributes $O(\|f_M\|_1^2)=O(\tau)$, and the Tate projection changes $f_M$ by a vector of norm
$O(1)$ uniformly in $M$ (moments $O(\sqrt\tau)$, Gram matrix $\asymp\tau$), which changes the pairing by $O(\sqrt M)$.
Hence, with the normalization of \eqref{eq:coer-cross}, $\sup_{\mathcal V_2}|\mathfrak q|/\|\cdot\|^2=\frac\pi2$ exactly, and
$\sup_{\mathcal V_k}|\mathfrak q|/\|\cdot\|^2\ge\frac\pi2$ for every $k\ge2$ and every $\tau$; the adjacent part of
$\mathfrak q$ is $\le\pi\cos\frac\pi{k+1}\|F\|^2$ by the path bound of Proposition~\ref{prop:coer-prime}, and the
non-adjacent part is not estimated here. (If the cross term is counted once for each \emph{ordered} pair, as in some
working notes, the quadratic coupling constant reads $\pi$ instead of $\frac\pi2$; the conclusion below is the same.)
\item[(b)] The value $\frac\pi2$ is a property of the essential spectrum of the Carleman operator $1/(x+y)$ (symbol
$\pi/\cosh\pi\sigma$ under the Mellin transform, \cite{Carleman1923}): no condition of finite codimension, such as the
Tate conditions, lowers it. Galerkin values of this norm ($0.74$ at degree $10$) are what polynomials of low degree see
and grow with the degree without saturating (numerical).
\item[(c)] With \eqref{eq:coer-A}, the split bound is positive at $T=\frac12\log5$ only for
$\log\frac1\tau>1.21+\tau+\frac\pi2+1.674$, i.e.\ $k\ge70$ ($\dim E_k\ge138$), against $k=25$ from
Theorem~\ref{thm:coer-B}; and the true coupling cost at the Galerkin minimizers of $\Gamma|_{\mathcal V_k}$ is only
$0.17$--$0.18$ (numerical): the extremals of the coupling norm are not the minimizers of the form. A
relative bound $\mathfrak q\le\theta\sum_j\Gamma(F_j)$ with $\theta<1$ is observed numerically ($\theta\approx0.22$) but is
not proved and is not needed: Theorem~\ref{thm:coer-B} is the relative bound with $\theta=1$ and an explicit constant.
\end{itemize}
\end{remark}

\begin{remark}[labels]\label{rem:coer-labels}
Lemmas~\ref{lem:coer-decomp}--\ref{lem:coer-poinc}, Theorems~\ref{thm:coer-A}, \ref{thm:coer-B}, \ref{thm:coer-schur},
Proposition~\ref{prop:coer-prime} and Corollary~\ref{cor:coer-L1} are proved, with the numerical inputs
$w(2)=2.1045$, the closed-form values of $C_B$, and the arithmetic of the table (double precision evaluations of
elementary expressions). The following are numerical and not proved: $\lambda_{\mathrm A}(\tau)=\log\frac1\tau-0.510$;
the uniform bound $C_B\le1.74$ on $(0,0.1]$; the Galerkin values of $c(T,k)$, $\lambda_S$, $\|K\|$, of the mass of the
near-minimizer in $E_k$, of the coupling cost at the minimizers and of the relative constant $\theta\approx0.22$.
The positivity of $S_{T,k}$ in the declared range beyond what
Theorem~\ref{thm:log5A} already certifies for $\lamA$ is open; the route of Remark~\ref{rem:coer-nogo} is closed.
\end{remark}


\section{Concluding remarks}\label{sec:concluding}

The results of this paper move the question of positivity from the primal side, where it is a spectral
computation, to the dual side, where at each window it is a finite-dimensional positivity problem with a
known meaning: a certificate is the explicit formula, its value is a lower bound on a flat-top smoothing
of the zeros, and its existence at every window is equivalent to the Riemann Hypothesis. They also
indicate what a proof of positivity for every $T$ would have to supply. No finite $\nu$ attains the dual
value (Theorem~\ref{thm:O}), and the formal limit of the certificates is the zero measure itself
(Remark~\ref{rem:RHcert}), so such a proof must be an approximation argument with an explicit rate
depending on the class of weights; by Theorem~\ref{thm:Vp} the exchange between escaped primal mass and
value is $\Theta(\log R)$, so the rate cannot be uniform over the feasible set. Since $\Aconst$ is the exact
supremum of the prime comb (Proposition~\ref{prop:env}(ii)), the interval $[0,2\pi e^{\Aconst}]$ cannot be
shortened by bounding the primes one at a time. Nor can one argue upward in $T$: $\lamP$ is
nonincreasing in $T$ (Corollary~\ref{cor:mono}, Lemma~\ref{lem:monoT}), new contacts enter as $T$ grows, and
by Corollary~\ref{cor:margin1} every certificate at a window $T\ge1$ has value at most
$4\cdot10^9Te^{-15.449T}$, so no $T$-uniform margin exists (Corollary~\ref{cor:margin2}(i)). Finally, by
Theorem~\ref{thm:Z} and Remark~\ref{rem:Cheb}, a prime-side proof of $f_w\ge0$ is a statement about the zeros
in the range $[0,\tau_{\mathrm{env}}]$, and by Corollary~\ref{cor:void} it cannot rest only on properties of the
primes that are preserved by finite modifications of the weights.

Heuristically (\S\ref{subsec:where}), everything above the envelope scale is handled by the tail lemmas,
everything below the resolution scale $2\pi e^{2T}$ is a finite computation about the first zeros, and the
interval between the two, where the prime comb is nearly aligned, is where a proof must say something
new about the primes. We formulate this as a problem.

\begin{problem}[arithmetic certificate]\label{prob:B}
Exhibit, for every $T$, a weight $w_T\in\Wcl$ for which the elementary inequality
$f_{w_T}\ge0$ on $[0,\tau_{\mathrm{env}}(T)]$ can be proved from the distribution of the primes
below $e^{2T}$ and the Diophantine structure of the phases $\tau\log p$, without reference to
the zeros.
\end{problem}
By Corollary~\ref{cor:RH}, a solution of Problem~\ref{prob:B} proves RH. Natural candidates for $w_T$ are the
canonical scale $s=\frac12$ of Proposition~\ref{prop:M} and weights adapted to the primes outside the
window (Remark~\ref{rem:RHcert}); the data of \S\ref{sec:cert} neither support nor exclude them.


\section{A sum--product identity and the first failure of positivity}\label{sec:firstcontact}

This section studies a hypothetical first window at which positivity of the primitive form fails.
We prove five things: a finite operator identity in which additive Poisson summation and the
convolution identity $\Lambda*1=\log$ occur together (Theorem~\ref{thm:PCfinite}); a faithful
quadratic observation of a test function by the positive odd poles of the Gamma factor
\eqref{eq:Gammaenergy}; exact identities for the exterior collars produced by compressing the
prime shifts to the window \eqref{eq:collar}; an exact balance identity satisfied by any vector in
the kernel at the first failure \eqref{eq:Qbalance}; and a strict contraction of the Euler channel
(Proposition~\ref{prop:GEgap}).  The contraction does not advance the logical closure: on the
radical, its proposed reverse inequality is equivalent to vanishing of the detector.  What is
missing is the direct arithmetic annihilation statement in Problem~\ref{prob:annihilate}, together
with the first-contact bridge discussed in Remark~\ref{rem:completecontact}.  Nothing in this
section extends the unconditional range stated in the Introduction.

\subsection{The first-contact reduction}

Throughout, $\mathfrak a_T(F,G)$ denotes the sesquilinear form of \eqref{eq:form} on its form
domain, $J F(t)=F(-t)$, and $\Pi$ is the orthogonal projection of $L^2(\IT)$ onto $\PT$. Since the
symbol $\sT$ is even, $\mathfrak a_T$ is $J$-invariant, and since $A(JF)=B(F)$, $J$ preserves $\PT$
and commutes with $\Pi$. The orthogonal complement of $\PT$ in $L^2(\IT)$ is the two-dimensional
\emph{Tate space} $\mathrm{span}\{e^{t/2},e^{-t/2}\}$ (restricted to $\IT$).

The form $\mathfrak a_T$ restricted to $\PT$ is closed and bounded below (its form norm is
equivalent to $\|F\|^2+\frac1{2\pi}\int\log(2+|\tau|)|\Fh|^2d\tau$, by Lemma~\ref{lem:mbounds}, as in
the proof of Lemma~\ref{lem:contT}). Let $H_T$ be the associated self-adjoint operator on $\PT$, so
that $\mathfrak a_T(F,G)=\langle H_TF,G\rangle$ for $F$ in the domain of $H_T$ and $G\in\PT$ in the
form domain, and $\lamA=\min\operatorname{spec}H_T$. The sublevel sets of $\mathfrak a_T$ on the
unit sphere are compact (proof of Theorem~\ref{thm:O}; on $\PT$ this is
\cite[Thm.~3]{Bombieri2000}), so $H_T$ has compact resolvent: its spectrum is discrete, its
eigenspaces are finite-dimensional, and $\lamA$ is attained. Because $H_T$ commutes with $J$, each
eigenspace splits into even and odd parts.

\begin{lemma}[continuity of the primitive bottom]\label{lem:contA}
The function $T\mapsto\lamA$ is continuous on $(0,\infty)$.
\end{lemma}

\begin{proof}
Fix $T$ and consider windows $T'\in[T/2,2T]$ and factors $r\in[\frac12,2]$ with $rT'\in[T/2,2T]$; all
constants depend only on $T$. Put $M_0=\mathsf A_{2T}-C_0$, so that $\mathfrak a_S+M_0\|\cdot\|^2\ge0$
for every $S\le2T$ (Lemma~\ref{lem:mbounds}).

\emph{Step 1.} Steps (i) and (ii) of the proof of Lemma~\ref{lem:contT} do not involve the polar term.
Taking $\delta=8T|1/r-1|$ and summing the comb over $n<e^{8T}$, which covers every window in
$[T/2,2T]$ and every contact that enters or leaves, they give $\eta(r)\to0$ as $r\to1$ with
$|\mathfrak a_{rT'}(U_rF)-\mathfrak a_{T'}(F)|\le\eta(r)\big(\mathfrak a_{T'}(F)+C\|F\|^2\big)$ for all
$F$ in the form domain of $\mathfrak a_{T'}$, uniformly in $T'$.

\emph{Step 2.} If $F\in\mathcal P_{T'}$, then $A(U_rF)=r^{1/2}\int F(s)\big(e^{rs/2}-e^{s/2}\big)ds$,
so $|A(U_rF)|\le\varepsilon(r)\|F\|$ with $\varepsilon(r)=r^{1/2}\|e^{rs/2}-e^{s/2}\|_{L^2(I_{2T})}\to0$,
and likewise for $B$.

\emph{Step 3.} Fix real $\phi_1,\phi_2\in C_c^\infty(I_{T/4})$ for which the matrix
$\big(A(\phi_j),B(\phi_j)\big)_{j}$ is invertible. By zero extension (proof of Lemma~\ref{lem:monoT}),
$\mathfrak a_S(\phi_i,\phi_j)$ does not depend on $S\ge T/4$. For $u=U_rF$ with $F\in\mathcal P_{T'}$,
$\|F\|=1$, let $G=c_1\phi_1+c_2\phi_2$ have the same Tate moments as $u$. Then
$v:=u-G\in\mathcal P_{rT'}$, $\|G\|\le C_2\varepsilon(r)$ and $\mathfrak a_{rT'}(G)\le C_3\varepsilon(r)^2$.

\emph{Step 4.} By the Cauchy--Schwarz inequality for the nonnegative form
$\mathfrak a_{rT'}+M_0\|\cdot\|^2$,
$|\mathfrak a_{rT'}(u,G)|\le(\mathfrak a_{rT'}(u)+M_0)^{1/2}(\mathfrak a_{rT'}(G)+M_0\|G\|^2)^{1/2}+M_0\|G\|$.
Hence, if $\mathfrak a_{T'}(F)\le V$ for a fixed $V$, then
$\mathfrak a_{rT'}(v)\le\mathfrak a_{T'}(F)+\omega(r)$ with $\omega(r)\to0$ uniformly in such $F$. Since
also $\mathfrak a_{rT'}(v)\ge-M_0\|v\|^2$ and $\big|\|v\|-1\big|\le C_2\varepsilon(r)$, the Rayleigh
quotient gives $\lambda_{\min}(A_{rT'}|\mathcal P_{rT'})\le\mathfrak a_{T'}(F)+\omega'(r)$ with
$\omega'(r)\to0$.

\emph{Step 5.} By Lemma~\ref{lem:monoT}, $\lamA$ is bounded on $[T/2,2T]$ between $C_0-\mathsf A_{2T}$
and $\lambda_{\min}(A_{T/2}|\mathcal P_{T/2})$; take $V$ one more than the latter. Applying Step 4 to
near-minimizers at $T$ gives $\limsup_{r\to1}\lambda_{\min}(A_{rT}|\mathcal P_{rT})\le\lamA$, and applying
it to near-minimizers at $rT$ with the factor $1/r$ gives the reverse inequality for the $\liminf$.
\end{proof}

\begin{lemma}[first-contact reduction]\label{lem:firstcontact}
If $\lamA<0$ for some $T$, then there
is a window $T_*\ge\log2$ and a nonzero $F\in\mathcal P_{T_*}$ such that
\[
 H_{T_*}\ge0,\qquad F\in\mathcal N_{T_*}:=\ker H_{T_*},\qquad JF=\varepsilon F,\quad \varepsilon\in\{\pm1\}.
\]
Conversely, if for no $T$ does $\ker H_T\ne\{0\}$ occur together with $H_T\ge0$ and $H_{T'}\not\ge0$
for all $T'>T$, then $\lamA\ge0$ for every $T$, and hence RH holds.
\end{lemma}

\begin{proof}
By \cite{Paper42}, $\lamA\ge0$ for $T\le\log2$. Let $T_*=\inf\{T:\lamA<0\}\ge\log2$. By
Lemma~\ref{lem:monoT}, $\lamA<0$ for every $T>T_*$ and $\lamA\ge0$ for every $T<T_*$; by
Lemma~\ref{lem:contA}, $\lambda_{\min}(A_{T_*}|\mathcal P_{T_*})=0$. By compactness of the resolvent this
value is an eigenvalue of $H_{T_*}$, and since $\ker H_{T_*}$ is $J$-invariant it contains a nonzero
vector of definite parity. For the converse, if no such window exists then the set
$\{T:\lamA<0\}$ is empty by the first part, so $\mathfrak a_T\ge0$ on $\PT$ for every $T$, which is
equivalent to RH \cite[Thm.~11.3]{Paper45}.
\end{proof}

\begin{remark}[the complete-form alternative]\label{rem:completecontact}
Lemma~\ref{lem:contT} does give an unconditional first-contact reduction for the complete bottom
$\lamW$: if it ever becomes negative, its first zero is attained by a nonzero vector in the kernel
of the nonnegative complete form.  The complete null vector need not be primitive, and its primitive component is governed by the
Schur complement including the polar block, not by $H_T$.  The collar balance
\eqref{eq:Qbalance} and the detector argument below have been proved for $\ker H_T$, not for that
complete Schur-complement kernel.  With Lemma~\ref{lem:contA} the primitive reduction is unconditional, and the complete-form
route is not needed below.
\end{remark}

We call a nonzero $F\in\ker H_T$ with $JF=\varepsilon F$, at a window where $H_T\ge0$, a
\emph{first-contact vector}, and $\ker H_T$ the \emph{radical}. Equivalently, $F\in\PT$ satisfies
$\mathfrak a_T(F,\Pi G)=0$ for every $G$ in the form domain, i.e.\ the weak equation
\begin{equation}\label{eq:localweak}
 \mathfrak a_T(F,G)=\langle r,G\rangle\qquad(G\ \text{in the form domain of}\ L^2(\IT))
\end{equation}
holds for a unique $r$ in the Tate space. More generally, for $F$ in the domain of $H_T$ we define
$r\in(\PT)^\perp$ by $\mathfrak a_T(F,G)=\langle H_TF+r,G\rangle$.

\subsection{Operators on the line}

Let $\tau_aF(t)=F(t+a)$ and $XF(t)=tF(t)$ on the Schwartz space $\mathcal S(\R)$ (the \emph{Schwartz
core}). Let $\mathcal M=m(D)$ be the Fourier multiplier with the symbol $m$ of \eqref{eq:symbol},
$(\mathcal MF)\,\widehat{}\,=m\Fh$; we call it the completed Gamma multiplier. Let
$u_\infty(\tau)=\pi^{-i\tau}\Gamma(\frac14+\frac{i\tau}2)/\Gamma(\frac14-\frac{i\tau}2)$, a
unimodular function with $u_\infty(-\tau)=\overline{u_\infty(\tau)}$ and, by Tate's local
functional equation, $\frac{d}{d\tau}\arg u_\infty=m$. Put
\[
 \mathscr F=u_\infty(D)\,J .
\]
This is the additive Fourier transform in logarithmic coordinates: if $F(t)=e^{t/2}f(e^t)$ with $f$
an even Schwartz function on $\R$, then $\mathscr FF(t)=e^{t/2}\widehat f(e^t)$. On the Schwartz
core
\[
 \mathscr F^*=\mathscr F^{-1}=\mathscr F,\qquad
 J\tau_aJ=\tau_{-a},\qquad X+\mathscr FX\mathscr F=\mathcal M .
\]
The first identity follows from $Ju_\infty(D)J=u_\infty(D)^*$. For the third, if $U=e^{i\theta}(D)$
with $\theta$ real and of polynomial growth with its first two derivatives, then
$U^*XU=X+\theta'(D)$ (differentiate $\widehat{XG}=-i\partial_\tau\widehat G$, or see the time-delay
identity $\langle UG,XUG\rangle=\langle G,XG\rangle+\frac1{2\pi}\int\theta'|\widehat G|^2$); with
$JX=-XJ$ this gives $\mathscr FX\mathscr F=-UXU^*=-X+m(D)$.

Two shift conventions occur. In the full-line identity we use $\tau_a$; in the compressed
moment calculations of \S\ref{subsec:collars} we use the right shift $V_bF(t)=F(t-b)$.

\subsection{The finite sum--product identity}

For integers $M,N\ge1$ define on the Schwartz core
\begin{align*}
 S_N&=\sum_{n\le N}n^{-1/2}\tau_{\log n},&
 K_N&=\sum_{n\le N}(\log n)n^{-1/2}\tau_{\log n},\\
 L_M&=\sum_{d\le M}\Lambda(d)d^{-1/2}\tau_{\log d},&
 A_M&=\mathcal M-(L_M+L_M^*),\\
 E_N&=S_N\mathscr F-JS_N,&
 H_{M,N}&=L_MS_N-K_N,
\end{align*}
and $\{X,E\}=XE+EX$. For the full integer comb, Poisson summation formally reads
$S_\infty\mathscr F=JS_\infty$; thus $E_N$ is the finite Poisson defect, and $H_{M,N}$ is the
incomplete divisor hyperbola.

\begin{theorem}[finite Poisson--Chebyshev identity]\label{thm:PCfinite}
On the Schwartz core,
\begin{equation}\label{eq:PCfinite}
 S_NA_M\mathscr F
 =\{X,E_N\}-L_M^*E_N-JH_{M,N}-H_{M,N}\mathscr F.
\end{equation}
If $M\ge N$, then
\[
 H_{M,N}=\sum_{k>N}h_{M,N}(k)k^{-1/2}\tau_{\log k},\qquad
 h_{M,N}(k)=\sum_{\substack{dm=k\\d\le M,\ m\le N}}\Lambda(d)\ge0.
\]
In particular $H_{M,N}$ has no coefficient at $k\le N$.
\end{theorem}
\begin{proof}
The commutator $[X,S_N]=-K_N$ and the three Fourier identities above give
\[
 S_N\mathcal M\mathscr F=JK_N+K_N\mathscr F+\{X,E_N\}.
\]
Insert $K_N=L_MS_N-H_{M,N}$, use $JL_M=L_M^*J$ and
$JS_N=S_N\mathscr F-E_N$, and subtract $S_N(L_M+L_M^*)\mathscr F$.
This proves \eqref{eq:PCfinite}. The coefficient of
$k^{-1/2}\tau_{\log k}$ in $L_MS_N$ is the displayed divisor sum. For
$k\le N\le M$ every divisor occurs, so $\sum_{d\mid k}\Lambda(d)=\log k$ cancels the
coefficient of $K_N$.
\end{proof}

The identity is specific to $\Lambda$: replacing $\Lambda$ by other coefficients $a(d)$ leaves, at
$k\le N$, the coefficient $\sum_{d\mid k}a(d)-\log k$. On the other hand the square of either side
of \eqref{eq:PCfinite} is invariant under $A_M\mapsto-A_M$, so squaring it cannot determine the sign
of the Weil form.

\begin{remark}[renormalization]\label{rem:Sren}
The operators $S_N$ have no limit as $N\to\infty$ on the critical half-density. The renormalized
operators are
\[
 S_N^{\rm ren}=S_N-C_N,\qquad
 C_Nh(t)=e^{-t/2}\int_{-\infty}^{t+\log N}e^{u/2}h(u)\,du ,
\]
defined on distributions modulo the homogeneous solution $e^{-t/2}$; $C_N$ is the continuum
counterpart of the integer comb. For smooth compactly supported primitive $h$ the discrepancy
$S_N^{\rm ren}h$ converges locally uniformly as $N\to\infty$, by the prime number theorem with
the classical unconditional error term; moving derivatives onto a compactly supported test
function extends the limit continuously from compactly supported primitive $L^2$ functions to
$\mathcal D'(\R)$. This weak limit does not control a Hilbert norm of the limit, the projection
$\PiT$ applied to it, or its Gamma residues: an exterior exponential tail can tend to zero in every
fixed Sobolev norm while retaining a fixed meromorphic residue. For this reason every quadratic
expression below is formed at finite $N$, before any limit.
\end{remark}

\subsection{A faithful observation by the Gamma poles}

Fix $\delta>0$ and put, for $m\ge1$,
\[
 a_m=2m+\tfrac12,\qquad
 b_m=\sqrt{a_m}\,e^{-a_m(T+\delta)},\qquad
 K_F(a)=\int_{-T}^T e^{-at}F(t)\,dt,\qquad
 (RF)_m=b_mK_F(a_m).
\]
Since $\|e^{-a_mt}\|^2_{L^2(\IT)}\le2Te^{2a_mT}$, the map $R:L^2(\IT)\to\ell^2$ is Hilbert--Schmidt,
and
\begin{equation}\label{eq:Gammaenergy}
 \|RF\|^2
 =\sum_{m\ge1}a_me^{-2a_m(T+\delta)}|K_F(a_m)|^2
 =\iint F(t)\overline{F(u)}\,\bigl[-(k_\Gamma^+)'(2T+2\delta+t+u)\bigr]\,dt\,du,
\end{equation}
where
\[
 k_\Gamma^+(v)=\sum_{m\ge1}e^{-a_mv}.
\]
This is a positive Hankel form, and it is faithful: $RF=0$ implies $F=0$. Indeed, after the
substitution $x=e^{-2t}$, vanishing of every $K_F(a_m)$ says that $x^{1/4}F(-\frac12\log x)\,x^{-1}$
is orthogonal to every power $x^m$, $m\ge1$, on the compact interval $[e^{-2T},e^{2T}]$. Since this
interval does not contain $0$, the functions $x\,p(x)$, $p$ a polynomial, are dense in its continuous
functions; hence $F=0$.

The numbers $K_F(a_m)$ are the residues at the positive odd poles of the Gamma factor. Write
$s=\frac12+z$, $\chi(s)=\pi^{s-1/2}\Gamma(\frac{1-s}2)/\Gamma(\frac s2)$, so that
$\zeta(s)=\chi(s)\zeta(1-s)$, and $d_\Gamma=-\chi'/\chi$. Near $s_m=2m+1$ one has
$\psi(\frac{1-s}2)=2/(s-s_m)+O(1)$, hence $\operatorname{Res}_{s=s_m}d_\Gamma=1$. For a
first-contact vector $F$ of parity $\varepsilon$ let $W^0_\varepsilon(s)$ be the corrected product
defect: the function $d_\Gamma(s)K_F(s-\frac12)$ minus the Laplace transforms of the two compactly
supported Euler histories of $F$ and minus two polar corrections supported outside $\IT$. We use only
two properties of it: its inverse transform restricts on $\IT$ to the Tate vector $r$ of
\eqref{eq:localweak}, and every term other than $d_\Gamma K_F$ is holomorphic at each $s_m$. Hence
\begin{equation}\label{eq:residue}
 \operatorname*{Res}_{s=s_m}\bigl[\zeta(s)W^0_\varepsilon(s)\bigr]
 =\zeta(s_m)K_F(a_m).
\end{equation}
Consequently holomorphy of $\zeta W^0_\varepsilon$ at every $s_m$, and vanishing of
\eqref{eq:Gammaenergy}, are each sufficient for $F=0$. Formula \eqref{eq:residue} places the
obstruction outside the window: corrections supported in a compact interval have entire Laplace
transforms and cannot cancel these residues.

\subsection{Compressed collars and the balance on the radical}\label{subsec:collars}

Let $P$ denote restriction to $\IT$ followed by extension by zero. For $N\ge2$ put
\[
 B_N=P\sum_{2\le n\le N}\Lambda(n)n^{-1/2}V_{\log n}P,\qquad
 d_{m,N}=\sum_{2\le n\le N}\frac{\Lambda(n)}{n^{a_m+1/2}},\qquad
 D_N=\operatorname{diag}(d_{m,N})_{m\ge1}.
\]
For $N\ge e^{2T}$, $B_N$ contains every shift that acts on $L^2(\IT)$, and
$\mathfrak a_T(F,G)=\mathfrak m_T(F,G)-\langle(B_N+B_N^*)F,G\rangle$, where
\[
 \mathfrak m_T(F,G)=\frac1{2\pi}\int_\R m(\tau)\,\Fh(\tau)\overline{\widehat G(\tau)}\,d\tau
\]
is the compressed Gamma form, i.e.\ the form of $\mathcal M$ on the form domain of $L^2(\IT)$.
Thus \eqref{eq:localweak} is the weak form of $(\mathcal M-B_N-B_N^*)F=r$ on $\IT$.
The exterior collar $\mathcal E_N:L^2(\IT)\to\ell^2$ is
\[
 (\mathcal E_NF)_m=b_m\sum_{2\le n\le N}\Lambda(n)n^{-(a_m+1/2)}
 \int_{\max(-T,T-\log n)}^T e^{-a_mu}F(u)\,du.
\]
The substitution $u=t-\log n$ gives the exact identity
\begin{equation}\label{eq:collar}
 RB_N=D_NR-\mathcal E_N.
\end{equation}
Primes with $n\ge e^{2T}$ do not contribute to $B_N$ but do contribute to $\mathcal E_N$: once
$N\ge e^{2T}$, one has $\mathcal E_{N'}-\mathcal E_N=(D_{N'}-D_N)R$ for $N'\ge N$, so these terms
cannot be discarded.

Let $N\ge e^{2T}$, let $R_L,D_{N,L},\mathcal E_{N,L}$ denote the truncations to $1\le m\le L$, and put
$G_L=R_L^*R_L$ and $S_L=G_L+JG_LJ$. The operator $S_L$ commutes with $J$ and its kernel is a finite
sum of exponentials on $\overline{\IT}$, so $\Pi S_LF$ lies in the form domain and is an admissible
test. Let $F$ be a first-contact vector with Tate vector $r$. Taking $G=S_LF$ in
\eqref{eq:localweak}, using $JB_NJ=B_N^*$, equal parity of $F$ and $S_LF$, and \eqref{eq:collar}, one
obtains
\begin{equation}\label{eq:Qbalance}
 0=2c_{N,L}(F)+Q_{N,L}(F),\qquad
 c_{N,L}(F)=\langle R_LF,D_{N,L}R_LF\rangle,
\end{equation}
where, with no further boundary term,
\begin{align}\label{eq:Qfull}
 Q_{N,L}(F)={}&2\Rel\langle R_LF,R_LB_N^*F\rangle
 -2\Rel\langle R_LF,\mathcal E_{N,L}F\rangle\notag\\
 &-\Rel\mathfrak m_T(F,S_LF)+\Rel\langle r,S_LF\rangle.
\end{align}
The backward term involves the truncated moments
\[
 (R_LB_N^*F)_m=b_m\sum_{2\le n<e^{2T}}
 \Lambda(n)n^{a_m-1/2}\int_{-T+\log n}^T e^{-a_mu}F(u)\,du;
\]
replacing them by full moments would remove a collar.

Identity \eqref{eq:Qbalance} carries no sign information. Indeed, for $F$ in the domain of $H_T$
with $r$ defined as above, the same computation gives
\begin{equation}\label{eq:Qrecombine}
 Q_{N,L}=-2c_{N,L}-\Rel\mathfrak a_T(F,\Pi S_LF),
\end{equation}
and for an eigenvector $H_TF=\lambda F$ of parity $\varepsilon$, since
$\langle F,S_LF\rangle=2\|R_LF\|^2$,
\begin{equation}\label{eq:Qspectral}
 Q_{N,L}(F)=-2\sum_{m\le L}(d_{m,N}+\lambda)|(R_LF)_m|^2.
\end{equation}
Thus $Q_{N,L}$ is negative on every eigenvector with $\lambda\ge0$ and $R_LF\ne0$. In particular, the
inequality $Q_{N,L}\ge0$ on the radical is equivalent to $R_LF=0$ there, and $Q_{N,L}$ admits no
representation as a sum of squares valid on all of $\PT$.

\subsection{The Gamma--Euler contraction and what it does not prove}

The Euler channel $D_N$ is a strict contraction, uniformly in $m$ and $N$.

\begin{proposition}[strict Gamma--Euler contraction]\label{prop:GEgap}
For all $m\ge1$ and $N\ge2$,
\[
 0<d_{m,N}\le q_0:=\frac{\log2}{4}+\frac1{16}<0.236<1.
\]
Consequently
\begin{equation}\label{eq:gapop}
 \Delta_{N,L}:=R_L^*(I-D_{N,L})R_L
 \ge(1-q_0)R_L^*R_L,
\end{equation}
and, keeping the collar,
\begin{equation}\label{eq:gapcollar}
 \langle F,\Delta_{N,L}F\rangle
 =\|R_LF\|^2-\Rel\langle R_LF,R_LB_NF\rangle
  -\Rel\langle R_LF,\mathcal E_{N,L}F\rangle.
\end{equation}
For a first-contact vector $F$,
\begin{equation}\label{eq:gapQ}
 2\langle F,\Delta_{N,L}F\rangle=2\|R_LF\|^2+Q_{N,L}(F).
\end{equation}
\end{proposition}
\begin{proof}
Since $a_m+\frac12=2m+1\ge3$, $\Lambda(n)\le\log n$, and $x^{-3}\log x$ is decreasing on $[2,\infty)$,
\[
 d_{m,N}\le\sum_{n\ge2}\frac{\log n}{n^3}
 \le\frac{\log2}{8}+\int_2^\infty\frac{\log x}{x^3}\,dx
 =\frac{\log2}{4}+\frac1{16}.
\]
Positivity holds because the term $n=2$ is present. Equation \eqref{eq:gapop} follows coordinatewise,
and \eqref{eq:gapcollar} follows by pairing \eqref{eq:collar} with $R_LF$. Finally, on the radical
\eqref{eq:Qbalance} gives $Q_{N,L}=-2\langle R_LF,D_{N,L}R_LF\rangle$.
\end{proof}

Letting $N\to\infty$, $d_m:=\lim_Nd_{m,N}=-\zeta'(s_m)/\zeta(s_m)$. Hence the residue vector
$x_m=\zeta(s_m)b_mK_F(a_m)$ and the Chebyshev jet $j_m=-\zeta'(s_m)b_mK_F(a_m)$ satisfy, with
$D=\operatorname{diag}(d_m)$,
\begin{equation}\label{eq:smallgain}
 j=D x,\qquad \|j\|\le q_0\|x\|.
\end{equation}
This is a quantitative identity, but it supplies no independent leverage toward vanishing.  For a
fixed $L$ injective on the radical and $F\in\ker H_T$, the following are equivalent:
\begin{equation}\label{eq:equivgate}
 R_LF=0
 \quad\Longleftrightarrow\quad
 \langle F,\Delta_{N,L}F\rangle\le0
 \quad\Longleftrightarrow\quad
 Q_{N,L}(F)\ge0.
\end{equation}
The first equivalence follows from \eqref{eq:gapop}; the second follows from
$Q_{N,L}=-2\langle R_LF,D_{N,L}R_LF\rangle$ and strict positivity of $D_{N,L}$.
Likewise, because $j=Dx$ with every diagonal entry of $D$ in $(0,q_0]$,
\[
 \|x\|^2\le\Rel\langle x,j\rangle\quad\Longleftrightarrow\quad x=0.
\]
Thus neither the finite return inequality nor its infinite-channel version is an intermediate
estimate: each is the desired conclusion in another notation.  Cauchy--Schwarz and $q_0<1$ merely
verify this equivalence.

\begin{problem}[Gamma-moment annihilation at first contact]\label{prob:annihilate}
Starting from the finite PC-ren identity, prove directly that for every primitive first-contact
vector $F$ there is a finite $L$, injective on the radical, for which
\begin{equation}\label{eq:annihilate}
 R_LF=0.
\end{equation}
The proof must use all three pieces absent from the formal radical equation: (i) minimality of the
contact window, not merely $H_TF=0$; (ii) an oriented causal or Poisson--Tate identity which changes
under reversal of the Weil sign; and (iii) $\Lambda*1=\log$ in the step which forces vanishing or
fixes the sign, not only in the preliminary contraction estimate.  It must retain both Fourier
orientations, the continuum polar term, the full exterior collars and the crossed Tate endpoint
before taking a limit.
\end{problem}

Condition (i) is a falsification test: by \eqref{eq:Qspectral}, reasoning which also applied to
positive eigenvalues would give the wrong conclusion.  Condition (ii) excludes norm-squaring of
\eqref{eq:PCfinite}, which is invariant under $A_M\mapsto-A_M$.  Condition (iii) retains the genuine
arithmetic advantage of this route: changing one von Mangoldt weight creates an interior divisor
defect, while Poisson and a functional equation alone also occur for Epstein/Davenport--Heilbronn
examples without an Euler product.

A finite $L$ with $R_L$ injective on $\ker H_T$ exists at each fixed window. The kernel is
finite-dimensional because $H_T$ has compact resolvent; the subspaces $\ker H_T\cap\ker R_L$ decrease
with $L$ and have zero intersection by faithfulness of $R$; a decreasing chain of subspaces of a
finite-dimensional space stabilizes, so one of them is $\{0\}$. The argument gives no bound on $L$
that is uniform in $T$.

\begin{remark}[remarks on moment annihilation]\label{rem:returnremarks}
The following shortcuts do not prove Problem~\ref{prob:annihilate}. The weak renormalized limit of
Remark~\ref{rem:Sren} cannot be squared, and it does not make the extraction of residues continuous.
The identity $Q_{N,L}=-2c_{N,L}$ on the radical cannot be combined with an inequality $Q_{N,L}\ge0$,
since that inequality is equivalent to the desired conclusion, and \eqref{eq:Qspectral} excludes a
proof of it valid on all of $\PT$. The variance of a sum labelled by divisor paths is not the norm of
the coherent sum over $k$, which contains Gram cross terms between different $k$. Additive phases
diagonalize the sum over $k$ by Parseval but do not factor through $de=k$, whereas multiplicative
characters factor through $de=k$ but replace the Gamma/Poisson operator by a twisted one, different
from the $H_T$ whose kernel is under study. Unitarity of $\mathscr F$ gives the contractive direction
only, and reversing it without a new oriented Green identity for the compressed polar and Tate terms
amounts to assuming \eqref{eq:annihilate}.
\end{remark}


\section{Open problems}\label{sec:repro}

\begin{enumerate}
\item Problem~\ref{prob:A}, first for one contact and then for two.
\item Whether the optimal multipliers of the primitive problem converge to $\beta=(2,0)$, that is,
whether $\lamP=\lamW$ at the windows considered (Remark~\ref{rem:Kprim}).
\item A certificate at a two-prime window $\frac12\log3<T\le\log2$. No certificate of positive value
is known there, and by Theorem~\ref{thm:Vp} the size of the class it requires cannot be predicted by
a rate depending on escaped mass alone.
\item A closed-form profile for the weights $a_j$ of Theorem~\ref{thm:C}.
\item Whether $\lamW>0$ unconditionally for some $T>\frac12\log5$, beyond the range of \cite{Zhu2026}
and of Theorem~\ref{thm:log5W}, in particular at a window where the contact $5$ has entered. By
Theorem~\ref{thm:KEP3}(a), a positive answer at $T$ gives the indeterminacy of $\mathrm{KEP}(T)$, and
indeterminacy gives $\lamW\ge0$; by Theorem~\ref{thm:KEP3}(b), the question for every $T$ is the
Riemann Hypothesis. By Theorem~\ref{thm:parityschur}, at a window with $\lamA>0$ the question is
whether $\mathfrak g_{\mathrm{ev}}(T)>0$ and $\mathfrak g_{\mathrm{odd}}(T)>0$, and by
Proposition~\ref{prop:onewindow} a positive answer at one window gives it at all smaller ones.
\item Whether the flat-top certificates of Theorem~\ref{thm:Zexh} can be chosen with $w_T$ depending
on $T$ only through the primes outside the window (Problem~\ref{prob:B}).
\item Positivity of the Schur matrix $S_{T,k}$ of Theorem~\ref{thm:coer-schur} at a window beyond
Theorem~\ref{thm:log5A}, and a certified lower bound for its entries (Remark~\ref{rem:coer-honest}(2)); the
semilocal split of Corollary~\ref{cor:sl-split} for $S=\{\infty,2,3\}$ and the absolute trace formula at
levels $\lambda\ne1$ (Remark~\ref{rem:sl-not}).
\item Gamma-moment annihilation at primitive first contact, Problem~\ref{prob:annihilate}.  The
return inequalities are equivalent to this conclusion by \eqref{eq:equivgate}, not intermediate
estimates.  Separately, close the first-contact bridge of Remark~\ref{rem:completecontact}: prove
continuity of $\lamA$, or transfer the construction to the continuous complete-form first contact.
\end{enumerate}

\emph{Data.} The certificate of Theorem~\ref{thm:C} (the $437$ atoms and weights as exact decimal
strings, the multipliers, the constants of the cover and the interval enclosures), the test kernels
of Theorems~\ref{thm:N} and \ref{thm:Np}, the coefficient vectors of Theorems~\ref{thm:w11} and
\ref{thm:nest}, the matrices of Propositions~\ref{prop:haynsworth} and \ref{prop:negplane}, and the
scripts that verify them accompany this paper. Each verification runs in at most a few minutes on one
core with Python~3, \texttt{mpmath}~1.3.0 and \texttt{scipy}~1.11.4. The certificates of
Theorems~\ref{thm:log5A} and \ref{thm:log5W} (the verification scripts for the primitive and the complete
form, their outputs with every certified and failed value of $\mu$, the run logs, and the check at
$T=\frac12\log3$) use \texttt{python-flint}~0.9.0 (Arb) and take about seven minutes per parity.



\appendix

\section{Enclosure methods}\label{app:encl}

This appendix describes the computer verifications used in Theorems~\ref{thm:C}, \ref{thm:N},
\ref{thm:Np}, \ref{thm:index}, \ref{thm:log5A} and \ref{thm:log5W}. All of them are carried out in
outward-rounded interval arithmetic (\texttt{mpmath}), except the last two, which use ball arithmetic
(\S\ref{subsec:encllog5}); the data (atoms, coefficient vectors, test kernels and interval endpoints)
accompany the paper.

\subsection{The interval enclosures of Theorems~\ref{thm:N} and \ref{thm:Np}}
\emph{The window $T=\frac12\log3$.} All three integrals are computed in interval arithmetic at $25$ digits
(\texttt{mpmath}), the only non-interval inputs being the rational $\tau_0$ and the decimal $q$, which
are data of the bound and not approximations. The archimedean multiplier is summed from
Theorem~\ref{thm:S}(i), $m(\tau)=C_0+\sum_{k\ge0}\tau^2/(a_k(\tau^2+4a_k^2))$, with
$C_0=\psi(\frac14)-\log\pi=-\gamma-\frac\pi2-3\log2-\log\pi$ by the Gauss digamma theorem, the first
$100$ terms in interval arithmetic and the tail bracketed by the convex midpoint rule: with $A=K-\frac14$
and $t(a)=\tau^2/(\tau^2a+4a^3)$, the function $t$ is positive, decreasing and convex on $(0,\infty)$
(because $2u'^2-u''u=2\tau^4+24\tau^2a^2+192a^4>0$ for $u=\tau^2a+4a^3$), the points $a_k=k+\frac14$
are the midpoints of unit cells tiling $[A,\infty)$, so
$\sum_{k\ge K}t(a_k)\in[\Lambda-E,\Lambda]$ with $\Lambda=\frac12\log(1+\tau^2/4A^2)=\int_A^\infty t$
and $E=\frac1{24}|t'(A)|=\frac1{24}\tau^2(\tau^2+12A^2)/(\tau^2A+4A^3)^2$, at most $2.4\cdot10^{-7}$
on the range used. The multipliers are evaluated from the closed form
$\int_0^{2T}e^{su}\cos\tau u\,du=[e^{2Ts}(s\cos2T\tau+\tau\sin2T\tau)-s]/(s^2+\tau^2)$ at
$s=\pm\frac12$ and its $s$-derivative, without quadrature. The integrals over $W$ are composite
Simpson with the classical remainder $Ph^4C_4/180$, the fourth derivatives being bounded by
$|m^{(j)}|\le2^{-j}j!\,\zeta(j+1,\frac14)$, $|(v_2\cos(\tau\log2))^{(j)}|\le v_2(\log2)^j$,
$|\gh0^{(j)}|\le2\cosh T(2T)^{j+1}/(j+1)$, $|\gh1^{(j)}|\le2\sinh T(2T)^{j+2}/(j+2)$ and the Bernstein
bound $|\varphi^{(j)}|\le(d(M-1))^j(\sum_k|q_k|)^2$, combined by the Leibniz rule; the resolutions are
chosen so that the remainder is $7\cdot10^{-5}$ for $a$, which needs only $10^{-4}$, and below
$10^{-8}$ for $b_0,b_1$, which are multiplied by $B$. The denominator $\int_W\varphi=P\|q\|^2$ is
exact. An independent evaluation by adaptive quadrature and a library digamma, outside the interval
arithmetic, lands inside all three enclosures. Steps $d$ with $2T/d\notin\N$ are not covered by
Theorem~\ref{thm:N}: each such $d$ needs its own enclosure with $M'=\lceil2T/d\rceil$.

\emph{The windows $T=0.68$ and $T=\log2$.} The archimedean multiplier, the constant $C_0$, the closed forms of
$\gh0,\gh1$ and the exact denominator $P\|q\|^2$ are as in Theorem~\ref{thm:N}; nonnegativity of
$\varphi$ is structural, since it is evaluated as a sum of two squares, so no property of the
particular $q$ is used. Two things change. First, Simpson's rule is replaced by composite
Gauss--Legendre of order $p$ on panels of width $H$, with the classical remainder
$|E_{\mathrm{panel}}|\le H^{2p+1}(p!)^4\big/\big[(2p+1)((2p)!)^3\big]\cdot\sup|f^{(2p)}|$ and the
derivative bounds of Theorem~\ref{thm:N} extended from order $4$ to order $2p$, the archimedean one
being $|m^{(j)}|\le2^{-j}j!\,\zeta(j+1,\frac14)\le2^{-j}j!\,(4^{j+1}+\zeta(j+1))$ for $j\ge1$ and
$|m|\le\max(|C_0|,m(\tau_{\max}))$ for $j=0$; the nodes are Newton-refined at $60$ digits, widened
before entering interval arithmetic, and guarded by the archived residuals
$\max_i|P_p(x_i)|$ and $|\sum_iw_i-2|$, both below $10^{-60}$ in every run; this guard ensures that the
computed rule is the one to which the remainder formula applies. Second, $q$ is optimised: rather than the bottom eigenvector of
$C_M(\tau_0,\beta^\star)$, which makes $a$ small and leaves $b$ where it falls, one minimises
$a+B(|b_0|+|b_1|)$ over the span of the three or four lowest eigenvectors, which drives
$|b_0|+|b_1|$ to the $10^{-16}$--$10^{-18}$ level while $a$ stays within a fraction of a percent of
the double-precision ceiling. Consequently the $\beta$-allowance $(|b_0|+|b_1|)B$ is below
$1.3\cdot10^{-11}$ in every row, four orders below the smallest $|a|$, and the binding error is the
quadrature on $a$, between $0\%$ and $6.7\%$ of $|a|$. The bound survives any $B\le4.7\cdot10^9$,
a factor $2.5\cdot10^5$ above the constants of Step~1, so only the qualitative content of
Theorem~\ref{thm:Bbeta} --- that some finite a priori box exists --- is used. At two values of $M$
the minimax over $\tau_0$ has two active period placements, and no single pair $(\tau_0,q)$ can then
drive $b$ to zero; a scan over the best period placements supplies one that closes, at the cost of
recovering only a few percent of the double-precision ceiling there. As in Theorem~\ref{thm:N},
steps $d$ with $2T/d\notin\N$ are not covered: each needs its own enclosure with
$M'=\lceil2T/d\rceil$.

\subsection{Verification on the bounded part}

\begin{proposition}[second-order cover]\label{prop:TAY}
Write $f=m+q+r$ with $q=-\sum_{n\in\NT}v_n\cos(\tau\log n)+\nuh$ ($\nu$ atomic) and
$r=\beta_0\gh0+\beta_1\gh1$, and put
$M_2=\sum_{n\in\NT}v_n\log^2n+2\sum_j|a_j|u_j^2+|\beta_0|H_0+|\beta_1|H_1$ with
$H_0=\int_{-2T}^{2T}u^2\cosh\frac u2\,du$, $H_1=\int_{-2T}^{2T}|u|^3\sinh\frac{|u|}2\,du$. Then
$\sup|q''+r''|\le M_2$, and on every cell $[x_0-\frac w2,x_0+\frac w2]\subset[0,\infty)$
\[
\min_{\text{cell}}f\ \ge\ m\!\big(x_0-\tfrac w2\big)+(q+r)(x_0)-|(q+r)'(x_0)|\,\tfrac w2-M_2\tfrac{w^2}8 .
\]
The same bound applies to each row $\nuh+W_k$ of Lemmas~\ref{lem:F1}, \ref{lem:hybrid} with
$m(x+kP)$ and $-|\beta_i|k_i(x+kP)$ evaluated at the left endpoint. A cell is certified when the
right side exceeds the target and is bisected otherwise; near a contact of $f$ the admissible width
is $\approx\sqrt{8c/M_2}$, so the number of evaluations grows like $c^{-1/2}$ rather than the
$c^{-1}$ of a Lipschitz cover.
\end{proposition}
\begin{proof}
Taylor's formula with integral remainder for $q+r\in C^2$, and monotonicity of $m$ and of
$-|\beta_i|k_i$.
\end{proof}

\subsection{Grids and the band-limited minimum}
The following lemma quantifies how far the minimum over a grid can be from the true infimum: on a
$\delta$-grid a band-limited function can hide a dip of size $\asymp U^2\delta^2$ between two adjacent
nodes.

\begin{lemma}[Bernstein grid bound]\label{lem:bern}
Let $f:\R\to\R$ be bounded, real, and band-limited to $[-U,U]$: a trigonometric sum
$f(\tau)=\sum_{\ell}c_\ell e^{i\omega_\ell\tau}$ with $|\omega_\ell|\le U$, or more generally a
bounded function whose distributional Fourier transform is supported in $[-U,U]$. Then for every
$\delta>0$ and every $x_0\in\R$,
\[
\inf_{\tau\in\R}f(\tau)\ \ge\ \inf_{k\in\Z}f(x_0+k\delta)\ -\ \tfrac18U^2\|f\|_\infty\delta^2
\ \ge\ \inf_{k\in\Z}f(x_0+k\delta)-\tfrac12U^2\|f\|_\infty\delta^2 .
\]
\end{lemma}
\begin{proof}
By Bernstein's inequality $\|f''\|_\infty\le U^2\|f\|_\infty=:M$. Put $g=f-\inf f\ge0$; $g$ is
bounded with $\|g''\|_\infty\le M$. For every $x$ and $t$, $0\le g(x+t)\le g(x)+g'(x)t+\frac M2t^2$,
and a nonnegative quadratic in $t$ has nonpositive discriminant, so $g'(x)^2\le2Mg(x)$. Fix
$\eta>0$ and $x$ with $g(x)\le\eta$; let $x_0+k\delta$ be a nearest grid point, so
$|x-(x_0+k\delta)|\le\delta/2$. Taylor's formula with integral remainder gives
\[
f(x_0+k\delta)\ \le\ f(x)+|f'(x)|\tfrac\delta2+\tfrac M2\Big(\tfrac\delta2\Big)^2
\ \le\ \inf f+\eta+\sqrt{2M\eta}\,\tfrac\delta2+\tfrac18M\delta^2 .
\]
Letting $\eta\downarrow0$ yields $\inf_kf(x_0+k\delta)\le\inf_\R f+\frac18M\delta^2$.
\end{proof}

For the certificate function
$f=\sT+\nuh+\beta\cdot\widehat g$ the archimedean term is not band-limited, and the verification of
Theorem~\ref{thm:C} therefore uses no grid at all: the exhaustive second-order cover of
Proposition~\ref{prop:TAY} is the same Taylor estimate with the \emph{exact} constant $M_2$ in
place of the Bernstein bound $U^2\|f\|_\infty$, and its $2\,378\,593$ evaluations leave $0$
unresolved intervals, so the infimum over $\R$ --- not over a finite set of nodes --- is what is
certified. Conversely, for the flat-top objects of \S\ref{sec:Z}, where $f_w$ \emph{is}
band-limited of type $U=\sup\supp w$ whenever $w$ is compactly supported, Lemma~\ref{lem:bern}
turns any grid evaluation into a two-sided statement: at $U=14$ and $\delta=10^{-2}$ the admissible dip between two nodes is
$\frac18U^2\|f\|_\infty\delta^2\approx2.5\cdot10^{-3}\|f\|_\infty$.

\subsection{The certificates of Theorem~\ref{thm:index}}\label{subsec:enclindex}
In this subsection $\mathrm{par}\in\{\mathrm{ev},\mathrm{odd}\}$, $\chi_{\mathrm{ev}}(t)=\cosh\frac t2$ and
$\chi_{\mathrm{odd}}(t)=\sinh\frac t2$ on $\IT$, and $L^2_{\mathrm{par}}(\IT)$ is the space of real functions
of that parity, so that a real $F\in L^2_{\mathrm{par}}(\IT)$ lies in $\PT$ if and only if
$\langle F,\chi_{\mathrm{par}}\rangle=0$ (\S\ref{subsec:arch}). The basis $e_j$, the operators $S_c$ and the
tails $\Lambda_k$ are those of Theorem~\ref{thm:P}. Every other letter is local to this subsection.

\begin{lemma}[bordered inertia]\label{lem:border}
Let $q$ be a real quadratic form on a real vector space $V$, $f$ a linear functional on
$V$ and $v_0\in V$ with $f(v_0)=1$ and $q(\cdot,v_0)$ finite. On $V\oplus\R$ put
$b(v,\xi)=q[v]+2\xi f(v)$. Then $n_-(b)=n_-(q|_{\ker f})+1$, $n_-$ being the largest dimension of a
subspace on which the form is negative definite.
\end{lemma}
\begin{proof}
Write $v=w+tv_0$ with $w\in\ker f$. Then $b=q[w]+2tq(w,v_0)+t^2q[v_0]+2\xi t$, and the invertible
substitution $\eta=\xi+q(w,v_0)$ gives $b=q[w]+(t^2q[v_0]+2t\eta)$: a $b$-orthogonal sum of
$q|_{\ker f}$ and a binary form of determinant $-1$, whose inertia is $(1,1)$.
\end{proof}

\begin{lemma}[multi-band minorant]\label{lem:multiband}
Let $0=\tau^{(0)}<\dots<\tau^{(N)}=\tau_N$, let $m^-$ be a lower bound for $m$ at the edges,
and let $M_a\ge\sup m''$ on the band $[\tau^{(a)},\tau^{(a+1)}]$, of width $\Delta_a$. Define the even
function $\underline m$ by $\underline m=C_0$ on $[0,\tau^{(1)}]$; on the band $a\ge1$, the chord through
$(\tau^{(a)},m^-(\tau^{(a)}))$ and $(\tau^{(a+1)},m^-(\tau^{(a+1)}))$ lowered by $\max(M_a,0)\Delta_a^2/8$;
and $\underline m=\underline m_\infty:=m^-(\tau_N)$ for $|\tau|\ge\tau_N$. Then $\underline m\le m$ on
$\R$, and $\Gamma_T\ge\underline\Gamma$ as forms, where
$\underline\Gamma(F)=\frac1{2\pi}\int\underline m|\Fh|^2$. If $m^-$ increases along the edges, then also
$\underline m\le\underline m_\infty$.
\end{lemma}
\begin{proof}
On $[0,\tau^{(1)}]$ and beyond $\tau_N$ use that $m$ is increasing with $m(0)=C_0$
(Lemma~\ref{lem:mbounds}(i)). On a band $a\ge1$ let $L_a$ be the chord through the exact values and
$h=m-L_a$: $h$ vanishes at the ends and $h''\le M_a$, so $h\ge0$ if $M_a\le0$, and otherwise
$h+\frac12M_a(\tau-\tau^{(a)})(\tau^{(a+1)}-\tau)$ is concave and vanishes at the ends, whence
$h\ge-M_a\Delta_a^2/8$; lowering the end values to $m^-$ lowers the chord. Every end value of a
lowered chord is at most $m^-$ at that edge, hence at most $m^-(\tau_N)$ when $m^-$ increases along the
edges, and the chords are affine. Finally $\Gamma_T-\underline\Gamma=\frac1{2\pi}\int(m-\underline m)|\Fh|^2\ge0$.
\end{proof}
In the verification, $m^-$ is the lower end of an interval enclosure of $C_0+\sum_ks/(a_k(a_k^2+s))$, $s=\tau^2/4$
(Theorem~\ref{thm:S}(i)), summed to $3000$ terms with the tail of the decreasing summand bracketed by the
integral test; its width is at most $7\cdot10^{-9}$, while consecutive edge values of $m$ differ by at
least $0.03$, so $m^-$ increases along the edges. The bound $M_a$ sums, over $k<3000$, the termwise
majorant $a_k(a_k^2-2s)/(2(a_k^2+s)^3)$ of the summand $a_k(a_k^2-3s)/(2(a_k^2+s)^3)$ of $m''$, which is
decreasing-then-increasing in $s$ and is therefore bounded by its values at the ends of an $s$-interval
containing the band, plus the positive tail $\sum_{k\ge K_0}(2a_k^3)^{-1}\le(4(K_0-\frac34)^2)^{-1}$,
$K_0=3000$. There are $28$ bands: band $0$ is $[0,0.15]$, twenty fixed edges reach $8$, and eight
geometric edges run from $9$ to $\tau_N=c_*/T$, $c_*=16$ (to six decimals); the largest offset is $0.05791$.

\begin{proposition}[the upper certificate]\label{prop:haynsworth}
Fix $\mathrm{par}$, integers $k_1<k_2$ ($22$ and $36$ in the verification) and $\underline m$ as in
Lemma~\ref{lem:multiband}, with $\underline m\le\underline m_\infty$, and put $c_*=T\tau_N$,
$\omega=\underline m_\infty-\min\underline m$ and $\mathcal M=Q_T(\underline m_\infty-\underline m)(D)Q_T$,
so that $\underline\Gamma(F)=\underline m_\infty\|F\|^2-\langle\mathcal MF,F\rangle$ and
$0\le\mathcal M\le\omega S_{c_*}$. Let $E$ be the span of the $e_j$, $j<k_1$, of parity $\mathrm{par}$,
$Y$ its orthogonal complement in $L^2_{\mathrm{par}}(\IT)$, $P_Y$ the orthogonal projection onto $Y$,
and $\mathfrak b(F,\xi)=\underline\Gamma(F)+2\xi\langle F,\chi_{\mathrm{par}}\rangle$ on
$L^2_{\mathrm{par}}(\IT)\oplus\R$.
\emph{(a)} On $Y$, $\mathfrak b\ge q_Y:=\underline m_\infty-\omega\Lambda_{k_1}(c_*)$.
\emph{(b)} Suppose $q_Y>0$, and let $\mathsf R$ be the symmetric matrix on $E\oplus\R$ with entries
\[
\mathsf R_{ij}=\underline\Gamma(e_i,e_j)-\frac{\langle P_Y\mathcal Me_i,P_Y\mathcal Me_j\rangle}{q_Y},\quad
\mathsf R_{i\xi}=\langle e_i,\chi_{\mathrm{par}}\rangle+\frac{\langle P_Y\mathcal Me_i,P_Y\chi_{\mathrm{par}}\rangle}{q_Y},\quad
\mathsf R_{\xi\xi}=-\frac{\|P_Y\chi_{\mathrm{par}}\|^2}{q_Y}.
\]
Then the Schur complement $\mathfrak s$ of $\mathfrak b|_Y$ in $\mathfrak b$ satisfies $\mathfrak s\ge\mathsf R$
and $n_-(\mathfrak b)=n_-(\mathfrak s)$.
\emph{(c)} If $\mathsf R$ is positive definite on a subspace of codimension $r+1$, then
$n_{\mathrm{par}}(T)\le r$.
\end{proposition}
\begin{proof}
(a) $\mathfrak b|_Y=\underline m_\infty-P_Y\mathcal MP_Y$, and
$\|P_Y\mathcal MP_Y\|\le\omega\|P_YS_{c_*}P_Y\|\le\omega\operatorname{tr}(P^\perp S_{c_*}P^\perp)
=\omega\sum_{j\ge k_1}\langle S_{c_*}e_j,e_j\rangle=\omega\Lambda_{k_1}(c_*)$, where $P^\perp$ removes every
$e_j$, $j<k_1$, of either parity: $Y$ lies in its range, the norm of a positive operator is at most its
trace, and the last equality is the trace identity of Theorem~\ref{thm:P}. (b) For $i<k_1$ and $y\in Y$,
$\mathfrak b(e_i,y)=\underline\Gamma(e_i,y)=-\langle P_Y\mathcal Me_i,y\rangle$, since $\langle e_i,y\rangle=0$,
and the border pairs with $y$ through $+\langle P_Y\chi_{\mathrm{par}},y\rangle$. The congruence
$(x,y)\mapsto(x,y+(\mathfrak b|_Y)^{-1}\mathfrak b_{YX}x)$ gives $\mathfrak b\cong\mathfrak s\oplus\mathfrak b|_Y$
with $\mathfrak b|_Y>0$, so $n_-(\mathfrak b)=n_-(\mathfrak s)$; and
$\mathfrak b_{XY}(\mathfrak b|_Y)^{-1}\mathfrak b_{YX}$ is at most $q_Y^{-1}$ times the Gram matrix of the
$Y$-columns $-P_Y\mathcal Me_i$ and $+P_Y\chi_{\mathrm{par}}$, which is $\mathfrak s\ge\mathsf R$.
(c) Since $\Gamma_T\ge\underline\Gamma$ and $\PT^{\mathrm{par}}=\ker\langle\cdot,\chi_{\mathrm{par}}\rangle$,
Lemma~\ref{lem:border} gives
$n_{\mathrm{par}}(T)\le n_-(\underline\Gamma|_{\PT^{\mathrm{par}}})=n_-(\mathfrak b)-1=n_-(\mathfrak s)-1\le n_-(\mathsf R)-1\le r$.
\end{proof}
\emph{The computation.} For $i-l$ even, $\widehat{e_i}\overline{\widehat{e_l}}=2T\sqrt{(2i+1)(2l+1)}\,
(-1)^{(l-i)/2}j_ij_l(\tau T)$ (proof of Theorem~\ref{thm:P}), so with $\underline m=\alpha_a+\beta_a|\tau|$
on the band $a$ and $I_p^{(il)}(c)=\int_0^cx^pj_ij_l\,dx$,
\[
\underline\Gamma(e_i,e_l)=\underline m_\infty\delta_{il}+\frac2\pi\sqrt{(2i+1)(2l+1)}\,(-1)^{\frac{l-i}2}
\Big[\sum_a\big(\alpha_a\,\Delta I^{(il)}_0+\tfrac{\beta_a}T\,\Delta I^{(il)}_1\big)-\underline m_\infty I^{(il)}_0(c_*)\Big],
\]
$\Delta$ denoting the increment over the band, and the entries with $i-l$ odd vanish. The moments
$I_p$ come from the power series of $j_ij_l$ with exact rational coefficients, the alternating tail being
bounded by its first omitted term once the index exceeds $2c$; $\langle e_l,\chi_{\mathrm{par}}\rangle$ comes
from the positive series of the modified spherical Bessel function $i_l(T/2)$ (valid for $T<1$), and
$\|\chi_{\mathrm{par}}\|^2=\sinh T\pm T$. The inner products over $Y$ are exact sums over $k_1\le l<k_2$ of
the given parity ($\mathcal M$ preserves parity) plus boxes from
$\|P^\perp_{k_2}\mathcal Me_i\|^2\le\rho\,\mathcal M_{ii}$, $\rho=\omega\Lambda_{k_2}(c_*)$, and
Cauchy--Schwarz. In the verification $\Lambda_{22}(16)\le7.2151\cdot10^{-6}$ and $\rho\le4.02\cdot10^{-20}$, so
$q_Y\ge\underline m_\infty-5\cdot10^{-5}>0$. For (c), the eigenvectors of the midpoint of $\mathsf R$,
read as decimal data, give a matrix $U_r$ omitting the $r+1$ lowest; interval Cholesky of
$U_r^{\mathsf T}\mathsf RU_r-\mu I$ over the whole interval box of $\mathsf R$ succeeds for a certified
$\mu>0$, which forces $U_r$ to be injective and $\mathsf R$ to be positive definite on its range.
Everything is interval arithmetic at $40$ digits, with no quadrature; the entry widths of $\mathsf R$ are
at most $1.7\cdot10^{-13}$. The first eigenvalue of $\mathsf R$ is always the hyperbolic direction of
the border (Lemma~\ref{lem:border}); $r$ counts the others. The verified cases are:
\begin{center}\small
\begin{tabular}{@{}lcccc@{}}
\toprule
$T$ & sector & $\underline m_\infty$ & certified & margin $\mu$\\
\midrule
$0.58$ & even & $1.479384$ & $n_{\mathrm{ev}}\le0$ & $0.015998$\\
$0.58$ & odd & $1.479384$ & $n_{\mathrm{odd}}\le0$ & $0.047516$\\
$\log2$ & odd & $1.301146$ & $n_{\mathrm{odd}}\le0$ & $0.041093$\\
$0.90$ & even & $1.039940$ & $n_{\mathrm{ev}}\le1$ & $0.185541$\\
$0.90$ & odd & $1.039940$ & $n_{\mathrm{odd}}\le1$ & $0.144840$\\
\bottomrule
\end{tabular}
\end{center}
The cases $\log2$ (even, $\le1$, $\mu=0.420225$) and $0.80$ (even and odd, $\le1$,
$\mu=0.298000$, $0.161624$) were also verified; they follow from the table by
Lemma~\ref{lem:zeroext}.

\begin{proposition}[an explicit negative plane]\label{prop:negplane}
Let $T<\operatorname{arccosh}\sqrt2=0.88137$ and let $F,H\in L^2_{\mathrm{par}}(\IT)$ be
explicit real polynomials, given by their Legendre coefficients, such that $\langle\cdot,\chi_{\mathrm{par}}\rangle$
does not vanish on $\mathrm{span}\{F,H\}$. If the matrix
\[
B_3=\begin{pmatrix}\Gamma_T[F]&\Gamma_T(F,H)&\langle F,\chi_{\mathrm{par}}\rangle\\
\Gamma_T(F,H)&\Gamma_T[H]&\langle H,\chi_{\mathrm{par}}\rangle\\
\langle F,\chi_{\mathrm{par}}\rangle&\langle H,\chi_{\mathrm{par}}\rangle&0\end{pmatrix}
\]
is negative definite on a two-dimensional subspace, then $\mathrm{span}\{F,H\}\cap\PT^{\mathrm{par}}$
contains an $F_0\ne0$ with $\Gamma_T[F_0]<0$, and $n_{\mathrm{par}}(T)\ge1$.
\end{proposition}
\begin{proof}
$B_3$ is the bordered form of $\Gamma_T|_{\mathrm{span}\{F,H\}}$ with border $\langle\cdot,\chi_{\mathrm{par}}\rangle$,
so Lemma~\ref{lem:border} in dimension two gives
$n_-(\Gamma_T|_{\mathrm{span}\{F,H\}\cap\PT})=n_-(B_3)-1\ge1$.
\end{proof}
In the verification $F$ is of degree $<18$, $H$ is the Legendre truncation of
$\chi_{\mathrm{par}}$, and $\langle H,\chi_{\mathrm{par}}\rangle\in[0.072,1.69]$. The quadratic values
$\Gamma_T[G]=C_0c_G(0)+\int_0^{2T}(c_G(0)-c_G)K+c_G(0)\int_{2T}^\infty K$ (Theorem~\ref{thm:A} with
$\NT=\emptyset$), for $G=F,H,F\pm H$, are enclosed exactly as in steps (a)--(e) of the proof of
Theorem~\ref{thm:w11}, with the prime atoms and the pole term omitted; step (c) is the range condition
$T<0.88137$. The off-diagonal entry is obtained by polarisation, as in Theorem~\ref{thm:nest}, and the
pairings with $\chi_{\mathrm{par}}$ from exact moment series.

\begin{lemma}[rounding in the polarisation]\label{lem:polround}
If $F\pm H$ are formed in double precision, the value of $\Gamma_T(F,H)$ computed from
the rounded combinations differs from the true one by at most $10^{-10}$, for $0.6\le T\le0.9$ and the
vectors used.
\end{lemma}
\begin{proof}
The rounded combinations are $F\pm H+\delta_\pm$ with Legendre coefficients $|\delta_{\pm,k}|\le3\cdot10^{-16}$
(all coefficients of $F\pm H$ are at most $2.2$ in modulus), so on $\IT$ one has
$\|\delta\|_\infty\le3\cdot10^{-14}$ and $\|\delta'\|_\infty\le4\cdot10^{-12}$ for degree $<18$. The
Dirichlet form of Proposition~\ref{prop:nodes}(i) with $K(u)\le u^{-1}+\frac12$ gives
$\Lev(c_\delta)\le\|\delta'\|_\infty^2T^2(2T+2T^2)+\frac12\|\delta\|_\infty^2\int_{\IT}\kappa\le10^{-22}$,
and then
$|\Gamma_T[v+\delta]-\Gamma_T[v]|\le|C_0|(2\|v\|\|\delta\|+\|\delta\|^2)+2\sqrt{\Lev(c_v)\Lev(c_\delta)}+\Lev(c_\delta)\le2\cdot10^{-10}$
with $\Lev(c_v)\le50$; divide by four.
\end{proof}
Two further effects are covered by the same inflation: the stored interval ends are rounded to double
precision, and the value of Euler's constant entering $C_0$ is correct only to about $10^{-51}$. The
verification therefore rebuilds $B_3$ from the stored enclosures, inflates every entry by $\pm10^{-9}$, certifies negative definiteness on
a two-plane over the inflated box by interval Cholesky with a certified margin, and re-imposes
$T<0.88137$. It certifies $n_{\mathrm{ev}}\ge1$ at $T=0.60$ (margin $0.006627$) and $n_{\mathrm{odd}}\ge1$
at $T=0.81$ (margin $0.005124$), which is what Theorem~\ref{thm:index} uses, and also $n_{\mathrm{ev}}\ge1$
at $\log2$ and $0.80$ and $n_{\mathrm{odd}}\ge1$ at $0.82$ (margins $0.147170$, $0.290056$, $0.012368$),
which follow from the first two by Lemma~\ref{lem:zeroext}. A run at $T=0.90$ lies outside the range
condition and is not used.

\begin{lemma}[zero extension]\label{lem:zeroext}
For $0<T<T'$, extension by zero maps $\PT^{\mathrm{par}}$ isometrically into
$\mathcal P^{\mathrm{par}}_{T'}$ with the same values of $\Gamma$. Hence $n_{\mathrm{par}}$ is
nondecreasing and $\lambda_{\mathrm{par}}$ nonincreasing in $T$, with no range condition at either window.
\end{lemma}
\begin{proof}
The norm, the Tate moments, the parity and $\cF$ are unchanged, and $\Gamma_T(F)=C_0\|F\|^2+\Lev(\cF)$
depends on $F$ only through $\cF$ on $\R$ (Theorem~\ref{thm:A}); $\Ainf$ does not depend on $T$
(Proposition~\ref{prop:arch}(i)).
\end{proof}

\begin{lemma}[a certificate with $r=0$ is strict]\label{lem:rzero}
If Proposition~\ref{prop:haynsworth}(c) holds with $r=0$ at $T$, then
$\lambda_{\mathrm{par}}(T)>0$.
\end{lemma}
\begin{proof}
Then $\underline\Gamma\ge0$ on $\PT^{\mathrm{par}}$. At the attained minimiser $F^*$ of
Lemma~\ref{lem:strict}, $\Gamma_T[F^*]-\underline\Gamma[F^*]=\frac1{2\pi}\int(m-\underline m)|\Fh^*|^2>0$,
because $m-\underline m=m-C_0=\varphi(\tau^2)>0$ on $0<|\tau|<\tau^{(1)}$ (Theorem~\ref{thm:S}(i)) and
$\Fh^*$ is entire and not identically zero.
\end{proof}


\subsection{The certificates at \texorpdfstring{$T=\frac12\log5$}{T = (1/2) log 5}}\label{subsec:encllog5}
The notation $\chi_{\mathrm{par}}$, $e_j$, $S_c$, $\Lambda_k$ is that of \S\ref{subsec:enclindex}, and
$\chi_{\mathrm{par}}=\varphi_{\mathrm{par}}$ of \S\ref{subsec:alignineq}. We work with real functions, which
suffices as in \S\ref{subsec:alignineq}. Unlike Proposition~\ref{prop:haynsworth}, the primes are kept:
the lower form is a frequency truncation of the full symbol $\sT$, and the operator subtracted from a
multiple of the identity is indefinite.

\emph{The lower form.} Fix $\tau_N>0$ and a number $s_\infty\le m(\tau_N)-\Aconst$. Since $m$ is increasing
with $m(0)=C_0$ and the comb is at least $-\Aconst$ (Lemma~\ref{lem:mbounds}(i),(iii)), the even function
$\underline s$ equal to $\sT$ on $|\tau|\le\tau_N$ and to $s_\infty$ on $|\tau|>\tau_N$ satisfies
$\underline s\le\sT$. Write $g=s_\infty-\underline s=(s_\infty-\sT)\mathbf 1_{|\tau|<\tau_N}=g_+-g_-$ with
$g_\pm\ge0$, and $\mathcal G_{(\pm)}=Q_Tg_{(\pm)}(D)Q_T$. Then
\[
\mathfrak a_T(F)\ \ge\ s_\infty\|F\|^2-\langle\mathcal GF,F\rangle,\qquad
0\le\mathcal G_\pm\le\omega_\pm S_{c_*},\quad c_*=T\tau_N,
\]
with $\omega_+=s_\infty-C_0+\Aconst$ and $\omega_-=m(\tau_N)+\Aconst-s_\infty$. For $\mu\in\R$ and
$c\in\{0,\frac12,-\frac12\}$ put, on $L^2_{\mathrm{par}}(\IT)\oplus\R$,
\[
\mathfrak b_{\mu,c}(F,\xi)=(s_\infty-\mu)\|F\|^2-\langle\mathcal GF,F\rangle+2\xi\,\ell_{\mathrm{par}}(F)+c\,\xi^2 .
\]

\begin{proposition}[the lower certificate]\label{prop:lowercert}
Let $k_1<k_2$, let $E$ be the span of the $e_j$, $j<k_1$, of parity $\mathrm{par}$, $Y$ its orthogonal
complement in $L^2_{\mathrm{par}}(\IT)$, and $P_Y$ the orthogonal projection onto $Y$.
\emph{(a)} On $Y$, $\mathfrak b_{\mu,c}\ge q_Y:=s_\infty-\mu-\omega_+\Lambda_{k_1}(c_*)$.
\emph{(b)} Suppose $q_Y>0$, and let $\mathsf R$ be the symmetric matrix on $E\oplus\R$ with
\[
\mathsf R_{ij}=(s_\infty-\mu)\delta_{ij}-\langle\mathcal Ge_i,e_j\rangle
-\frac{\langle P_Y\mathcal Ge_i,P_Y\mathcal Ge_j\rangle}{q_Y},\quad
\mathsf R_{i\xi}=\langle e_i,\chi_{\mathrm{par}}\rangle+\frac{\langle P_Y\mathcal Ge_i,P_Y\chi_{\mathrm{par}}\rangle}{q_Y},\quad
\mathsf R_{\xi\xi}=c-\frac{\|P_Y\chi_{\mathrm{par}}\|^2}{q_Y}.
\]
Then $n_-(\mathfrak b_{\mu,c})\le n_-(\mathsf R)$, and the same holds for any symmetric $\mathsf R'\le\mathsf R$.
\emph{(c)} If $n_-(\mathsf R')\le1$ with $c=0$, then $\mathfrak a_T\ge\mu$ on the unit sphere of
$\PT^{\mathrm{par}}$. If $n_-(\mathsf R')=0$ with $c=\frac12$ and $\mathrm{par}=\mathrm{odd}$, or $n_-(\mathsf R')\le1$ with
$c=-\frac12$ and $\mathrm{par}=\mathrm{ev}$, then $W_T\ge\mu\|\cdot\|^2$ on $\mathcal D^{\mathrm{par}}_T$.
\end{proposition}
\begin{proof}
(a) As in Proposition~\ref{prop:haynsworth}(a): $-P_Y\mathcal GP_Y\ge-P_Y\mathcal G_+P_Y$ and
$\|P_Y\mathcal G_+P_Y\|\le\omega_+\operatorname{tr}(P^\perp S_{c_*}P^\perp)=\omega_+\Lambda_{k_1}(c_*)$.
(b) For $i<k_1$ and $y\in Y$, $\mathfrak b_{\mu,c}(e_i,y)=-\langle P_Y\mathcal Ge_i,y\rangle$ and the border pairs
with $y$ through $+\langle P_Y\chi_{\mathrm{par}},y\rangle$. The congruence of the proof of
Proposition~\ref{prop:haynsworth}(b) gives $\mathfrak b_{\mu,c}\cong\mathfrak s\oplus\mathfrak b_{\mu,c}|_Y$ with
$\mathfrak b_{\mu,c}|_Y\ge q_Y>0$, so $n_-(\mathfrak b_{\mu,c})=n_-(\mathfrak s)$, and the subtracted term of the Schur
complement $\mathfrak s$ is at most $q_Y^{-1}$ times the Gram matrix of the $Y$-columns
$-P_Y\mathcal Ge_i$ and $P_Y\chi_{\mathrm{par}}$; hence $\mathfrak s\ge\mathsf R\ge\mathsf R'$.
(c) Put $\mathfrak q_\mu(F)=\mathfrak b_{\mu,0}(F,0)$, so that $\mathfrak a_T-\mu\|\cdot\|^2\ge\mathfrak q_\mu$. For
$c=0$, Lemma~\ref{lem:border} gives $n_-(\mathfrak q_\mu|_{\ker\ell_{\mathrm{par}}})=n_-(\mathfrak b_{\mu,0})-1=0$. For
$c=\frac12$, $\mathfrak b_{\mu,1/2}=\mathfrak q_\mu-2\ell_{\mathrm{par}}^2+\frac12(\xi+2\ell_{\mathrm{par}})^2$, so
$\mathfrak q_\mu-2\ell_{\mathrm{odd}}^2\ge0$, and $W_T=\mathfrak a_T-2\ell_{\mathrm{odd}}^2$ on odd $F$
(Theorem~\ref{thm:parityschur}(o)). For $c=-\frac12$,
$\mathfrak b_{\mu,-1/2}=\mathfrak q_\mu+2\ell_{\mathrm{par}}^2-\frac12(\xi-2\ell_{\mathrm{par}})^2$; the substitution
$\eta=\xi-2\ell_{\mathrm{par}}(F)$ is a congruence, so
$n_-(\mathfrak q_\mu+2\ell_{\mathrm{ev}}^2)=n_-(\mathfrak b_{\mu,-1/2})-1=0$, and $W_T=\mathfrak a_T+2\ell_{\mathrm{ev}}^2$ on even
$F$. The statements pass from $L^2_{\mathrm{par}}(\IT)$ to the form domain because $\mathfrak q_\mu$ is bounded.
\end{proof}

\emph{The computation.} For $i-l$ even,
$\langle\mathcal Ge_i,e_l\rangle=\frac{2T}\pi\sqrt{(2i+1)(2l+1)}\,(-1)^{(l-i)/2}\int_0^{\tau_N}g(\tau)j_i(\tau T)j_l(\tau T)\,d\tau$
(proof of Theorem~\ref{thm:P}), and the entries with $i-l$ odd vanish. The integral is computed by composite
Gauss--Legendre quadrature with $46$ nodes on panels of length $\frac12$, with the remainder bound for
functions analytic in a Bernstein ellipse \cite[Thm.~19.3]{Trefethen2013}, $\rho=3.2$. The ellipse lies
within distance $0.439<\frac12$ of the panel, while the poles of $\psi(\frac14\pm\frac{i\tau}2)$ lie at
$|\Iml\tau|\ge\frac12$. The modulus on the ellipse is bounded with $|j_n(z)|\le e^{|\Iml z|}$,
$|\cos(z\log n)|\le\cosh(\Iml z\log n)$ and, for $\psi$, the recurrence $\psi(w)=\psi(w+8)-\sum_{k<8}(w+k)^{-1}$
together with Binet's bound $|\psi(u)|\le\log|u|+\frac\pi2+\frac1{2\Rel u}+\frac1{12(\Rel u)^2}$ for
$\Rel u\ge8$. The resulting remainder, $1.5\cdot10^{-43}$ per entry (below $2\cdot10^{-42}$ under either
point-count convention for the rule), is added as a ball to every entry. The inner products over $Y$ are
exact sums over $k_1\le l<k_2$ of the given parity plus tail boxes: by Cauchy--Schwarz for the positive
operators $\mathcal G_\pm$, $\|P^\perp_{k_2}\mathcal G_\pm e\|^2\le\|P^\perp_{k_2}\mathcal G_\pm P^\perp_{k_2}\|
\langle\mathcal G_\pm e,e\rangle\le\omega_\pm^2\Lambda_{k_2}(c_*)$, so
$\|P^\perp_{k_2}\mathcal Ge\|\le(\omega_++\omega_-)\Lambda_{k_2}(c_*)^{1/2}$, and
$\|P^\perp_{k_2}\chi_{\mathrm{par}}\|^2=\sinh T\pm T-\sum_{l<k_2}\langle e_l,\chi_{\mathrm{par}}\rangle^2$ with
$\langle e_l,\chi_{\mathrm{par}}\rangle=\sqrt{(2l+1)/2T}\,2T\,i_l(T/2)$. The tails $\Lambda_k$ are bounded by
$\Lambda^*_k$ of Corollary~\ref{cor:P1}. Everything is midpoint--radius ball arithmetic at $640$ bits
(Arb \cite{Johansson2017}, through \texttt{python-flint}), including the digamma and Bessel values and the
Gauss--Legendre nodes and weights; every Bessel ball is asserted to have radius below $10^{-60}$, and
comparisons are ball predicates.

\emph{Inertia.} In $\mathsf R$ the term $\|P_Y\chi_{\mathrm{par}}\|^2/q_Y$ is of order $10^{-190}$; it is replaced
by an exact upper bound $\delta'$, which gives $\mathsf R'\le\mathsf R$. When $n_-(\mathsf R')=0$ is to be shown,
interval Cholesky of $\mathsf R'$ is applied. When $n_-(\mathsf R')\le1$ is to be shown, a double-precision
Householder matrix $H$ that sends the midpoint of the column $(\mathsf R'_{i\xi})_i$ to a multiple of $e_0$ is
used as an exact matrix; $\|H^{\mathsf T}H-I\|_\infty<1$ is verified in balls, so $\mathrm{diag}(H,1)$ is a
congruence. In the new coordinates the $2\times2$ block $B$ on $(e_0,\xi)$ has $\det B<0$, certified, so its
inertia is $(1,1)$, and by Haynsworth's inertia additivity \cite{Haynsworth1968}
$n_-(\mathsf R')=1+n_-(S)$ for the Schur complement $S$ of $B$, whose positivity is verified by interval
Cholesky. The only non-rigorous objects are $H$ and the midpoints, which enter only through this exact
congruence.

\emph{Data.} $T=\frac12\log5$, $\tau_N=240$, $c_*=193.13$, $k_1=280$, $k_2=330$ ($\dim E=140$ in each parity).
From $m(\tau_N)=3.64276\ldots$: $s_\infty=0.700787608\ldots$, $\omega_+=9.0150$, $\omega_-=5.8839$,
$\Lambda^*_{k_1}(c_*)\le4.4\cdot10^{-20}$, $\Lambda^*_{k_2}(c_*)\le3.9\cdot10^{-70}$, $q_Y\ge0.70079$, tail
box $(\omega_++\omega_-)\Lambda^*_{k_2}(c_*)^{1/2}\le2.9\cdot10^{-34}$; the entry radii of $\mathsf R$ are at
most $4.4\cdot10^{-41}$. The verified cases are:
\begin{center}\small
\begin{tabular}{@{}llccc@{}}
\toprule
form & sector & certified $\mu$ & minimal pivot at that $\mu$ & fails at\\
\midrule
$\mathfrak a_T|\PT$ ($c=0$) & even & $2.5\cdot10^{-12}$ & $8.47\cdot10^{-3}$ & $3\cdot10^{-12}$\\
$\mathfrak a_T|\PT$ ($c=0$) & odd & $3.0\cdot10^{-10}$ & $2.03\cdot10^{-2}$ & $4\cdot10^{-10}$\\
$W_T$ ($c=-\frac12$) & even & $10^{-18}$ & $4.47\cdot10^{-4}$ & $10^{-16}$\\
$W_T$ ($c=\frac12$) & odd & $10^{-16}$ & $8.14\cdot10^{-14}$ & $10^{-14}$\\
\bottomrule
\end{tabular}
\end{center}
The radius of the smallest pivot in the last row is below $10^{-24}$. A failed Cholesky is not evidence of
negativity. A double-precision eigenvalue computation of the midpoint of $S$ at $\mu=0$ gives
$2.53\cdot10^{-12}$ (even) and $3.62\cdot10^{-10}$ (odd) for the primitive form, so the certified bounds are
within a small factor of the bottom of the lower form. As a check, the same procedure at
$T=\frac12\log3$ ($\tau_N=40$, $k_1=50$) certifies $\lambda_{\min}(A_T|\PT^{\mathrm{ev}})\ge6.3\cdot10^{-4}$,
consistent with Theorem~\ref{thm:C}. Each run takes about seven minutes per parity on one core.


\section{The Li coefficients}\label{app:li}

This appendix places the Li--Keiper criterion in the language of this paper. Its kernels have
nonnegative Fourier transform, which Corollary~\ref{cor:nontermwise} forbids for flat-top weights, and
they fall outside the flat-top class for a sharper reason as well.

\begin{corollary}[the Li kernel is outside the flat-top class]\label{cor:Li}
For $n\ge1$ let $\lambda_n=\sum_\rho\big[1-(1-1/\rho)^n\big]$, the sum taken in the symmetric pairing
of $\rho$ with $1-\rho$. Write $\rho=\frac12+i\gamma$ and $B(\gamma)=(2i\gamma-1)/(2i\gamma+1)$. Then
$\lambda_n=\sum_\rho\Phi_n(\gamma_\rho)$ with
\[
\Phi_n(\gamma)=1-\tfrac12\big(B(\gamma)^n+B(\gamma)^{-n}\big),
\]
an even rational function with poles of order $n$ at $\gamma=\pm\frac i2$ and no others, satisfying
$\Phi_n(\gamma)=1-\cos\big(n\theta(\gamma)\big)\ge0$ for real $\gamma$, where
$\theta(\gamma)=2\arctan\frac1{2\gamma}$. Moreover $\Phi_n=\widehat{w_n}$ with
\[
w_n(u)=-\tfrac12e^{-|u|/2}L_n'(|u|),
\]
$L_n$ the $n$-th Laguerre polynomial. For every $n\ge1$ and every $T>0$, $w_n\notin\Wcl$: its transform
is nonnegative, which Corollary~\ref{cor:nontermwise} forbids for any flat-top weight; its decay is
exactly $e^{-|u|/2}$, i.e.\ $\delta=0$ in Definition~\ref{def:W}; and it is nowhere locally constant,
so it has no flat top on any window. In the sharp form of the second reason,
$\widehat{w_n}(\pm\frac i2)=\infty$ and the associated prime sum
$\sum_m\Lambda(m)m^{-1}P_n(\log m)$ diverges, so $w_n$ is not an admissible test function for the
explicit formula at all: $\lambda_n$ is the zero-sum half of a Weil functional, not a Weil-functional
evaluation in the sense of Theorem~\ref{thm:Z}.
\end{corollary}
\begin{proof}
$1-1/\rho=B(\gamma)$ and $1-1/(1-\rho)=B(\gamma)^{-1}$, so pairing $\rho$ with $1-\rho$ halves the sum
into $\Phi_n$. $B$ is a M\"obius map with a simple pole at $\frac i2$ and a simple zero at $-\frac i2$,
and $B(\gamma)\to1$ as $|\gamma|\to\infty$ with
$\frac12(B^n+B^{-n})=1-n^2/(2\gamma^2)+O(\gamma^{-4})$, so $\Phi_n\in L^1(\R)\cap C(\R)$ and the
pairing sum converges absolutely. For real $\gamma$, $1-1/\rho=(\gamma+\frac i2)^2/|\gamma+\frac i2|^2$
has modulus $1$ and argument $\theta(\gamma)=2\arg(\gamma+\frac i2)$ for $\gamma>0$, the case
$\gamma<0$ following by evenness, whence
$\Phi_n=1-\cos n\theta\ge0$. Fourier inversion of the rational $\Phi_n\in L^1$ by contour closure
collects the single $n$-fold residue at $\pm\frac i2$ and gives $w_n(u)=P_n(|u|)e^{-|u|/2}$ with
$\deg P_n=n-1$, and $P_n=-\frac12L_n'$ by the generating identity for the Laguerre polynomials. The
three exclusions are Corollary~\ref{cor:nontermwise}, Definition~\ref{def:W} and the definition of a
flat top; for the last statement, $\Phi_n$ has a pole of order $n$ at $\pm\frac i2$, the boundary of
the strip of admissibility, and $\sum_{m\le X}\Lambda(m)(\log m)^j/m\asymp(\log X)^{j+1}$ diverges for
every $j\ge0$.
\end{proof}
In this sense the Li criterion is termwise:
$\Phi_n\ge0$ pointwise on $\R$, so positivity of $\lambda_n$ needs no cancellation across the zeros,
and Corollary~\ref{cor:nontermwise} says that no member of $\Wcl$ can be used that way.

\begin{remark}[the Li coefficients do not see collectivity]\label{rem:LiSC}
The analogue of Theorem~\ref{thm:SC} in the Li coordinates of Corollary~\ref{cor:Li} is false. Define
$\lambda_n^{(S)}$ by deleting from the arithmetic side all contacts but $2$ and $3$, i.e.\ the
archimedean term of $\lambda_n$ minus the finite explicit sum over the prime powers of $2$ and $3$.
Then $\lambda_n^{(S)}$ is eventually positive \emph{unconditionally} --- the sub-comb side alone takes
the values $-1.24,+0.63,-0.31,-1.64$ at $n=1,10,10^2,10^3$, bounded, against an archimedean term of
order $n\log n$ --- because the Li weights carry $m^{-1}$, which is subcritical, while the weights
$v_n=2\Lambda(n)n^{-1/2}$ of Theorem~\ref{thm:SC} are critical. Under the Riemann Hypothesis
$\lambda_n^{(S)}$ differs from $\lambda_n$ by $o(n\log n)$, by the Bombieri--Lagarias and Coffey
asymptotics \cite{BombieriLagarias1999,Coffey2005}. So the collectivity of positivity --- that no
proper sub-comb of mass below $-C_0$ carries the sign --- is invisible in those coordinates. In the
computed range the arithmetic content of $\lambda_n$ is a lower-order correction to a growing
archimedean main term.
\end{remark}

\section{Numerical observations}\label{app:diag}

This appendix collects numerical values computed in double precision (Galerkin values, linear
programs, grid evaluations). They are upper bounds or approximations, not certified values, and no
proof in the paper depends on them.

\begin{itemize}
\item At $T=\frac12\log3$, a Galerkin compression with $K=120$ modes gives the upper bound
$\lamA\le9.893\cdot10^{-4}$; the value $8.8577\cdot10^{-5}$ of Theorem~\ref{thm:C} is about $9\%$ of it.
\item At the two-prime windows ($B_*=7.621$), with the Galerkin upper bound at $T=0.58$, the radius of
Theorem~\ref{thm:P3}(ii) is $R(\tfrac12,\tfrac12)=8.98\cdot10^5$, $R(\tfrac12,\tfrac14)=1.02\cdot10^{13}$
and $R(\tfrac1{10},\tfrac1{10})=2.70\cdot10^{32}$; these are smaller than the envelope threshold of
Proposition~\ref{prop:env} for the certificates actually built.
\item The linear program with free $\beta$ at a two-contact window returned $\beta_0=2.0400$,
$\beta_1=-0.1478$.
\item At $T=\frac12\log3$, lattice orders $M=8,12,16,19,24$ gave linear-program values
$-7.1\cdot10^{-3}$, $-4.4\cdot10^{-3}$, $+2.2\cdot10^{-4}$, $+3.4\cdot10^{-4}$, $+1.7\cdot10^{-4}$ on
their own $\tau$-grids. These grid values are unreliable: a coarse-grid value $+1.7\cdot10^{-4}$
corresponded to a true infimum of $-2.1\cdot10^{-2}$. After re-optimization in $128$ blocks on a grid
$10^3$ times finer, the linear-program value and the true infimum agreed to $2\cdot10^{-9}$.
\item In the verification of Theorem~\ref{thm:C}, the measured discrepancy between the double-precision
and the interval evaluators at the four fine-grid minima and at $60$ random points was
$5.4\cdot10^{-13}$, consistent with the a priori bound $10^{-9}$.
\item For $\frac12\log3<T\le\log2$ the Galerkin upper bound is $2.16\cdot10^{-4}$ at $T=0.58$. For
larger $T$ the computed bounds fall below $10^{-4}$ and do not stabilize in the basis size
(Remark~\ref{rem:galerkin}); being upper bounds, they are consistent with both a positive and a
negative infimum.
\item At $T=0.58$, a lattice certificate built as in \S\ref{sec:cert} ($M=24$, $414$ atoms, $\tau_1=30$,
$N=2$, $128$-block re-optimization, exhaustive second-order cover with no unresolved cell) has
infimum $-3.1964\cdot10^{-3}$. Adding the canonical column $\nu_{\frac12,0}$ (Lemma~\ref{lem:hybrid})
did not improve it within the available solver accuracy, and a clustered placement of the atoms was
worse by a factor $4$. Varying the linear-program tolerance, the grid, the verifier and the tail lemma
did not change this; what remains is the size of the certificate class.
\item The optimized certificates of the primitive problem have polar multiplier $\beta_0$ close to $2$, and $\beta_1$ drifts when a second contact is added; by Theorem~\ref{thm:Zexh}(ii) this drift measures the correction in $\beta$. The computed certificates show no visible arithmetic structure, presumably because the class of admissible $w$ is large and the optimizer selects an arbitrary member.
\item Test-kernel ceilings at $T=\frac12\log3$ (double precision, two independent quadratures agreeing to six figures). With the $\beta$ of Theorem~\ref{thm:C}, $\lambda_{\min}(C_8)=-1.33\cdot10^{-2}$ and $\lambda_{\min}(C_{24})=4.15\cdot10^{-4}$, tending to about $4.18\cdot10^{-4}$ as $M\to\infty$. After the concave maximization over $\beta$, $\operatorname{Ceil}(8)=-5.4\cdot10^{-3}$, $\operatorname{Ceil}(12)=7.6\cdot10^{-5}$ and $\operatorname{Ceil}(24)=5.9\cdot10^{-4}$; the certificate of Theorem~\ref{thm:C} realizes about $15\%$ of its class ceiling.
\item The same eigenvalue computation at a two-prime window decides, before any linear program is run, whether a given lattice resolution can reach $c>0$ there.
\item The linear-program value at $T=\frac12\log3$, $M=8$ is $-7.1\cdot10^{-3}$, below the certified ceiling $-5.70\cdot10^{-3}$ of Theorem~\ref{thm:N}.
\item At $T=\log2$ the nodes of the bottom eigenvector of $\Ainf$ on $\PT$ lie at $\pm0.3152$ (one nodal pair).
\item Galerkin value $\laminfP(\frac12\log3)\approx+0.083$; with the primes switched on the bottom drops to about $9.8\cdot10^{-4}$ (Galerkin upper bound $9.89\cdot10^{-4}$), consistent with the certificate of Theorem~\ref{thm:C} having $\beta\approx(2,0)$.
\item At $T=\frac12\log3$ a quadrature of the exact form gives $\Ainf(F^{\mathrm{odd}}_T)=-0.0268769$, against the bound $U(\frac12\log3)=-0.0200526$ of Proposition~\ref{prop:archneg}.
\item Galerkin values in an orthonormal Legendre basis of $L^2(\IT)$ give the variational upper bounds $T^*_\infty\le0.3716$ and $T^*_{\infty,\mathrm P}\le0.5958$; the second lies inside the proved bracket of Corollary~\ref{cor:Tarch}.
\item At $T=\log2$ the Galerkin bottom of $\Ainf$ on $\PT$ is about $-0.154$, and on that direction the prime form supplies about $+0.188$.
\item The triangle bound $\lamA\ge\laminf-\|\PiT\|$ at $T=\frac12\log3$ evaluates (with Galerkin values) to about $-0.92$. Measured Galerkin norms of $\PiT$ are $0.49$ on the one-contact range and $0.69$--$0.90$ on the two-prime range; the jump is the entry of the contact $n=3$.
\item Evaluating $\Lambda_k$ of Theorem~\ref{thm:P} by the trapezoid rule on a $2\cdot10^5$-point grid at $\tau_1=6.4222$, $6.4222$, $6.3576$ gives the counts $5$, $6$, $7$ at $T=\frac12\log3,\log2,0.80$ (margins $\Lambda_k/\theta\approx0.76$--$0.86$), slightly sharper than Corollary~\ref{cor:P3}; correspondingly the brackets on $L^2(\IT)$ would read $[2,5]$, $[2,6]$, $[3,7]$.
\item At $T=\log2$ and the optimal $\tau_1$, $\lambda_4(S_{T,\tau_1})=2.18\cdot10^{-2}$ against $\theta=3.87\cdot10^{-3}$, so the exact count in Proposition~\ref{prop:NC} is $5$ (Remark~\ref{rem:P4}).
\item A Galerkin computation with $N=60$ finds exactly $12$ competing even modes at $T=\log2$, the thirteenth at level $2.2568>\Aconst$.
\item In unpublished numerical work of the author, the primal minimizer at window $T$ clears the zeros up to a height $H_\star\approx2\pi e^{2T+1}$, which is of the order of the resolution scale $2\pi e^{2T}$ of \S\ref{subsec:where}.
\item At the two windows for which certificates were computed, the last sign change of $\sT$ lies at $10.5$ and $46.1$, against the envelope thresholds $\tau_{\mathrm{env}}=16.7$ and $59.5$.
\item Minimized over $k$, the right side of \eqref{eq:Bmargin} exceeds $1$ for $T<1.430$ and falls below the value $8.8577\cdot10^{-5}$ of Theorem~\ref{thm:C} only at $T=2.054$ (numerical evaluation of the explicit expression).
\item At $T=0.58,0.62,0.65$ and $\log2$ the pivot value of \S\ref{subsec:schur} is $w_{11}=0.0924,0.0626,0.0415,0.0342$, far above the (unconverged, Remark~\ref{rem:galerkin}) Galerkin bottom of $W_T|\PT$.
\item At the same four windows the crude sufficient condition $\lambda_2-\|\PiT|_V\|-c^2/w_{11}>0$, $c=\|b\|$, of Proposition~\ref{prop:schur} fails by $1.07$, $1.19$, $1.32$, $1.49$.
\item By parity the Schur correction acts only in the even sector, where the first available archimedean level is the third; numerically the ground state of the Schur form $S$ lies there, with $99.46\%$ of its mass on the third archimedean eigenvector at $\log2$. The condition $\lambda_3-c^2/w_{11}>0$, in which the prime form is dropped, fails by $0.064$, $0.101$, $0.154$, $0.310$ at $T=0.58,0.62,0.65,\log2$.
\item In the two-prime computations of \S\ref{subsec:twoprime}, $\lamA$ falls by three orders of magnitude between $T=0.58$ and $T=\log2$.

\end{itemize}


\begin{remark}\label{rem:Zcheck}
Both sides of the identity \eqref{eq:fw}--\eqref{eq:Z} have been evaluated numerically and
independently for $C^2$ flat-top weights at eight windows in $[\frac12\log3,1.20]$, the zero side
using the first $10^5$ ordinates of Odlyzko's table and the prime side the elementary right-hand
side of \eqref{eq:fw}; the two agree to between $2.6\cdot10^{-7}$ and $2.8\cdot10^{-6}$. At
$T=\frac12\log3$ the truncations of the zero side to the first $10$, $10^2$ and $10^3$ ordinates
leave residuals $4.0\cdot10^{-4}$, $6.0\cdot10^{-6}$ and $1.2\cdot10^{-7}$, the last at the floor of
the double-precision evaluation; the orders of magnitude are the ones Lemma~\ref{lem:smooth}(iii)
predicts at $k=2$.
\end{remark}

\begin{remark}\label{rem:comblp}
The comb-vanishing class has been optimised over. Taking for $w$ the compactly supported $C^2$
cubic splines on a lattice of $48$ cells with $U=6$, matched to $1$ across the window edge, with the
comb equalities imposed exactly and $\beta$ free, the linear program over $[0,200]$ returns
negative minima at $T=\frac12\log3$, $0.58$, $0.62$ and $\log2$, namely $-1.2\cdot10^{-2}$,
$-1.7\cdot10^{-2}$, $-5.2\cdot10^{-2}$ and $-8.9\cdot10^{-2}$ in double precision, the zero-side and
prime-side evaluations agreeing to $10^{-5}$ at each argument of the minimum. For the resulting
weights $\tau_{\mathrm{stop}}<200$ and $L_w>0$ beyond (the constants $D(w)$ being quadratures with
a margin of about $30\%$ before the closure fails, and the comb residual $E_{\mathrm{at}}\le2\cdot10^{-14}$),
so for these weights the failure is a statement about all of $\R$ and not about an interval chosen
for convenience. The binding minimum sits at $\tau\approx4$--$8$, whereas the same program without
the comb conditions has its minimum near $\tau\approx43$: the small-$\tau$ failure is induced by
the comb conditions. One structural observation: since
consecutive prime powers are spaced $\approx e^{-u}$, once the lattice step exceeds that spacing
every added cell is consumed by the conditions inside it, so comb-vanishing forces an effectively
bounded support whatever $U$ is taken to be.
\end{remark}

The measured ceilings at $T=\frac12\log3$ (double precision, an unrestricted
supremum over $\beta$ on a finite grid of $\tau_0$) rise from $6.8\cdot10^{-5}$ at $M=12$ through
$5.9\cdot10^{-4}$ at $M=24$ to $9.1\cdot10^{-4}$ at $M=96$, with the maximising $\beta$ staying near
$(2.00,-0.02)$, well inside the box; the Galerkin upper bound is $9.893\cdot10^{-4}$. No rate is
asserted: Theorem~\ref{thm:R} proves convergence and, unconditionally, gives $\Delta\cdot\mu(|\tau|>P/2)=O(1)$
only, so the rate is part of Problem~\ref{prob:A}.

At the two-prime windows the same eigenvalue has been computed in double precision. At $T=0.58$ the ceiling is capped at about
$2.0\cdot10^{-4}$, stable to two significant figures under refinement of the $\tau_0$ grid; the
negative value $-3.20\cdot10^{-3}$ reached there by the optimiser (\S\ref{subsec:twoprime}) is
therefore a property of the optimiser and not of the class, and a certificate at that window is not
excluded by the ceiling. At $T=0.68$ and $T=\log2$ the ceiling at $M=24$ is negative, and the interval
enclosure of Theorem~\ref{thm:Np} turns that into a class impossibility at both windows, for every
$M\le24$.

The test-kernel ceilings of Proposition~\ref{prop:TK} for the archimedean form
alone, $\mathrm{Ceil}_\infty(M,T)$, approach $\laminfP$ from below at both $T=\frac12\log3$ and
$T=\log2$, with a deficit $\approx0.28\,M^{-1.2}$ over $8\le M\le64$, essentially independent of $T$
for $M\ge24$. This is a statement about the \emph{archimedean} program, where $\beta=0$, the polar
block is absent and $G$ vanishes to order $r$ at $2T$, so there is no edge and sampling really is
$O(M^{-r})$; it must therefore not be compared with the edge rate $\Theta(M^{-1})$ of
Theorem~\ref{thm:AL}, which is a statement about the full program. In the full program no ceiling
has been computed. Whether $D\asymp M^{-1}$ in Problem~\ref{prob:A} is open.

Two observations on of the matrix $(\pi_{ij})$ in the even archimedean basis at
$T=\log2$. First, what makes the $2$-plane of Theorem~\ref{thm:nest} positive is a cancellation
between the two prime atoms: $\pi_{13}=+0.1867\ (\text{atom }2)-0.1410\ (\text{atom }3)=+0.0457$, a
$76\,\%$ cancellation; with the atom at $\log2$ alone the determinant would be $-0.033$ and the form
already indefinite in two dimensions. Second, $\PiT$ mixes the archimedean modes at full strength ---
the diagonal is not dominant, and the near-zero bottom of the Galerkin form appears at the
\emph{last} even mode of every truncation, at $N=18$ and at $N=36$ alike. The obstruction is therefore
not captured by any fixed number of archimedean modes, which is both why Remark~\ref{rem:galerkin} is
forced and why a tower of subspaces such as $E_k$ cannot converge to a statement about the form.

\begin{remark}\label{rem:nodal}
At $T=\log2$, with $F^*=e_1$ and $e_3$ the two even vectors of Theorem~\ref{thm:nest}, the overlap
$I(a)=I_{F^*,e_3}(a)$ vanishes at $a=0.2128$, $0.6087$ and $1.0225$, and the two contacts straddle
the third of these: $\log2=0.6931<1.0225<1.0986=\log3$. That straddle is the $76\,\%$ cancellation
of the preceding paragraph --- the two atoms computed from Lemma~\ref{lem:N} are $+0.1866$ and
$-0.1410$, agreeing with the matrix entries above to three digits, which is the quadrature level.
The forced-sign threshold of Proposition~\ref{prop:G} is $a^*=T+z_e=1.1717$, and $\log3$ misses it
by $6.65\,\%$: one contact more, or a window a few per cent wider, and both atoms would be forced
positive, with no cancellation to make. The same straddle holds at $T=0.62$ and $T=0.65$, where the
third zero of $I$ sits at $0.9139$ and $0.9584$ and $\log3$ is still inside the window. Finally the
nodes of $F^*$, at $\pm0.3152$, are forced and not accidental: the two Tate constraints
$A(F)=B(F)=0$ both read $\int F(t)\cosh\frac t2\,dt=0$ for even $F$, and $\cosh>0$, so every even
primitive $F$ changes sign. Numerically, then, the positivity of the $2$-plane at $\log2$ is
governed by the position of a zero of an overlap function, which is determined by the node of the
archimedean ground state.
\end{remark}

\begin{remark}\label{rem:signrule}
The nodes decide the large entries of the prime matrix and only those. Refine
$(\log n-T,T)$ by the nodes of $e_j$ and by the nodes of $e_i$ shifted by $\log n$; on each cell
the integrand of Lemma~\ref{lem:N} has a constant sign, readable from the nodes alone, and predict
the sign of $\PiT(e_i,e_j)$'s $n$-th atom from the longest cell (Proposition~\ref{prop:G} is the
case of a single cell, where the prediction is a theorem). Measured against the exact overlaps on
the $8\times8$ even block at $T=\log2$, the rule predicts $100\,\%$ of the entries with
$|I|\ge0.10$, which carry $89$--$96\,\%$ of the total mass $\sum|I|$, and fails only on the small
near-cancelling entries (mean $|I|$ at a failure $0.014$--$0.028$ against $0.07$--$0.11$ over the
whole matrix); the overall rate is $59$--$80\,\%$. In the same computation the $k$-th even
eigenvector has exactly $k$ nodes in $(0,T)$ for $k\le8$, with strict interlacing of consecutive
modes --- a Sturm pattern whose base case is the lower bound of
Proposition~\ref{prop:nodes}(ii)--(iii) and which is otherwise a conjecture, the oscillation theory
that would prove it being unavailable by Remark~\ref{rem:PF2}. Lemma~\ref{lem:N} reproduces the
prime matrix to $1.7\cdot10^{-5}$.
\end{remark}


\begin{remark}\label{rem:bandtable}
The negative set of $\sT$ at the two regimes of Proposition~\ref{prop:align}, from bisection on a grid of
$4\cdot10^5$ points in double precision. For $\NT=\{2,3\}$ its components in $\tau\ge0$ are
$(0,2.278)$, $(4.816,7.272)$, $(9.666,11.966)$, $(16.251,18.494)$, $(27.558,28.938)$, $(45.149,46.129)$,
so $\tau_{\mathrm{last}}\approx46.129$ (the $46.13$ of Proposition~\ref{prop:env}); for $\NT=\{2,3,4\}$ the same
six within $0.6$, then $(62.616,63.777)$ and $(108.495,108.950)$, so $\tau_{\mathrm{last}}\approx108.950$. The
components beyond $\tau_a$ sit at the pairs $(k,l)=(3,5)$, $(5,8)$ for $\{2,3\}$, with $2^l/3^k=1.185$, $1.053$,
and additionally $(7,11)$, $(12,19)$ for $\{2,3,4\}$, with $0.936$, $0.9865$; the fractions $8/5$ and $19/12$
carrying the last components are convergents of $\log3/\log2$, and $5/3$, $11/7$ are intermediate
fractions. The traces of Proposition~\ref{prop:negtrace}(i) are $\operatorname{tr}\mathrm{Neg}_T=5.233$ at $T=\log2$
(with $\sup\sigma_-=7.621=B_*$) and $7.039$ at $T=\frac12\log5$, against $\operatorname{tr}M_-=3.116$ for the
archimedean part alone at $\log2$; the bound $\lamA\ge-5.233$ they give is useless for the sign. Finally,
the test function of Theorem~\ref{thm:margin}, evaluated by quadrature of the symbol (no zeros enter), has
Rayleigh quotient smaller than the proved bound $\mathcal B(T,k)$ by six to ten orders of magnitude over
$\frac12\log3\le T\le1.5$: the proof pays for the sinc-power bound at every zero, for the counting bound
of Lemma~\ref{lem:Ncount}, and for the weight $\tau^8$.
\end{remark}

\begin{remark}\label{rem:odd080}
Theorem~\ref{thm:index} leaves $n_{\mathrm{odd}}(0.80)\in\{0,1\}$, hence $n_-(\Ainf|\mathcal P_{0.80})\in\{1,2\}$;
whether $(\mathrm{H1})(0.80)$ holds, equivalently whether $T^*_{\mathrm{odd}}>0.80$, is open. The upper
certificate of Proposition~\ref{prop:haynsworth} at $T=0.80$ in the odd sector stops at $r=1$: the second
eigenvalue of the midpoint of $\mathsf R$ is $-0.0047$. With $77$ bands instead of $28$ it is $-0.0043$,
and with the bands extended to $\tau_N=30$ ($c_*=24$, $k_1=30$, $k_2=44$) it is $-0.0019$; so the residual
is set by the constant minorant $\underline m_\infty$ beyond $\tau_N$ rather than by the band grid
(, three runs). Galerkin bottoms of the odd sector, which are upper bounds on
$\lambda_{\mathrm{odd}}$ only up to the floating-point constraint, are $+0.0684$, $+0.0288$, $+0.0033$ at
$T=0.75$, $0.78$, $0.80$ and $-0.0092$ at $T=0.81$; they place the odd crossing in $(0.80,0.81)$ and prove
neither sign at $0.80$. Likewise the Galerkin upper bound $0.5958$ on $T^*_{\mathrm{ev}}=T^*_{\infty,\mathrm P}$
(\S\ref{subsec:arch}) is ; the certified bracket is $[0.58,0.60]$.
\end{remark}

\begin{remark}\label{rem:polarcapdiag}
These are float64 Legendre--Galerkin values of the polar-capacity constants of Theorem~\ref{thm:parityschur}.
A Galerkin $\mathfrak g_{\mathrm{par}}$ is a minimum over a subspace, hence an \emph{upper} bound on the true constant.
A negative value is therefore evidence of a true negative, up to rounding and quadrature error. A positive
value is no evidence at all. At $T=\frac12\log3$ with the primes deleted, $\mathfrak g_{\mathrm{ev}}\approx-0.451$ and
$\mathfrak g_{\mathrm{odd}}\approx-20.08$, with primitive sector bottoms $\approx+0.0825$ and $+0.3816$, so both channels
fail. For $\zeta$ at the same window, $\mathfrak g_{\mathrm{ev}}\approx8.6\cdot10^{-8}$ and $\mathfrak g_{\mathrm{odd}}\approx1.35\cdot10^{-3}$,
stable over $16$, $24$ and $32$ modes. There $S^{\mathrm{Sch}}_T\approx\mathrm{diag}(-1.99999991,+2.00135)\approx-J$:
the primitive Schur complement nearly cancels the polar block. In the family $a_\theta$ of
Corollary~\ref{cor:convfail}(ii) the complete form fails for $\theta$ below $\theta_*\approx0.99991$, and its
even channel fails already at $\theta=1.001$ ($\mathfrak g_{\mathrm{ev}}\approx-4.9\cdot10^{-5}$). With the primes
deleted, the odd and the even channel of the complete form cross zero at $T\approx0.372$ and $T\approx0.44$. The
positivity of $\mathfrak g_{\mathrm{ev}},\mathfrak g_{\mathrm{odd}}$ for $\zeta$ at $T\le0.8$ rests on \cite{Zhu2026} and
Theorem~\ref{thm:parityschur}(c), not on these numbers.
\end{remark}

The labels of these measurements, moved here from the table of \S\ref{sec:repro}:

\begin{thebibliography}{99}
\bibitem{BoasKac1945} R.~P.~Boas and M.~Kac, \emph{Inequalities for Fourier transforms of positive
functions}, Duke Math.~J. \textbf{11} (1945), 189--206.
\bibitem{BombieriLagarias1999} E.~Bombieri and J.~C.~Lagarias, \emph{Complements to Li's criterion
for the Riemann Hypothesis}, J.~Number Theory \textbf{77} (1999), 274--287.
\bibitem{Bombieri2000} E.~Bombieri, \emph{Remarks on Weil's quadratic functional in the theory of
prime numbers, I}, Atti Accad.~Naz.~Lincei Cl.~Sci.~Fis.~Mat.~Natur.~Rend.~Lincei (9) Mat.~Appl.
\textbf{11} (2000), no.~3, 183--233.
\bibitem{Carleman1923} T.~Carleman, \emph{Sur les \'equations int\'egrales singuli\`eres \`a noyau r\'eel et sym\'etrique}, Uppsala Universitets \AA rsskrift, 1923.
\bibitem{Coffey2005} M.~W.~Coffey, \emph{Toward verification of the Riemann hypothesis:
application of the Li criterion}, Math.~Phys.~Anal.~Geom. \textbf{8} (2005), 211--255.
\bibitem{CohnElkies2003} H.~Cohn and N.~Elkies, \emph{New upper bounds on sphere packings I},
Ann.~of Math. \textbf{157} (2003), 689--714.
\bibitem{ConnesConsani2021} A.~Connes and C.~Consani, \emph{Weil positivity and trace formula, the archimedean place}, Selecta Math.\ (N.S.) \textbf{27} (2021), no.~4, Paper No.~77; arXiv:2006.13771.
\bibitem{CCquasi2020} A.~Connes and C.~Consani, \emph{Quasi-inner functions and local factors}, arXiv:2008.10974 (2020).
\bibitem{CCM2023} A.~Connes, C.~Consani and H.~Moscovici, \emph{Zeta zeros and prolate wave operators}, arXiv:2310.18423 (2023).
\bibitem{CCM2024} A.~Connes, C.~Consani and H.~Moscovici, \emph{On $q$-series and the moment problem associated to local factors}, arXiv:2403.01247 (2024).
\bibitem{CvS2025} A.~Connes and W.~D.~van Suijlekom, \emph{Quadratic forms, real zeros and echoes of the spectral action}, arXiv:2511.23257 (2025).
\bibitem{DavenportHeilbronn1936} H.~Davenport and H.~Heilbronn, \emph{On the zeros of certain
Dirichlet series}, J.~London Math.~Soc. \textbf{11} (1936), 181--185 and 307--312.
\bibitem{DLMF} NIST Digital Library of Mathematical Functions, \url{https://dlmf.nist.gov}, \S5.9.
\bibitem{Groskin2026} A.~Groskin, \emph{A finite Guinand--Weil dictionary and archimedean tail
order for the truncated Weil quadratic form}, arXiv:2607.02828 (2026).
\bibitem{Haynsworth1968} E.~V.~Haynsworth, \emph{Determination of the inertia of a partitioned Hermitian
matrix}, Linear Algebra Appl.\ \textbf{1} (1968), 73--81.
\bibitem{HLP} G.~H.~Hardy, J.~E.~Littlewood and G.~P\'olya, \emph{Inequalities}, 2nd ed., Cambridge Univ.~Press, 1952.
\bibitem{Katznelson} Y.~Katznelson, \emph{An Introduction to Harmonic Analysis}, 3rd ed.,
Cambridge Univ.~Press, 2004.
\bibitem{IK} H.~Iwaniec and E.~Kowalski, \emph{Analytic Number Theory}, Amer.~Math.~Soc.
Colloq.~Publ.~53, 2004.
\bibitem{Johansson2017} F.~Johansson, \emph{Arb: efficient arbitrary-precision midpoint-radius interval
arithmetic}, IEEE Trans.\ Comput.\ \textbf{66} (2017), 1281--1292.
\bibitem{Karlin1968} S.~Karlin, \emph{Total Positivity, Vol.~I}, Stanford Univ.~Press, 1968.
\bibitem{Kim2026} T.~Kim, Y.~Hong, M.~Kim, S.~Choi, J.~Jang and M.~Kim, \emph{A numerical realization of Suzuki's Weil-quadratic-form operator: the archimedean spectral law, its universality, and an operator form of Weil's positivity criterion}, arXiv:2607.24830 (2026).
\bibitem{Krein1940} M.~G.~Kre\u{\i}n, \emph{Sur le probl\`eme du prolongement des fonctions
hermitiennes positives et continues}, Dokl.~Akad.~Nauk SSSR \textbf{26} (1940), 17--22.
\bibitem{PlattTrudgian2021} D.~J.~Platt and T.~S.~Trudgian, \emph{The Riemann hypothesis is true up to
$3\cdot10^{12}$}, Bull.\ London Math.\ Soc.\ \textbf{53} (2021), 792--797.
\bibitem{Poineau2010} J.~Poineau, \emph{La droite de Berkovich sur $\Z$}, Ast\'erisque \textbf{334} (2010).
\bibitem{Rockafellar} R.~T.~Rockafellar, \emph{Convex Analysis}, Princeton Univ.~Press, 1970.
\bibitem{Rosser1941} J.~B.~Rosser, \emph{Explicit bounds for some functions of prime numbers}, Amer.\ J.\
Math.\ \textbf{63} (1941), 211--232.
\bibitem{RosserYoheSchoenfeld1968} J.~B.~Rosser, J.~M.~Yohe and L.~Schoenfeld, \emph{Rigorous computation and
the zeros of the Riemann zeta-function}, Proc.\ IFIP Congress 1968, 70--76.
\bibitem{Sasvari1994} Z.~Sasv\'ari, \emph{Positive Definite and Definitizable Functions},
Akademie Verlag, Berlin, 1994.
\bibitem{Schoenberg1951} I.~J.~Schoenberg, \emph{On P\'olya frequency functions.~I. The
totally positive functions and their Laplace transforms}, J.~Analyse Math. \textbf{1} (1951), 331--374.
\bibitem{Simon2005} B.~Simon, \emph{Trace Ideals and Their Applications}, 2nd ed., Math.\ Surveys Monogr.\ \textbf{120}, Amer.~Math.~Soc., 2005.
\bibitem{SSV2012} R.~L.~Schilling, R.~Song and Z.~Vondra\v{c}ek, \emph{Bernstein Functions:
Theory and Applications}, 2nd ed., de Gruyter, 2012.
\bibitem{SongVondracek2003} R.~Song and Z.~Vondra\v{c}ek, \emph{Potential theory of subordinate
killed Brownian motion in a domain}, Probab.~Theory Related Fields \textbf{125} (2003), 578--592.
\bibitem{Suzuki2023jlms} M.~Suzuki, \emph{Aspects of the screw function corresponding to the
Riemann zeta function}, J.~Lond.~Math.~Soc.~(2) \textbf{107} (2023), 1846--1888.
\bibitem{Suzuki2023screwline} M.~Suzuki, \emph{The screw line of the Riemann zeta-function and
its applications}, arXiv:2209.04658.
\bibitem{Suzuki2025hilbert} M.~Suzuki, \emph{On the Hilbert space derived from the Weil
distribution}, arXiv:2301.00421.
\bibitem{Titchmarsh} E.~C.~Titchmarsh, \emph{The Theory of the Riemann
Zeta-Function}, 2nd ed., revised by D.~R.~Heath-Brown, Oxford Univ.~Press, 1986.
\bibitem{Trefethen2013} L.~N.~Trefethen, \emph{Approximation Theory and Approximation Practice},
SIAM, Philadelphia, 2013.
\bibitem{Paper42} D.~A.~Trejo Pizzo, \emph{An arithmetic Lefschetz programme over $\Z$: a metrized
bivariant correspondence square, its determinant pairing and its corner trace}, preprint, 2026.
\bibitem{Paper45} D.~A.~Trejo Pizzo, \emph{An explicit formula and a comparison theorem for the canonical trace on a correspondence
algebra of the idele class group}, preprint, 2026.
\bibitem{Weil52} A.~Weil, \emph{Sur les ``formules explicites'' de la th\'eorie des nombres
premiers}, Comm.~S\'em.~Math.~Univ.~Lund (1952), 252--265.
\bibitem{Yoshida1992} H.~Yoshida, \emph{On Hermitian forms attached to zeta functions}, in:
\emph{Zeta Functions in Geometry}, Adv.~Stud.~Pure Math.~\textbf{21}, Kinokuniya, Tokyo, 1992, 281--325.
\bibitem{Zhu2026} X.~Zhu, \emph{Weil positivity in compact windows: a finite reduction, certified
two-sided bounds, and a Landau--Widom decay law}, arXiv:2608.24827 (2026).
\end{thebibliography}

\end{document}
